% concatenated from main-v4.tex
%----------------------------------------------------------
% DOCUMENTCLASS & ENCODING
%----------------------------------------------------------
\documentclass[11pt,reqno]{amsart}

\usepackage[margin=1in]{geometry}
\usepackage{microtype}                     % Improves typesetting
\usepackage{subcaption}
%----------------------------------------------------------
% MATH PACKAGES
%----------------------------------------------------------
\usepackage{amsmath,amssymb,amsfonts,amsthm} % Core AMS math
\usepackage{mathrsfs,euscript}             % Script fonts
\usepackage{mathtools}
\usepackage{xcolor}
%----------------------------------------------------------
% GRAPHICS
%----------------------------------------------------------
\usepackage{graphicx}                      % Include graphics

%----------------------------------------------------------
% HYPERLINKS & CROSS-REFERENCES
%----------------------------------------------------------
\usepackage{hyperref}
\usepackage[nameinlink]{cleveref}          % Smart cross-referencing

%----------------------------------------------------------
% THEOREM ENVIRONMENTS
%----------------------------------------------------------
\theoremstyle{plain}
\newtheorem{theorem}{Theorem}[section]
\newtheorem{lemma}[theorem]{Lemma}
\newtheorem{proposition}[theorem]{Proposition}
\newtheorem{corollary}[theorem]{Corollary}
\newtheorem{conjecture}[theorem]{Conjecture}

\theoremstyle{definition}
\newtheorem{definition}[theorem]{Definition}
\newtheorem{remark}[theorem]{Remark}

\numberwithin{equation}{section}

%----------------------------------------------------------
% CUSTOM MATH OPERATORS
%----------------------------------------------------------
\DeclareMathOperator{\tr}{tr}
\DeclareMathOperator{\Id}{Id}

%----------------------------------------------------------
% SHORTCUT COMMANDS
%----------------------------------------------------------
\newcommand{\paren}[1]{\left(#1\right)}
\newcommand{\jump}[1]{\left[\!\left[ #1\right]\!\right]}
\newcommand{\D}[2]{\frac{d#1}{d#2}}
\newcommand{\PD}[2]{\frac{\partial#1}{\partial#2}}
\newcommand{\p}{\partial}
\newcommand{\PDD}[3]{\frac{\partial^{#1}{#2}}{\partial{#3}^{#1}}}
\newcommand{\lap}[1]{\Delta #1}
\newcommand{\at}[2]{\left. #1 \right|_{#2}}
\newcommand{\grad}[1]{\nabla #1}
\newcommand{\mb}[1]{\mathbf{#1}}
\newcommand{\mc}[1]{\mathcal{#1}}
\newcommand{\wh}[1]{\widehat{#1}}
\newcommand{\wt}[1]{\widetilde{#1}}
\newcommand{\bm}[1]{\boldsymbol{#1}}
\newcommand{\abs}[1]{\left\lvert #1 \right\rvert}
\newcommand{\norm}[1]{\left\lVert #1 \right\rVert}
\newcommand{\inprod}[2]{\left( #1,#2 \right)}
\newcommand{\dual}[2]{\left\langle #1,#2 \right\rangle}
\newcommand{\mbs}{{\mathbb{S}^2}}
\newcommand{\mcu}{{\mc{U}}}
\newcommand{\mbr}{{\mathbb{R}^2}}
\newcommand{\starnorm}[1]{\left| #1 \right|_*}
\newcommand{\circlenorm}[1]{\left| #1 \right|_\circ}
\newcommand{\vph}{\varphi}
\newcommand{\eqdef}{\overset{\mbox{\tiny{def}}}{=}}
\newcommand{\hsdual}[3]{ \left\langle #1,#2 \right\rangle_{H^{-#3}_{\overline S}(\Gamma), H^{#3}(S)}  }
\newcommand{\hgdual}[3]{ \left\langle #1,#2 \right\rangle_{H^{-#3}(\Gamma), H^{#3}(\Gamma)}  }

\newcommand{\R}{\mathbb{R}}
\newcommand{\C}{\mathbb{C}}
\newcommand{\N}{\mathbb{N}}
\newcommand{\Z}{\mathbb{Z}}
\newcommand{\T}{\mathbb{T}}
\newcommand{\thk}{\theta_\kappa}
\newcommand{\tk}{\bm\tau_\kappa}
\newcommand{\nk}{\bm n_\kappa}
\newcommand{\xk}{\bm X_\kappa}

\newcommand{\tka}[1]{\bm\tau_{\kappa_{#1}}}
\newcommand{\thka}[1]{\theta_{\kappa_{#1}}}
\newcommand{\nka}[1]{\bm n_{\kappa_{#1}}}
\newcommand{\xka}[1]{\bm X_{\kappa_{#1}}}

\newcommand{\ci}[1]{C^{#1}(I)}
\newcommand{\hi}[1]{H^{#1}(I)}
\newcommand{\thi}[1]{\wt H^{#1}(I)}
\newcommand{\hr}[1]{H^{#1}(\R)}


\title[Open Inextensible Filaments in Stokes Flow]
{Open Inextensible Filaments in Planar Stokes Flow: Well-Posedness, Endpoint Asymptotics, and Straightening}

\author{Han Zhou}
\address{Department of Mathematics, University of Pennsylvania, Philadelphia, PA 19104, USA}
\email{hzhou24@sas.upenn.edu}

\subjclass[2020]{Primary 35Q35, 76D07; Secondary 74F10, 35B40, 35B65.}
\keywords{Open inextensible filament, Stokes flow, supported Sobolev spaces, Wiener--Hopf factorization, endpoint asymptotics, energy threshold, exponential straightening.}

\date{\today}

\begin{document}

\begin{abstract}
We study an inextensible open filament with free ends in a planar Stokes fluid. The system reduces to a third-order nonlocal curvature equation coupled to an elliptic equation for the tension. We prove local well-posedness for nearly critical initial data in supported Sobolev spaces $\widetilde H^s$, $-1/2<s\le0$, satisfying the arc-chord condition. For positive times, we prove improved Sobolev regularity and derive a $d^{3/2}$-type expansion at each free end using Wiener--Hopf factorization, where $d$ denotes the distance to that endpoint. We further prove global existence and exponential convergence to a straight filament for sufficiently small initial data and for finite-energy initial data satisfying $E(0)<\pi^2/4$. The finite-energy result follows from an energy identity and a geometric estimate relating the bending energy to the arc-chord constant. More generally, any finite-time breakdown must be accompanied by loss of the arc-chord condition, while every global solution either converges exponentially to a straight filament or has arc-chord constants tending to zero along a sequence of times tending to infinity.
\end{abstract}

\maketitle

% \setcounter{tocdepth}{1}
\tableofcontents

\section{Introduction}



We consider a free-ended, inextensible elastic filament immersed in a two-dimensional Stokes flow.
The filament is modeled as a zero-width interface with two endpoints. It exerts singular forces derived from the bending energy on the surrounding fluid, whose flow moves and deforms the filament.
Similar setups have been used to model the motion of fibers or threads in viscous fluids, including cilia and flagella~\cite{TornbergShelley2004,EtienneLoheacSaramito2010Inextensible,BouzarthLaytonYoung2011Cellular,LiLayton2012OpenInterface}.

\subsection{Problem formulation}
\label{subsec:prob-formulation}
Let $\Gamma(t)$ denote the open filament, and suppose that the surrounding domain $\mathbb{R}^2\setminus\overline{\Gamma(t)}$ is occupied by a Stokes fluid.
The fluid velocity $\bm u$ and pressure $p$ satisfy the Stokes equations
\begin{equation}\label{eqn:intro-stokes-bulk}
    \Delta\bm u-\grad p=0,
    \quad \grad\cdot\bm u=0
    \quad\text{in }\R^2\setminus\overline{\Gamma(t)}.
\end{equation}
We impose the far-field normalization $\bm u(\bm x,t)\to0$ and
$p(\bm x,t)\to0$ as $\abs{\bm x}\to\infty$. Let
$I=(-1,1)$ and let $\bm X:\overline I\times[0,T]\to\R^2$ be a material
parametrization of the filament
$\Gamma(t)=\bm X(I,t)$. The unit tangent and clockwise unit normal are
\begin{equation}
    \bm\tau = \p_s\bm X, \quad \bm n  = Q_{\frac{\pi}{2}}^{-1}\bm \tau, \quad Q_{\alpha} =
    \begin{bmatrix}
        \cos\alpha & -\sin\alpha\\
        \sin\alpha & \cos\alpha
    \end{bmatrix},
\end{equation}
where $Q_\alpha$ denotes counterclockwise rotation through the angle $\alpha$.
We label the side toward which $\bm n$ points by $-$ and the opposite side
by $+$. For the corresponding traces $\phi^\pm$ on $\Gamma(t)$, we define the jump by
$\jump{\phi}=\phi^+-\phi^-$.
Across the filament, the velocity is continuous and the jump in
hydrodynamic traction balances the elastic bending force:
\begin{equation}\label{eqn:intro-stokes-interface}
    \jump{\bm u}=0,
    \quad
    \jump{\Sigma\bm n}
    =\bm F(\bm X)
    \quad\text{on }\Gamma(t),
\end{equation}
where $\bm F$ is the elastic force density. The fluid stress $\Sigma$ and the symmetric velocity gradient
$\grad_S\bm u$ are defined by
\begin{equation*}
    \Sigma=2\grad_S\bm u-p\Id,
    \quad
    \grad_S\bm u=\frac12\paren{\grad\bm u+\grad\bm u^T}.
\end{equation*}
The filament is transported by the fluid and satisfies the inextensibility
constraint:
\begin{equation}\label{eqn:intro-kinematic}
    \p_t\bm X(s,t)=\bm u(\bm X(s,t),t),
    \quad \abs{\p_s\bm X(s,t)}=1
    \quad\text{for }s\in\overline I,
\end{equation}
where endpoint velocities are understood through the continuous trace.
The initial filament is arclength-parametrized:
\begin{equation}\label{eqn:intro-initial}
    \bm X(s,0)=\bm X_0(s),
    \quad \abs{\p_s\bm X_0(s)}=1.
\end{equation}
To specify the bending force and boundary conditions, we introduce the bending energy:
\begin{equation}
    \mc E[\bm X]
    =\frac{1}{2}\int_I |\p_s^2\bm X|^2\,ds.
\end{equation}
The tension $\lambda$ is a Lagrange multiplier that enforces the inextensibility constraint and plays a role analogous to that of the pressure $p$ in the Stokes equations.
The constrained functional is
\begin{equation}
    \mc A[\bm X,\lambda]
    =\frac{1}{2}\int_I
    \left(|\p_s^2\bm X|^2
    +\lambda\paren{|\p_s\bm X|^2-1}\right)\,ds.
\end{equation}
Its first variation in the direction $\bm Y$ is
\begin{equation}
\begin{aligned}
    &
    \at{\dfrac{d}{d\varepsilon}
    \mc A[\bm X+\varepsilon\bm Y,\lambda]}{\varepsilon=0}
    = \int_I \paren{\p_s^4\bm X
    -\p_s\paren{\lambda\p_s\bm X}}\cdot\bm Y\,ds \\
    &\quad
    + \at{\paren{\p_s^2\bm X\cdot\p_s\bm Y}}{\p I}
    + \at{\paren{\paren{\lambda\p_s\bm X-\p_s^3\bm X}
    \cdot\bm Y}}{\p I}.
\end{aligned}
\end{equation}
The elastic force density is therefore the negative variational derivative
with respect to $\bm X$:
\begin{equation}\label{eqn:bending-force}
\begin{aligned}
    \bm F(\bm X(s))=\p_s\paren{\lambda\p_s\bm X  -\p_s^3\bm X }.
\end{aligned}
\end{equation}
The boundary terms give the free-end conditions
\begin{equation}\label{eqn:intro-free-end}
    \p_s^2\bm X=0,
    \quad \lambda\p_s\bm X-\p_s^3\bm X=0
    \quad\text{at }s=\pm1.
\end{equation}
Since $\p_t\abs{\p_s\bm X}^2=2\p_s\bm X\cdot\p_s\bm u$, the initial
condition $\abs{\p_s\bm X_0}=1$ is preserved when
$\grad_\Gamma\cdot\bm u=\bm\tau\cdot\p_s\bm u=0$ on $\Gamma$.
For smooth solutions, testing~\eqref{eqn:intro-stokes-bulk} against the fluid velocity
and using the free-end conditions~\eqref{eqn:intro-free-end} formally gives
\begin{equation}\label{eqn:energy-law}
    \frac{d}{dt} \mc E[\bm X(t)]
    = - 2\int_{\R^2} \abs{\grad_S \bm u}^2 \,d\bm x.
\end{equation}

The constraint~\eqref{eqn:intro-kinematic} preserves the arclength parametrization, making the curvature $\kappa$ a natural variable for describing the evolving shape.
The curvature formulation is derived in Subsection~\ref{subsec:curvature-formulation}.
Write the tangent as
$\bm\tau=(\cos\theta,\sin\theta)$, set $\kappa=\p_s\theta$, and introduce the
modified tension $\sigma=\lambda+\kappa^2$.  Eliminating $\bm u$ and $p$ with boundary integral representations reduces the coupled system to the curvature equation and the tension equation:
\begin{equation}\label{eqn:intro-schematic-system}
    \p_t\kappa=\Lambda\kappa+\p_s\wt N(\kappa,\sigma),
    \quad
    T_\kappa\sigma=F(\kappa),
\end{equation}
where
\begin{equation}\label{eqn:intro-principal-operators}
    \Lambda=\frac14\p_s^2\mc H\p_s,
    \quad
    T_\kappa=\frac14\p_s\mc H+R_\kappa.
\end{equation}
Here, $\mc H$ is the finite Hilbert transform on the interval $I$, $\wt N$ and $F$ collect the nonlinear terms arising from the geometry and tension, and $R_\kappa$ is the
remainder operator in the tension equation.
The free-end conditions~\eqref{eqn:intro-free-end} become
\begin{equation}\label{eqn:cur-bc}
    \kappa=0,\quad \p_s\kappa=0,\quad \sigma=0
    \quad\text{on }\p I.
\end{equation}
Solving the second equation in~\eqref{eqn:intro-schematic-system} determines $\sigma=\sigma(\kappa)$ and hence
the reduced nonlinearity
$N(\kappa)\coloneqq\wt N(\kappa,\sigma(\kappa))$.






% The free endpoints shape the analysis in three ways. The arclength
% constraint leads to a nonlocal tension equation on an interval, with
% positive first-order principal part $\frac14\p_s\mc H$. To establish
% invertibility of the tension operator, we prove injectivity using the Stokes
% energy identity for an open-arc single-layer field. The third-order
% principal part of the curvature evolution, realized on supported Sobolev
% spaces, has a boundary obstruction reflected in a characteristic $d^{3/2}$
% endpoint profile. Finally, the free-end conditions force every equilibrium
% to be straight. Together with the energy law and a geometric chord estimate,
% this rigidity provides control of the long-time dynamics.



% For arc-chord admissible initial curvature in
% $\widetilde H^s(I)$, $-\frac12<s\le0$, we prove local well-posedness in the
% continuity classes stated below, with the time-weighted smoothing associated
% with the third-order evolution. At positive times, solutions admit a precise
% $d^{3/2}$ endpoint expansion. Once the initial orientation and centroid are
% fixed, we reconstruct the filament, tension, velocity, and pressure, which
% satisfy the original coupled equations pointwise for positive times in the
% sense specified below. We also prove global existence and exponential decay
% for small data near the critical regularity, together with a
% continuation criterion and an energy gap for general data.  In particular,
% every initial curvature satisfying
% $\norm{\kappa_0}_{L^2(I)}<\pi/\sqrt2$ generates a global solution that
% straightens exponentially.

% \subsection{Related work}
% Much of the rigorous PDE theory for immersed elastic interfaces has been developed for closed interfaces. 
% The closest related results concern closed inextensible filaments, for which 
% the tension equation and the resulting evolution equation were
% studied in
% \cite{KuoLaiMoriRodenberg2023Tension,
% GarciaJuarezKuoMori2025Inextensible}. The absence of endpoints in that setting 
% permits a Fourier characterization of the operators and their inverses. A fixed-point argument in the tangent angle formulation gives local well-posedness in subcritical H\"{o}lder spaces.

% For extensible interfaces, the Peskin problem describes motion driven by elastic stretching forces, without bending or an inextensibility
% constraint. Local well-posedness and stability near
% circular equilibria were established in
% \cite{MoriRodenbergSpirn2019Peskin,LinTong2019StokesIB}. Subsequent work addresses
% nonlinear tension laws, rough or critical data, viscosity contrast, and
% three-dimensional interfaces
% \cite{CameronStrain2024CriticalPeskin,ChenNguyen2023Bdot,
% GarciaJuarezHaziot2025GeneralTension,ChenHuNguyen2026Schauder,
% GarciaJuarezMoriStrain2023ViscosityContrast,
% GarciaJuarezKuoMoriStrain2025Peskin3D,
% GarciaJuarezHaziotKuoMoriZhou2026Peskin3DStability}.
% Tong and Wei~\cite{TongWei2024GeometricPeskin} study the geometry of the evolving string and prove global existence under a condition on its initial shape, without a smallness assumption on the stretching.

% Open filaments also arise when one end is anchored to a rigid wall and the other is free, as in models of cilia.
% Li and Layton~\cite{LiLayton2012OpenInterface} develop a numerical method for Stokes flow driven by an open immersed interface and apply it to a beating cilium attached to a wall.
% The wall introduces hydrodynamic interactions in addition to the effects of the filament endpoints.
% In a related geometry, Kandalam and Spirn~\cite{KandalamSpirn2026AnchoredPeskin} study the anchored Peskin problem in a half-plane, with both endpoints attached to a no-slip wall, and prove local well-posedness and instantaneous smoothing.
% The filament considered here has two free endpoints and evolves in the whole plane.


% The main analytical difficulty comes from the open-interface geometry of the Stokes problem, which can produce endpoint singularities even when the geometry and data are smooth.
% Related singularities occur in screen problems in scattering theory~\cite{Hayashi1973OpenBoundary,Stephan1987ScreenProblems} and crack problems in continuum mechanics~\cite{DuduchavaWendland1995CrackProblems,CostabelDaugeDuduchava2003CrackAsymptotics}.
% Here, the Stokes field is coupled to the filament evolution, so its singular behavior also affects the regularity of the filament.
% In contrast to the closed-curve case~\cite{GarciaJuarezKuoMori2025Inextensible}, parabolic smoothing need not give smoothness up to the endpoints.
% Epstein, Mori, and Zhou~\cite{EpsteinMoriZhou2026BulkSurface} study an elliptic bulk--surface coupled Laplace system on a fixed open arc.
% They reduce the system to an integro-differential equation on the arc and use Wiener--Hopf analysis to derive endpoint asymptotics.

% This endpoint behavior is more transparent in the curvature formulation.
% On the supported Sobolev spaces used here, the operators $\frac{1}{4}\p_s\mc H$ and $\Lambda$ are the restricted fractional Laplacians $\frac{1}{4}\paren{-\Delta}^{1/2}$ and $-\frac{1}{4}\paren{-\Delta}^{3/2}$, respectively: they act on the functions supported on $\overline{I}$, and their outputs are then restricted to $I$.
% The zero exterior condition allows fractional-power endpoint profiles that limit Sobolev regularity even for smooth data.
% Boundary regularity for fractional powers between zero and one is studied by Ros-Oton and Serra~\cite{RosOtonSerra2014BoundaryRegularity}.
% Grubb~\cite{Grubb2015FractionalLaplacians} develops a boundary theory that also covers higher fractional powers, while Abatangelo et al.~\cite{AbatangeloEtAl2019HalfSpace} give explicit solution formulas and boundary estimates for positive powers of the Laplacian in a half-space.
% For the curvature equation, the leading endpoint profile is $d^{3/2}$, where $d$ is the distance to the endpoint; we isolate this term below using Wiener--Hopf factorization, following the framework of~\cite{eskin1981boundary}.


% Free-ended filaments and rods also appear in geometric-flow and reduced
% hydrodynamic models, including curve straightening, classical
% elastohydrodynamics, discrete-link models, and PDE models for immersed rods
% \cite{Oelz2011CurveStraightening,MoriOhm2023Elastohydrodynamics,
% AlougesLefebvreLepotLevillainMoreau2025NLink,
% AlbrittonOhm2025RodsInFlows}. Related nonlocal and slender-body models are
% studied in
% \cite{TornbergShelley2004,MoriOhmSpirn2020SBT,
% MoriOhmSpirn2020SBTFreeEnds,Ohm2026NonlocalCurve,
% Ohm2026FreeBoundary3D,Ohm2025SlenderBodyFreeBoundary}. 
% Unlike the finite-radius filaments considered in three-dimensional slender-body theory, the present model treats the filament as a zero-width open interface in a planar Stokes fluid.

\subsection{Related work}

Rigorous PDE theory for immersed elastic interfaces has largely concerned
closed interfaces. For closed inextensible filaments, the tension and
evolution equations are studied in
\cite{KuoLaiMoriRodenberg2023Tension,
GarciaJuarezKuoMori2025Inextensible}.
Without endpoints, the operators and their inverses admit a Fourier
characterization; a fixed-point argument in the tangent angle gives local
well-posedness in subcritical H\"{o}lder spaces.

The Peskin problem concerns extensible interfaces driven by stretching,
without bending or inextensibility. Local well-posedness and stability
near circular equilibria are proved in
\cite{MoriRodenbergSpirn2019Peskin,LinTong2019StokesIB}.
Subsequent work treats nonlinear tension laws, rough or critical data, viscosity contrast, and three-dimensional interfaces
\cite{CameronStrain2024CriticalPeskin,ChenNguyen2023Bdot,
GarciaJuarezHaziot2025GeneralTension,ChenHuNguyen2026Schauder,
GarciaJuarezMoriStrain2023ViscosityContrast,
GarciaJuarezKuoMoriStrain2025Peskin3D,
GarciaJuarezHaziotKuoMoriZhou2026Peskin3DStability}.
A geometric condition on the initial configuration gives global existence
without a smallness assumption on stretching in
\cite{TongWei2024GeometricPeskin}.

For open interfaces, a numerical Stokes method applied to a beating cilium
attached to a rigid wall is developed in
\cite{LiLayton2012OpenInterface}.
Local well-posedness and instantaneous smoothing for a half-plane problem
with both endpoints attached to a no-slip wall are proved in
\cite{KandalamSpirn2026AnchoredPeskin}.
The wall adds hydrodynamic interactions absent from the present whole-plane
problem with two free endpoints.

Free-ended filaments and rods also appear in geometric-flow and reduced hydrodynamic models, including curve straightening, classical elastohydrodynamics, discrete-link models, and PDE models for immersed rods
\cite{Oelz2011CurveStraightening,MoriOhm2023Elastohydrodynamics,
AlougesLefebvreLepotLevillainMoreau2025NLink,
AlbrittonOhm2025RodsInFlows}.
Related nonlocal and slender-body models are studied in
\cite{TornbergShelley2004,MoriOhmSpirn2020SBT,
MoriOhmSpirn2020SBTFreeEnds,Ohm2026NonlocalCurve,
Ohm2026FreeBoundary3D,Ohm2025SlenderBodyFreeBoundary}.
Unlike the finite-radius filaments considered in three-dimensional slender-body theory, the present model treats the filament as a zero-width open interface in a planar Stokes fluid.

The main analytical difficulty specific to the open-interface geometry is the behavior near the endpoints. Singularities can occur there even when the geometry and data are smooth. Related phenomena arise for screen problems in scattering theory
\cite{Hayashi1973OpenBoundary,Stephan1987ScreenProblems}
and for crack problems in continuum mechanics
\cite{DuduchavaWendland1995CrackProblems,CostabelDaugeDuduchava2003CrackAsymptotics}.
In the present problem, the Stokes flow is coupled to the filament evolution, so the endpoint behavior of the fluid also enters the regularity for the filament. In contrast to the closed-curve setting
\cite{GarciaJuarezKuoMori2025Inextensible},
parabolic smoothing does not in general give smoothness up to the free endpoints. A related elliptic bulk--surface system on a fixed open arc was studied in
\cite{EpsteinMoriZhou2026BulkSurface},
where reduction to an integro-differential equation and Wiener--Hopf analysis were used to derive endpoint asymptotics.

This endpoint behavior is particularly transparent in the principal operators~\eqref{eqn:intro-principal-operators}. On the supported Sobolev spaces used here, the operators $\frac{1}{4}\p_s\mc H$ and $\Lambda$ are the restricted fractional Laplacians
$\frac{1}{4}\paren{-\Delta}^{1/2}$ and
$-\frac{1}{4}\paren{-\Delta}^{3/2}$, respectively: they act on functions supported on $\overline{I}$, and their outputs are then restricted to $I$. The zero exterior condition allows fractional-power endpoint profiles that limit Sobolev regularity even for smooth data. Boundary regularity for fractional powers between zero and one is studied in
\cite{RosOtonSerra2014BoundaryRegularity},
while higher fractional powers are treated in
\cite{Grubb2015FractionalLaplacians}.
Explicit solution formulas and boundary estimates for positive powers of the Laplacian in a half-space are given in
\cite{AbatangeloEtAl2019HalfSpace}.
For the curvature equation considered here, we isolate this term in~\eqref{eqn:main-endpoint-split} by Wiener--Hopf factorization, following the framework of
\cite{eskin1981boundary}.


\subsection{Main results}
We define the supported Sobolev space
$\widetilde H^s(I)=\overline{C_c^\infty(I)}^{\,H^s(\R)}$, with the inherited $H^s(\R)$ norm. Let $\theta_\kappa$ be the tangent angle satisfying
$\p_s\theta_\kappa=\kappa$ and $\theta_\kappa(-1)=0$, and set
$\bm X_\kappa(s)=\int_{-1}^s(\cos\theta_\kappa,\sin\theta_\kappa)(\eta)\,d\eta$.
The arc-chord constant is defined as
\begin{equation}\label{eqn:intro-arc-chord}
    \abs{\kappa}_*
    =\inf_{s,\eta\in I,\ s\ne\eta}
    \frac{\abs{\bm X_\kappa(s)-\bm X_\kappa(\eta)}}{\abs{s-\eta}}.
\end{equation}
We assume the arc-chord condition $\abs{\kappa_0}_*>0$, which excludes degeneration and self-intersection of the filament.
We also introduce $\omega_\Lambda
    =\sup\left\{\operatorname{Re}z:z\in\sigma(\Lambda)\right\}<0.$ 
The function spaces and curvature formulation are described in detail in
Section~\ref{sec:formulation}.

\begin{theorem}[Local well-posedness]\label{thm:main-local}
Fix $0<\delta\le\frac12$ and let the initial datum
$\kappa_0\in\thi{-\frac12+\delta}$ satisfy $\abs{\kappa_0}_*>0$.
There exist $T=T(\kappa_0)>0$ and a unique mild solution satisfying
\begin{equation}\label{eqn:main-local-class}
    \kappa \in C\paren{[0,T];\thi{-\frac12+\delta}} \cap C\paren{(0,T];\thi{\frac32}}, \quad \sup_{0<t\le
    T}t^{\frac{2-\delta}{3}} \norm{\kappa(t)}_{\hr{\frac32}} <\infty.
\end{equation}
For this $T$, there are $r_0,C>0$, depending only on
$\delta,\kappa_0$, such that each $\vph_0\in\thi{-\frac12+\delta}$ with
$\norm{\vph_0-\kappa_0}_{\hr{-\frac12+\delta}}\le r_0$ has a unique solution
in the class~\eqref{eqn:main-local-class}, with
\begin{equation}\label{eqn:main-local-lipschitz}
\begin{aligned}
    &\sup_{0\le t\le T}
    \norm{\kappa(t;\vph_0)-\kappa(t;\kappa_0)}_{
        \hr{-\frac12+\delta}}\\
    &\quad+
    \sup_{0<t\le T}t^{\frac{2-\delta}{3}}
    \norm{\kappa(t;\vph_0)-\kappa(t;\kappa_0)}_{\hr{\frac32}}
    \le C\norm{\vph_0-\kappa_0}_{\hr{-\frac12+\delta}}.
\end{aligned}
\end{equation}
\end{theorem}

The tension operator $T_\kappa:\thi{\frac12}\to\hi{-\frac12}$ is a
compact perturbation of $\frac14\p_s\mc H$. Uniqueness from the Stokes
energy identity and the Fredholm alternative give invertibility of
$T_\kappa$ in~\eqref{eqn:ten-det}, eliminating $\sigma$ from~\eqref{eqn:intro-schematic-system}.
We then solve the reduced curvature equation~\eqref{eqn:kap_evo} using the
analytic-semigroup smoothing estimate~\eqref{eqn:lambda-semigroup-est}
and a fixed-point argument in the time-weighted supported-space
norm~\eqref{eqn:XT-norm-intro}. The fixed-point neighborhood retains a
common arc-chord lower bound, while nonlinear commutator and boundary
integral estimates control both equations.

At positive times, the curvature gains Sobolev regularity and admits the
following endpoint expansion.

\begin{theorem}[Positive-time regularity and endpoint asymptotics]
\label{thm:main-regularity}
Let $\kappa$ be the solution of Theorem~\ref{thm:main-local}.
For every $0<\tau<T$ and $0\le\varepsilon<\frac12$,
\begin{equation}\label{eqn:main-positive-time-regularity}
    \kappa\in C\paren{[\tau,T];\thi{\frac32+\varepsilon}}
    \cap C^1\paren{[\tau,T];\hi{-\frac32+\varepsilon}}.
\end{equation}
Moreover, let $\chi_\pm$ be smooth cutoffs with disjoint supports, equal
to one near $s=\pm1$, respectively. For each $t\in(0,T]$, there exist
$c_\pm(t)\in\R$ and
$\kappa_r(t)\in\thi{\frac52+\varepsilon}$ such that $\kappa$ has the form
\begin{equation}\label{eqn:main-endpoint-split}
    \kappa(s,t)
    =c_-(t)\chi_-(s)d_-(s)^{\frac32}
     +c_+(t)\chi_+(s)d_+(s)^{\frac32}
     +\kappa_r(s,t),
\end{equation}
where $d_-(s)=(s+1)_+$ and $d_+(s)=(1-s)_+$, with $r_+=\max\{r,0\}$.
\end{theorem}

Since $d^{3/2}\in H^{2-\eta}$ for every $\eta>0$ but $d^{3/2}\notin H^2$, each nonzero coefficient $c_\pm(t)$ is the precise obstruction to $H^2$ regularity at that endpoint.

We bootstrap the spatial regularity at positive times and use the Duhamel
formula to prove time differentiability. With a further gain in regularity,
we write the curvature equation as a supported third-order elliptic problem.
After localizing near the endpoints, we apply Wiener--Hopf factorization
to identify the leading $d^{3/2}$ free-end profile of the restricted operator.

This positive-time regularity suffices to reconstruct the physical variables
and verify the original Stokes--filament equations pointwise in the regions
specified below.

\begin{corollary}[Reconstruction of a classical solution]\label{cor:main-reconstruction}
Let $\kappa$ be as in Theorem~\ref{thm:main-regularity}.
Given an initial tangent angle $\theta_{*,0}\in\R$ and
centroid $\overline{\bm X}_0\in\R^2$, $\kappa$ uniquely
determines a reconstruction $(\bm X,\bm u,p,\lambda)$ solving
the Stokes--filament system
\eqref{eqn:intro-stokes-bulk}--\eqref{eqn:intro-kinematic} and \eqref{eqn:intro-free-end}
classically for $t\in(0,T]$. The initial filament $\bm X_0$ is determined
by $\kappa_0$, $\theta_{*,0}$, and $\overline{\bm X}_0$.
The fluid fields satisfy the prescribed far-field normalization.
The bulk equations hold pointwise off $\overline{\Gamma(t)}$, the
interface conditions hold pointwise along the filament
interior, and the free-end conditions hold at $s=\pm1$.
The kinematic equation and inextensibility constraint hold
throughout $\overline I$, with the velocity at the endpoints
understood through its continuous trace.
\end{corollary}

% The angle and centroid equations reconstruct $\bm X$, while the
% single-layer potential reconstructs $(\bm u,p)$ and the tension equation
% recovers $\lambda$.  Positive-time Sobolev estimates for $\bm X$ and
% $\lambda$, followed by Sobolev embedding and standard Stokes layer-potential
% regularity, show that the original coupled equations hold classically away
% from the two endpoints and that the free-end conditions hold pointwise.

For $0<a<b$ in the interval of existence, the reconstructed solution satisfies the energy identity:
\begin{equation}\label{eqn:intro-energy-identity}
    \mc E(b)
    +2\int_a^b\int_{\R^2}
        \abs{\grad_S\bm u(\bm x,t)}^2\,d\bm x\,dt
    =\mc E(a),
    \quad \mc E(t)=\frac12\norm{\kappa(t)}_{L^2(I)}^2.
\end{equation}
If $\kappa_0\in L^2(I)$, then~\eqref{eqn:intro-energy-identity} extends to
$a=0$, with $\mc E(0)=\frac12\norm{\kappa_0}_{L^2(I)}^2$.
The $L^2$ control in~\eqref{eqn:intro-energy-identity} is stronger than the
critical norm and supplies the bound needed for continuation.

The next theorem gives small-data global existence in part~(1).
Part~(2) concerns continuation and the alternatives between straightening
and arc-chord degeneration, with the energy gap~\eqref{eqn:main-energy-gap}
distinguishing the two cases. Part~(3) gives an explicit finite-energy threshold.

\begin{theorem}[Global well-posedness and long-time behavior]
\label{thm:main-global}
Fix $0<\delta\le\frac12$ and let
$\kappa_0\in\thi{-\frac12+\delta}$ satisfy $\abs{\kappa_0}_*>0$.
Let $\kappa$ be the unique maximal continuation of the solution in
Theorem~\ref{thm:main-local}, with lifetime $\tau_{\mathrm{max}}$.
\begin{enumerate}
    \item \emph{Nearly critical small data.}
    For each $\omega\in(\omega_\Lambda,0)$, there are $M_0,C>0$, depending
    only on $\delta,\omega$, such that, if
    $\norm{\kappa_0}_{\hr{-\frac12+\delta}}\le M_0$, then the solution
    is global and
    \begin{equation}\label{eqn:main-global-decay}
        \norm{\kappa(t)}_{\hr{-\frac12+\delta}} +\min\left\{t^{\frac{2-\delta}{3}},1\right\}
        \norm{\kappa(t)}_{\hr{\frac32}} \le Ce^{t\omega} \norm{\kappa_0}_{\hr{-\frac12+\delta}}, \quad t>0.
    \end{equation}
    The same smallness condition implies $\abs{\kappa_0}_*>0$.

    \item \emph{Continuation and long-time alternatives.}
    Let $\mc E_{\mathrm{ac}}=\frac{\pi^2}{4}$. The limiting energy exists and satisfies
    \begin{equation}\label{eqn:main-energy-gap}
        \mc E_*:=\lim_{t\uparrow\tau_{\mathrm{max}}}\mc E(t)
        \in\{0\}\cup[\mc E_{\mathrm{ac}},\infty),
    \end{equation}
    where the limit is taken as $t\to\infty$ when
    $\tau_{\mathrm{max}}=\infty$.
    If $\tau_{\mathrm{max}}<\infty$, then
    $\mc E_*\ge\mc E_{\mathrm{ac}}$ and
    \begin{equation*}
        \lim_{t\uparrow\tau_{\mathrm{max}}}\abs{\kappa(t)}_*=0.
    \end{equation*}
    In particular, $\mc E_*=0$ implies global existence.

    Every global solution satisfies exactly one of the following alternatives:
    \begin{enumerate}
        \item If $\mc E_*=0$, then $\kappa(t)\to0$ in $L^2(I)$ and
        $\abs{\kappa(t)}_*\to1$. For each
        $\omega\in(\omega_\Lambda,0)$, there is $T_\omega>0$ such that
        \begin{equation*}
            \norm{\kappa(T_\omega+t)}_{L^2(I)}
            +\min\{t^{1/2},1\}
            \norm{\kappa(T_\omega+t)}_{\hr{\frac32}}
            \le Ce^{t\omega}\norm{\kappa(T_\omega)}_{L^2(I)},
            \quad t>0,
        \end{equation*}
        where $C$ depends only on $\omega$. In particular, the filament
        straightens exponentially modulo rigid motions.

        \item If $\mc E_*>0$, then $\mc E_*\ge\mc E_{\mathrm{ac}}$ and
        \begin{equation*}
            \liminf_{t\to\infty}\abs{\kappa(t)}_*=0.
        \end{equation*}
    \end{enumerate}

    \item \emph{Finite-energy threshold.}
    If $\kappa(t_0)\in L^2(I)$ and
    $\norm{\kappa(t_0)}_{L^2(I)}<\pi/\sqrt2$ for some
    $t_0\in[0,\tau_{\mathrm{max}})$, then the solution is global,
    $\mc E_*=0$, and
    \begin{equation}\label{eqn:main-post-threshold-arc-chord}
        \abs{\kappa(t)}_*
        \ge\cos\sqrt{\mc E(t_0)}>0,
        \quad t\ge t_0.
    \end{equation}
    Thus the exponential straightening in part~\emph{(2)} applies.
    In particular, every $\kappa_0\in L^2(I)$ with
    $\norm{\kappa_0}_{L^2(I)}<\pi/\sqrt2$ generates such a global solution.
\end{enumerate}
\end{theorem}

For small data, the semigroup estimate~\eqref{eqn:lambda-semigroup-est}
and quadratic bound~\eqref{eqn:global-nonlinear-bounds} give a contraction
uniform in time in the exponentially weighted space~\eqref{eqn:Y-intro};
see~\eqref{eqn:global-map-bounds}.

For general data, the $L^2$ bound away from $t=0$
in~\eqref{eqn:intro-energy-identity} and Theorem~\ref{thm:main-local} allow continuation unless the arc-chord
condition is lost. The total-turning estimate~\eqref{eqn:turning-chord-bound}
and its energy consequence~\eqref{eqn:energy-arc-chord-bound} exclude
this degeneration below $\mc E_{\mathrm{ac}}=\pi^2/4$.
For a global solution with a uniform positive arc-chord bound, we use
compactness of time translates to obtain a limiting solution with constant
energy. This solution has zero Stokes dissipation, and the rigidity argument
using the free-end conditions~\eqref{eqn:intro-free-end} forces it to be
straight. Thus $\mc E_*<\mc E_{\mathrm{ac}}$ implies
$\mc E_*=0$, so no positive limiting energy lies below the threshold.
When the energy tends to zero, we apply the small-data result at a later
time to obtain exponential straightening.

The reconstructed filament then converges as follows:
\begin{corollary}[Convergence to a straight filament]
\label{cor:main-straightening}
Let $\bm X$ and $\theta$ be reconstructed from a solution in
Theorem~\ref{thm:main-global} with $\mc E_*=0$.
There exist $\theta_\infty\in\R$ and
$\overline{\bm X}_\infty\in\R^2$ such that for every
$\omega\in(\omega_\Lambda,0)$ there is $C_\omega>0$ satisfying
\begin{equation*}
    |\bar\theta(t)-\theta_\infty|
    +|\overline{\bm X}(t)-\overline{\bm X}_\infty|
    +\norm{\bm X(\cdot,t)-\bm X_\infty}_{C^1(\overline I)}
    \le C_\omega e^{\omega t},\quad t\ge0.
\end{equation*}
Here $\bm X_\infty(s)=\overline{\bm X}_\infty
+s(\cos\theta_\infty,\sin\theta_\infty)$.
\end{corollary}

\begin{remark}
The rescaling $\kappa_\rho(s,t)=\rho\kappa(\rho s,\rho^3t)$ leaves the homogeneous $\dot H^{-1/2}$ norm invariant when the interval is rescaled accordingly. The local well-posedness and small-data results are nearly critical, since they hold at subcritical regularities arbitrarily close to the scaling-critical index.
\end{remark}

\subsection{Organization of the paper}
The remainder of the paper is organized as follows.
Section~\ref{sec:formulation} derives the curvature formulation and introduces
the supported Sobolev spaces and semigroup estimates. It also proves the
half-line Wiener--Hopf estimates and their interval counterparts, and
collects the commutator and boundary integral estimates.
Section~\ref{sec:tension} solves the tension constraint and estimates the
tension and its differences. Section~\ref{sec:local-theory} derives the
curvature nonlinearity estimates and proves local well-posedness.
Section~\ref{sec:regularity-theory} proves positive-time regularity and the
endpoint expansion. It then uses this regularity to reconstruct a solution of
the original Stokes--filament system. Section~\ref{sec:global-theory} uses
small-data stability, energy dissipation, and geometric control to prove the
continuation criterion, energy gap, and long-time alternatives.

The technical estimates for the nonlinear operators are collected in
Appendix~\ref{app:operator-estimates}.


\section{Preliminaries}
\label{sec:formulation}
\label{sec:prelim}

Subsection~\ref{subsec:curvature-formulation} derives the curvature--tension
system and isolates its principal operators.
Subsections~\ref{subsec:notation}--\ref{subsec:linear-theory} establish the
supported Sobolev and geometric framework, then prove coercivity and
semigroup estimates for the principal operators.
The Wiener--Hopf estimates in Subsection~\ref{subsec:endpoint-factorization}
give the endpoint expansions used later at the free ends.
Subsections~\ref{subsec:commutator-estimates}--\ref{sec:divided-diff}
control the lower-order nonlinear terms through commutator and
boundary integral estimates.

\subsection{Curvature formulation}
\label{subsec:curvature-formulation}
We derive the curvature--tension system from
\eqref{eqn:intro-stokes-bulk}--\eqref{eqn:intro-initial} and~\eqref{eqn:intro-free-end}.
Similar reductions have been used for related curve evolution
problems~\cite{Oelz2011CurveStraightening,MoriOhm2023Elastohydrodynamics,
Ohm2026NonlocalCurve}.
Let $G$ be the two-dimensional Stokeslet:
\begin{equation}\label{eqn:stokeslet-definition}
    G_{ij}(\bm x) = \frac{1}{4\pi}\paren{-\log|\bm x|\delta_{ij} + \frac{x_i x_j}{|\bm x|^2}}, \quad i,j = 1,2.
\end{equation}
Writing $\Delta\bm X=\bm X(s,t)-\bm X(\eta,t)$, the filament evolution takes
the form
\begin{equation}\label{eqn:filament-boundary-evolution}
    \p_t\bm X=\int_I G(\Delta\bm X)\p_\eta\paren{\lambda\p_\eta\bm X  -\p_\eta^3\bm X}\,d\eta,\quad \abs{\p_s\bm X} \equiv 1,
\end{equation}
with initial condition $\bm X(\cdot,0)=\bm X_0(\cdot)$ and boundary
conditions~\eqref{eqn:intro-free-end}.
We write
\begin{equation}
    \bm\tau(s)=\bm e_r(\theta(s))
    =\paren{\cos\theta(s),\sin\theta(s)}^T.
\end{equation}
Set $\kappa=\p_s\theta$ and $\sigma=\kappa^2+\lambda$. Then
\begin{equation}
    \lambda\p_s\bm X-\p_s^3\bm X
    =
    \paren{\lambda+\kappa^2}\bm\tau+\p_s\kappa\,\bm n
    =
    \sigma\bm\tau+\p_s\kappa\,\bm n .
\end{equation}
Since $\abs{\p_s\bm X}=1$ and $\p_s^2\bm X=-\kappa\bm n$, the bending
energy is
\begin{equation}\label{eqn:curvature-energy}
    \mc E[\bm X]
    =\frac12\norm{\kappa}_{L^2(I)}^2
    =\frac12\norm{\kappa}_{L^2(\R)}^2,
\end{equation}
where the $L^2(\R)$ norm for $\kappa$ is taken after zero extension from $I$.
For a time-dependent reconstructed curve, we write
$\mc E(t)\coloneqq\mc E[\bm X(t)]$.

The curvature determines the filament up to a rotation and a
translation. Set
\begin{equation}\label{eqn:normalized-geometry}
\begin{aligned}
    \thk(s)&=J\kappa(s)\coloneqq\int_{-1}^s\kappa(\eta)\,d\eta,
    &\tk(s)&=\paren{\cos\thk(s),\sin\thk(s)}^T,\\
    \nk(s)&=Q_{\frac{\pi}{2}}^{-1}\tk(s),
    &\xk(s)&=\int_{-1}^s\tk(\eta)\,d\eta.
\end{aligned}
\end{equation}
Since $\p_s\theta=\kappa$,
$\theta(s)=\theta_*+\thk(s)$, where the constant
$\theta_*=\theta(-1)$ is the tangent angle at the left endpoint. Hence
\begin{equation}
    \Delta\bm X=Q_{\theta_*}\Delta\xk,
    \quad
    \bm\tau=Q_{\theta_*}\tk,
    \quad
    \bm n=Q_{\theta_*}\nk.
\end{equation}
Here and below,
$\Delta\xk=\xk(s)-\xk(\eta)$.

In angle variables, the evolution equation and the inextensibility constraint
take the form
\begin{subequations}\label{eqn:theta-eqn}
\begin{align}
    \p_t\theta &= - \bm n(s)\cdot\p_s\int_I G(\Delta\bm X)\p_\eta\paren{ \p_\eta^2\theta\bm n+\sigma\bm \tau } \,d\eta ,\\
     0&= \bm\tau(s) \cdot\p_s\int_I G(\Delta\bm X)\p_\eta\paren{ \p_\eta^2\theta \bm n + \sigma\bm \tau } \,d\eta.
\end{align}
\end{subequations}
Differentiating the first equation in~\eqref{eqn:theta-eqn} with respect to $s$ and using the rotational
invariance of the Stokeslet~\eqref{eqn:stokeslet-definition}, $G(Q\bm x)=QG(\bm x)Q^T$, gives the
curvature--tension system
\begin{subequations}\label{eqn:kappa-eqn}
\begin{align}
    \p_t\kappa
    &= -\p_s\paren{\nk(s)\cdot\p_s\int_I
    G(\Delta\xk)\p_\eta\paren{
    \p_\eta\kappa\nk+\sigma\tk}\,d\eta},\\
    0
    &= \tk(s)\cdot\p_s\int_I
    G(\Delta\xk)\p_\eta\paren{
    \p_\eta\kappa\nk+\sigma\tk}\,d\eta,
\end{align}
\end{subequations}
subject to the endpoint conditions~\eqref{eqn:cur-bc}.
The system~\eqref{eqn:kappa-eqn} is independent of
$\theta_*$.

The remaining rotation and translation are recovered after solving
\eqref{eqn:kappa-eqn}: the mean angle
$\bar\theta=\frac12\int_I\theta(\eta)\,d\eta$ and centroid
$\overline{\bm X}=\frac12\int_I\bm X(\eta)\,d\eta$ satisfy
\begin{equation}\label{eqn:theta-bar}
    \p_t\bar\theta
    = -\frac{1}{2}\int_I \bm n(s)\cdot\p_s\int_I
    G(\Delta\bm X)\p_\eta\paren{ \p_\eta\kappa\bm n+\sigma\bm \tau }
    \,d\eta \,ds,
\end{equation}
and
\begin{equation}\label{eqn:X-bar}
    \p_t \overline{\bm X}
    = \frac{1}{2}\int_I \int_I
    G(\Delta\bm X)\p_\eta\paren{ \p_\eta\kappa\bm n+\sigma\bm \tau }
    \,d\eta ds.
\end{equation}
An initial tangent angle $\theta_{*,0}$ at the left endpoint and
centroid $\overline{\bm X}_0$ give the initial values
\begin{equation}\label{eqn:reconstruction-initial-means}
    \bar\theta(0)
    =\theta_{*,0}
    +\frac12\int_I(1-\eta)\kappa_0(\eta)\,d\eta,
    \quad
    \overline{\bm X}(0)=\overline{\bm X}_0.
\end{equation}
At each time, the tangent angle at the left endpoint is
\begin{equation}\label{eqn:reconstruction-endpoint-angle}
    \theta_*=\bar\theta-\frac12\int_I(1-\eta)\kappa(\eta)\,d\eta.
\end{equation}
The curve is then reconstructed from
\begin{equation}\label{eqn:curve-reconstruction}
    \bm X(s)
    =\overline{\bm X}
    -\frac12\int_I(1-\eta)\bm e_r(\theta(\eta))\,d\eta
    +\int_{-1}^s\bm e_r(\theta(\eta))\,d\eta.
\end{equation}

Returning to the curvature--tension system~\eqref{eqn:kappa-eqn}, we separate
the straight-filament singular part, which yields $\frac14\p_s\mc H$
and $\Lambda=\frac14\p_s^2\mc H\p_s$, from the lower-order geometric
remainder estimated in Subsection~\ref{sec:divided-diff}. Define
\begin{equation}\label{eqn:R-kernel}
    r(s,\eta;\kappa) = \log\abs{\frac{\Delta\xk}{s-\eta} } \Id - \frac{ \Delta\xk \otimes \Delta\xk }{|\Delta\xk |^2}.
\end{equation}
Componentwise,
\begin{equation}\label{eqn:stokeslet-splitting}
    G_{ij}(\Delta\xk)
    =-\frac{1}{4\pi}\log\abs{s-\eta}\,\delta_{ij}
    -\frac{1}{4\pi}r_{ij}(s,\eta;\kappa).
\end{equation}
Let $\mc H_{\R}$ denote the whole-line Hilbert transform and write
$\mc H=\mc H_I$ for its restriction to $I$ on functions $f$ supported in $I$:
$\mc H_I f=(\mc H_{\R}f)|_I$. For $f\in C_c^\infty(I)$,
\begin{equation}\label{eqn:finite-hilbert-definition}
    \mc Hf(s) = \frac{1}{\pi}\operatorname{p.v.} \int_I \frac{f(\eta)}{s-\eta} \,d\eta.
\end{equation}
We define the commutator by
\begin{equation}
    [A,B] f \coloneqq A(Bf)-B(Af).
\end{equation}
For $\bm b=(b_1,b_2)$ and $\bm q=(q_1,q_2)$, we use the componentwise
commutator and its contraction:
\begin{equation*}
    [A,\bm b]g\coloneqq([A,b_j]g)_{j=1}^2,
    \quad
    [A,\bm b]\cdot\bm q\coloneqq\sum_{j=1}^2[A,b_j]q_j.
\end{equation*}


The tension equation in~\eqref{eqn:kappa-eqn} is equivalent to
\begin{equation}\label{eqn:ten-det}
    T_\kappa\sigma=F(\kappa),
\end{equation}
where $T_\kappa$ is the tension operator
in~\eqref{eqn:intro-principal-operators} and $F(\kappa)$ is its curvature
forcing. The regular part $R_\kappa$ and the forcing are
\begin{equation}
\begin{aligned}
    R_\kappa\sigma
    &= \frac{1}{4}\tk \cdot [\p_s \mc H, \tk]\sigma
    - \frac{1}{4\pi} \tk\cdot\p_s\int_I
    \p_\eta r(s,\eta;\kappa)\sigma(\eta)\tk(\eta)\,d\eta,
\end{aligned}
\label{eqn:Rop-def}
\end{equation}
and
\begin{equation}\label{eqn:F-kappa}
\begin{aligned}
    F(\kappa)
    &= \frac{1}{4}[\p_s\mc H, \tk]\cdot(\p_s\kappa\nk)
    + \frac{1}{4\pi}\tk\cdot\p_s\int_I
    \p_\eta r(s,\eta;\kappa)\p_\eta\kappa\nk(\eta)\,d\eta.
\end{aligned}
\end{equation}

Before eliminating the tension, the curvature nonlinearity in~\eqref{eqn:intro-schematic-system} is
\begin{equation}
\begin{aligned}
    \wt N(\kappa,\sigma)
    &= -\frac14[\p_s\mc H,\nk]\cdot
    \paren{\p_s\kappa\nk+\sigma\tk}\\
    &\quad -\frac{1}{4\pi}\nk(s)\cdot\p_s\int_I
    \p_\eta r(s,\eta;\kappa)
    \paren{\p_\eta\kappa\nk+\sigma\tk}\,d\eta.
\end{aligned}
\label{eqn:N-def}
\end{equation}
Section~\ref{sec:tension} proves that the tension constraint~\eqref{eqn:ten-det}
can be solved for $\sigma=\sigma(\kappa)$. Substituting this solution defines
$N(\kappa)=\wt N(\kappa,\sigma(\kappa))$ and reduces
\eqref{eqn:kappa-eqn} to a single nonlocal parabolic equation:
\begin{equation}\label{eqn:kap_evo}
    \p_t\kappa=\Lambda\kappa+\p_s\paren{N(\kappa)}.
\end{equation}

Given an initial datum $\kappa_0$, define the Duhamel map by
\begin{equation}\label{eqn:mild-def}
    (\Psi_{\kappa_0}\kappa)(t)
    \coloneqq
    e^{t\Lambda}\kappa_0
    +\int_0^t e^{(t-\tau)\Lambda}
    \p_s\paren{N(\kappa(\tau))}\,d\tau.
\end{equation}
The precise solution class and the definition of a mild solution are given
in Subsection~\ref{subsec:local-proof}, where the weighted space
$X_T^{(\delta)}$ is introduced.

\subsection{Supported Sobolev spaces and geometric estimates}
\label{subsec:notation}

\emph{Supported spaces and endpoint conditions.}
The Fourier transform on $\R$ is defined by
\begin{equation}
    \wh f(\xi) = \frac{1}{\sqrt{2\pi}}\int_{\R} e^{-is\xi}f(s)\,ds.
\end{equation}
Write $\langle\xi\rangle=(1+\abs{\xi}^2)^{1/2}$. For $s\in\R$, define
\begin{equation}
    H^s(\R)
    \coloneqq
    \left\{f\in\mathcal S'(\R):
    \langle\xi\rangle^s\wh f(\xi)\in L^2(\R)\right\},
\end{equation}
with the norm
\begin{equation}
    \norm{f}_{H^s(\R)}^2
    =
    \int_{\R}\langle\xi\rangle^{2s}\abs{\wh f(\xi)}^2\,d\xi .
\end{equation}

Let $I\subset\R$ be an open interval, bounded or unbounded, and let $p_I$
denote restriction to $I$:
\begin{equation}
    p_I:\mathcal D'(\R)\longrightarrow\mathcal D'(I),
    \quad p_IU\coloneqq U|_I.
\end{equation}
Define the restriction space $H^s(I)$ by
\begin{equation}
    H^s(I)
    \coloneqq
    \left\{u\in\mathcal D'(I):
    u=p_IU\text{ for some }U\in H^s(\R)\right\},
\end{equation}
equipped with the norm
\begin{equation}
    \norm{u}_{H^s(I)}
    \coloneqq
    \inf\left\{\norm{U}_{H^s(\R)}:
    U\in H^s(\R),\ p_IU=u\right\}.
\end{equation}
Restriction is bounded, and every $u\in H^s(I)$ has an extension
$U\in H^s(\R)$ with $\norm{U}_{H^s(\R)}\le C\norm{u}_{H^s(I)}$.
Define the supported Sobolev space by
\begin{equation}\label{eqn:supported-space-definition}
    \wt H^s(I)
    \coloneqq
    \overline{C_c^\infty(I)}^{\,H^s(\R)},
\end{equation}
with the inherited $H^s(\R)$ norm.
Under this identification, zero extension is bounded from
$\wt H^s(I)$ into $H^s(\R)$.
The space of Sobolev distributions supported in $\overline I$ is
\begin{equation}
    H^s_{\overline I}(\R)
    \coloneqq
    \left\{u\in H^s(\R):\operatorname{supp}u\subset\overline I\right\}.
\end{equation}
These definitions extend to a general domain $\Omega \subset \R^n$, with
$\wt H^s(\Omega)\subset H^s_{\overline \Omega}(\R^n)$.
If $\Omega$ is a $C^0$ domain, then
$\wt H^s(\Omega)= H^s_{\overline \Omega}(\R^n)$;
see~\cite[Chapter~3]{McLean2000}.
In particular, for an interval,
\begin{equation}
    \wt H^s(I)=H^s_{\overline I}(\R),\quad s\in\R.
\end{equation}
With $L^2(I)$ as pivot space, we have the canonical dualities
\begin{equation}
    \paren{\wt H^s(I)}'=H^{-s}(I),
    \quad
    \paren{H^s(I)}'=\wt H^{-s}(I).
\end{equation}
The pairing $\dual{\cdot}{\cdot}$ extends the $L^2$ inner product.
For $0<s<1$, define the Slobodeckij seminorm by
\begin{equation}
    [u]_{H^s(I)}^2
    \coloneqq
    \int_I\int_I
    \frac{\abs{u(x)-u(y)}^2}{\abs{x-y}^{1+2s}}\,dx\,dy,
\end{equation}
and analogously on $\R$. For $s=1$, we set
$[u]_{H^1(I)}=\norm{\p_su}_{L^2(I)}$.
The quantity $\norm{u}_{L^2(I)}^2+[u]_{H^s(I)}^2$ defines an equivalent squared norm on
$H^s(I)$; see~\cite[Chapter~3]{McLean2000}.

The supported space $\wt H^s(I)$ must be distinguished from $H_0^s(I)$ at half-integer indices.
For $s\ge0$, set
$H_0^s(I)=\overline{C_c^\infty(I)}^{\,H^s(I)} \supset \wt H^s(I)$. If $s\notin\{\frac12,\frac32,\frac52,\ldots\}$, then $\wt H^s(I) = H_0^s(I)$.
For $0\le s<\frac12$, these spaces also coincide with $H^s(I)$. For instance, $\wt L^2(I) = L^2(I)$.
For $s>\frac12$, $\wt H^s(I)$ and $H^s(I)$ differ since Sobolev embedding implies that elements of
$\wt H^s(I)$ have zero endpoint values, whereas elements of $H^s(I)$ need not.
At the half-integer indices, the space $\wt H^s(I)$ imposes additional
endpoint integrability. In particular, $\wt H^{\frac12}(I)$ coincides
with the Lions--Magenes space $H_{00}^{\frac12}(I)$, whose zero boundary
condition cannot be interpreted in the ordinary trace sense.
Likewise, $u\in\wt H^{\frac32}(I)$ satisfies
$\p_su\in\wt H^{\frac12}(I)$, but an ordinary endpoint trace of
$\p_su$ need not exist. For details, see~\cite[Chapter~3]{McLean2000}
and~\cite[Chapter~33]{Tartar2007}.
These supported conditions encode the free-end conditions before
positive-time regularity gives ordinary endpoint traces.

For an integer $m\ge0$ and $0<\alpha<1$, the H\"older space
$\ci{m,\alpha}$ consists of functions $f\in\ci{m}$ with finite norm
\begin{align*}
    \norm{f}_{\ci{m,\alpha}}
    \coloneqq\sum_{j=0}^{m}\norm{\p_s^j f}_{L^\infty(I)}
        +[\p_s^m f]_{\ci{0,\alpha}}, \quad 
    [f]_{\ci{0,\alpha}}
    \coloneqq\sup_{x\ne y\in I}
        \frac{\abs{f(x)-f(y)}}{\abs{x-y}^{\alpha}}.
\end{align*}
We identify $\ci{m,\alpha}$ with $C^{m,\alpha}(\overline I)$ by
continuous extension to the endpoints. For vector-valued functions,
absolute values denote Euclidean norms.

From this point onward, we set $I=(-1,1)$ and collect useful results on these spaces.

% \begin{corollary}\label{point_mult_sobo}
% Let $f\in H^r(I)$ and $g\in\thi{s}$ for $r > \frac{1}{2}$ and $r \ge |s|$. It follows that $fg \in \thi{s}$ and
% \begin{equation}
%     \norm{fg}_{H^s(\R)} \le C \norm{g}_{H^r(I)}\norm{f}_{H^s(\R)} .
% \end{equation}
% Let $f\in L^2(I)$ and $g\in\thi{\frac{1}{2}}$. It follows that $fg \in \thi{-\varepsilon}$ for every $\varepsilon<0$ and
% \begin{equation}
%     \norm{fg}_{H^{-\varepsilon}(\R)} \le C \norm{g}_{L^2(I)}\norm{f}_{H^\frac{1}{2}(\R)} ,
% \end{equation}
% \end{corollary}

\emph{Interpolation and multiplication.}
We first record how differentiation acts on supported spaces.

\begin{lemma}\label{lem:wt-der}
For every $s\in\R$, differentiation defines a bounded map $\p_s:\wt H^s(I)\to\wt H^{s-1}(I).$
\end{lemma}
\begin{proof}
If $f\in\wt H^s(I)$, then its distributional derivative remains supported
in $\overline I$. The Fourier multiplier estimate
\begin{equation}
    \norm{\p_s f}_{H^{s-1}(\R)} \le \norm{f}_{H^s(\R)}
\end{equation}
therefore gives $\p_s f\in\wt H^{s-1}(I)$ and proves boundedness.
\end{proof}

We write $(Y_0,Y_1)_{\theta,2}$ for real interpolation. The following
identity identifies the interpolation spaces used in the semigroup estimates.

\begin{lemma}\label{lem:hs-intp}
Let $s>0$. With equivalent norms,
\begin{equation}\label{eqn:hs-intp}
    \paren{\hi{-s},\thi{s}}_{\theta,2}
    =
    \begin{cases}
        \thi{s(2\theta-1)}, & \frac12<\theta<1,\\
        L^2(I), & \theta=\frac12,\\
        \hi{-s(1-2\theta)}, & 0<\theta<\frac12.
    \end{cases}
\end{equation}
The supported spaces carry the inherited norms from~\eqref{eqn:supported-space-definition}, including at half-integer indices.
\end{lemma}
\begin{proof}
The canonical duality and dense inclusions
$\thi{s}\hookrightarrow L^2(I)\hookrightarrow\hi{-s}$ give, by
\cite[Lemma~B.10]{McLean2000},
\begin{equation*}
    \paren{\hi{-s},\thi{s}}_{\frac12,2}=L^2(I).
\end{equation*}
Reiteration \cite[Theorem~B.6]{McLean2000} reduces the remaining cases
to interpolation between $H^{-s}(I)$ and $L^2(I)$ with parameter
$2\theta$, or between $L^2(I)$ and $\wt H^s(I)$ with parameter
$2\theta-1$. The restriction-space and supported-space identities in
\cite[Theorems~B.8 and~B.9]{McLean2000} then give \eqref{eqn:hs-intp}.
\end{proof}

The multiplication estimate below is used in the composition bounds for
the reconstructed tangent and in the nonlinear operator estimates.

\begin{lemma}\label{point_mult_sobo}
Let $s,s_1,s_2\in\R$ satisfy one of the following conditions:
\begin{enumerate}
    \item $s\ge0$, $s_i\ge s$ for $i=1,2$, and
    $s_1+s_2-s>\frac12$;
    \item $s<0$, $s_i\ge s$ for $i=1,2$, $\min\{s_1,s_2\}<0$,
    $s_1+s_2\ge0$, and $s_1+s_2-s>\frac12$;
    \item $s<0$, $s_i\ge s$ for $i=1,2$, $\min\{s_1,s_2\}\ge0$,
    $s_1+s_2>0$, and $s_1+s_2-s>\frac12$.
\end{enumerate}
Then multiplication extends to bounded bilinear maps
\begin{align*}
    \hi{s_1}\times\thi{s_2}&\longrightarrow\thi{s},\\
    \hi{s_1}\times\hi{s_2}&\longrightarrow\hi{s},
\end{align*}
with the estimates
\begin{align}
    \norm{fg}_{\hr{s}}
    &\le C\norm{f}_{\hi{s_1}}\norm{g}_{\hr{s_2}},
    &&f\in\hi{s_1},\quad g\in\thi{s_2},\label{eqn:product-supported}\\
    \norm{fg}_{\hi{s}}
    &\le C\norm{f}_{\hi{s_1}}\norm{g}_{\hi{s_2}},
    &&f\in\hi{s_1},\quad g\in\hi{s_2},\label{eqn:product-restriction}
\end{align}
where $C=C(s,s_1,s_2)$.
\end{lemma}
\begin{proof}
Under the stated conditions, the whole-line multiplication theorem
\cite[Theorems~7.3, 8.1, and~8.2]{Behzadan2021} gives
\begin{equation*}
    \norm{FG}_{\hr{s}}
    \le C\norm{F}_{\hr{s_1}}\norm{G}_{\hr{s_2}}.
\end{equation*}
Let $F\in\hr{s_1}$ extend $f$. For $g\in\thi{s_2}$, choose
$g_n\in C_c^\infty(I)$ converging to $g$ in $\hr{s_2}$.
Each product $Fg_n$ belongs to $\thi{s}$ and depends only on the
restriction of $F$ to $I$. The whole-line estimate shows that
$Fg_n\to Fg$ in $\hr{s}$. Since $\thi{s}$ is closed, this defines
$fg\in\thi{s}$ independently of the extension $F$. Taking the infimum
over extensions gives~\eqref{eqn:product-supported}.

For $g\in\hi{s_2}$, also extend $g$ to $G\in\hr{s_2}$ and define
$fg=p_I(FG)$. The whole-line product is local, so its restriction is
independent of both extensions. Boundedness of restriction and the
whole-line estimate give~\eqref{eqn:product-restriction} after taking the infimum
over extensions of $f$ and $g$.
\end{proof}

\emph{Reconstruction and arc-chord control.}
The primitive estimate below controls the tangent angle $J\kappa=\theta_\kappa$;
Lemmas~\ref{tau-n-est} and~\ref{lem:Xk-holder} then give Sobolev bounds
for the tangent and H\"older bounds for the reconstructed curve.

\begin{lemma}\label{lem:theta-kappa}
For $s\ge -1$ and $\kappa\in\thi{s}$, one has
\begin{equation}\label{eqn:theta-Hs+1}
    \norm{\theta_\kappa}_{\hi{s+1}} \le C \norm{\kappa}_{\hr{s}}.
\end{equation}
\end{lemma}
\begin{proof}
Since $C_c^\infty(I)$ is dense in $\thi{s}$, it is enough to prove the
estimate for smooth, compactly supported functions. Define
\begin{equation}
    (Jf)(x) \coloneqq \int_{-1}^x f(y)\,dy,
    \quad
    (J^*f)(x) \coloneqq \int_x^1 f(y)\,dy.
\end{equation}
Then $(Jf)'=f$, $(J^*f)'=-f$, and $\theta_\kappa=J\kappa$.

We first establish the two endpoint estimates. The pointwise bound
\begin{equation}
    \abs{Jf(x)} \le \int_I \abs{f(y)}\,dy
\end{equation}
and the identity $(Jf)'=f$ give
\begin{equation}\label{eqn:Jf-1}
    \norm{Jf}_{\hi{1}} \le C\norm{f}_{L^2(I)}.
\end{equation}
The same estimate holds for $J^*$. Moreover, Fubini's theorem shows that
$J^*$ is the adjoint of $J$. Thus, for $g\in L^2(I)$,
\begin{equation}
    \abs{\dual{Jf}{g}}
    = \abs{\dual{f}{J^*g}}
    \le \norm{f}_{\hr{-1}}\norm{J^*g}_{\hi{1}}
    \le C\norm{f}_{\hr{-1}}\norm{g}_{L^2(I)}.
\end{equation}
Taking the supremum over $g$ yields
\begin{equation}\label{eqn:Jf-2}
    \norm{Jf}_{L^2(I)} \le C \norm{f}_{\hr{-1}}.
\end{equation}

Consequently, $J:\thi{-1}\to L^2(I)$ and
$J:L^2(I)\to \hi{1}$ are bounded. Interpolation gives
\begin{equation}
    \norm{Jf}_{\hi{s+1}} \le C\norm{f}_{\hr{s}},
    \quad -1\le s\le 0.
\end{equation}

For the nonnegative range, let $\ell\ge0$ be an integer. Since
$\partial_x^{j+1}Jf=\partial_x^j f$, \eqref{eqn:Jf-1} implies
\begin{equation}
    \norm{Jf}_{\hi{\ell+1}}^2 = \norm{Jf}_{L^2(I)}^2 + \sum_{j=0}^{\ell}\norm{\partial_x^{j+1}Jf}_{L^2(I)}^2 \le
    C \norm{f}_{L^2(I)}^2 + \sum_{j=0}^{\ell}\norm{\partial_x^j f}_{L^2(I)}^2 \le C \norm{f}_{\hi{\ell}}^2.
\end{equation}
Interpolating between integer orders therefore gives
\begin{equation}
    \norm{Jf}_{\hi{s+1}} \le C\norm{f}_{\hi{s}},
    \quad s\ge 0.
\end{equation}
Applying this to $f=\kappa$ and using $\norm{\kappa}_{\hi{s}}\le C\norm{\kappa}_{\hr{s}}$
proves~\eqref{eqn:theta-Hs+1} for $s\ge0$, while the negative range follows
from the first interpolation argument.
\end{proof}

Since $\xk'=\bm e_r(\theta_\kappa)$, the angle estimate
\eqref{eqn:theta-Hs+1} gives the following tangent bounds.

\begin{lemma}\label{tau-n-est}
Let $-\frac{1}{2}<s\le 0$ and $\kappa\in \thi{s}$. Then
\begin{align}
    [\xk']_{\hi{s+1}} &\le C\norm{\kappa}_{\hr{s}},\\
    \norm{\xk''}_{\hi{s}} &\le C\paren{1+\norm{\kappa}_{\hr{s}}}\norm{\kappa}_{\hr{s}},\\
    \norm{\xk'}_{\hi{s+1}} &\le C\paren{1+\norm{\kappa}_{\hr{s}}}.\label{eqn:tangent-Sobolev-bound}
\end{align}
Moreover, if $\kappa_i\in \thi{s}$ satisfy
$\norm{\kappa_i}_{\hr{s}}\le M$ for $i=1,2$, then
\begin{align}
    \norm{\xka{1}'-\xka{2}'}_{\hi{s+1}}
    &\le C\paren{1+M} \norm{\kappa_1-\kappa_2}_{\hr{s}},\\
    \norm{\xka{1}''-\xka{2}''}_{\hi{s}}
    &\le C\paren{1+M} \norm{\kappa_1-\kappa_2}_{\hr{s}}.
\end{align}
\end{lemma}

The proof is given in Appendix~\ref{app:geometry-estimates}.


\begin{lemma}\label{lem:Xk-holder}
Let $m\in \mathbb N\cup\{0\}$, $\delta>0$, and
$s=m-\frac{1}{2}+\delta$. Suppose that
$0<\alpha<\min\{\delta,1\}$ and $\kappa\in \thi{s}$. Then
\begin{equation}\label{eqn:Xk-holder-higher}
    \norm{\xk}_{\ci{m+1,\alpha}} \le C\paren{1+\norm{\kappa}_{\hr{s}}}^{m+1}.
\end{equation}
When $m=0$, we have the sharper seminorm estimate
\begin{equation}\label{eqn:Xk-holder}
    [\xk']_{\ci{0,\alpha}} \le C \norm{\kappa}_{\hr{s}}.
\end{equation}
Moreover, if $\kappa_i\in \thi{s}$ satisfy
$\norm{\kappa_i}_{\hr{s}}\le M$ for $i=1,2$, then
\begin{equation}\label{eqn:Xk-holder-higher-lip}
    \norm{\xka{1}-\xka{2}}_{\ci{m+1,\alpha}}
    \le C\paren{1+M}^{m+1}\norm{\kappa_1-\kappa_2}_{\hr{s}}.
\end{equation}
\end{lemma}

The proof is given in Appendix~\ref{app:geometry-estimates}.

For a parametrized curve $\bm X$, define the arc-chord constant by
\begin{equation}\label{eqn:arc-chord-intro}
    \abs{\bm X}_*
    \coloneqq
    \inf_{s,\eta\in I,\ s\neq\eta}
    \frac{\abs{\bm X(s)-\bm X(\eta)}}{\abs{s-\eta}}.
\end{equation}
We restrict attention to curves satisfying the arc-chord condition
$\abs{\bm X}_*>0$. For the curve $\xk$ reconstructed from $\kappa$, write
$\abs{\kappa}_*=\abs{\xk}_*$. Since $\abs{\xk'}=1$, the chord length
is bounded by the arclength:
\begin{equation}
    \abs{\xk(s)-\xk(\eta)}\le\abs{s-\eta},
    \quad s,\eta\in I.
\end{equation}
Thus $0\le\abs{\kappa}_*\le1$.

Fix $0<\delta\le\frac12$. The admissible curvatures form the set
\begin{equation}\label{eqn:U-intro}
    \mc U_\delta
    \coloneqq
    \left\{
    \kappa\in\thi{-\frac12+\delta}:\abs{\kappa}_*>0
    \right\}.
\end{equation}

Throughout the estimates, $C_{M,m}$ denotes a positive constant depending
on the norm bound $M$ and arc-chord lower bound $m$ specified in the
statement. It is locally bounded for $M\ge0$, $m>0$, nondecreasing in $M$
and $m^{-1}$, and may change from line to line. Dependence on fixed
regularity exponents is suppressed. Factors expressing vanishing at zero
are kept explicit.

To formulate the local theory, let $J\kappa=\thk$ be the
integration operator introduced above and define
\begin{equation}
    C_{J,\delta}
    \coloneqq
    \sup_{0\neq\kappa\in\thi{-\frac12+\delta}}
    \frac{\norm{J\kappa}_{\ci{0}}}
    {\norm{\kappa}_{\hr{-\frac12+\delta}}}.
\end{equation}
Lemma~\ref{lem:theta-kappa} and Sobolev embedding show that
$C_{J,\delta}<\infty$.
For $\kappa_0\in\mc U_\delta$, the following neighborhood gives a common
low-norm curvature bound and a common arc-chord lower bound for reconstructed
curves and their convex combinations, as required in the difference estimates:

\begin{equation}\label{eqn:Bkappa0}
\mc B(\kappa_0) \coloneqq \left\{\kappa\in \thi{-\frac{1}{2}+\delta}: \norm{\kappa-\kappa_0}_{\hr{-\frac{1}{2}+\delta}}\le \frac{1}{2C_{J,\delta}}\abs{\kappa_0}_* \right\}.
\end{equation}

\begin{lemma}\label{lem:set-ok0}
For $\kappa_0\in \mc U_\delta$, let
\begin{equation}
    M
    = \norm{\kappa_0}_{\hr{-\frac{1}{2}+\delta}}
    +\frac{1}{2C_{J,\delta}}\abs{\kappa_0}_*,
    \quad
    m = \frac{1}{2}\abs{\kappa_0}_*.
\end{equation}
If $\kappa_1,\kappa_2 \in \mc B(\kappa_0)$, then
\begin{align}\label{eqn:neighborhood-geometry}
    \norm{\kappa_i}_{\hr{-\frac{1}{2}+\delta}} \le M,
    \quad
    \abs{\rho\xka{1} + \paren{1-\rho}\xka{2}}_* \ge m,
\end{align}
for $i=1,2$ and every $\rho\in[0,1]$.
Hence $\mc B(\kappa_0)\subset\mc U_\delta$.
\end{lemma}
\begin{proof}
We have
\begin{equation}
    \norm{\kappa_i}_{\hr{-\frac{1}{2}+\delta}}
    \le \norm{\kappa_0}_{\hr{-\frac{1}{2}+\delta}}
    + \frac{1}{2C_{J,\delta}}\abs{\kappa_0}_*
    =M.
\end{equation}
For $\bm X,\bm Y\in\ci{1}$, the arc-chord constant satisfies
\begin{equation}\label{eqn:arc-chord-C1-continuity}
\begin{aligned}
    &\abs{\abs{\bm X}_* - \abs{\bm Y}_*}
    = \abs{ \inf_{s\neq t}\abs{\frac{\bm X(s)-\bm X(t)}{s-t}}
    - \inf_{s\neq t}\abs{\frac{\bm Y(s)-\bm Y(t)}{s-t}} } \\
    &\le \sup_{s\neq t}\abs{\abs{\frac{\bm X(s)-\bm X(t)}{s-t}}
    - \abs{\frac{\bm Y(s)-\bm Y(t)}{s-t}} }
    \le \sup_{s\neq t}\abs{\frac{(\bm X-\bm Y)(s)-(\bm X-\bm Y)(t)}{s-t}} \\
    &= \sup_{s\neq t}\abs{\int_0^1
    (\bm X-\bm Y)'(t+\lambda(s-t))\,d\lambda}
    \le \norm{\bm X'-\bm Y'}_{\ci{0}}.
\end{aligned}
\end{equation}
For $i=1,2$, Lemma~\ref{lem:theta-kappa} gives
\begin{equation}
\begin{aligned}
    &\norm{\xka{i}'-\xka{0}'}_{\ci{0}} = \sup_{s\in I} \abs{\bm e_r(\thka{i}(s)) - \bm e_r(\thka{0}(s))} \le \sup_{s\in I} \abs{\thka{i}(s)-\thka{0}(s)} \\
    & \le C_{J,\delta}\norm{\kappa_i-\kappa_0}_{\hr{-\frac{1}{2}+\delta}} \le \frac{1}{2}\abs{\kappa_0}_*.
\end{aligned}
\end{equation}
For $\rho\in[0,1]$, the preceding arc-chord estimate gives
\begin{equation}
\begin{aligned}
&\abs{\abs{\rho\xka{1}+(1-\rho)\xka{2}}_* - \abs{\kappa_0}_* }
\le \rho\norm{\xka{1}'-\xka{0}'}_{\ci{0}}
+(1-\rho)\norm{\xka{2}'-\xka{0}'}_{\ci{0}} \\
&\le C_{J,\delta} \rho\norm{\kappa_1-\kappa_0}_{\hr{-\frac{1}{2}+\delta}}
+C_{J,\delta}(1-\rho)\norm{\kappa_2-\kappa_0}_{\hr{-\frac{1}{2}+\delta}}
\le \frac{1}{2}\abs{\kappa_0}_*.
\end{aligned}
\end{equation}
Hence
\begin{equation}
    \abs{\rho\xka{1} + \paren{1-\rho}\xka{2}}_*
    \ge \frac{1}{2}\abs{\kappa_0}_*=m.
\end{equation}
\end{proof}

\subsection{Linear operators and semigroup estimates}
\label{subsec:linear-theory}
\label{subsec:interval-regularity}

The order-one operator $\p_s\mc H$ is the elliptic principal part of the
tension equation, while the positive order-three operator $-\Lambda$
governs curvature dissipation and smoothing. We prove coercivity and
invertibility for both operators, then derive the analytic-semigroup
estimates used in Sections~\ref{sec:local-theory}--\ref{sec:global-theory}.
We begin with boundedness of the finite Hilbert transform and the
fractional Poincar\'e inequality needed for coercivity.

\begin{lemma}\label{Hf-bounded}
For every $s\in\R$, the finite Hilbert transform defines a bounded map
$\mc H:\wt H^s(I)\to H^s(I)$. More precisely,
\begin{equation}\label{eqn:hilbert-bound}
    \norm{\mc Hf}_{H^s(I)} \le C\norm{f}_{H^s(\R)}
\end{equation}
for every $f\in\wt H^s(I)$.
\end{lemma}
\begin{proof}
The whole-line Hilbert transform $\mc H_{\R}$ has bounded Fourier symbol
$-i\operatorname{sgn}(\xi)$. Restricting $\mc H_{\R}f$ to $I$ gives
\begin{equation}
    \norm{\mc Hf}_{H^s(I)}
    \le \norm{\mc H_{\R}f}_{H^s(\R)}
    \le C\norm{f}_{H^s(\R)}.
\end{equation}
\end{proof}

\begin{lemma}\label{lem:frac-poincare}
There exists $C>0$ such that every $f\in\thi{\frac{1}{2}}$ satisfies
\begin{equation}
    \norm{f}_{L^2(\R)}\le C[f]_{\hr{\frac{1}{2}}}.
\end{equation}
\end{lemma}
\begin{proof}
For $f\in C_c^\infty(I)$, the interaction between $I$ and its complement
gives
\begin{equation}
    [f]_{\hr{\frac{1}{2}}}^2 \ge 2\int_I\abs{f(y)}^2 \left(\int_{-\infty}^{-1}\frac{dx}{\abs{x-y}^2}
    +\int_1^\infty\frac{dx}{\abs{x-y}^2}\right)dy =4\int_I\frac{\abs{f(y)}^2}{1-y^2}\,dy \ge
    2\norm{f}_{L^2(\R)}^2.
\end{equation}
The result follows by density.
\end{proof}

The Fourier symbol of $\p_s\mc H$ and the preceding Poincar\'e inequality
give coercivity of the tension principal part.

\begin{lemma}\label{ds-H-biject}
The operator
$\p_s\mc H:\thi{\frac{1}{2}}\to H^{-\frac{1}{2}}(I)$ is bounded and
bijective.
\end{lemma}
\begin{proof}
Lemmas~\ref{lem:wt-der} and~\ref{Hf-bounded}, together with Plancherel's
theorem, give boundedness:
\begin{equation}
\begin{aligned}
    \norm{\p_s\mc Hf}_{H^{-\frac{1}{2}}(I)}^2
    \le \norm{\mc H_{\R}f}_{\hr{\frac{1}{2}}}^2  =\int_\R\paren{1+\abs{\xi}^2}^{\frac12}
    \abs{\wh{\mc H_{\R}f}(\xi)}^2\,d\xi
    =\norm{f}_{\hr{\frac{1}{2}}}^2.
\end{aligned}
\end{equation}

For the lower bound, symmetrizing the kernel gives
\begin{equation}\label{eqn:dsh-coer}
\begin{aligned}
    \dual{\p_s\mc Hf}{f}
    =\frac{1}{2\pi}\int_\R\int_\R
    \frac{\paren{f(s)-f(\eta)}^2}{\paren{s-\eta}^2}\,d\eta ds =\frac{1}{2\pi}[f]_{\hr{\frac{1}{2}}}^2.
\end{aligned}
\end{equation}
By Lemma~\ref{lem:frac-poincare}, the homogeneous seminorm on the right
controls the full $H^{1/2}(\R)$ norm. The associated bilinear form is
therefore bounded and coercive on $\thi{\frac{1}{2}}$. The Lax--Milgram
theorem gives bijectivity and the bounded inverse
$(\p_s\mc H)^{-1}:H^{-\frac{1}{2}}(I)\to\thi{\frac{1}{2}}$.
\end{proof}

The same coercivity argument, applied to $\p_s f$, gives the corresponding
result for the order-three curvature operator.

\begin{lemma}\label{ds-Lam-biject}
The operator
$-\Lambda=-\frac14\p_s^2\mc H\p_s:
\thi{\frac{3}{2}}\to H^{-\frac{3}{2}}(I)$ is bounded and bijective.
\end{lemma}
\begin{proof}
Lemmas~\ref{ds-H-biject} and~\ref{lem:wt-der} give
\begin{equation}
\begin{aligned}
    \norm{\Lambda f}_{H^{-\frac{3}{2}}(I)}^2
    &\le \frac{1}{16}\norm{\p_s \mc H\p_s f}_{\hi{-\frac{1}{2}}}^2
    \le \frac{1}{16}\norm{\p_s f}_{\hr{\frac{1}{2}}}^2
    \le \frac{1}{16}\norm{f}_{\hr{\frac{3}{2}}}^2.
\end{aligned}
\end{equation}

For coercivity, integration by parts and~\eqref{eqn:dsh-coer}, applied to
$\p_s f$, give
\begin{equation}\label{eqn:lam-coer}
    \dual{-\Lambda f}{f} = \frac{1}{4}\int_\R f(s)(-\p_s^2\mc H_{\R}\p_s f)(s)\,ds = \frac{1}{4}\int_\R \p_s
    f(s)(\p_s\mc H_{\R}\p_s f)(s)\,ds = \frac{1}{8\pi}[\p_s f]_{\hr{\frac{1}{2}}}^2.
\end{equation}
Lemma~\ref{lem:frac-poincare} controls $\norm{\p_s f}_{L^2(\R)}$, and the
usual Poincar\'e inequality applies because
$f\in\thi{\frac{3}{2}}\hookrightarrow H_0^1(I)$. Consequently,
$\norm{f}_{H^1(\R)}$ is controlled by the right-hand side of
\eqref{eqn:lam-coer}. We obtain
$\dual{-\Lambda f}{f}\ge C\norm{f}_{\hr{\frac{3}{2}}}^2$.
The associated bilinear form is therefore bounded and coercive on
$\thi{\frac{3}{2}}$, so the Lax--Milgram theorem gives bijectivity.
\end{proof}

The coercivity bounds~\eqref{eqn:dsh-coer} and~\eqref{eqn:lam-coer} give equivalent norms on the supported spaces and
their duals. For either
$B=\p_s\mc H$ on $V=\thi{\frac12}$ or
$B=-\Lambda$ on $V=\thi{\frac32}$, boundedness and coercivity give
\begin{equation}\label{eqn:coercive-norm-equivalences}
\begin{aligned}
    c\norm{u}_V^2
    &\le \dual{Bu}{u}\le C\norm{u}_V^2,
    &&u\in V,\\
    c\norm{h}_{V'}^2
    &\le \dual{h}{B^{-1}h}\le C\norm{h}_{V'}^2,
    &&h\in V'.
\end{aligned}
\end{equation}
The second follows by writing $h=Bu$ and using the boundedness of $B$ and
$B^{-1}$.

We use \cite[Chapter~2]{Lunardi1995} for semigroup estimates for
$\Lambda$ on the spaces

\begin{equation}
    X=\hi{-\frac32},
    \quad D(\Lambda)=\thi{\frac32}.
\end{equation}
The domain is dense in $X$, and the endpoint conditions~\eqref{eqn:cur-bc}
are imposed in the supported Sobolev-space sense. By Lemma~\ref{ds-Lam-biject},
its graph norm is equivalent to the $H^{3/2}(\R)$ norm.
At the interpolation endpoints, we use the convention
$(X,D(\Lambda))_{0,2}=X$ and
$(X,D(\Lambda))_{1,2}=D(\Lambda)$.

Using symmetry and coercivity, we equip $X$ with an equivalent Hilbert
norm in which $-\Lambda$ is positive and self-adjoint. The following
proposition then gives the semigroup bound~\eqref{eqn:lambda-semigroup-est}
used in the evolution estimates.

\begin{proposition}\label{prop:lambda-semigroup}
The realization $\Lambda:D(\Lambda)\subset X\to X$ is sectorial and
generates an analytic semigroup $e^{t\Lambda}$ on $X$. Its spectral bound
satisfies
\begin{equation}
    \omega_\Lambda
    \coloneqq\sup\{\operatorname{Re}z:z\in\sigma(\Lambda)\}<0.
\end{equation}
For $\omega\in(\omega_\Lambda,0)$, $n\in\{0,1\}$, and
$\alpha,\beta\in[0,1]$, with $\alpha\le\beta$ when $n=0$, one has
\begin{equation}\label{eqn:lambda-semigroup-est}
    \norm{\Lambda^n e^{t\Lambda}f}_{(X,D(\Lambda))_{\beta,2}}
    \le C_{n,\alpha,\beta,\omega}e^{\omega t}
    \max\{t^{-n+\alpha-\beta},1\}
    \norm{f}_{(X,D(\Lambda))_{\alpha,2}}
\end{equation}
for every $f\in(X,D(\Lambda))_{\alpha,2}$ and $t>0$.
Moreover, for $0<\alpha<1$ and
$f\in(X,D(\Lambda))_{\alpha,2}$,
\begin{equation}\label{eqn:lambda-semigroup-cont}
    \lim_{t\downarrow0}
    \norm{e^{t\Lambda}f-f}_{(X,D(\Lambda))_{\alpha,2}}=0.
\end{equation}
\end{proposition}
\begin{proof}
Write $B=-\Lambda$. The symmetry and coercivity in
\eqref{eqn:lam-coer} and \eqref{eqn:coercive-norm-equivalences} show that
\begin{equation}\label{eqn:x-lambda-norm}
    (f,g)_{X_{-\Lambda}}
    \coloneqq\dual{f}{B^{-1}g}_{X,D(\Lambda)}
\end{equation}
defines an inner product on $X$ whose induced norm is equivalent to the
original norm.
On this Hilbert space, $B^{-1}$ is bounded and satisfies
\begin{equation*}
    (B^{-1}f,g)_{X_{-\Lambda}}
    =(B^{-1}f,B^{-1}g)_{L^2(I)},
    \quad f,g\in X.
\end{equation*}
Thus $B^{-1}$ is positive, self-adjoint, and injective, with dense range
$D(\Lambda)$. Its inverse $B$ is therefore positive and self-adjoint
with domain exactly $D(\Lambda)$. The spectral theorem gives
\begin{equation*}
    \sigma(\Lambda)\subset(-\infty,-\mu],
    \quad
    \mu=\norm{B^{-1}}_{\mc L(X_{-\Lambda})}^{-1}>0,
\end{equation*}
hence $\omega_\Lambda<0$, sectoriality, and analytic generation.

The estimates of
\cite[Propositions~2.2.2, 2.2.9, and~2.3.1]{Lunardi1995}, with $p=2$,
give \eqref{eqn:lambda-semigroup-est} for $0<\alpha,\beta<1$.
The endpoint cases follow from the base-space derivative estimates and
graph-norm bounds by interpolation, using
$\Lambda e^{t\Lambda}f=e^{t\Lambda}\Lambda f$ for $f\in D(\Lambda)$.
The graph-norm equivalence and~\eqref{eqn:hs-intp} identify the
interpolation spaces with the stated Sobolev spaces. Continuity at $t=0$
follows from \cite[Proposition~2.2.8]{Lunardi1995} with $p=2$.
\end{proof}

\subsection{Wiener--Hopf factorization and endpoint regularity}
\label{subsec:endpoint-factorization}

For an integer $k\ge1$, let $A_k$ be the Fourier multiplier on $\R$
with symbol $(1+\xi^2)^{k/2}$. We factor $A_k$ into two operators that
preserve support on opposite half-lines, then apply $\Pi_+$ to select the
positive-support component and obtain an explicit half-line solution formula.
For $|\varepsilon|<\frac12$, forcing in
$H^{-k/2+\varepsilon}(\R_+)$ gives a solution in
$\wt H^{k/2+\varepsilon}(\R_+)$.
If the forcing has one additional order of Sobolev regularity, the inverse
Fourier transform of the factored forcing lies in $H^{1+\varepsilon}(\R)$
and has a trace at the origin. Differentiating its
half-line truncation produces a point mass there. The contribution of this
point mass to the solution is $c\mathbf1_{[0,\infty)}e^{-x}x^{k/2}$, with
remainder in $\wt H^{k/2+1+\varepsilon}(\R_+)$.
After localization to $I$, $k=1$ governs the tension endpoint behavior,
while $k=3$ gives the $d^{3/2}$ curvature profile in
Theorem~\ref{thm:main-regularity}.

Write $a_k$ for the Fourier symbol of $A_k$ and $A_{k,\pm}$ for the multipliers
with symbols $a_{k,\pm}$, where
\begin{equation}\label{eqn:orderk-symbol-factorization}
    a_k(\xi)
    =(1+\xi^2)^{\frac{k}{2}}
    =a_{k,+}(\xi)a_{k,-}(\xi),
    \quad
    a_{k,+}(\xi)=(\xi-i)^{\frac{k}{2}},
    \quad
    a_{k,-}(\xi)=(\xi+i)^{\frac{k}{2}}.
\end{equation}
The factor $a_{k,+}$ is the boundary value of the holomorphic function
$(z-i)^{k/2}=\exp(\frac{k}{2}\log(z-i))$ in the lower half-plane, with
$-\pi<\arg(z-i)<0$. Similarly, $a_{k,-}$ is the boundary value of
$(z+i)^{k/2}$ in the upper half-plane, with $0<\arg(z+i)<\pi$.
For real $\xi$, the arguments of $\xi-i$ and $\xi+i$ are opposite and
both moduli equal $(1+\xi^2)^{1/2}$. Thus their $k/2$ powers multiply
to exactly $(1+\xi^2)^{k/2}$, including when $k$ is odd.
The reciprocal factors are holomorphic in the same respective half-planes.

With the Fourier convention of Subsection~\ref{subsec:notation}, their
support properties follow directly from the Laplace-transform identity
\begin{equation}\label{eq:halfline-factor-kernel}
    \mathcal F^{-1}\!\bigl[(\xi\mp i)^{-a}\bigr](x)
    =\frac{\sqrt{2\pi}\,e^{\pm i\pi a/2}}{\Gamma(a)}
    \mathbf1_{\R_\pm}(x)e^{-|x|}|x|^{a-1},
    \quad a>0.
\end{equation}
Here $\Gamma(a)$ is Euler's gamma function.
Consequently $A_{k,\pm}^{-1}$ preserve support in $\overline{\R_\pm}$.
Writing $(\xi\mp i)^{k/2}=(\xi\mp i)^N(\xi\mp i)^{-(N-k/2)}$
with an integer $N>k/2$ proves the same assertion for $A_{k,\pm}$.
The multiplier bounds extend these identities to all Sobolev orders:
for either sign, $A_{k,\pm}:H^r(\R)\to H^{r-k/2}(\R)$ is an isomorphism
whose inverse raises the order by $k/2$. Both $A_{k,\pm}$ and
$A_{k,\pm}^{-1}$ preserve support in $\overline{\R_\pm}$.

For $g\in L^2(\R)$, define the half-line projections
\begin{equation}\label{eq:Pi-support-projections}
    \Pi_\pm g
    \coloneqq
    \mathcal F\!\left(
    \mathbf1_{\R_\pm}\mathcal F^{-1}g
    \right).
\end{equation}
Multiplication by $\mathbf1_{\R_\pm}$ is bounded on $H^s(\R)$ for
$|s|<\frac12$. Consequently, $\Pi_\pm$ are bounded on
$\mathcal F(H^s(\R))$ in this range and project onto
$\mathcal F(H^s_{\overline{\R_\pm}}(\R))$. For negative $s$, the
projections are understood by continuous extension. In this range,
$\Pi_++\Pi_-=1$ and $\Pi_+\Pi_-=\Pi_-\Pi_+=0$, giving the direct sum
\begin{equation*}
    H^s(\R)=H^s_{\overline{\R_+}}(\R)
    \oplus H^s_{\overline{\R_-}}(\R).
\end{equation*}
The intersection is trivial: a distribution supported only at the origin
is a finite linear combination of derivatives of $\delta_0$, and no
nonzero such combination belongs to $H^s(\R)$ for $s>-\frac12$. Thus $\Pi_+$ acts as the identity on the positive-support
component and annihilates the negative-support component.
Consider
\begin{equation}\label{eq:A-orderk-halfline}
    p_+ A_k\psi_+ = f \quad \text{on } \R_+,
\end{equation}
where $p_+$ is restriction to $\R_+$ and $\psi_+$ is the zero extension:
$\psi_+(x)=\psi(x)$ for $x>0$ and $\psi_+(x)=0$ for $x<0$.

\begin{lemma}\label{lem:A-orderk-halfline}
Let $k\ge 1$ be an integer and let $|\varepsilon|<\frac12$.
The constants in the estimates below depend only on $k,\varepsilon$.
\begin{enumerate}
    \item If $f\in H^{-\frac{k}{2}+\varepsilon}(\R_+)$, then the equation
    \eqref{eq:A-orderk-halfline} admits a unique solution
    $\psi_+\in H^{\frac{k}{2}+\varepsilon}_{\overline{\R_+}}(\R)$, and
    \begin{equation}
        \norm{\psi_+}_{H^{\frac{k}{2}+\varepsilon}(\R)}
        \le C \norm{f}_{H^{-\frac{k}{2}+\varepsilon}(\R_+)}.
    \end{equation}
    \item If $f\in H^{-\frac{k}{2}+1+\varepsilon}(\R_+)$, then the solution admits the decomposition
    \begin{equation}
        \psi_+(x)
        = c\,\mathbf1_{[0,\infty)}(x)e^{-x}x^{\frac{k}{2}} + \psi_r(x),
    \end{equation}
    where $c$ depends on $f$ and
    $\psi_r\in H^{\frac{k}{2}+1+\varepsilon}_{\overline{\R_+}}(\R)$. Moreover,
    \begin{equation}
        \norm{\psi_r}_{H^{\frac{k}{2}+1+\varepsilon}(\R)}
        + |c|
        \le C \norm{f}_{H^{-\frac{k}{2}+1+\varepsilon}(\R_+)}.
    \end{equation}
\end{enumerate}
\end{lemma}
\begin{proof}
\emph{Solution formula.}
Let $\wt f\in H^{-\frac{k}{2}+\varepsilon}(\R)$ be an extension of $f$ such that
$\norm{\wt f}_{H^{-\frac{k}{2}+\varepsilon}(\R)}
\le 2 \norm{f}_{H^{-\frac{k}{2}+\varepsilon}(\R_+)}$.
For a solution of~\eqref{eq:A-orderk-halfline}, set
\begin{equation}
    \psi_-=\wt f-A_k\psi_+
    \in H^{-\frac{k}{2}+\varepsilon}_{\overline{\R_-}}(\R).
\end{equation}
Taking Fourier transforms gives
\begin{equation}
    a_{k,+}(\xi)a_{k,-}(\xi)\wh{\psi_+}(\xi) + \wh{\psi_-}(\xi)
    = \wh{\wt f}(\xi).
\end{equation}
Multiplication by $a_{k,-}^{-1}(\xi)$ gives
\begin{equation}\label{eq:A-orderk-wh}
    a_{k,+}(\xi)\wh{\psi_+}(\xi)
    + a_{k,-}^{-1}(\xi)\wh{\psi_-}(\xi)
    = a_{k,-}^{-1}(\xi)\wh{\wt f}(\xi).
\end{equation}
Since $a_{k,-}^{-1}$ has order $-\frac{k}{2}$, the right-hand side belongs to
$\mathcal{F}(H^\varepsilon(\R))$. The support properties following~\eqref{eq:halfline-factor-kernel} give
$a_{k,+}\wh{\psi_+}\in \mathcal{F}(H^\varepsilon_{\overline{\R_+}}(\R))$
and
$a_{k,-}^{-1}\wh{\psi_-}\in \mathcal{F}(H^\varepsilon_{\overline{\R_-}}(\R))$.
Because $\Pi_+:\mathcal{F}(H^\varepsilon(\R))\to\mathcal{F}(H^\varepsilon(\R))$
is bounded for $|\varepsilon|<\frac12$, applying $\Pi_+$ to
\eqref{eq:A-orderk-wh} yields
$a_{k,+}(\xi)\wh{\psi_+}(\xi)
= \Pi_+\!\bigl[a_{k,-}^{-1}\wh{\wt f}\bigr](\xi)$.
Hence
\begin{equation}\label{eq:A-orderk-repr}
    \wh{\psi_+}(\xi)
    = a_{k,+}^{-1}(\xi)\Pi_+\!\bigl[a_{k,-}^{-1}\wh{\wt f}\bigr](\xi).
\end{equation}

\emph{Existence and uniqueness.}
Conversely, \eqref{eq:A-orderk-repr} defines a solution: if
$g=a_{k,-}^{-1}\wh{\wt f}$, then
$\wh{\wt f}-a_k\wh{\psi_+}=a_{k,-}\Pi_-g$, and the inverse Fourier
transform of the right-hand side is supported in $\overline{\R_-}$.
Thus
$\wh{\psi_+}(\xi)\in
\mathcal{F}(H^{\frac{k}{2}+\varepsilon}_{\overline{\R_+}}(\R))$ and
\begin{equation}
    \norm{\psi_+}_{H^{\frac{k}{2}+\varepsilon}(\R)}
    \le C \norm{\wt f}_{H^{-\frac{k}{2}+\varepsilon}(\R)}
    \le C \norm{f}_{H^{-\frac{k}{2}+\varepsilon}(\R_+)}.
\end{equation}
For homogeneous data, take $\wt f=0$ in
\eqref{eq:A-orderk-wh}. Applying $\Pi_+$ gives
$a_{k,+}\wh{\psi_+}=0$, hence $\psi_+=0$. This proves uniqueness in
the stated supported Sobolev class and also shows that
\eqref{eq:A-orderk-repr} is independent of the chosen extension.

\emph{Extraction of the endpoint term.}
For part (2), assume
$f\in H^{-\frac{k}{2}+1+\varepsilon}(\R_+)$ and choose an extension
$\wt f\in H^{-\frac{k}{2}+1+\varepsilon}(\R)$ such that
\begin{equation}
    \norm{\wt f}_{H^{-\frac{k}{2}+1+\varepsilon}(\R)}
    \le 2\norm{f}_{H^{-\frac{k}{2}+1+\varepsilon}(\R_+)}.
\end{equation}
Then the same representation~\eqref{eq:A-orderk-repr} holds and
\begin{equation}
    g\coloneqq a_{k,-}^{-1}\wh{\wt f}
    \in\mathcal F(H^{1+\varepsilon}(\R)).
\end{equation}
Let $h=\mathcal F^{-1}g$. The extra regularity of the forcing gives
$h\in H^{1+\varepsilon}(\R)$, so $h(0)$ is well defined since
$1+\varepsilon>\frac12$; this trace generates the explicit $x^{k/2}$
boundary term. Define
\begin{equation}\label{eq:Pi-prime-definition}
    \ell_0(g)
    \coloneqq
    -\frac{i}{\sqrt{2\pi}}h(0)
    =
    -\frac{i}{\sqrt{2\pi}}
    \left(\mathcal F^{-1}g\right)(0).
\end{equation}
The distributional identity
\begin{equation}
    \p_x\paren{\mathbf1_{\R_+}h}
    =\mathbf1_{\R_+}\p_xh+h(0)\delta_0
\end{equation}
and the conventions $\mathcal F(\p_xh)=i\xi\wh h$ and
$\mathcal F\delta_0=(2\pi)^{-1/2}$ give
\begin{equation*}
    i(\xi-i)\Pi_+g
    =i\Pi_+\!\left[(\xi-i)g\right]+\frac{h(0)}{\sqrt{2\pi}}.
\end{equation*}
Dividing by $i(\xi-i)$ yields
\begin{equation}\label{eq:Pi-first-order-expansion}
    \Pi_+g
    =(\xi-i)^{-1}\ell_0(g)
    +(\xi-i)^{-1}\Pi_+\!\left[(\xi-i)g\right].
\end{equation}
The term $(\xi-i)^{-1}\ell_0(g)$ comes from $h(0)\delta_0$ and will
determine the coefficient $c$.
Here $\Pi_+$ is applied to $(\xi-i)g\in\mathcal F(H^\varepsilon(\R))$,
not as a bounded projection on $\mathcal F(H^{1+\varepsilon}(\R))$.
Thus, with $b_+(\xi)=\xi-i$, the representation
\eqref{eq:A-orderk-repr} becomes
\begin{equation}\label{eq:A-orderk-expansion}
    \wh{\psi_+}(\xi)
    = a_{k,+}^{-1}(\xi)b_+^{-1}(\xi)
      \ell_0\!\bigl[a_{k,-}^{-1}\wh{\wt f}\bigr]
    + a_{k,+}^{-1}(\xi)b_+^{-1}(\xi)
      \Pi_+\!\bigl[b_+a_{k,-}^{-1}\wh{\wt f}\bigr](\xi).
\end{equation}
The second term in~\eqref{eq:A-orderk-expansion} belongs to $\mathcal{F}(H^{\frac{k}{2}+1+\varepsilon}_{\overline{\R_+}}(\R))$ because
$b_+a_{k,-}^{-1}$ has order $1-\frac{k}{2}$ and
$\Pi_+$ is bounded on $\mathcal F(H^\varepsilon(\R))$. Hence
\begin{equation}
    \psi_r
    \coloneqq \mathcal{F}^{-1}\!\left(
    a_{k,+}^{-1}b_+^{-1}
    \Pi_+\!\bigl[b_+a_{k,-}^{-1}\wh{\wt f}\bigr]\right)
    \in H^{\frac{k}{2}+1+\varepsilon}_{\overline{\R_+}}(\R).
\end{equation}
For the first term, note that
$a_{k,+}^{-1}(\xi)b_+^{-1}(\xi)
=(\xi-i)^{-\frac{k}{2}-1}$.
By \eqref{eq:halfline-factor-kernel} with $a=k/2+1$, its inverse
Fourier transform is
$\sqrt{2\pi}e^{i\pi(k/2+1)/2}/\Gamma(k/2+1)$ times
$\mathbf1_{[0,\infty)}(x)e^{-x}x^{\frac{k}{2}}$.
Hence the first term in
\eqref{eq:A-orderk-expansion} has the form
$c\,\mathbf1_{[0,\infty)}(x)e^{-x}x^{\frac{k}{2}}$,
and the trace and multiplier bounds give
\begin{equation}
    |c|
    +\norm{\psi_r}_{H^{\frac{k}{2}+1+\varepsilon}(\R)}
    \le C\norm{\mathcal F^{-1}\!\left(
    a_{k,-}^{-1}\wh{\wt f}\right)}_{H^{1+\varepsilon}(\R)}
    \le C \norm{f}_{H^{-\frac{k}{2}+1+\varepsilon}(\R_+)}.
\end{equation}
The coefficient $c$ is intrinsic. By~\eqref{eq:halfline-factor-kernel},
the Fourier transform of $\mathbf1_{[0,\infty)}e^{-x}x^{k/2}$ has modulus
a nonzero constant times $\langle\xi\rangle^{-k/2-1}$, so this profile
does not belong to $H^{k/2+1+\varepsilon}(\R)$ for
$\varepsilon>-\frac12$. Two such decompositions of the same solution
therefore have the same coefficient and remainder.
\end{proof}

Lemma~\ref{lem:A-orderk-halfline} is the half-line model. For an equation on the interval $I$, localization at the two endpoints of $I$ gives the following result.

\begin{lemma}\label{lem:A-orderk-interval}
Let $k\ge 1$ be an integer, let $0\le \varepsilon<\frac12$, and let
$I=(-1,1)$. Suppose $u\in H^{\frac{k}{2}}_{\overline I}(\R)$ satisfies
\begin{equation}\label{eq:A-orderk-interval}
    p_I A_k u = f \quad \text{on } I,
\end{equation}
with $f\in H^{-\frac{k}{2}+\varepsilon}(I)$.
The constants in these estimates depend only on $k$, $\varepsilon$, and
the fixed cutoffs.
\begin{enumerate}
    \item Then $u\in H^{\frac{k}{2}+\varepsilon}_{\overline I}(\R)$ and
\begin{equation}\label{eq:A-orderk-interval-est}
    \norm{u}_{H^{\frac{k}{2}+\varepsilon}(\R)}
    \le C\paren{
    \norm{f}_{H^{-\frac{k}{2}+\varepsilon}(I)}
    + \norm{u}_{H^{\frac{k}{2}}(\R)}} .
\end{equation}
    \item If in addition $f\in H^{-\frac{k}{2}+1+\varepsilon}(I)$, then
    \begin{equation}\label{eq:A-orderk-interval-split}
        u(s)=c_-\chi_-(s)e^{-(s+1)}(s+1)_+^{\frac{k}{2}}
        +c_+\chi_+(s)e^{-(1-s)}(1-s)_+^{\frac{k}{2}}+u_r(s),
    \end{equation}
    where $\chi_\pm$ are real-valued smooth cutoffs supported near the
    respective endpoints and equal to one there, with supports disjoint
    from the opposite endpoint, $c_\pm\in\mathbb C$, and
    $u_r\in H^{\frac{k}{2}+1+\varepsilon}_{\overline I}(\R)$.
    The coefficients are real when $u$ is real-valued. Moreover,
    \begin{equation}\label{eq:A-orderk-interval-split-est}
        |c_-|+|c_+|+\norm{u_r}_{H^{\frac{k}{2}+1+\varepsilon}(\R)}
        \le C\paren{
        \norm{f}_{H^{-\frac{k}{2}+1+\varepsilon}(I)}
        +\norm{u}_{H^{\frac{k}{2}}(\R)}} .
    \end{equation}
\end{enumerate}
\end{lemma}

\begin{proof}
We localize at the left endpoint and translate to the half-line, where
Lemma~\ref{lem:A-orderk-halfline} applies. Reflection treats the right
endpoint; the whole-line $A_k$-multiplier estimate treats the interior. We then
combine these three pieces.
The cutoffs introduce commutators, for which we use

\begin{equation}\label{eq:A-orderk-commutator}
    \norm{[A_k,\phi]v}_{H^{r-k+1}(\R)}
    \le C_{k,r,\phi}\norm{v}_{H^r(\R)},
    \quad r\in\R,\quad \phi\in C_c^\infty(\R).
\end{equation}
Indeed, its Fourier kernel is
$(2\pi)^{-1/2}(\langle\xi\rangle^k-\langle\eta\rangle^k)
\wh\phi(\xi-\eta)$. The mean-value estimate for the symbol, Peetre's
inequality, and rapid decay of $\wh\phi$ give
\eqref{eq:A-orderk-commutator} by Young's inequality. Thus the
commutator has order $k-1$, although it need not have compactly supported
output.

Let $\wt f\in H^{-\frac{k}{2}+\varepsilon}(\R)$ be an extension of $f$ with
$\norm{\wt f}_{H^{-\frac{k}{2}+\varepsilon}(\R)}
\le C\norm{f}_{H^{-\frac{k}{2}+\varepsilon}(I)}$. Choose real-valued smooth cutoffs
$\vph_-$, $\vph_+$, and $\vph_0$ such that $\vph_-$ is supported in a small
neighborhood of $s=-1$, $\vph_+$ is supported in a small neighborhood of
$s=1$, $\vph_0$ is compactly supported in $I$, and
$\vph_-+\vph_++\vph_0=1$ on $I$. Take $\vph_\pm=1$ near the respective
endpoint, with support disjoint from the other endpoint. Set
$\chi_\pm=\vph_\pm$ in the statement.

First consider the left endpoint. Shift $s=-1$ to the origin by setting
$x=s+1$, and define $u_-(x)=\vph_-(s)u(s)$ for $x>0$ and $u_-(x)=0$ for
$x<0$. Under this change of variables,
$\operatorname{supp}\vph_-\cap I$ lies in $\R_+$. Since $\vph_-$
vanishes near and to the right of $s=1$, the identity
$\vph_-A_ku=\vph_-\wt f$ holds on $s>-1$. Thus the localized equation,
with all terms on the right translated to the $x$ variable, is
\begin{equation}\label{eq:A-orderk-left-local}
    p_+ A_k u_-
    = p_+\paren{\vph_-\wt f + [A_k,\vph_-]u}.
\end{equation}
The first term satisfies
\begin{equation}
    \norm{p_+\vph_-\wt f}_{H^{-\frac{k}{2}+\varepsilon}(\R_+)}
    \le C\norm{f}_{H^{-\frac{k}{2}+\varepsilon}(I)}.
\end{equation}
For the commutator term, \eqref{eq:A-orderk-commutator} with $r=k/2$
gives
\begin{equation}
    \norm{p_+[A_k,\vph_-]u}_{H^{-\frac{k}{2}+\varepsilon}(\R_+)}
    \le \norm{[A_k,\vph_-]u}_{H^{-\frac{k}{2}+\varepsilon}(\R)}
    \le C\norm{u}_{H^{\frac{k}{2}}(\R)}.
\end{equation}
Here we used $\varepsilon<1$ to embed $H^{1-k/2}(\R)$ into
$H^{-k/2+\varepsilon}(\R)$. Thus \eqref{eq:A-orderk-left-local} has
right-hand side in $H^{-k/2+\varepsilon}(\R_+)$.
Lemma~\ref{lem:A-orderk-halfline}(1), with the stated $\varepsilon$, gives
a solution $w_-\in H^{\frac{k}{2}+\varepsilon}_{\overline{\R_+}}(\R)$
of \eqref{eq:A-orderk-left-local}. Since both $u_-$ and $w_-$ belong to
$H^{\frac{k}{2}}_{\overline{\R_+}}(\R)$ and solve the same equation,
uniqueness in Lemma~\ref{lem:A-orderk-halfline}(1) with $\varepsilon=0$
gives $u_-=w_-$. Hence
\begin{equation}\label{eq:A-orderk-left-est}
    \norm{\vph_-u}_{H^{\frac{k}{2}+\varepsilon}(\R)}
    \le C\paren{
    \norm{f}_{H^{-\frac{k}{2}+\varepsilon}(I)}
    + \norm{u}_{H^{\frac{k}{2}}(\R)}} .
\end{equation}

At the right endpoint, reflect by $x=1-s$. Since the symbol
$(1+\xi^2)^{k/2}$ is even, the same argument gives
\begin{equation}\label{eq:A-orderk-right-est}
    \norm{\vph_+u}_{H^{\frac{k}{2}+\varepsilon}(\R)}
    \le C\paren{
    \norm{f}_{H^{-\frac{k}{2}+\varepsilon}(I)}
    + \norm{u}_{H^{\frac{k}{2}}(\R)}} .
\end{equation}

For the interior piece, $\vph_0$ is compactly supported in $I$, so we work
on the whole line:
\begin{equation}\label{eq:A-orderk-int-local}
    A_k(\vph_0u)=\vph_0\wt f + [A_k,\vph_0]u .
\end{equation}
As in the proof of \eqref{eq:A-orderk-left-est},
\begin{equation}
    \norm{\vph_0\wt f + [A_k,\vph_0]u}_{H^{-\frac{k}{2}+\varepsilon}(\R)}
    \le C\paren{
    \norm{f}_{H^{-\frac{k}{2}+\varepsilon}(I)}
    + \norm{u}_{H^{\frac{k}{2}}(\R)}} .
\end{equation}
The whole-line $A_k$-multiplier estimate then gives
\begin{equation}\label{eq:A-orderk-int-est}
    \norm{\vph_0u}_{H^{\frac{k}{2}+\varepsilon}(\R)}
    \le C\paren{
    \norm{f}_{H^{-\frac{k}{2}+\varepsilon}(I)}
    + \norm{u}_{H^{\frac{k}{2}}(\R)}} .
\end{equation}
Combining \eqref{eq:A-orderk-left-est}, \eqref{eq:A-orderk-right-est}, and
\eqref{eq:A-orderk-int-est}, and using
$u=\vph_-u+\vph_+u+\vph_0u$, proves part (1).

For part (2), the Sobolev gain from part (1) improves the commutator
terms and allows extraction of the endpoint profiles.
Assume $f\in H^{-\frac{k}{2}+1+\varepsilon}(I)$ and choose the
extension $\wt f\in H^{-\frac{k}{2}+1+\varepsilon}(\R)$. By part (1),
$u\in H^{\frac{k}{2}+\varepsilon}_{\overline I}(\R)$. Hence the commutator
terms in \eqref{eq:A-orderk-left-local} and \eqref{eq:A-orderk-int-local}
satisfy, by \eqref{eq:A-orderk-commutator} with $r=k/2+\varepsilon$,
\begin{equation}
\begin{aligned}
    \norm{p_+[A_k,\vph_-]u}_{H^{-\frac{k}{2}+1+\varepsilon}(\R_+)}
    +\norm{[A_k,\vph_0]u}_{H^{-\frac{k}{2}+1+\varepsilon}(\R)}
    \le C\norm{u}_{H^{\frac{k}{2}+\varepsilon}(\R)} .
\end{aligned}
\end{equation}
Together with \eqref{eq:A-orderk-interval-est}, this is bounded by
$C\paren{\norm{f}_{H^{-\frac{k}{2}+1+\varepsilon}(I)}
+\norm{u}_{H^{\frac{k}{2}}(\R)}}$.
Thus the right-hand side of the left endpoint equation
\eqref{eq:A-orderk-left-local} belongs to
$H^{-\frac{k}{2}+1+\varepsilon}(\R_+)$. Applying
Lemma~\ref{lem:A-orderk-halfline}(2) gives
\begin{equation}
    \vph_-u
    =c_-\mathbf1_{[0,\infty)}(x)e^{-x}x^{\frac{k}{2}}+u_{-,r},
    \quad u_{-,r}\in H^{\frac{k}{2}+1+\varepsilon}_{\overline{\R_+}}(\R),
\end{equation}
with
\begin{equation}
    |c_-|+\norm{u_{-,r}}_{H^{\frac{k}{2}+1+\varepsilon}(\R)}
    \le C\paren{
    \norm{f}_{H^{-\frac{k}{2}+1+\varepsilon}(I)}
    +\norm{u}_{H^{\frac{k}{2}}(\R)} } .
\end{equation}
Returning to the variable $s$ gives the singular term
$c_-e^{-(s+1)}(s+1)_+^{\frac{k}{2}}$. Redefine the left remainder by
\begin{equation*}
    u_{-,r}(s)=\vph_-(s)u(s)
    -c_-\vph_-(s)e^{-(s+1)}(s+1)_+^{\frac{k}{2}}.
\end{equation*}
This adds
\begin{equation*}
    c_-(1-\vph_-(s))e^{-(s+1)}(s+1)_+^{k/2}
\end{equation*}
to the translated half-line remainder. Since $1-\vph_-$ vanishes near
$s=-1$, this correction is smooth across that endpoint; it vanishes for
$s<-1$ and decays exponentially with all derivatives as $s\to+\infty$.
It belongs to every $H^r(\R)$, with norm bounded by a constant times
$|c_-|$, so the same Sobolev estimate holds for the redefined remainder.
Moreover, both terms in this definition are supported in
$\overline I$; hence
$u_{-,r}\in H^{k/2+1+\varepsilon}_{\overline I}(\R)$.
The right endpoint, after the
reflection $x=1-s$, gives
$c_+\vph_+(s)e^{-(1-s)}(1-s)_+^{\frac{k}{2}}$ and a remainder
$u_{+,r}\in H^{\frac{k}{2}+1+\varepsilon}_{\overline I}(\R)$ with the same
estimate, defined in the same way.
The whole-line $A_k$-multiplier estimate applied to
\eqref{eq:A-orderk-int-local} yields
$\vph_0u\in H^{\frac{k}{2}+1+\varepsilon}(\R)$ with the same bound.
Setting
$u_r=u_{-,r}+u_{+,r}+\vph_0u$ proves
\eqref{eq:A-orderk-interval-split} and
\eqref{eq:A-orderk-interval-split-est}.
Localizing the difference of two decompositions near each endpoint and
using the nonmembership of the half-line profile in the remainder space
shows that $c_\pm$ are unique. For real-valued $u$, complex conjugation
gives another decomposition with the same real cutoffs and profiles, so
uniqueness implies $c_\pm\in\R$. This applies in particular to the
curvature expansion in Theorem~\ref{thm:main-regularity}.
\end{proof}

\subsection{Commutator estimates}
\label{subsec:commutator-estimates}

Commutators of $\p_s\mc H$ with the frame fields appear
in~\eqref{eqn:Rop-def}--\eqref{eqn:N-def}.
Lemmas~\ref{dsH_f_comm}--\ref{lem:tn-comm} control them at the low
regularity used for the tension equation and local well-posedness.
Lemma~\ref{lem:comm-psH} gives the higher-order estimate needed for
positive-time regularity and reconstruction.

\begin{lemma}\label{dsH_f_comm}
Let $0\le\varepsilon<\alpha<\delta\le\frac12$,
$f\in H^{\frac12+\delta}(I)$, and $g\in\wt H^{\frac12}(I)$. Then, with
$C=C(\varepsilon,\alpha,\delta)$,
\begin{equation}\label{eqn:commutator-low-bound}
    \norm{[\p_s\mc H,f]g}_{\hi{-\frac12+\varepsilon}}
    \le C\paren{[f]_{\ci{0,\alpha}}
    +\norm{\p_s f}_{\hi{-\frac12+\delta}}}
    \norm{g}_{H^{\frac12}(\R)}.
\end{equation}
\end{lemma}
\begin{proof}
We separate the commutator into a multiplication term and a Hilbert-transform
commutator:
\begin{equation}
    [\p_s\mc H,f]g
    =\p_s f\,\mc Hg+\p_s\paren{[\mc H,f]g}.
\end{equation}
For $\vph\in\thi{\frac12-\varepsilon}$, duality and
Lemmas~\ref{point_mult_sobo} and~\ref{Hf-bounded} give
\begin{equation}
\begin{aligned}
    \abs{\dual{\p_s f\mc H g}{\vph}}
    &\le C\norm{\vph}_{H^{\frac12-\varepsilon}(\R)}
    \norm{\p_s f}_{\hi{-\frac12+\delta}}
    \norm{\mc Hg}_{H^{\frac12}(I)} \\
    &\le C\norm{\vph}_{H^{\frac12-\varepsilon}(\R)}
    \norm{\p_s f}_{\hi{-\frac12+\delta}}
    \norm{g}_{H^{\frac12}(\R)}.
\end{aligned}
\end{equation}
Taking the supremum over $\vph$ gives
\begin{equation}\label{eqn:dsHf1}
    \norm{\p_s f\,\mc Hg}_{\hi{-\frac12+\varepsilon}}
    \le C\norm{\p_s f}_{\hi{-\frac12+\delta}}
    \norm{g}_{H^{\frac12}(\R)}.
\end{equation}

It remains to estimate $[\mc H,f]g$. Set
\begin{equation}
    q=\frac12+\varepsilon,\quad \beta=\delta-\varepsilon>0.
\end{equation}
Since $[\mc H,f]$ is unchanged when a constant is added to $f$, there exists a constant $c\in\R$ and a bounded extension
$F\in\hr{\frac12+\delta}$ of $f-c$ satisfying
\begin{equation}
    \norm{F}_{\hr{\frac12+\delta}}
    \le C\paren{[f]_{\ci{0,\alpha}}
    +\norm{\p_s f}_{\hi{-\frac12+\delta}}}.
\end{equation}
Identifying $g$ with its zero extension, we have
\begin{equation}
    [\mc H,f]g=\left.[\mc H_{\R},F]g\right|_I.
\end{equation}

The Fourier representation of the whole-line commutator is
\begin{equation}
    \wh{[\mc H_{\R},F]g}(\xi)
    =c_0\int_\R
    \paren{\operatorname{sgn}(\xi)-\operatorname{sgn}(\eta)}
    \wh F(\xi-\eta)\wh g(\eta)\,d\eta.
\end{equation}
The difference of the signs vanishes unless $\xi\eta<0$. This cancellation
implies $\abs{\xi}\le\abs{\xi-\eta}$ on the support of the integrand, so
\begin{equation}
    \langle\xi\rangle^q
    \abs{\wh{[\mc H_{\R},F]g}(\xi)}
    \le C\int_\R
    \langle\xi-\eta\rangle^q
    \abs{\wh F(\xi-\eta)}\abs{\wh g(\eta)}\,d\eta.
\end{equation}
Define $U$ and $V$ by
\begin{equation}
    \wh U(\xi)=\langle\xi\rangle^q\abs{\wh F(\xi)},
    \quad
    \wh V(\xi)=\abs{\wh g(\xi)}.
\end{equation}
Plancherel's theorem and the whole-line multiplication estimate
\cite[Theorem~7.3]{Behzadan2021}, applied with $s=0$, $s_1=\beta$, and
$s_2=\frac12$, give
\begin{equation}
\begin{aligned}
    \norm{[\mc H_{\R},F]g}_{\hr{q}}
    &\le C\norm{UV}_{L^2(\R)}
    \le C\norm{U}_{\hr{\beta}}\norm{V}_{\hr{\frac12}} \\
    &=C\norm{F}_{\hr{q+\beta}}\norm{g}_{\hr{\frac12}}
    =C\norm{F}_{\hr{\frac12+\delta}}\norm{g}_{\hr{\frac12}}.
\end{aligned}
\end{equation}
The multiplication estimate applies because $\beta>0$. Restriction to $I$
and the bound for $F$ now give
\begin{equation}\label{eqn:dsHf2}
    \norm{\p_s[\mc H,f]g}_{\hi{-\frac12+\varepsilon}}
    \le C\norm{[\mc H,f]g}_{\hi{\frac12+\varepsilon}}
    \le C\paren{[f]_{\ci{0,\alpha}}
    +\norm{\p_s f}_{\hi{-\frac12+\delta}}}
    \norm{g}_{\hr{\frac12}}.
\end{equation}
Combining \eqref{eqn:dsHf1} and~\eqref{eqn:dsHf2} proves the claim.
\end{proof}


Combining~\eqref{eqn:commutator-low-bound} with the reconstruction estimates
specializes Lemma~\ref{dsH_f_comm} to the tangent and normal fields.
\begin{lemma}\label{lem:tn-comm}
Let $0<\varepsilon<\delta\le\frac12$,
$\kappa\in\thi{-\frac12+\delta}$, and $f\in\thi{\frac12}$. Then
\begin{equation}
    \norm{[\p_s\mc H,a]f}_{\hi{-\frac12+\varepsilon}}
    \le C\paren{1+\norm{\kappa}_{\hr{-\frac12+\delta}}}
    \norm{\kappa}_{\hr{-\frac12+\delta}}
    \norm{f}_{\hr{\frac12}},\quad a\in\{\tk,\nk\}.
\end{equation}
If $\delta=\frac12$ and $\kappa\in\wt L^2(I)$, then the sharper estimate
\begin{equation}
    \norm{[\p_s\mc H,a]f}_{\hi{-\frac12+\varepsilon}}
    \le C\norm{\kappa}_{L^2(\R)}\norm{f}_{\hr{\frac12}},
    \quad a\in\{\tk,\nk\},
\end{equation}
holds for every $f\in\thi{\frac12}$.
\end{lemma}
\begin{proof}
Choose $\alpha\in(\varepsilon,\delta)$. We apply
Lemma~\ref{dsH_f_comm} with the multipliers $\tk$ and $\nk$.
Since $\tk=\xk'$, Lemma~\ref{lem:Xk-holder} with $m=0$ gives
\begin{equation}
    [\tk]_{\ci{0,\alpha}}\le C\norm{\kappa}_{\hr{-\frac12+\delta}},
\end{equation}
and Lemma~\ref{tau-n-est} gives
\begin{equation}
    \norm{\p_s\tk}_{\hi{-\frac12+\delta}}
    =\norm{\xk''}_{\hi{-\frac12+\delta}}
    \le C\paren{1+\norm{\kappa}_{\hr{-\frac12+\delta}}}
        \norm{\kappa}_{\hr{-\frac12+\delta}}.
\end{equation}
Since $\nk=Q_{\frac{\pi}{2}}^{-1}\tk$, the same two estimates hold for
$\nk$. Substitution into Lemma~\ref{dsH_f_comm} proves the first estimate.

For $\delta=\frac12$, the Frenet identities
$\p_s\tk=-\kappa\nk$ and $\p_s\nk=\kappa\tk$, together with
$\abs{\tk}=\abs{\nk}=1$, yield
\begin{equation}
    [\tk]_{\ci{0,\alpha}}+[\nk]_{\ci{0,\alpha}}
    +\norm{\p_s\tk}_{L^2(I)}+\norm{\p_s\nk}_{L^2(I)}
    \le C\norm{\kappa}_{L^2(\R)}.
\end{equation}
Applying Lemma~\ref{dsH_f_comm} once more gives the $L^2$-based estimate.
\end{proof}


\begin{lemma}\label{lem:comm-psH}
Let $0\le\varepsilon<\delta\le\frac12$,
$f\in H^{\frac52+\delta}(I)$, and
$g\in\thi{\frac12+\varepsilon}$. Then
\begin{equation}\label{eq:comm-psH-bound}
    \norm{[\p_s\mc H,f]g}_{\hi{\frac12+\varepsilon}}
    \le C\norm{f}_{H^{\frac52+\delta}(I)}
    \norm{g}_{\hr{\frac12+\varepsilon}},
\end{equation}
where $C$ depends only on $\delta$ and $\varepsilon$.
\end{lemma}

The proof is given in Appendix~\ref{app:commutator-estimates}.

\subsection{Boundary integral estimates}
\label{sec:divided-diff}

We control the Stokes remainder by deriving pointwise and finite-difference
kernel bounds, then converting them into Sobolev mapping estimates.
Lemma~\ref{LM-kernel-est} gives the kernel bounds, and
Lemma~\ref{K-op} supplies the abstract weakly singular operator estimate
used in the mapping and difference bounds of
Lemmas~\ref{LM-op-est}--\ref{lem:LM-Lip}.
Lemmas~\ref{lem:Sop-est}--\ref{lem:Sop-Lip} specialize these bounds to
$S_\kappa$.
Finally, Lemma~\ref{lem:stokes-layer-open-arc} supplies the Stokes trace
estimate and energy identity.

For $s,\eta,s+h\in I$ and $h\in\R$, define
\begin{align}
    \Delta f(s,\eta)&\coloneqq f(s)-f(\eta), \\
    T_{s,h}f(s,\eta)&\coloneqq f(s+h,\eta), \\
    \Delta_{s,h}f(s,\eta)&\coloneqq f(s+h,\eta)-f(s,\eta).
\end{align}
For $s,\eta\in I$ with $s\ne\eta$, let
\begin{equation}\label{eqn:LM-kernel}
    \ell(s,\eta;\bm Z)
    \coloneqq\log\abs{\frac{\Delta\bm Z}{s-\eta}},
    \quad
    m_{ij}(s,\eta;\bm Z)
    \coloneqq\frac{\Delta Z_i\Delta Z_j}{\abs{\Delta\bm Z}^2}.
\end{equation}
The associated operators are
\begin{align}
    (L_{\bm Z}f)(s)
    &\coloneqq\int_I\p_\eta \ell(s,\eta;\bm Z)f(\eta)\,d\eta,
    \label{eqn:F-def}\\
    (M_{\bm Z,ij}f)(s)
    &\coloneqq\int_I\p_\eta m_{ij}(s,\eta;\bm Z)f(\eta)\,d\eta,
    \label{eqn:M-def}\\
    (S_{\kappa,ij}f)(s)
    &\coloneqq\bigl(\delta_{ij}L_{\bm X_\kappa}f
    -M_{\bm X_\kappa,ij}f\bigr)(s) = \int_I\p_\eta r_{ij}(s,\eta;\kappa)f(\eta)\,d\eta,\label{eqn:Sop-definition}
\end{align}
for $i,j=1,2$, where $r_{ij}$ is given in~\eqref{eqn:R-kernel}.
For vector-valued inputs, write $S_\kappa=(S_{\kappa,ij})_{i,j=1}^2$, so that
\begin{equation*}
    (S_\kappa\bm f)(s)
    =\int_I\p_\eta r(s,\eta;\kappa)\bm f(\eta)\,d\eta.
\end{equation*}


The divided-difference estimates below are open interval versions of those in
\cite[Section~2]{MoriRodenbergSpirn2019Peskin}.
\begin{lemma}\label{LM-kernel-est}
Let $\alpha\in(0,1)$ and
$\bm Z=(Z_1,Z_2)^T\in\ci{1,\alpha}$ satisfy $\abs{\bm Z}_*>0$.
Then
\begin{align}
    \abs{\p_s\ell(s,\eta;\bm Z)},\quad
    \abs{\p_\eta \ell(s,\eta;\bm Z)}
    &\le C\abs{\bm Z}_*^{-1}[\bm Z']_{\ci{0,\alpha}}
    \abs{s-\eta}^{\alpha-1},\label{eqn:LM-1}\\
    \abs{\p_sm_{ij}(s,\eta;\bm Z)},\quad
    \abs{\p_\eta m_{ij}(s,\eta;\bm Z)}
    &\le C\abs{\bm Z}_*^{-1}[\bm Z']_{\ci{0,\alpha}}
    \abs{s-\eta}^{\alpha-1}.
    \label{eqn:LM-2}
\end{align}
If $s,\eta,s+h\in I$ and $\abs{s-\eta}\ge2\abs h$, then
\begin{align}
    \abs{\Delta_{s,h}\paren{\p_\eta \ell(s,\eta;\bm Z)}},\quad
    \abs{\Delta_{s,h}\paren{\p_s\ell(s,\eta;\bm Z)}}
    &\le C\abs{\bm Z}_*^{-2}\norm{\bm Z'}_{\ci{0}}[\bm Z']_{\ci{0,\alpha}}\abs h^\alpha\abs{s-\eta}^{-1},
    \label{eqn:LM-3}\\
    \abs{\Delta_{s,h}\paren{\p_\eta m_{ij}(s,\eta;\bm Z)}},\quad
    \abs{\Delta_{s,h}\paren{\p_sm_{ij}(s,\eta;\bm Z)}}
    &\le C\abs{\bm Z}_*^{-2}\norm{\bm Z'}_{\ci{0,\alpha}}[\bm Z']_{\ci{0,\alpha}}\abs h^\alpha\abs{s-\eta}^{-1}.
    \label{eqn:LM-4}
\end{align}
\end{lemma}

The proof is given in Appendix~\ref{app:boundary-integral-estimates}.


We use the following general mapping estimate for integral operators with weakly singular kernels.
\begin{lemma}\label{K-op}
Let $0\le\varepsilon<\alpha<1$.
For $f\in\wt L^2(I)$, define
\begin{equation}
    K f  (s) = \int_I k(s,\eta) f(\eta)\,d\eta.
\end{equation}
Suppose that the pointwise bound
\begin{equation}\label{eqn:kernel-pointwise-hypothesis}
    \abs{k(s,\eta)} \le a_1 \abs{s-\eta}^{\alpha-1}
\end{equation}
holds for all $s,\eta\in I$ with $s\ne\eta$, and that
\begin{equation}\label{eqn:kernel-increment-hypothesis}
    \abs{k(s+h,\eta)-k(s,\eta)}
    \le a_2\abs{h}^\alpha\abs{s-\eta}^{-1}
\end{equation}
whenever $s,\eta,s+h\in I$, $h\in\R$, and
$\abs{s-\eta}\ge2\abs h$.
Then there exists a constant $C$, independent of $f,a_1,a_2$, such that
\begin{equation}\label{eqn:kernel-mapping-bound}
    \norm{K f}_{\hi{\varepsilon}} \le C \paren{a_1+a_2}\norm{f}_{L^2(\R)}.
\end{equation}
\end{lemma}
The proof combines the Schur test with the difference characterization of
$H^\varepsilon(I)$ and is given in Appendix~\ref{app:boundary-integral-estimates}.


\begin{lemma}\label{LM-op-est}
Let $0\le\varepsilon<\alpha\le\frac12$ and
$\bm Z=(Z_1,Z_2)^T\in\ci{1,\alpha}$ satisfy $\abs{\bm Z}_*>0$.
Then, for every $r\in[0,1]$ and $f\in\thi{r}$,
\begin{align}
    \norm{L_{\bm Z}f}_{H^{r+\varepsilon}(I)}
    &\le C\abs{\bm Z}_*^{-2}\norm{\bm Z'}_{\ci{0}}[\bm Z']_{\ci{0,\alpha}}\norm{f}_{H^r(\R)}, \\
    \norm{M_{\bm Z,ij}f}_{H^{r+\varepsilon}(I)}
    &\le C\abs{\bm Z}_*^{-2}\norm{\bm Z'}_{\ci{0,\alpha}}[\bm Z']_{\ci{0,\alpha}}\norm{f}_{H^r(\R)}.
\end{align}
\end{lemma}
\begin{proof}
We prove the endpoint estimates $r=0$ and $r=1$ and then interpolate.
For $f\in\wt L^2(I)$, \eqref{eqn:LM-1}--\eqref{eqn:LM-4} show that
the kernels $\p_\eta\ell$ and $\p_\eta m_{ij}$ satisfy~\eqref{eqn:kernel-pointwise-hypothesis}--\eqref{eqn:kernel-increment-hypothesis}. Lemma~\ref{K-op} therefore gives
\begin{align}
    \norm{L_{\bm Z}f}_{\hi{\varepsilon}}
    &\le C\abs{\bm Z}_*^{-2}\norm{\bm Z'}_{\ci{0}}[\bm Z']_{\ci{0,\alpha}}\norm{f}_{L^2(\R)},\\
    \norm{M_{\bm Z,ij}f}_{\hi{\varepsilon}}
    &\le C\abs{\bm Z}_*^{-2}\norm{\bm Z'}_{\ci{0,\alpha}}[\bm Z']_{\ci{0,\alpha}}\norm{f}_{L^2(\R)}.
\end{align}

For the $r=1$ endpoint, let $f\in\wt H^1(I)$. Integration by parts gives
\begin{align}
    \p_s(L_{\bm Z}f)(s)
    &=-\int_I\p_s\ell(s,\eta;\bm Z)\p_\eta f(\eta)\,d\eta,\\
    \p_s(M_{\bm Z,ij}f)(s)
    &=-\int_I\p_sm_{ij}(s,\eta;\bm Z)\p_\eta f(\eta)\,d\eta.
\end{align}
The boundary terms vanish because $f\in\wt H^1(I)=H_0^1(I)$. Applying
Lemmas~\ref{K-op} and~\ref{LM-kernel-est} to the kernels on the right,
and combining the resulting derivative estimates with the $r=0$ bounds,
gives
\begin{align}
    \norm{L_{\bm Z}f}_{\hi{1+\varepsilon}}
    &\le C\abs{\bm Z}_*^{-2}\norm{\bm Z'}_{\ci{0}}[\bm Z']_{\ci{0,\alpha}}\norm{f}_{H^1(\R)},\\
    \norm{M_{\bm Z,ij}f}_{\hi{1+\varepsilon}}
    &\le C\abs{\bm Z}_*^{-2}\norm{\bm Z'}_{\ci{0,\alpha}}[\bm Z']_{\ci{0,\alpha}}\norm{f}_{H^1(\R)}.
\end{align}
Finally, the interpolation identities
$\paren{\wt L^2(I),\wt H^1(I)}_{r,2}=\wt H^r(I)$ and
$\paren{H^\varepsilon(I),H^{1+\varepsilon}(I)}_{r,2}
=H^{r+\varepsilon}(I)$ prove the result.
\end{proof}


% \begin{lemma}[LM Lipschitz ]

% \end{lemma}
% \begin{proof}
% \begin{equation}
% \begin{aligned}
%     &\abs{\p_s \ell(s,\eta;\bm Y) - \p_s \ell(s,\eta;\bm X)} = \abs{ \int_0^1 \frac{d}{d\eta} \p_s \ell(s,\eta;\bm X + \eta(\bm Y-\bm X)) \,d\eta }\\
%     \le \int_0^1 \abs{ \p_{\bm X}\paren{\p_s \ell(s,\eta;\bm X + \eta(\bm Y-\bm X))} \paren{\bm Y-\bm X}} \, d\eta \\
%     \le C \int_0^1 \frac{\norm{\paren{\bm X + \eta(\bm Y-\bm X)}'}_{\ci{0,\alpha}}}{\abs{\bm X + \eta(\bm Y-\bm X)}_*^2} d\eta\, \norm{(\bm Y- \bm X)'}_{\ci{0,\alpha}}\abs{s-\eta}^{\alpha-1}.
% \end{aligned}
% \end{equation}
% \end{proof}

\begin{lemma}\label{lem:LM-Lip}
Let $0\le\varepsilon<\alpha\le\frac12$ and
$\bm X,\bm Y\in\ci{1,\alpha}$ satisfy
$\norm{\bm X'}_{\ci{0,\alpha}},\norm{\bm Y'}_{\ci{0,\alpha}}\le M$ and
$\abs{\rho\bm X+(1-\rho)\bm Y}_*\ge m$ for every $\rho\in[0,1]$, where
$M,m>0$.
% Let
% \begin{equation}
%     \ell(s,\eta;\bm X) = \log\abs{\frac{\Delta\bm X}{s-\eta}} , \quad m_{ij}(s,\eta;\bm X) = \frac{\Delta X_i \Delta X_j}{\abs{\Delta\bm X}^2}.
% \end{equation}
% and
% \begin{align}
%     &L_{\bm X} f (s) = \int_I \p_\eta \ell(s,\eta;\bm X) f(\eta) \,d\eta,\\
%     &M_{\bm X,ij} f (s) = \int_I \p_\eta m_{ij}(s,\eta;\bm X) f(\eta)\,d\eta,
% \end{align}
Then, for every $f\in\thi{s}$ with $s\in[0,1]$,
\begin{align}
    \norm{L_{\bm X}f-L_{\bm Y}f}_{\hi{s+\varepsilon}}
    &\le C_{M,m}
    \norm{\bm X'-\bm Y'}_{\ci{0,\alpha}}\norm{f}_{\hr{s}},
    \label{eqn:Lop-Lip}\\
    \norm{M_{\bm X,ij}f-M_{\bm Y,ij}f}_{\hi{s+\varepsilon}}
    &\le C_{M,m}
    \norm{\bm X'-\bm Y'}_{\ci{0,\alpha}}\norm{f}_{\hr{s}}.
    \label{eqn:Mop-Lip}
\end{align}
Here $L_{\bm Z}$ and $M_{\bm Z,ij}$ are defined in
\eqref{eqn:F-def} and~\eqref{eqn:M-def}.
\end{lemma}

The proof is given in Appendix~\ref{app:boundary-integral-estimates}.


We now specialize the geometric mapping and difference estimates to the
reconstructed filament $\xk$, obtaining the corresponding bounds for
$S_\kappa$.

\begin{lemma}\label{lem:Sop-est}
Let $0\le\varepsilon<\delta\le\frac12$ and
$\kappa\in\thi{-\frac12+\delta}$ satisfy $\abs{\kappa}_*>0$.
Then, for every $s\in[0,1]$ and $f\in\thi{s}$,
\begin{equation}\label{eqn:Sop-bound}
    \norm{S_{\kappa,ij}f}_{\hi{s+\varepsilon}}
    \le C\abs{\kappa}_*^{-2}
    \norm{\kappa}_{\hr{-\frac12+\delta}}
    \paren{1+\norm{\kappa}_{\hr{-\frac12+\delta}}}
    \norm{f}_{\hr{s}}.
\end{equation}
\end{lemma}
\begin{proof}
Set $\alpha=\frac12(\varepsilon+\delta)$, so that
$\varepsilon<\alpha<\delta$. The reconstruction estimate
\eqref{eqn:Xk-holder} gives $\xk\in\ci{1,\alpha}$ and
\begin{equation}
    [\xk']_{\ci{0,\alpha}} \le C\norm{\kappa}_{\hr{-\frac12+\delta}}, \quad \norm{\xk'}_{\ci{0}} \le
    C\paren{1+\norm{\kappa}_{\hr{-\frac12+\delta}}}.
\end{equation}
Since $S_{\kappa,ij}=\delta_{ij}L_{\xk}-M_{\xk,ij}$ by
\eqref{eqn:Sop-definition}, the result follows from
Lemma~\ref{LM-op-est}.
\end{proof}


\begin{lemma}\label{lem:Sop-Lip}
Let $0\le\varepsilon<\delta\le\frac12$ and
$\kappa_1,\kappa_2\in\thi{-\frac12+\delta}$ satisfy
$\norm{\kappa_i}_{\hr{-\frac12+\delta}}\le M$ for $i=1,2$. Assume also that
$\abs{\rho\xka{1}+(1-\rho)\xka{2}}_*\ge m$ for every $\rho\in[0,1]$.
Then, for every $s\in[0,1]$ and $f\in\thi{s}$,
\begin{equation}\label{eqn:Sop-difference}
    \norm{S_{\kappa_1,ij}f-S_{\kappa_2,ij}f}
    _{\hi{s+\varepsilon}}
    \le C_{M,m}
    \norm{\kappa_1-\kappa_2}_{\hr{-\frac12+\delta}}
    \norm{f}_{\hr{s}}.
\end{equation}
The dependence of $C_{M,m}$ on $\delta$ and $\varepsilon$ is suppressed.
\end{lemma}
\begin{proof}
Set $\alpha=\frac12(\varepsilon+\delta)$.
% Since $\abs{\kappa}_* = \abs{\xk}_*\le \norm{\bm X_\kappa'}_{\ci{0}}\le C\paren{1 + \norm{\kappa_0}_{\hr{-\frac{1}{2}+\delta}}}$, there holds that
% \begin{equation}
%     \norm{\kappa_i}_{\hr{-\frac{1}{2}+\delta}} \le \norm{\kappa_0}_{\hr{-\frac{1}{2}+\delta}} + \frac{1}{2C_{J,\delta}}\abs{\kappa_0}_* \le C \paren{1 + \norm{\kappa_0}_{\hr{-\frac{1}{2}+\delta}}}.
% \end{equation}
Lemma~\ref{lem:Xk-holder}, with $m=0$, gives
\begin{equation}
    \norm{\bm X_{\kappa_i}'}_{\ci{0,\alpha}}
    \le C\paren{1+\norm{\kappa_i}_{\hr{-\frac12+\delta}}}
    \le C_{M,m}.
\end{equation}
% For any $\bm X,\bm Y\in \ci{1}$ with $\abs{\bm X}_*, \abs{\bm Y}_*>0$, we have
% \begin{equation}
% \begin{aligned}
%     &\abs{\abs{\bm X}_* - \abs{\bm Y}_*} = \abs{ \inf_{s\neq t}\abs{\frac{\Delta \bm X}{s-\eta}} - \inf_{s\neq t}\abs{\frac{\Delta \bm Y}{s-\eta}} } \le \sup_{s\neq t}\abs{\abs{\frac{\Delta \bm X}{s-\eta}} - \abs{\frac{\Delta \bm Y}{s-\eta}} } \\
%     & \le \sup_{s\neq t} \abs{ \frac{\Delta \paren{\bm X-\bm Y}}{s-\eta} } = \sup_{s\neq t} \abs{ \int_0^1 \paren{\bm X-\bm Y}'(t + \eta (s-\eta)) \,d\eta } \le \norm{\bm X'-\bm Y'}_{\ci{0}}.
% \end{aligned}
% \end{equation}
% Since we have, for $i=1,2$,
% \begin{equation}
% \begin{aligned}
%     &\norm{\xka{i}'-\xka{0}'}_{\ci{0}} = \sup_{s\in I} \abs{\bm e_r(\thka{i}(s)) - \bm e_r(\thka{0}(s))} \le \sup_{s\in I} \abs{\thka{i}(s)-\thka{0}(s)} \\
%     & \le C_{J,\delta}\norm{\kappa_i-\kappa_0}_{\hr{-\frac{1}{2}+\delta}} \le \frac{1}{2}\abs{\kappa_0}_*.
% \end{aligned}
% \end{equation}
% for $\eta\in[0,1]$, we have
% \begin{equation}
% \begin{aligned}
% &\abs{\abs{\eta\xka{1}+(1-\eta)\xka{2}}_* - \abs{\kappa_0}_* } \le \eta\norm{\xka{1}'-\xka{0}'}_{\ci{0}} + (1-\eta)\norm{\xka{2}'-\xka{0}'}_{\ci{0}} \\
% & \le C_{J,\delta} \eta\norm{\kappa_1-\kappa_0}_{\hr{-\frac{1}{2}+\delta}} + C_{J,\delta}(1-\eta)\norm{\kappa_2-\kappa_0}_{\hr{-\frac{1}{2}+\delta}} \le \frac{1}{2}\abs{\kappa_0}_*.
% \end{aligned}
% \end{equation}
% This leads to $\abs{\eta\xka{1}+(1-\eta)\xka{2}}_* \ge \frac{1}{2}\abs{\kappa_0}_*$.
% We set
% \begin{equation}
%     m = \frac{1}{2}\abs{\kappa_0}_*, \quad M = 1 + M.
% \end{equation}
Lemma~\ref{lem:LM-Lip} controls the operator difference by the difference of
the reconstructed tangents, while Lemma~\ref{lem:Xk-holder} bounds the latter
in terms of $\kappa_1-\kappa_2$. Thus
\begin{equation}
\begin{aligned}
    &\norm{S_{\kappa_1,ij}f-S_{\kappa_2,ij}f} _{\hi{s+\varepsilon}} \le \norm{L_{\bm X_{\kappa_1}}f-L_{\bm X_{\kappa_2}}f} _{\hi{s+\varepsilon}} \\
    &\quad+ \norm{M_{\bm X_{\kappa_1},ij}f -M_{\bm X_{\kappa_2},ij}f}_{\hi{s+\varepsilon}} \\
    &\quad\le C_{M,m} \norm{\bm X_{\kappa_1}'-\bm X_{\kappa_2}'}_{\ci{0,\alpha}} \norm{f}_{\hr{s}} \le C_{M,m}
    \norm{\kappa_1-\kappa_2}_{\hr{-\frac12+\delta}} \norm{f}_{\hr{s}}.
\end{aligned}
\end{equation}
\end{proof}

The next lemma constructs the finite-energy Stokes solution and gives its
common velocity trace, traction jump, uniform trace estimate, and energy
identity. We use these properties in the tension uniqueness and global
energy arguments.

\begin{lemma}[Stokes single-layer potential and energy identity]\label{lem:stokes-layer-open-arc}\label{lem:weak-stokes-open-arc}
Let $0<\delta\le\frac12$, $\kappa\in\mc U_\delta$, and
$\Gamma=\xk(I)$, with closure $\overline{\Gamma}=\xk(\overline I)$. For
$\bm g\in\wt H^{\frac12}(I)$, set
$\bm f=\p_s\bm g\in\wt H^{-\frac12}(I)$.
The Stokes single-layer potential with density $\bm f$ defines a solution
$(\bm u,p)$ on $\R^2\setminus\overline{\Gamma}$ with
$\bm u\in H^1_{\rm loc}(\R^2)$, $\grad\bm u,p\in L^2(\R^2)$, and
$\bm u(\bm x),p(\bm x)\to0$ as $|\bm x|\to\infty$.
For smooth supported data,
\begin{equation}\label{eqn:stokes-single-layer}
    \bm u(\bm x)=\int_I G(\bm x-\xk(\eta))\bm f(\eta)\,d\eta.
\end{equation}
The velocity has a common parametrized trace
$\bm v=(\bm u|_\Gamma)\circ\xk\in H^{\frac12}(I)$ on the two sides, and
$\jump{\Sigma\bm n}\circ\xk=\bm f$ on $I$ with the convention
in~\eqref{eqn:intro-stokes-interface}.
If $\norm{\kappa}_{\hr{-\frac12+\delta}}\le M$ and
$\abs{\kappa}_*\ge m>0$, then
\begin{equation}\label{eqn:stokes-layer-bound}
    \norm{\bm v}_{\hi{\frac12}}
    \le C_{M,m}\norm{\bm g}_{\hr{\frac12}}.
\end{equation}
Moreover,
\begin{equation}\label{eqn:weak-stokes-energy}
    2\int_{\R^2\setminus\overline{\Gamma}}\abs{\grad_S\bm u}^2\,d\bm x
    =\dual{\bm f}{\bm v}
    =-\dual{\p_s\bm v}{\bm g},
\end{equation}
where the pairings are the
$\wt H^{-\frac12}(I)$--$H^{\frac12}(I)$ and
$H^{-\frac12}(I)$--$\wt H^{\frac12}(I)$ dualities, respectively.
\end{lemma}
\begin{proof}
Fix $0<\alpha<\beta<\delta$. Lemmas~\ref{tau-n-est}
and~\ref{lem:Xk-holder}, together with $\abs{\kappa}_*>0$, show that
$\overline{\Gamma}$ is a regular embedded $C^{1,\beta}$ arc. Extend $\xk$ slightly past each endpoint along its tangent,
choosing the extensions sufficiently short that they are disjoint from each
other and from the original arc except at their attachment points.
The two new endpoints can then be joined by a smooth simple arc, disjoint
from the extended arc except at its endpoints and with matching oriented
endpoint tangents. Such a connection can be routed through the complement
to a sufficiently large enclosing circle and completed outside that circle.
This gives a regular $C^{1,\alpha}$ Jordan curve $\wt\Gamma$.
Choose a regular periodic parametrization $\wt{\bm X}$ on $\mathbb T$
that agrees with $\xk$ on $\overline I$. Periodizing the zero extension of
$\bm g$ gives $\wt{\bm g}\in H^{\frac12}(\mathbb T)$. Define the density
on $\wt\Gamma$ by
\begin{equation*}
    \dual{\wt{\bm f}}{\bm\phi}_{\wt\Gamma}
    =\dual{\p_s\wt{\bm g}}{\bm\phi\circ\wt{\bm X}}_{\mathbb T}.
\end{equation*}
Then $\wt{\bm f}\in H^{-\frac12}(\wt\Gamma)$ is supported on $\overline{\Gamma}$
and has zero resultant, since a periodic derivative annihilates constants.
For this fixed parametrization, its norm is bounded by
$C\norm{\bm g}_{\hr{\frac12}}$.

The closed-curve theory in
\cite[Propositions~7.1--7.3 and~9.2--9.3]{SayasSelgas2014} gives
$(\bm u,p)=(S_u\wt{\bm f},S_p\wt{\bm f})$ with the stated regularity,
traces, jump relation, and single-layer representation. Since the density
vanishes on the added portion, the weak Stokes equations hold across it;
interior regularity therefore gives a solution on $\R^2\setminus\overline{\Gamma}$.
The zero resultant cancels the leading far-field terms in the kernels
\cite[(9.1) and~(9.4)--(9.5)]{SayasSelgas2014}, giving
$\bm u(\bm x)=O(|\bm x|^{-1})$ and
$|\grad\bm u(\bm x)|+|p(\bm x)|=O(|\bm x|^{-2})$.

Testing the closed-curve weak formulation
\cite[(7.6)]{SayasSelgas2014} with $\bm u$ gives
$2\norm{\grad_S\bm u}_{L^2}^2
=\dual{\wt{\bm f}}{\bm u|_{\wt\Gamma}}_{\wt\Gamma}$.
The support of $\wt{\bm f}$ reduces this to the canonical interval
pairing $\dual{\bm f}{\bm v}$, and periodic distributional integration
by parts gives $-\dual{\p_s\bm v}{\bm g}$.
This proves~\eqref{eqn:weak-stokes-energy} without endpoint terms.

The trace estimate follows from the closed-curve mapping properties;
Appendix~\ref{app:boundary-integral-estimates} proves uniformity of its
constant by compactness of the admissible family of arcs.
Bounded differentiation and~\eqref{eqn:weak-stokes-energy} also give uniform
bounds for $\p_s\bm v$ in $H^{-\frac12}(I)$ and $\grad_S\bm u$ in $L^2$.
\end{proof}


\section{The tension equation}
\label{sec:tension}

We prove solvability of the tension equation~\eqref{eqn:ten-det}, derive uniform bounds for
the inverse of $T_\kappa$, and estimate the tension.

\subsection{Invertibility of the tension operator}
\label{subsec:tension-existence}

The estimates in Subsections~\ref{subsec:commutator-estimates}
and~\ref{sec:divided-diff} give compactness of the remainder in $T_\kappa$,
while the Stokes energy identity~\eqref{eqn:weak-stokes-energy} gives uniqueness. The Fredholm alternative
then yields solvability for each admissible curvature.

\begin{proposition}[Solvability of the tension equation]\label{prop:main-tension}
Fix $0<\delta\le\frac12$ and let $\kappa\in\mc U_\delta$. Then
\begin{equation}
    T_\kappa:
    \thi{\frac12}\longrightarrow\hi{-\frac12}
\end{equation}
has a bounded inverse. Consequently, for every
$F\in\hi{-\frac12}$, the equation
\begin{equation}
    T_\kappa\sigma=F
\end{equation}
has a unique solution $\sigma\in\thi{\frac12}$ satisfying
\begin{equation}\label{eqn:c-kappa}
    \norm{\sigma}_{\hr{\frac12}}
    \le C_\kappa\norm{F}_{\hi{-\frac12}},
    \quad
    C_\kappa
    \coloneqq\norm{T_\kappa^{-1}}_{
    \hi{-\frac12}\mapsto\thi{\frac12}}.
\end{equation}
\end{proposition}

The compact-perturbation argument uses the following gain in the remainder
operator.

\begin{lemma}\label{R-est}
Let $0\le \varepsilon < \delta \le \frac{1}{2}$.
Let $\kappa\in\mc U_\delta$ and $\sigma\in \thi{\frac{1}{2}}$. Then
\begin{equation}\label{eqn:Rsigma}
\norm{R_\kappa \sigma}_{\hi{-\frac{1}{2}+\varepsilon}} \le C\abs{\kappa}_*^{-2} \paren{1+\norm{\kappa}_{\hr{-\frac{1}{2}+\delta}}}^4\norm{\kappa}_{\hr{-\frac{1}{2}+\delta}}\norm{\sigma}_{\hr{\frac{1}{2}}}.
\end{equation}
\end{lemma}
\begin{proof}
Set $\alpha=\frac{1}{2}\paren{\varepsilon+\delta}$.
Lemmas~\ref{point_mult_sobo}, \ref{dsH_f_comm},
and~\ref{tau-n-est} give the commutator bound
\begin{equation}
\begin{aligned}
    &\norm{\tk\cdot[\p_s\mc H, \tk]\sigma}_{\hi{-\frac{1}{2}+\varepsilon}} \le C\norm{\tk}_{\hi{\frac{1}{2}+\delta}} \paren{[\tk]_{\ci{0,\alpha}}+\norm{\p_s\tk}_{\hi{-\frac{1}{2}+\delta}}} \norm{\sigma}_{\hr{\frac{1}{2}}}\\
    & \le C\paren{1+\norm{\kappa}_{\hr{-\frac{1}{2}+\delta}}}^2 \norm{\kappa}_{\hr{-\frac{1}{2}+\delta}} \norm{\sigma}_{\hr{\frac{1}{2}}}.
\end{aligned}
\end{equation}
For the Stokes remainder, Lemmas~\ref{lem:Sop-est}
and~\ref{point_mult_sobo} give
\begin{equation}
\begin{aligned}
    &\norm{\tk\cdot\p_s S_\kappa(\sigma\tk)}_{\hi{-\frac{1}{2}+\varepsilon}}
    \le C \norm{\tk}_{\hi{\frac{1}{2}+\delta}}^2 \abs{\kappa}_*^{-2}\paren{1+\norm{\kappa}_{\hr{-\frac{1}{2}+\delta}}}\norm{\kappa}_{\hr{-\frac{1}{2}+\delta}} \norm{\sigma}_{\hr{\frac{1}{2}}} \\
    % \le C \frac{\norm{\tk}_{\ci{0,\alpha}}[\tk]_{\ci{0,\alpha}}}{\abs{\xk}_*^2} \norm{\tk}_{\hi{\frac{1}{2}+\delta}}^2\norm{\sigma}_{\hr{\frac{1}{2}}} \\
    &\le C\abs{\kappa}_*^{-2} \paren{1+\norm{\kappa}_{\hr{-\frac{1}{2}+\delta}}}^3\norm{\kappa}_{\hr{-\frac{1}{2}+\delta}}\norm{\sigma}_{\hr{\frac{1}{2}}}.
\end{aligned}
\end{equation}
Finally,
\begin{equation}
    \abs{\kappa}_*
    \le \norm{\tk}_{\ci{0}}
    \le C\norm{\tk}_{\hi{\frac{1}{2}+\delta}}
    \le C\paren{1+\norm{\kappa}_{\hr{-\frac{1}{2}+\delta}}},
\end{equation}
so the two bounds combine to give \eqref{eqn:Rsigma}.
\end{proof}

\iffalse
Old uniqueness argument kept for comparison.
\begin{lemma}\label{ten-unique-old}
Let $\sigma_1,\sigma_2$ be two solutions to the tension determination problem~\eqref{eqn:ten-det} with respect to the right-hand sides $F_1,F_2$. It follows that $\sigma_1=\sigma_2$.
\end{lemma}
\begin{proof}
We consider the homogeneous equation of~\eqref{eqn:ten-det}, i.e. $F=0$.
Let $\bm u(\bm x) = \int_I G(\bm x- \xk) \p_\eta \paren{\sigma\tk}\,d\eta$.
It follows that $\tk\cdot\p_s\bm u = 0$.
It satisfies the Stokes equation~\eqref{eqn:intro-stokes-bulk} with interface conditions
\begin{equation}
    \jump{\bm u} = 0, \quad \jump{\Sigma\bm n} = \p_s\paren{\sigma\tk}.
\end{equation}
Multiplying by $\bm u$ and integrate by parts, we have
\begin{equation}
    \int_{\R^2} \abs{\grad_S\bm u}^2\,d\bm x = \int_I \p_s\paren{\sigma\tk}\cdot\bm u\,ds = \int_I \sigma \tk\cdot\p_s\bm u\,ds = 0.
\end{equation}
Hence, $\grad_S\bm u=0$ everywhere. Since $\Gamma$ is open and the domain $\R\setminus\Gamma$ is connected, the pressure $p$ is a constant everywhere and $\jump{p}=0$ on $\Gamma$. Using the interface condition again, we have
\begin{equation}
    0 = -\jump{p\bm n} = \p_s\paren{\sigma\tk} =  \p_s\sigma\tk -\sigma\kappa\nk.
\end{equation}
Taking inner product with $\tk$ on both sides leads to $\p_s \sigma=0$, so $\sigma$ is a constant on $\Gamma$. Due to the boundary condition $\sigma=0$, we conclude that $\sigma=0$.
\end{proof}
\fi

\begin{lemma}\label{ten-unique}
Let $0<\delta\le\frac12$, let
$\kappa\in\mc U_\delta$, and let
$\sigma\in\thi{\frac{1}{2}}$ solve the homogeneous tension equation
\begin{equation}\label{eqn:tension-homogeneous}
    T_\kappa\sigma=0.
\end{equation}
Then $\sigma=0$. Thus the tension equation~\eqref{eqn:ten-det} has at
most one solution in $\thi{\frac{1}{2}}$ for each right-hand side
$F\in H^{-\frac{1}{2}}(I)$.
\end{lemma}
\begin{proof}
The homogeneous constraint gives zero tangential strain and hence zero
Stokes dissipation; closed-curve single-layer coercivity then forces the
density to vanish.
Let $(\bm u,p)$ be the finite-energy Stokes solution associated with
$\bm f=\p_s(\sigma\tk)$ and represented by the single-layer potential in
Lemma~\ref{lem:stokes-layer-open-arc}, and set
$\bm v=(\bm u|_\Gamma)\circ\xk$. For smooth $\sigma$, the Stokeslet
decomposition~\eqref{eqn:R-kernel}--\eqref{eqn:stokeslet-splitting},
together with the definition~\eqref{eqn:Rop-def}, gives
\begin{equation}\label{eqn:tension-tangential-id}
    \tk\cdot\p_s\bm v
    =
    -T_\kappa\sigma.
\end{equation}
Indeed, the straight-filament contribution is
$-\frac{1}{4}\p_s\mc H\sigma$, while the geometric remainder and the
commutator terms are exactly those in \eqref{eqn:Rop-def}.
Both sides of
\eqref{eqn:tension-tangential-id} are continuous as maps
$\thi{\frac{1}{2}}\to\hi{-\frac{1}{2}}$: the left-hand side by
Lemma~\ref{lem:stokes-layer-open-arc}, and the right-hand side by
Lemma~\ref{R-est}. Density therefore extends~\eqref{eqn:tension-tangential-id} to
$\sigma\in\thi{\frac{1}{2}}$. Equation~\eqref{eqn:tension-homogeneous} then gives
\begin{equation}
    \tk\cdot\p_s\bm v=0.
\end{equation}
By Lemma~\ref{point_mult_sobo},
$\bm g=\sigma\tk\in\wt H^{\frac12}(I)$, and
\eqref{eqn:weak-stokes-energy} gives
\begin{equation*}
    2\norm{\grad_S\bm u}_{L^2}^2
    =-\dual{\tk\cdot\p_s\bm v}{\sigma}=0.
\end{equation*}
Let $V$ be the single-layer trace operator on the closed extension
from the proof of Lemma~\ref{lem:stokes-layer-open-arc}. For the extended
density, \eqref{eqn:weak-stokes-energy} reads
$\dual{\wt{\bm f}}{V\wt{\bm f}}_{\wt\Gamma}=0$.
The closed-curve coercivity modulo the normal density
\cite[Proposition~7.4]{SayasSelgas2014} implies
$\wt{\bm f}=c\wt{\bm n}$. Since $\wt{\bm f}$ vanishes on the added
portion of $\wt\Gamma$, $c=0$. Thus $\bm f=\p_s(\sigma\tk)=0$.

It follows that $\sigma\tk$ is a constant vector on $I$. The only constant
vector whose zero extension belongs to $H^{\frac12}(\R)$ is zero. Hence
$\sigma\tk=0$, and $\abs{\tk}=1$ gives $\sigma=0$.

If $\sigma_1$ and $\sigma_2$ have the same right-hand side, their difference
solves~\eqref{eqn:tension-homogeneous}, so $\sigma_1=\sigma_2$.
\end{proof}
 

\begin{proof}[Proof of
Proposition~\ref{prop:main-tension}]
Choose $\varepsilon\in(0,\delta)$.  The embedding
$\hi{-\frac{1}{2}+\varepsilon}\hookrightarrow\hi{-\frac{1}{2}}$ is
compact.  Lemma~\ref{R-est} therefore shows that
\begin{equation}
    R_\kappa:\thi{\frac{1}{2}}\to\hi{-\frac{1}{2}}
\end{equation}
is compact.  By Lemma~\ref{ds-H-biject},
\begin{equation}
    T_\kappa:
    \thi{\frac{1}{2}}\to\hi{-\frac{1}{2}}
\end{equation}
is Fredholm of index zero.  Its kernel is trivial by
Lemma~\ref{ten-unique}; hence the Fredholm alternative gives a bounded
inverse, proving existence and uniqueness for~\eqref{eqn:ten-det}.
Moreover,
\begin{equation}\label{eqn:sigma-H12}
    \norm{\sigma}_{\hr{\frac{1}{2}}}
    \le C_\kappa \norm{F}_{\hi{-\frac{1}{2}}},
\end{equation}
where $C_\kappa<\infty$ is defined in~\eqref{eqn:c-kappa}.
\end{proof}

Lemma~\ref{R-Lip} gives a difference estimate for $R_\kappa$.
We use it in Lemma~\ref{lem:inv-op-cont} to prove continuity of $T_\kappa^{-1}$,
then combine this continuity with compactness in Lemma~\ref{lem:c-kappa}
to obtain a uniform inverse bound.

\begin{lemma}\label{R-Lip}
Let $0\le \varepsilon<\delta\le \frac{1}{2}$.
Let $\kappa_1,\kappa_2 \in \thi{-\frac{1}{2}+\delta}$ satisfy
\begin{equation*}
    \norm{\kappa_i}_{\hr{-\frac{1}{2}+\delta}}\le M,
    \quad i=1,2,
    \quad
    \abs{\rho \xka{1} + \paren{1-\rho}\xka{2}}_*\ge m,
    \quad \rho\in[0,1].
\end{equation*}
Then, for every $\sigma\in\thi{\frac12}$,
\begin{equation}\label{eqn:R-curvature-difference}
    \norm{\paren{R_{\kappa_1} - R_{\kappa_2}} \sigma}_{\hi{-\frac{1}{2}+\varepsilon}} \le C_{M,m} \norm{\kappa_1-\kappa_2}_{\hr{-\frac{1}{2}+\delta}} \norm{\sigma}_{\hr{\frac{1}{2}}}.
\end{equation}
\end{lemma}

The proof is given in Appendix~\ref{app:nonlinear-estimates}.

\begin{lemma}[Continuous dependence on the curvature]
\label{lem:inv-op-cont}
The map
\begin{equation}
    \kappa\longmapsto
    T_\kappa^{-1}
\end{equation}
is continuous from $\mc U_\delta$ into
$\mc L(\hi{-\frac{1}{2}},\thi{\frac{1}{2}})$.
\end{lemma}
\begin{proof}
Lemma~\ref{lem:set-ok0} shows that $\mc U_\delta$ is open. Fix
$\kappa_0\in\mc U_\delta$ and set
\begin{equation}
    M
    =\norm{\kappa_0}_{\hr{-\frac{1}{2}+\delta}}
    +\frac{1}{2C_{J,\delta}}\abs{\kappa_0}_*,
    \quad
    m = \frac{1}{2}\abs{\kappa_0}_*.
\end{equation}
For $0<r\le \abs{\kappa_0}_*/(2C_{J,\delta})$, we define
\begin{equation}
    B_r(\kappa_0) = \{\kappa\in\thi{-\frac{1}{2}+\delta}:\norm{\kappa-\kappa_0}_{\hr{-\frac{1}{2}+\delta}}\le r\}\subset\mc U_\delta.
\end{equation}
For $\kappa\in B_r(\kappa_0)$, set
\begin{equation}
    \Delta R=T_\kappa-T_{\kappa_0}
    =R_\kappa-R_{\kappa_0}.
\end{equation}
By Proposition~\ref{prop:main-tension},
$T_{\kappa_0}:\thi{\frac{1}{2}}\to\hi{-\frac{1}{2}}$ is boundedly invertible.
Then
\begin{equation}
    T_\kappa^{-1}-T_{\kappa_0}^{-1} =\paren{\paren{I+T_{\kappa_0}^{-1}\Delta R}^{-1}-I} T_{\kappa_0}^{-1}
    =-\paren{I+T_{\kappa_0}^{-1}\Delta R}^{-1} T_{\kappa_0}^{-1}\Delta R T_{\kappa_0}^{-1}.
\end{equation}
Lemmas~\ref{lem:set-ok0} and~\ref{R-Lip} give
\begin{equation}
    \norm{\Delta R}_{\thi{\frac{1}{2}}\mapsto\hi{-\frac{1}{2}}}
    \le K_0\norm{\kappa-\kappa_0}_{\hr{-\frac{1}{2}+\delta}}
    \le K_0 r,
\end{equation}
where $K_0=C_{M,m}$. Choose $r$ small enough that
\begin{equation}
    \norm{T_{\kappa_0}^{-1}\Delta R}_{\thi{\frac{1}{2}}\mapsto\thi{\frac{1}{2}}} \le
    \norm{T_{\kappa_0}^{-1}}_{\hi{-\frac{1}{2}}\mapsto\thi{\frac{1}{2}}} \norm{\Delta R}_{\thi{\frac{1}{2}}\mapsto\hi{-\frac{1}{2}}} \le C_{\kappa_0}K_0 r < 1.
\end{equation}
The Neumann-series formula then gives
\begin{equation}
    \norm{T_\kappa^{-1}-T_{\kappa_0}^{-1}}_{\hi{-\frac{1}{2}}\to\thi{\frac{1}{2}}}
    \le \frac{C_{\kappa_0}^2K_0r}{1-C_{\kappa_0}K_0r}.
\end{equation}
\end{proof}

\begin{lemma}[Uniform bound for the inverse]\label{lem:c-kappa}
Let $0<\delta\le\frac{1}{2}$ and let $M,m>0$. Define
\begin{equation}
    \mc K_{M,m}
    \coloneqq
    \left\{
    \kappa\in\thi{-\frac{1}{2}+\delta}:
    \norm{\kappa}_{\hr{-\frac{1}{2}+\delta}}\le M,
    \abs{\kappa}_*\ge m
    \right\}.
\end{equation}
Then the inverse constants $C_\kappa$ defined in~\eqref{eqn:c-kappa} are
uniformly bounded on $\mc K_{M,m}$:
\begin{equation}\label{eqn:inverse-uniform-bound}
    \sup_{\kappa\in\mc K_{M,m}} C_\kappa \le C_{M,m}<\infty.
\end{equation}
\end{lemma}
\begin{proof}
Choose $\delta_0\in(0,\delta)$. Since the embedding
$\thi{-\frac{1}{2}+\delta}\hookrightarrow\thi{-\frac{1}{2}+\delta_0}$ is
compact, $\mc K_{M,m}$ is precompact in
$\thi{-\frac{1}{2}+\delta_0}$.
Let $\kappa_n\in\mc K_{M,m}$ and suppose
$\kappa_n\to\kappa$ in $\thi{-\frac{1}{2}+\delta_0}$. The bounded map
$J:\thi{-\frac{1}{2}+\delta_0}\to C^0(I)$ gives
$J\kappa_n\to J\kappa$ and
$\bm X_{\kappa_n}'\to\bm X_{\kappa}'$ in $C^0(I)$. Hence
\begin{equation}
    \abs{\abs{\kappa_n}_*-\abs{\kappa}_*}
    \le \norm{\bm X_{\kappa_n}'-\bm X_{\kappa}'}_{\ci{0}}\to0.
\end{equation}
Thus $\abs{\kappa}_*\ge m$, so the closure of $\mc K_{M,m}$ in
$\thi{-\frac{1}{2}+\delta_0}$ remains inside the admissible set
$\{\kappa\in\thi{-\frac{1}{2}+\delta_0}:\abs{\kappa}_*>0\}$.

The proof of Lemma~\ref{lem:inv-op-cont}, with $\delta$ replaced by
$\delta_0$, shows that
$\kappa\mapsto T_\kappa^{-1}$ is continuous
on this closure in the operator norm
$\mc L(\hi{-\frac{1}{2}},\thi{\frac{1}{2}})$. Thus
$\kappa\mapsto C_\kappa$ is continuous on this compact set and bounded on
$\mc K_{M,m}$.
\end{proof}

\subsection{Tension estimates}
\label{subsec:nonlinear-estimates}

We first estimate the forcing $F(\kappa)$ in~\eqref{eqn:F-kappa} and apply
the inverse of $T_\kappa$ to control $\sigma$.
Comparison with $A_1$ gives higher tension regularity
through~\eqref{eq:A-orderk-interval-est}, after which we derive the forcing
and tension difference bounds.

\begin{lemma}\label{F-est}
Let $0\le\varepsilon<\delta\le\frac12$.
Let $\kappa\in\thi{\frac{3}{2}}$ with $\abs{\kappa}_*>0$. Then
\begin{equation}\label{eqn:forcing-bound}
\begin{aligned}
    \norm{F(\kappa)}_{\hi{-\frac{1}{2}+\varepsilon}} \le C \abs{\kappa}_*^{-2}\paren{1+\norm{\kappa}_{\hr{-\frac{1}{2}+\delta}}}^4\norm{\kappa}_{\hr{-\frac{1}{2}+\delta}}\norm{\kappa}_{\hr{\frac{3}{2}}}.
\end{aligned}
\end{equation}

\end{lemma}
\begin{proof}
Set $\alpha=\frac12(\varepsilon+\delta)$.
The commutator and Stokes-remainder estimates used in
Lemma~\ref{R-est} bound the two contributions in~\eqref{eqn:F-kappa}. Using
$\abs{\kappa}_*\le
C\paren{1+\norm{\kappa}_{\hr{-\frac{1}{2}+\delta}}}$, we obtain
\begin{equation}
\begin{aligned}
    \norm{F(\kappa)}_{\hi{-\frac{1}{2}+\varepsilon}} &\le C\paren{[\tk]_{\ci{0,\alpha}}+\norm{\p_s\tk}_{\hi{-\frac{1}{2}+\delta}}}\norm{\nk}_{\hi{\frac{1}{2}+\delta}}\norm{\p_s\kappa}_{\hr{\frac{1}{2}}} \\
    & + C\abs{\kappa}_*^{-2}\norm{\tk}_{\ci{0,\alpha}}[\tk]_{\ci{0,\alpha}}\norm{\tk}_{\hi{\frac{1}{2}+\delta}}\norm{\nk}_{\hi{\frac{1}{2}+\delta}}\norm{\p_s\kappa}_{\hr{\frac{1}{2}}} \\
    & \le C\abs{\kappa}_*^{-2}\paren{1+\norm{\kappa}_{\hr{-\frac{1}{2}+\delta}}}^4\norm{\kappa}_{\hr{-\frac{1}{2}+\delta}}\norm{\kappa}_{\hr{\frac{3}{2}}}.
\end{aligned}
\end{equation}
\end{proof}

\begin{lemma}[Tension estimate]\label{sigma-est-F}
Let $0<\delta\le\frac12$.
Let $\kappa\in\thi{\frac{3}{2}}$ with $\abs{\kappa}_*>0$, and let
$F(\kappa)$ be given by~\eqref{eqn:F-kappa}. There exists a unique solution
$\sigma\in\thi{\frac{1}{2}}$ to~\eqref{eqn:ten-det} satisfying
\begin{equation}\label{eqn:tension-bound}
    \norm{\sigma}_{\hr{\frac{1}{2}}} \le C C_\kappa \abs{\kappa}_*^{-2}\paren{1+\norm{\kappa}_{\hr{-\frac{1}{2}+\delta}}}^4\norm{\kappa}_{\hr{-\frac{1}{2}+\delta}}\norm{\kappa}_{\hr{\frac{3}{2}}},
\end{equation}
where $C_\kappa < \infty$ is given in~\eqref{eqn:c-kappa}.
\end{lemma}
\begin{proof}
Lemma~\ref{F-est}, with $\varepsilon=0$, gives~\eqref{eqn:forcing-bound} in the form
\begin{equation}
    \norm{F(\kappa)}_{\hi{-\frac12}}
    \le
    C\abs{\kappa}_*^{-2}
    \paren{1+\norm{\kappa}_{\hr{-\frac12+\delta}}}^4
    \norm{\kappa}_{\hr{-\frac12+\delta}}
    \norm{\kappa}_{\hr{\frac32}} .
\end{equation}
Proposition~\ref{prop:main-tension} gives the unique solution
of~\eqref{eqn:ten-det}; its inverse bound~\eqref{eqn:sigma-H12} yields
\eqref{eqn:tension-bound}:
\begin{equation}
    \norm{\sigma}_{\hr{\frac12}}
    \le C_\kappa\norm{F(\kappa)}_{\hi{-\frac12}}.
\end{equation}
\end{proof}

\begin{lemma}[Higher regularity of tension]\label{sigma-est-F-higher}
Let $0 < \varepsilon<\delta\le \frac{1}{2}$.
Let $\kappa\in\thi{\frac{3}{2}}$, $\abs{\kappa}_* > 0$, and let $\sigma$ be the
unique solution to~\eqref{eqn:ten-det}. Then $\sigma\in \thi{\frac{1}{2}+\varepsilon}$.
If $\norm{\kappa}_{\hr{-\frac12+\delta}}\le M$ and $\abs{\kappa}_*\ge m>0$, then
\begin{equation}\label{eqn:tension-higher-bound}
\begin{aligned}
    \norm{\sigma}_{\hr{\frac{1}{2}+\varepsilon}}
    \le C_{M,m}(1+C_\kappa)\norm{\kappa}_{\hr{\frac{3}{2}}}.
\end{aligned}
\end{equation}
\end{lemma}
\begin{proof}
We compare $\p_s\mc H_{\R}$ with $A_1$, whose difference has order $-1$,
and apply Lemma~\ref{lem:A-orderk-interval} to gain regularity.
Recall that $A_1$ is the Fourier multiplier on $\R$ with symbol
\begin{equation}
    a_1(\xi) = (1+\xi^2)^{\frac{1}{2}}.
\end{equation}
By Lemma~\ref{lem:A-orderk-interval}(1) with $k=1$, applied to $\sigma$,
which is supported in $\overline I$,
\begin{equation}\label{eqn:sigma-A-priori}
    \norm{\sigma}_{\hr{\frac{1}{2}+\varepsilon}}
    \le C\left(
    \norm{p_I A_1\sigma}_{\hi{-\frac{1}{2}+\varepsilon}}
    +\norm{\sigma}_{\hr{\frac{1}{2}}}
    \right).
\end{equation}
From~\eqref{eqn:ten-det}, with the remainder defined in~\eqref{eqn:Rop-def},
\begin{equation}
    \p_s\mc H \sigma = -4R_\kappa\sigma + 4F(\kappa),
\end{equation}
hence, on $I$,
\begin{equation}\label{eqn:sigma-principal-decomposition}
    p_I A_1\sigma
    =p_I\paren{A_1-\p_s\mc H_{\R}}\sigma
    -4R_\kappa\sigma+4F(\kappa).
\end{equation}
The operator $A_1-\p_s\mc H_{\R}$ is a Fourier multiplier with symbol
\begin{equation}
    (1+\xi^2)^{\frac{1}{2}}-\abs{\xi} = \frac{1}{\sqrt{1+\xi^2}+\abs{\xi}} \sim \abs{\xi}^{-1},
    \quad \text{as } \abs{\xi}\to \infty.
\end{equation}
Thus the difference is an operator of order $-1$, and
\begin{equation}\label{eqn:sigma-lower-order-bound}
    \norm{p_I\paren{A_1-\p_s\mc H_{\R}}\sigma}_{\hi{-\frac{1}{2}+\varepsilon}}
    \le C\norm{\sigma}_{\hr{-\frac{3}{2}+\varepsilon}}
    \le C\norm{\sigma}_{\hr{\frac{1}{2}}} .
\end{equation}
Applying~\eqref{eqn:sigma-lower-order-bound} and Lemma~\ref{R-est} to
the first two terms in~\eqref{eqn:sigma-principal-decomposition}, we obtain
\begin{equation}
    \norm{p_I A_1\sigma}_{\hi{-\frac{1}{2}+\varepsilon}} \le{} C_{M,m}\norm{\sigma}_{\hr{\frac{1}{2}}} +
    C\norm{F(\kappa)}_{\hi{-\frac{1}{2}+\varepsilon}}.
\end{equation}
Lemmas~\ref{F-est} and~\ref{sigma-est-F}, together with
\eqref{eqn:sigma-A-priori}, then give
\begin{equation}
    \norm{\sigma}_{\hr{\frac{1}{2}+\varepsilon}}
    \le C_{M,m}(1+C_\kappa)\norm{\kappa}_{\hr{\frac{3}{2}}}.
\end{equation}
\end{proof}

% \begin{equation}
%     F(\kappa) = \frac{1}{4}[\p_s\mc H, \bm \tau_\kappa]\cdot (\p_s \kappa \nk) + \frac{1}{4\pi}\tk\cdot\p_s \int_I\p_\eta r(s,\eta;\kappa)\p_\eta\kappa\nk(\eta) \,d\eta.
% \end{equation}

The next estimate controls the change in forcing with the curvature.

\begin{lemma}\label{F-Lip-kappa}
Let $0\le \varepsilon<\delta\le \frac{1}{2}$.
Let $\kappa_1,\kappa_2 \in \thi{\frac{3}{2}}$ satisfy
\begin{equation*}
    \norm{\kappa_i}_{\hr{-\frac{1}{2}+\delta}}\le M,
    \quad i=1,2,
    \quad
    \abs{\rho \xka{1} + \paren{1-\rho}\xka{2}}_*\ge m,
    \quad \rho\in[0,1].
\end{equation*}
Then
\begin{equation}\label{eqn:forcing-difference}
\begin{aligned}
    \norm{F(\kappa_1)-F(\kappa_2)}_{\hi{-\frac{1}{2}+\varepsilon}} &\le C_{M,m}
    \norm{\kappa_1-\kappa_2}_{\hr{-\frac{1}{2}+\delta}} \times \paren{ \norm{\kappa_1}_{\hr{\frac{3}{2}}}
    +\norm{\kappa_2}_{\hr{\frac{3}{2}}} } \\
    &\quad+ C_{M,m}M \norm{\kappa_1-\kappa_2}_{\hr{\frac{3}{2}}}.
\end{aligned}
\end{equation}
\end{lemma}

The proof is provided in Appendix~\ref{app:nonlinear-estimates}.

\begin{lemma}[Lipschitz continuity of tension with respect to curvature]\label{sigma-Lip-F}
Let $0\le \varepsilon<\delta\le \frac{1}{2}$.
Let $\kappa_1,\kappa_2 \in \thi{\frac{3}{2}}$ satisfy
\begin{equation*}
    \norm{\kappa_i}_{\hr{-\frac{1}{2}+\delta}}\le M,
    \quad i=1,2,
    \quad
    \abs{\rho \xka{1} + \paren{1-\rho}\xka{2}}_*\ge m,
    \quad \rho\in[0,1].
\end{equation*}
Set $C'=(1+C_{\kappa_1})(1+C_{\kappa_2})$.
Let $\sigma_i$ be the solution to~\eqref{eqn:ten-det} with right-hand side
$F(\kappa_i)$, $i=1,2$.
Then
\begin{equation}\label{eqn:tension-difference}
\begin{aligned}
    \norm{\sigma_1 - \sigma_2}_{\hr{\frac{1}{2}}} &\le C_{M,m}C'
    \norm{\kappa_1-\kappa_2}_{\hr{-\frac{1}{2}+\delta}} \times \paren{ \norm{\kappa_1}_{\hr{\frac{3}{2}}}
    +\norm{\kappa_2}_{\hr{\frac{3}{2}}} } \\
    &\quad+ C_{M,m}C'M \norm{\kappa_1-\kappa_2}_{\hr{\frac{3}{2}}}.
\end{aligned}
\end{equation}
\end{lemma}
\begin{proof}
Subtracting~\eqref{eqn:ten-det} for $\kappa_1$ and $\kappa_2$ gives
\begin{equation}\label{eqn:tension-difference-equation}
    T_{\kappa_1}\paren{\sigma_1-\sigma_2}
    =F(\kappa_1)-F(\kappa_2)
    -\paren{R_{\kappa_1}-R_{\kappa_2}}\sigma_2.
\end{equation}
Applying~\eqref{eqn:sigma-H12} to~\eqref{eqn:tension-difference-equation}
and bounding the right-hand side
by~\eqref{eqn:R-curvature-difference}, \eqref{eqn:forcing-difference},
and~\eqref{eqn:tension-bound}, we obtain
\begin{equation}
\begin{aligned}
    &\norm{\sigma_1-\sigma_2}_{\hr{\frac{1}{2}}} \le C_{M,m}C'
    \norm{\kappa_1-\kappa_2}_{\hr{-\frac{1}{2}+\delta}} \paren{ \norm{\kappa_1}_{\hr{\frac{3}{2}}}
    +\norm{\kappa_2}_{\hr{\frac{3}{2}}} } \\
    &\quad+ C_{M,m}C'M \norm{\kappa_1-\kappa_2}_{\hr{\frac{3}{2}}} + C_{M,m}C' M
    \norm{\kappa_1-\kappa_2}_{\hr{-\frac{1}{2}+\delta}} \norm{\kappa_2}_{\hr{\frac{3}{2}}} \\
    &\le C_{M,m}C' \norm{\kappa_1-\kappa_2}_{\hr{-\frac{1}{2}+\delta}} \paren{
    \norm{\kappa_1}_{\hr{\frac{3}{2}}} +\norm{\kappa_2}_{\hr{\frac{3}{2}}} } + C_{M,m}C'M
    \norm{\kappa_1-\kappa_2}_{\hr{\frac{3}{2}}}.
\end{aligned}
    % &+C\frac{\paren{1+M^6}}{m^3}\norm{\kappa_1-\kappa_2}_{\hr{-\frac{1}{2}+\delta}}C_{\kappa_2} \frac{\paren{1+M^4}}{m^2}\norm{\kappa_2}_{\hr{-\frac{1}{2}+\delta}}\norm{\kappa_2}_{\hr{\frac{3}{2}}} \\
    % & \le C \paren{C_{\kappa_1}+C_{\kappa_2}} \frac{\paren{1+M^{10}}}{m^5}\norm{\kappa_1-\kappa_2}_{\hr{\frac{3}{2}}}\paren{\norm{\kappa_1}_{\hr{\frac{3}{2}}} + \norm{\kappa_2}_{\hr{\frac{3}{2}}}}.
\end{equation}
\end{proof}


\section{Local well-posedness}
\label{sec:local-theory}

We prove Theorem~\ref{thm:main-local} using the tension estimates to bound
$N(\kappa)$ and a weighted Duhamel estimate to obtain a contraction in
$X_T$. The spatial neighborhood $\mc B(\kappa_0)$ retains a uniform
arc-chord lower bound throughout the construction. Uniqueness gives
consistency under time restart and under changes of the regularity exponent.
We use these consistency properties to define a common maximal lifespan
and to propagate the local Lipschitz estimate along the solution.

% \subsection{Estimates for the curvature nonlinearity}
% \label{subsec:curvature-nonlinearity}

\subsection{Local well-posedness at fixed regularity}
\label{subsec:local-proof}
The operator estimates from Section~\ref{sec:formulation} and tension
bounds from Section~\ref{sec:tension} give the two nonlinear estimates
needed below: a bound for $N(\kappa)$ to preserve the fixed-point ball,
and a difference bound for contraction and continuous dependence.

\begin{lemma}\label{lem:N-est}
Let $0 \le \varepsilon < \delta \le \frac{1}{2}$.
Let $M,m>0$ and let $\kappa\in\thi{\frac{3}{2}}$ satisfy
$\norm{\kappa}_{\hr{-\frac12+\delta}}\le M$ and $\abs{\kappa}_*\ge m$. Then
\begin{equation}\label{eqn:N-pointwise-bound}
    \norm{N(\kappa)}_{\hi{-\frac{1}{2}+\varepsilon}}
    \le C_{M,m}
    \norm{\kappa}_{\hr{-\frac{1}{2}+\delta}}
    \norm{\kappa}_{\hr{\frac{3}{2}}} .
\end{equation}
Let $\kappa_1,\kappa_2\in\thi{\frac{3}{2}}$ satisfy
$\norm{\kappa_i}_{\hr{-\frac{1}{2}+\delta}}\le M$ for $i=1,2$. Suppose
$\abs{\rho\xka{1}+\paren{1-\rho}\xka{2}}_*\ge m$ whenever
$\rho\in[0,1]$.
Then
\begin{equation}\label{eqn:N-pointwise-difference}
\begin{aligned}
    \norm{N(\kappa_1) - N(\kappa_2)}_{\hi{-\frac{1}{2}+\varepsilon}} &\le C_{M,m} \norm{\kappa_1-\kappa_2}_{\hr{-\frac{1}{2}+\delta}} \paren{\norm{\kappa_1}_{\hr{\frac{3}{2}}}+\norm{\kappa_2}_{\hr{\frac{3}{2}}}} \\
    & + C_{M,m}M \norm{\kappa_1-\kappa_2}_{\hr{\frac{3}{2}}}.
\end{aligned}
\end{equation}
\end{lemma}
\begin{proof}[Proof of the first estimate]
We control the frame and tension, then the force primitive $\bm g$
entering $N(\kappa)$.
Set $\alpha=\frac{1}{2}\paren{\varepsilon+\delta}$, so that
$\varepsilon<\alpha<\delta$, and write
$K=\norm{\kappa}_{\hr{-\frac{1}{2}+\delta}}\le M$.
Lemma~\ref{lem:c-kappa} bounds the inverse constant $C_\kappa$ by $C_{M,m}$.
Lemmas~\ref{tau-n-est} and~\ref{lem:Xk-holder} give
\begin{equation}
    [\nk]_{\ci{0,\alpha}}
    +\norm{\p_s\nk}_{\hi{-\frac{1}{2}+\delta}}
    \le C_{M,m}K,\quad
    \norm{\nk}_{\hi{\frac{1}{2}+\delta}}
    +\norm{\tk}_{\hi{\frac{1}{2}+\delta}}
    \le C_{M,m}.
\end{equation}
Let $\sigma=\sigma(\kappa)$ and set
$\bm g=\p_s\kappa\nk+\sigma\tk$. We use fixed bounded Sobolev
extensions of the frame fields to $\mathbb R$; $\bm g$ remains supported
in $\overline I$. Lemmas~\ref{lem:wt-der}, \ref{point_mult_sobo},
and~\ref{sigma-est-F} give $\bm g\in\thi{\frac12}$ and
\begin{equation}
    \norm{\bm g}_{\hr{\frac12}}
    \le C_{M,m}
    \paren{
        \norm{\kappa}_{\hr{\frac{3}{2}}}
        +\norm{\sigma}_{\hr{\frac{1}{2}}}
    }
    \le C_{M,m}\norm{\kappa}_{\hr{\frac32}}.
\end{equation}
Lemma~\ref{dsH_f_comm} applies with this choice of $\alpha$, including when
$\varepsilon=0$. For the Stokes remainder, apply Lemma~\ref{lem:Sop-est}
with $s=\frac12$, then differentiate and use Lemma~\ref{point_mult_sobo}
to multiply by $\nk$:
\begin{equation}
\begin{aligned}
    \norm{[\p_s\mc H,\nk]\cdot\bm g}_{\hi{-\frac12+\varepsilon}}
    &\le C_{M,m}K\norm{\bm g}_{\hr{\frac12}},\\
    \norm{\nk\cdot\p_s(S_\kappa\bm g)}_{\hi{-\frac12+\varepsilon}}
    &\le C_{M,m}K\norm{\bm g}_{\hr{\frac12}}.
\end{aligned}
\end{equation}
Substituting into~\eqref{eqn:N-def} gives
\begin{equation}
    \norm{N(\kappa)}_{\hi{-\frac12+\varepsilon}}
    \le C_{M,m}K\norm{\kappa}_{\hr{\frac32}}.
\end{equation}
The proof of~\eqref{eqn:N-pointwise-difference} is given in
Appendix~\ref{app:nonlinear-estimates}.
\end{proof}


Fix $0<\delta\le\frac12$ and recall $\mc U_\delta$
and $\mc B(\kappa_0)$ from
Subsection~\ref{subsec:notation}. The third-order semigroup gains
$2-\delta$ derivatives from $\thi{-1/2+\delta}$ to $\thi{3/2}$ with
singularity $t^{-(2-\delta)/3}$, which motivates the time weight below.
Define the function space
\begin{equation}\label{eqn:XT-intro}
    X_T
    \coloneqq
    \left\{
    \kappa\in C\paren{[0,T];\thi{-\frac12+\delta}}
    \cap C\paren{(0,T];\thi{\frac32}}:
    \norm{\kappa}_{X_T}<\infty
    \right\},
\end{equation}
with a time-weighted norm
\begin{equation}\label{eqn:XT-norm-intro}
    \norm{\kappa}_{X_T}
    =
    \sup_{0\le t\le T}
    \norm{\kappa(t)}_{\hr{-\frac12+\delta}}
    +
    \sup_{0<t\le T}
    t^{\frac{2-\delta}{3}}
    \norm{\kappa(t)}_{\hr{\frac32}}.
\end{equation}
We write $X_T^{(\delta)}$ when specifying the exponent.

Fix $0<\delta\le\frac12$ and $\kappa_0\in\mc U_\delta$, where
$\mc U_\delta$ is defined in~\eqref{eqn:U-intro}.
A \emph{mild solution} on $[0,T]$ is a function $\kappa$ in the space
$X_T^{(\delta)}$ defined by
\eqref{eqn:XT-intro}--\eqref{eqn:XT-norm-intro}, such that
$\kappa(0)=\kappa_0$, $\kappa(t)\in\mc U_\delta$ for every $t\in[0,T]$,
and $\kappa=\Psi_{\kappa_0}\kappa$ on $[0,T]$, with the time integral
in~\eqref{eqn:mild-def} understood as a Bochner integral in $H^{-3/2}(I)$.
A function $\kappa$ is a global mild solution if its restriction to every
finite interval $[0,T]$ is a mild solution.

\begin{proposition}[Local well-posedness for a fixed exponent]
\label{prop:local-fixed-exponent}
Let $\kappa_0\in\mc U_\delta$. There exist $T,r_0,C>0$, depending only on
$\delta,\kappa_0$, such that for every $\vph_0\in\thi{-\frac12+\delta}$ with
$\norm{\vph_0-\kappa_0}_{\hr{-\frac12+\delta}}\le r_0$, there exists a unique mild
solution $\kappa(\cdot;\vph_0)\in X_T$. Moreover,
\begin{equation}\label{eqn:local-fixed-lipschitz}
    \norm{\kappa(\cdot;\vph_0)-\kappa(\cdot;\kappa_0)}_{X_T}
    \le C\norm{\vph_0-\kappa_0}_{\hr{-\frac12+\delta}}.
\end{equation}
\end{proposition}
\begin{proof}
Fix $\varepsilon=\delta/2$ and write $a=(2-\delta)/3$.

\emph{1. Choice of the fixed-point set and preservation of the geometry.}
We center the ball at the linear trajectory $e^{t\Lambda}\kappa_0$.
Strong continuity keeps this trajectory near $\kappa_0$ for short time.
The radius of the fixed-point ball then ensures that every trajectory remains
in $\mc B(\kappa_0)$, where Lemma~\ref{lem:set-ok0} controls the geometry. Define
\begin{equation}\label{eqn:OT-intro}
    \mc O_T
    \coloneqq
    \left\{
    \kappa\in X_T:
    \norm{\kappa-e^{\cdot\Lambda}\kappa_0}_{X_T}
    \le \frac{\abs{\kappa_0}_*}{4C_{J,\delta}}
    \right\}.
\end{equation}
For the neighborhood~\eqref{eqn:Bkappa0}, strong
continuity~\eqref{eqn:lambda-semigroup-cont} gives $T_1\in(0,1]$ such that
\begin{equation}\label{eqn:local-linear-closeness}
    \sup_{0\le t\le T_1}
    \norm{e^{t\Lambda}\kappa_0-\kappa_0}_{\hr{-\frac12+\delta}}
    \le \frac{\abs{\kappa_0}_*}{4C_{J,\delta}}.
\end{equation}
By~\eqref{eqn:OT-intro}, \eqref{eqn:XT-norm-intro}, and
\eqref{eqn:local-linear-closeness}, for $0<T\le T_1$ and $\kappa\in\mc O_T$,
\begin{equation}
    \begin{aligned}
        \norm{\kappa(t)-\kappa_0}_{\hr{-\frac12+\delta}}
        &\le \norm{\kappa(t)-e^{t\Lambda}\kappa_0}_{\hr{-\frac12+\delta}}
        +\norm{e^{t\Lambda}\kappa_0-\kappa_0}_{\hr{-\frac12+\delta}}\\
        &\le \frac{\abs{\kappa_0}_*}{4C_{J,\delta}}
        +\frac{\abs{\kappa_0}_*}{4C_{J,\delta}}
        =\frac{\abs{\kappa_0}_*}{2C_{J,\delta}},
        \quad 0\le t\le T,
    \end{aligned}
\end{equation}
so $\kappa(t)\in\mc B(\kappa_0)$. Set the low-norm and arc-chord bounds
\begin{equation}
    M=\norm{\kappa_0}_{\hr{-\frac12+\delta}}
      +\frac{\abs{\kappa_0}_*}{2C_{J,\delta}},
    \quad
    m=\frac12\abs{\kappa_0}_*.
\end{equation}
By Lemma~\ref{lem:set-ok0},
$\norm{\kappa(t)}_{\hr{-\frac12+\delta}}\le M$ and
$\abs{\kappa(t)}_*\ge m$.
Since $\kappa_i(t)\in\mc B(\kappa_0)$ for any
$\kappa_1,\kappa_2\in\mc O_T$, \eqref{eqn:neighborhood-geometry} gives the arc-chord bound
for $\rho\xka{1}(t)+(1-\rho)\xka{2}(t)$, $0\le\rho\le1$.
This also verifies the geometric hypothesis of the difference estimate
\eqref{eqn:N-pointwise-difference} in Lemma~\ref{lem:N-est}.

\emph{2. Nonlinear and Duhamel estimates.}
Equations~\eqref{eqn:lambda-semigroup-est} and~\eqref{eqn:OT-intro} then give
\begin{equation}\label{eqn:local-ball-norm}
    \norm{\kappa}_{X_T}
    \le \norm{e^{\cdot\Lambda}\kappa_0}_{X_T}
       +\frac{\abs{\kappa_0}_*}{4C_{J,\delta}}
    \le C(1+M).
\end{equation}
Applying~\eqref{eqn:N-pointwise-bound} at each positive time and using
\eqref{eqn:XT-norm-intro} and~\eqref{eqn:local-ball-norm} gives
\begin{equation*}
    t^a\norm{N(\kappa(t))}_{\hi{-\frac12+\varepsilon}}
    \le C_{M,m}M\,t^a\norm{\kappa(t)}_{\hr{\frac32}}
    \le C_{M,m}.
\end{equation*}
For the difference, the two terms in~\eqref{eqn:N-pointwise-difference}
carry the same weight:
\begin{equation*}
\begin{aligned}
    t^a\norm{N(\kappa_1(t))-N(\kappa_2(t))}_{\hi{-\frac12+\varepsilon}}
    &\le C_{M,m}\norm{\kappa_1(t)-\kappa_2(t)}_{\hr{-\frac12+\delta}}
        \sum_{i=1}^2 t^a\norm{\kappa_i(t)}_{\hr{\frac32}}\\
    &\quad+C_{M,m}M t^a
        \norm{\kappa_1(t)-\kappa_2(t)}_{\hr{\frac32}}\\
    &\le C_{M,m}\norm{\kappa_1-\kappa_2}_{X_T}.
\end{aligned}
\end{equation*}
Uniformly for $T\le T_1$, we therefore have
\begin{equation}\label{eqn:local-nonlinear-bounds}
\begin{aligned}
    \sup_{0<t\le T}t^a
    \norm{N(\kappa(t))}_{\hi{-\frac12+\varepsilon}}
    &\le C_{M,m},\\
    \sup_{0<t\le T}t^a
    \norm{N(\kappa_1(t))-N(\kappa_2(t))}_{\hi{-\frac12+\varepsilon}}
    &\le C_{M,m}\norm{\kappa_1-\kappa_2}_{X_T}.
\end{aligned}
\end{equation}
Equation~\eqref{eqn:N-pointwise-difference} also gives
$N(\kappa)\in C((0,T];\hi{-\frac12+\varepsilon})$, since
$\kappa\in C((0,T];\thi{\frac32})$.

\emph{Duhamel estimate.}
For the integral term in~\eqref{eqn:mild-def}, we prove
$\mc Df\in X_T$ and $\norm{\mc Df}_{X_T}\le CT^{\varepsilon/3}F_T$.
Here $f\in C((0,T];\hi{-\frac12+\varepsilon})$ and
\begin{equation*}
    F_T=\sup_{0<r\le T}r^a\norm{f(r)}_{\hi{-\frac12+\varepsilon}}<\infty.
\end{equation*}
Set
\begin{equation}\label{eqn:local-D-definition}
    (\mc Df)(t)
    =\int_0^t e^{(t-r)\Lambda}\p_s f(r)\,dr.
\end{equation}
Write $X=H^{-3/2}(I)$ and $D(\Lambda)=\thi{3/2}$.
The forcing satisfies
$\p_s f\in H^{-3/2+\varepsilon}(I)=(X,D(\Lambda))_{\varepsilon/3,2}$;
the targets $\thi{-1/2+\delta}$ and $\thi{3/2}$ correspond to indices
$(1+\delta)/3$ and $1$, with zero extension bounded at the lower index.
Proposition~\ref{prop:lambda-semigroup}, with $\alpha=\varepsilon/3$
and these target indices, yields from~\eqref{eqn:lambda-semigroup-est}
\begin{equation*}
\begin{aligned}
    \norm{(\mc Df)(t)}_{\hr{-\frac12+\delta}}
    &\le C\int_0^t
    (t-r)^{-\frac{1+\delta-\varepsilon}{3}}
    \norm{f(r)}_{\hi{-\frac12+\varepsilon}}\,dr,\\
    \norm{(\mc Df)(t)}_{\hr{\frac32}}
    &\le C\int_0^t
    (t-r)^{-1+\frac{\varepsilon}{3}}
    \norm{f(r)}_{\hi{-\frac12+\varepsilon}}\,dr.
\end{aligned}
\end{equation*}
Using $\norm{f(r)}_{\hi{-\frac12+\varepsilon}}\le F_T r^{-a}$ and
substituting $u=r/t$ gives
\begin{equation*}
\begin{aligned}
    \int_0^t (t-r)^{-\frac{1+\delta-\varepsilon}{3}}r^{-a}\,dr
    &=t^{\frac{\varepsilon}{3}}
    \int_0^1 (1-u)^{-\frac{1+\delta-\varepsilon}{3}}u^{-a}\,du,\\
    t^a\int_0^t (t-r)^{-1+\frac{\varepsilon}{3}}r^{-a}\,dr
    &=t^{\frac{\varepsilon}{3}}
    \int_0^1 (1-u)^{-1+\frac{\varepsilon}{3}}u^{-a}\,du.
\end{aligned}
\end{equation*}
Both integrals are finite: $a<1$ controls $u=0$, and
$(1+\delta-\varepsilon)/3<1$ and $\varepsilon>0$ control $u=1$.
Both components of~\eqref{eqn:XT-norm-intro} gain $t^{\varepsilon/3}$;
taking their suprema gives the small factor $T^{\varepsilon/3}$:
\begin{equation}\label{eqn:local-duhamel-est}
    \norm{\mc Df}_{X_T}
    \le C T^{\frac{\varepsilon}{3}}
    \sup_{0<t\le T}t^a\norm{f(t)}_{\hi{-\frac12+\varepsilon}}.
\end{equation}
The constant depends on $\delta,\varepsilon$, but not on $T\le T_1$.

Beyond~\eqref{eqn:local-duhamel-est}, membership in $X_T$ requires the
time continuity in~\eqref{eqn:XT-intro}. For continuity in $\thi{3/2}$
at $t_0\in(0,T]$, fix $0<\eta<t_0/4$ and, for $t\in(0,T]$ with
$\abs{t-t_0}<\eta/2$, split at $t_0-\eta$:
\begin{equation*}
    (\mc Df)(t)
    =\int_0^{t_0-\eta}e^{(t-r)\Lambda}\p_s f(r)\,dr
    +\int_{t_0-\eta}^t e^{(t-r)\Lambda}\p_s f(r)\,dr.
\end{equation*}
In the first integral, $t-r\ge\eta/2$, so analyticity and the integrable
bound $\norm{f(r)}_{\hi{-\frac12+\varepsilon}}\le F_T r^{-a}$ give
continuity in $\thi{\frac32}$ by dominated convergence.
For the second integral, $r\ge t_0-\eta>t_0/2$, so
\begin{equation*}
\begin{aligned}
    \norm{\int_{t_0-\eta}^t e^{(t-r)\Lambda}\p_s f(r)\,dr}_{\hr{\frac32}}
    &\le C F_T t_0^{-a}
    \int_{t_0-\eta}^t(t-r)^{-1+\frac{\varepsilon}{3}}\,dr\\
    &\le C F_T t_0^{-a}\eta^{\frac{\varepsilon}{3}}.
\end{aligned}
\end{equation*}
This bound also holds at $t_0$; letting $t\to t_0$ and then
$\eta\downarrow0$ proves continuity there.
Continuity at zero follows from the lower-norm estimate
\begin{equation*}
    \norm{(\mc Df)(t)}_{\hr{-\frac12+\delta}}
    \le C F_T t^{\frac{\varepsilon}{3}}\longrightarrow0
    \quad\text{as }t\downarrow0.
\end{equation*}
With $(\mc Df)(0)=0$, we have $\mc Df\in X_T$.

\emph{3. Contraction and dependence on the initial data.}
For $\kappa,\kappa_1,\kappa_2\in\mc O_T$, definitions
\eqref{eqn:mild-def} and~\eqref{eqn:local-D-definition} give
\begin{equation*}
    \begin{aligned}
        \Psi_{\vph_0}\kappa-e^{\cdot\Lambda}\kappa_0
        &=e^{\cdot\Lambda}(\vph_0-\kappa_0)+\mc D N(\kappa),\\
        \Psi_{\vph_0}\kappa_1-\Psi_{\vph_0}\kappa_2
        &=\mc D\paren{N(\kappa_1)-N(\kappa_2)}.
    \end{aligned}
\end{equation*}
Equation~\eqref{eqn:lambda-semigroup-est}, together with
\eqref{eqn:local-nonlinear-bounds} and~\eqref{eqn:local-duhamel-est}, gives
\begin{equation}\label{eqn:local-map-bounds}
\begin{aligned}
    \norm{\Psi_{\vph_0}\kappa-e^{\cdot\Lambda}\kappa_0}_{X_T}
    &\le C_{M,m}\paren{
        T^{\frac{\varepsilon}{3}}
        +\norm{\vph_0-\kappa_0}_{\hr{-\frac12+\delta}}
    },\\
    \norm{\Psi_{\vph_0}\kappa_1-\Psi_{\vph_0}\kappa_2}_{X_T}
    &\le C_{M,m}T^{\frac{\varepsilon}{3}}
        \norm{\kappa_1-\kappa_2}_{X_T}.
\end{aligned}
\end{equation}
Choose $T\le T_1$ and $r_0>0$ sufficiently small such that
\begin{equation}\label{eqn:local-T-r-choice}
    C_{M,m}\paren{T^{\frac{\varepsilon}{3}}+r_0}
    \le\frac{\abs{\kappa_0}_*}{4C_{J,\delta}},
    \quad
    C_{M,m}T^{\frac{\varepsilon}{3}}\le\frac12.
\end{equation}
By~\eqref{eqn:local-map-bounds}, the first inequality in
\eqref{eqn:local-T-r-choice} ensures that $\Psi_{\vph_0}$ maps
$\mc O_T$ from~\eqref{eqn:OT-intro} into itself; the second gives
a contraction constant at most $1/2$.
Both hold for every $\vph_0$ in the stated neighborhood, with the same
choice of $T$. The Banach fixed-point theorem
gives a mild solution with initial value $\vph_0$.

For two fixed points with initial data $\vph_0,\psi_0$ in this neighborhood,
subtracting~\eqref{eqn:mild-def} gives
\begin{equation*}
\begin{aligned}
    \kappa(\cdot;\vph_0)-\kappa(\cdot;\psi_0)
    &=e^{\cdot\Lambda}(\vph_0-\psi_0)\\
    &\quad+\mc D\paren{N(\kappa(\cdot;\vph_0))-N(\kappa(\cdot;\psi_0))}.
\end{aligned}
\end{equation*}
The same contraction estimate~\eqref{eqn:local-map-bounds}--\eqref{eqn:local-T-r-choice} yields
\begin{equation}
\begin{aligned}
    &\norm{\kappa(\cdot;\vph_0)-\kappa(\cdot;\psi_0)}_{X_T}\\
    &\quad\le C\norm{\vph_0-\psi_0}_{\hr{-\frac12+\delta}}
    +\frac12
    \norm{\kappa(\cdot;\vph_0)-\kappa(\cdot;\psi_0)}_{X_T}.
\end{aligned}
\end{equation}
Absorbing the last term and taking $\psi_0=\kappa_0$ proves
\eqref{eqn:local-fixed-lipschitz}.

\emph{4. Uniqueness among all mild solutions.}
The contraction gives uniqueness only within $\mc O_T$
in~\eqref{eqn:OT-intro}. We show that every mild solution with the same
initial datum enters such a ball on a sufficiently short interval.
Let $\kappa_1,\kappa_2\in X_T$ be mild solutions with the same initial
value $f\in\mc U_\delta$. By continuity, both lie in
$\mc B(f)$ near zero.
Equations~\eqref{eqn:N-pointwise-bound} and
\eqref{eqn:local-duhamel-est} give, for sufficiently small $h>0$,
\begin{equation}
    \norm{\kappa_i-e^{\cdot\Lambda}f}_{X_h}
    \le C h^{\frac{\varepsilon}{3}}\norm{\kappa_i}_{X_h}
    \longrightarrow0,
    \quad i=1,2,
\end{equation}
where $C$ is independent of $h$. The convergence follows from
$\norm{\kappa_i}_{X_h}\le\norm{\kappa_i}_{X_T}$. For small $h$, both
solutions lie in the ball~\eqref{eqn:OT-intro} with $\kappa_0=f$ and $T=h$, so they
agree near zero. Set
\begin{equation*}
    T_*=\sup\{t\in[0,T]:\kappa_1=\kappa_2\text{ on }[0,t]\}.
\end{equation*}
If $T_*<T$, then $T_*>0$ and the semigroup property applied to
\eqref{eqn:mild-def} gives the same formula after shifting time by $T_*$,
with initial value
$g=\kappa_1(T_*)=\kappa_2(T_*)$.
The definition of a mild solution gives $g\in\mc U_\delta$, so $g$
is admissible initial data for the local uniqueness argument.
For $0<h\le T-T_*$, the shifted
solutions $u_i(t)=\kappa_i(T_*+t)$ belong to $X_h$. Indeed, they are
continuous in $\thi{\frac32}$ up to the new initial time, and
\begin{equation*}
    \sup_{0<t\le h}t^a\norm{u_i(t)}_{\hr{\frac32}}
    \le h^a\sup_{T_*\le r\le T_*+h}\norm{\kappa_i(r)}_{\hr{\frac32}}
    <\infty.
\end{equation*}
The uniqueness argument at zero, applied with initial value $g$,
gives $u_1=u_2$ for small $h$. Their equality extends beyond
$T_*$, a contradiction.
\end{proof}

\subsection{Maximal lifespan and continuous dependence}
\label{subsec:lifespan}

The next proposition gives consistency under time restart and under
change of regularity exponent. These allow us to continue local
solutions and compare their maximal lifespans.

\begin{proposition}[Consistency of local solutions]
\label{lem:restart-uniqueness}
Fix $0<\delta\le\frac12$. Let $\kappa$ be a mild solution on $[0,T]$.
For $t_0\in[0,T)$, let $\vph$ be the local fixed-point solution from
Proposition~\ref{prop:local-fixed-exponent} with initial value
$\vph(0)=\kappa(t_0)$.
Then
\begin{equation}\label{eqn:restart-identity}
    \vph(t)=\kappa(t_0+t)
\end{equation}
throughout their common interval of existence.

In addition, let $0<\delta_1\le\delta_2\le\frac12$, and let
$f\in\thi{-\frac12+\delta_2}$ satisfy $\abs{f}_*>0$. If $\vph_i$ is the
local fixed-point solution with initial value $f$ constructed at exponent
$\delta_i$, then $\vph_1=\vph_2$ throughout their common interval of
existence.
\end{proposition}
\begin{proof}
Set $\kappa^{[t_0]}(t)=\kappa(t_0+t)$. The semigroup property and the Duhamel
formula~\eqref{eqn:mild-def} for $\kappa$ give
\begin{equation}\label{eqn:shifted-duhamel}
    \kappa^{[t_0]}(t)
    =e^{t\Lambda}\kappa(t_0)
    +\int_0^t e^{(t-r)\Lambda}
    \p_s\paren{N(\kappa^{[t_0]}(r))}\,dr.
\end{equation}
The shifted solution inherits the arc-chord bound of $\kappa$ and belongs to
$X_h$ on every common interval $[0,h]$; for $t_0>0$, this membership follows
from positive-time continuity in $\thi{\frac32}$. Uniqueness in
Proposition~\ref{prop:local-fixed-exponent} gives \eqref{eqn:restart-identity}.

For consistency across exponents, set $a_i=(2-\delta_i)/3$, $i=1,2$.
The Sobolev embedding
$\thi{-\frac12+\delta_2}\hookrightarrow\thi{-\frac12+\delta_1}$ and the identity
\begin{equation}
    t^{a_1}\norm{\vph_2(t)}_{\hr{\frac32}}
    =t^{a_1-a_2}
    \paren{t^{a_2}\norm{\vph_2(t)}_{\hr{\frac32}}}
\end{equation}
give $X_T^{(\delta_2)}\hookrightarrow X_T^{(\delta_1)}$ on every finite
interval, since $a_1\ge a_2$. Both solutions satisfy~\eqref{eqn:mild-def} with initial value $f$, so uniqueness at exponent $\delta_1$ in
Proposition~\ref{prop:local-fixed-exponent} gives equality.
\end{proof}

\begin{proof}[Proof of Theorem~\ref{thm:main-local}]
We first identify a common maximal lifespan, then propagate the local
Lipschitz estimate along the reference trajectory to the chosen time $T$.

\emph{Independence of the lifespan.}
For each admissible exponent $\delta$,
Propositions~\ref{prop:local-fixed-exponent} and~\ref{lem:restart-uniqueness}
give a maximal solution $\kappa_\delta$ with lifespan $\tau_\delta$.
If $0<\delta_1\le\delta_2\le\frac12$ and
$\kappa_0\in\thi{-\frac12+\delta_2}$, the two solutions agree on their
common interval. The $\delta_2$ solution lies in the
$\delta_1$ solution class, so $\tau_{\delta_2}\le\tau_{\delta_1}$.
If this inequality is strict, choose
$0<a<\tau_{\delta_2}<b<\tau_{\delta_1}$.
On $[0,a]$, equality with $\kappa_{\delta_2}$ gives the initial continuity
and weighted bound at exponent $\delta_2$. On $[a,b]$, positive-time
regularity $\kappa_{\delta_1}\in C([a,b];\thi{3/2})$ gives the same
properties, since the time weight is comparable to a constant on $[a,b]$.
Thus $\kappa_{\delta_1}\in X_b^{(\delta_2)}$. It satisfies
\eqref{eqn:mild-def} and the arc-chord condition in~\eqref{eqn:U-intro},
so it extends $\kappa_{\delta_2}$ beyond its maximal lifespan, a
contradiction. All admissible exponents therefore have the
same lifespan $\tau_{\mathrm{max}}(\kappa_0)$.
Choose a finite $T=T(\kappa_0)\in(0,\tau_{\mathrm{max}}(\kappa_0))$.
This choice is independent of $\delta$, and the solution satisfies
\eqref{eqn:main-local-class} on $[0,T]$ for every admissible exponent.

\emph{Lipschitz dependence on the initial data.}
Fix $\delta$ and put $a=(2-\delta)/3$. We propagate the local Lipschitz
estimate to time $T$ by covering the compact reference trajectory
$\mc K=\{\kappa(t):0\le t\le T\}\subset\mc U_\delta$ with finitely
many local well-posedness neighborhoods and restarting finitely many times.
At each point of $\mc K$, Proposition~\ref{prop:local-fixed-exponent}
gives a ball of initial data with a common existence time, and its proof
gives a Lipschitz estimate for any two initial data in that ball.
The balls with half these radii still cover $\mc K$. Choose a finite
subcover and let $r>0$ be the minimum of these half-radii, bounding
the admissible perturbation of the initial data. Choose a uniform restart time $0<h\le T$
so that $2h$ is no larger than any local existence time, and a common
Lipschitz constant $L\ge1$. Then
\begin{equation}\label{eqn:uniform-local-lipschitz}
    \norm{\kappa(\cdot;q)-\kappa(\cdot;g)}_{X_{2h}^{(\delta)}}
    \le L\norm{q-g}_{\hr{-\frac12+\delta}}
\end{equation}
whenever $g\in\mc K$, $q\in\thi{-\frac12+\delta}$, and
$\norm{q-g}_{\hr{-\frac12+\delta}}\le r$.
Indeed, $g$ lies in one of the smaller balls, so $q$ and $g$ both lie
in the corresponding ball of full radius. The quantities $h,r,L$ may
depend on $\delta,\kappa_0,T$.

Construct the nearby solution by restarting at intervals of length $h$.
Set $n=\lceil T/h\rceil$ and choose $r_0=rL^{-n}$, so that up to
$n$ factors of $L$ keep the perturbation within radius $r$. For
\begin{equation*}
    d_0=\norm{\vph_0-\kappa_0}_{\hr{-\frac12+\delta}}\le r_0,
\end{equation*}
let $u(t)=\kappa(t;\vph_0)$ be the local solution. At each restart,
\eqref{eqn:uniform-local-lipschitz} increases the bound on the difference
in $\hr{-\frac12+\delta}$ by at most a factor $L$. Thus, inductively,
\begin{equation*}
    \norm{u(jh)-\kappa(jh)}_{\hr{-\frac12+\delta}}
    \le L^j d_0\le r,
    \quad j=0,\ldots,n-1.
\end{equation*}
Each local solution exists for time $2h$, and uniqueness
identifies the solutions on their overlaps. They therefore extend $u$
to $[0,T]$, with
\begin{equation*}
    \sup_{0\le t\le T}\norm{u(t)-\kappa(t)}_{\hr{-\frac12+\delta}}
    \le L^n d_0.
\end{equation*}

A restart at $t_0$ controls the high norm with the local weight
$(t-t_0)^a$, whereas $X_T$ uses $t^a$. On $(0,h]$,
\eqref{eqn:uniform-local-lipschitz} gives
\begin{equation*}
    \sup_{0<t\le h}t^a\norm{u(t)-\kappa(t)}_{\hr{\frac32}}
    \le Ld_0.
\end{equation*}
On $[jh,(j+1)h]\cap[0,T]$, use the solution restarted one interval
earlier, at $(j-1)h$. Then $h\le t-(j-1)h\le2h$, so the ratio of
global to local weights is at most $(T/h)^a$:
\begin{equation*}
\begin{aligned}
    t^a\norm{u(t)-\kappa(t)}_{\hr{\frac32}}
    &\le \left(\frac{t}{t-(j-1)h}\right)^a L^j d_0\\
    &\le \left(\frac{T}{h}\right)^a L^j d_0,
    \quad j=1,\ldots,n-1.
\end{aligned}
\end{equation*}
The lower-norm estimate and these weighted estimates give
\begin{equation*}
    \norm{u-\kappa}_{X_T^{(\delta)}}
    \le C_{\delta,\kappa_0,T}
    \norm{\vph_0-\kappa_0}_{\hr{-\frac12+\delta}}.
\end{equation*}
Since $T$ depends only on $\kappa_0$, this proves
\eqref{eqn:main-local-lipschitz} with $r_0,C$ depending only on
$\delta,\kappa_0$. Local uniqueness gives uniqueness on the whole interval.
\end{proof}


\section{Positive-time regularity and reconstruction}
\label{sec:regularity-theory}

We prove Theorem~\ref{thm:main-regularity} and
Corollary~\ref{cor:main-reconstruction}. We estimate $N(\kappa)$ at higher order
and restart at positive times to improve Sobolev regularity. We then apply
Lemma~\ref{lem:A-orderk-interval} to the supported third-order elliptic equation
to obtain the $d^{3/2}$ endpoint terms. For reconstruction, we prove a small
additional Sobolev gain with a time-integrable bound near $t=0$, then
reconstruct $(\bm X,\bm u,p,\lambda)$ and verify the original Stokes--filament system.

\subsection{Sobolev regularity and endpoint asymptotics}

\label{subsec:regularity-proof}

We first estimate the nonlinearity in a higher-order Sobolev space.
\begin{lemma}\label{lem:N-higher-est}
Let $0<\varepsilon<\delta< \frac{1}{2}$.
Let $L,m>0$ and let $\kappa\in\thi{\frac{3}{2}+\delta}$ satisfy
$\norm{\kappa}_{\hr{\frac32+\delta}}\le L$ and $\abs{\kappa}_*\ge m$.
Let $\sigma$ be the solution of~\eqref{eqn:ten-det}. Then
$N(\kappa)\in \hi{\frac{1}{2}+\varepsilon}$ and
\begin{equation}\label{eqn:N-higher-est}
    \norm{N(\kappa)}_{\hi{\frac{1}{2}+\varepsilon}}
    \le C_{L,m}
    \norm{\kappa}_{\hr{\frac{3}{2}+\delta}} .
\end{equation}
\end{lemma}
\begin{proof}
The lower Sobolev norm is also bounded by $L$, so
\eqref{eqn:inverse-uniform-bound} gives $C_\kappa\le C_{L,m}$.
Constants independent of the arc-chord bound are denoted by $C_L$.
Fix $\alpha=(\varepsilon+\delta)/2$ and set
$\bm q=\p_s\kappa\,\nk+\sigma\tk$. Before estimating the two terms in
$N(\kappa)$, we control the frame, the tension, and then $\bm q$.
By~\eqref{eqn:theta-Hs+1},
$\norm{\theta_\kappa}_{\hi{\frac{5}{2}+\delta}}
\le C\norm{\kappa}_{\hr{\frac{3}{2}+\delta}}$.
Set $r=\frac12+\delta\in(\frac12,1)$ and write $\theta=\theta_\kappa$.
The chain rule bounds $\sin\theta$ and $\cos\theta$ in $H^1(I)$,
and hence in $H^r(I)$, by $C_L$. Since $H^r(I)$ is an algebra by
Lemma~\ref{point_mult_sobo} and $\theta',\theta''\in H^r(I)$, the identities
\begin{align*}
    \p_s^2(\sin\theta)
    &=\cos\theta\,\theta''-\sin\theta\,(\theta')^2,\\
    \p_s^2(\cos\theta)
    &=-\sin\theta\,\theta''-\cos\theta\,(\theta')^2
\end{align*}
bound both second derivatives in $H^r(I)$ by $C_L$. Thus both
compositions are bounded in $H^{r+2}(I)$, giving
\begin{equation}\label{eqn:tau-n-high}
    \norm{\tk}_{\hi{\frac{5}{2}+\delta}}
    +\norm{\nk}_{\hi{\frac{5}{2}+\delta}}
    \le C_L .
\end{equation}
Lemma~\ref{lem:Xk-holder}, with $m=2$, also gives
\begin{equation}\label{eqn:X-high-C}
    \norm{\xk}_{\ci{3,\alpha}}
    \le C_L .
\end{equation}
The higher tension estimate
\eqref{eqn:tension-higher-bound} gives
\begin{equation}\label{eqn:sigma-high-for-N}
    \norm{\sigma}_{\hr{\frac{1}{2}+\varepsilon}}
    \le C_{L,m}
        \norm{\kappa}_{\hr{\frac{3}{2}+\delta}} .
\end{equation}

For $\bm q$, Lemma~\ref{point_mult_sobo} and
\eqref{eqn:tau-n-high}--\eqref{eqn:sigma-high-for-N} give
\begin{equation}\label{eqn:q-high}
\begin{aligned}
    &\norm{\bm q}_{\hr{\frac{1}{2}+\varepsilon}} \le \norm{\p_s\kappa\,\nk}_{\hr{\frac{1}{2}+\varepsilon}} +
    \norm{\sigma\tk}_{\hr{\frac{1}{2}+\varepsilon}} \\
    &\le C\norm{\p_s\kappa}_{\hr{\frac{1}{2}+\delta}} \norm{\nk}_{\hi{\frac{1}{2}+\varepsilon}} +
    C\norm{\sigma}_{\hr{\frac{1}{2}+\varepsilon}} \norm{\tk}_{\hi{\frac{1}{2}+\varepsilon}} \le C_{L,m}
    \norm{\kappa}_{\hr{\frac{3}{2}+\delta}} .
\end{aligned}
\end{equation}
In particular,
\begin{equation}\label{eqn:q-L2-high}
    \norm{\bm q}_{L^2(\R)}
    \le C_{L,m}
        \norm{\kappa}_{\hr{\frac{3}{2}+\delta}} .
\end{equation}

For the commutator term in~\eqref{eqn:N-def} with $\sigma=\sigma(\kappa)$,
Lemma~\ref{lem:comm-psH}, \eqref{eqn:tau-n-high}, and~\eqref{eqn:q-high} give
\begin{equation}\label{eqn:N-comm-high}
    \norm{[\p_s\mc H,\nk]\cdot\bm q}_{\hi{\frac{1}{2}+\varepsilon}} \le C\norm{\nk}_{\hi{\frac{5}{2}+\delta}}
    \norm{\bm q}_{\hr{\frac{1}{2}+\varepsilon}} \le C_{L,m} \norm{\kappa}_{\hr{\frac{3}{2}+\delta}} .
\end{equation}

For the Stokes-remainder term, write
$r_{ij}(s,\eta;\kappa)=\phi_{ij}(V(s,\eta))$, where
$V(s,\eta)=\frac{\xk(s)-\xk(\eta)}{s-\eta}$ and
$\phi_{ij}(v)=\frac{1}{2}\log(\abs{v}^2)\delta_{ij}-\frac{v_i v_j}{|v|^2}$.
The arc-chord bound and $|\p_s\xk|=1$ give $m\le|V|\le1$, including
on the diagonal by continuity. Thus the divided differences stay in a
compact subset of $\R^2\setminus\{0\}$, where $\phi_{ij}$ is smooth.
Together with $\xk\in\ci{3,\alpha}$ from~\eqref{eqn:X-high-C} and
$\varepsilon<\alpha$, this verifies the hypotheses of
Lemma~\ref{lem:K-higher-order} with $k=3$. Applied componentwise, it gives
\begin{equation}\label{eqn:R-kernel-high}
    \norm{\int_I \p_\eta r(s,\eta;\kappa)\bm q(\eta)\,d\eta}_{\hi{2+\varepsilon}}
    \le C_{L,m}
        \norm{\bm q}_{L^2(\R)} .
\end{equation}
After differentiating once in $s$,
\begin{equation}\label{eqn:dR-kernel-high}
\begin{aligned}
    &\norm{\p_s\int_I \p_\eta r(s,\eta;\kappa)\bm q(\eta)\,d\eta}_{\hi{\frac{1}{2}+\varepsilon}} \le C\norm{\int_I
    \p_\eta r(s,\eta;\kappa)\bm q(\eta)\,d\eta}_{\hi{\frac{3}{2}+\varepsilon}}  \\
    &\le C\norm{\int_I \p_\eta r(s,\eta;\kappa)\bm q(\eta)\,d\eta}_{\hi{2+\varepsilon}} \le C_{L,m}
    \norm{\bm q}_{L^2(\R)} \le C_{L,m} \norm{\kappa}_{\hr{\frac{3}{2}+\delta}} ,
\end{aligned}
\end{equation}
where the last two inequalities use \eqref{eqn:X-high-C} and
\eqref{eqn:q-L2-high}. By~\eqref{eqn:Sop-definition},
\begin{equation}
    \p_s S_\kappa\bm q
    =
    \p_s\int_I\p_\eta r(s,\eta;\kappa)\bm q(\eta)\,d\eta .
\end{equation}
Multiplying by $\nk$ and using Lemma~\ref{point_mult_sobo},
\eqref{eqn:tau-n-high}, and \eqref{eqn:dR-kernel-high} gives
\begin{equation}\label{eqn:N-rem-high}
    \norm{\nk\cdot\p_s S_\kappa\bm q}_{\hi{\frac{1}{2}+\varepsilon}} \le C\norm{\nk}_{\hi{\frac{1}{2}+\varepsilon}}
    \norm{\p_s S_\kappa\bm q}_{\hi{\frac{1}{2}+\varepsilon}} \le C_{L,m} \norm{\kappa}_{\hr{\frac{3}{2}+\delta}} .
\end{equation}
Combining \eqref{eqn:N-comm-high} and \eqref{eqn:N-rem-high} in the formula
\eqref{eqn:N-def} proves~\eqref{eqn:N-higher-est}.
\end{proof}

\begin{proof}[Proof of Theorem~\ref{thm:main-regularity}]
Fix $0<\tau<T$ and $0\le\varepsilon<\frac12$, and set
$t_0=\tau/4$ and $t_1=\tau/2$.

\emph{1. Positive-time spatial Sobolev regularity.}
Restarting away from $t=0$ provides $\thi{3/2}$ data for the parabolic
estimates and a uniform positive arc-chord bound on the remaining
compact time interval. Choose
\begin{equation*}
    \varepsilon<\eta<\eta_1<\frac12.
\end{equation*}
Here $\varepsilon$ is the final gain and $\eta$ an intermediate spatial
gain; $\eta_1>\eta$ makes the convolution exponent below greater than $-1$.
The continuity in~\eqref{eqn:XT-intro}, the condition $\kappa(t)\in\mc U_\delta$ from~\eqref{eqn:U-intro}, and~\eqref{eqn:arc-chord-C1-continuity} give finite bounds
\begin{equation*}
    B=\sup_{t_0\le r\le T}\norm{\kappa(r)}_{\hr{\frac32}},
    \quad m=\inf_{t_0\le r\le T}\abs{\kappa(r)}_*>0.
\end{equation*}
Constants in this step may depend on $B,m,T-t_0$ and the fixed
exponents, but are independent of $t,r$.
Choose an auxiliary exponent $\delta_1\in(\eta_1,\frac12]$.
The $\thi{3/2}$ bound controls the $\thi{-1/2+\delta_1}$ norm, so
Lemma~\ref{lem:N-est} applies uniformly with this auxiliary exponent:
\begin{equation}\label{eqn:pN-low}
    \p_sN(\kappa)\in
    L^\infty\paren{[t_0,T];\hi{-\frac32+\eta_1}}.
\end{equation}
For $t\in(t_0,T]$, \eqref{eqn:shifted-duhamel} reads
\begin{equation}\label{eqn:mild-reg}
    \kappa(t)=e^{(t-t_0)\Lambda}\kappa(t_0)
    +\int_{t_0}^t e^{(t-r)\Lambda}\p_sN(\kappa(r))\,dr.
\end{equation}
For forcing in $H^{-3/2+\eta_1}(I)$ and target $H^{-3/2+\eta}(I)$,
Proposition~\ref{prop:lambda-semigroup} with $n=1$ gives the exponent
$-1+(\eta_1-\eta)/3$ in~\eqref{eqn:lambda-semigroup-est}, while the
initial datum in $D(\Lambda)$ gives $-\eta/3$:
\begin{equation}\label{eqn:Lam-32}
\begin{aligned}
    \norm{\Lambda\kappa(t)}_{\hi{-\frac32+\eta}}
    &\le C(t-t_0)^{-\eta/3}\norm{\kappa(t_0)}_{\hr{\frac32}}\\
    &\quad+C\int_{t_0}^t
        (t-r)^{-1+(\eta_1-\eta)/3}
        \norm{\p_sN(\kappa(r))}_{\hi{-\frac32+\eta_1}}\,dr\\
    &\le C(t-t_0)^{-\eta/3}\norm{\kappa(t_0)}_{\hr{\frac32}}\\
    &\quad+C\int_{t_0}^t
        (t-r)^{-1+(\eta_1-\eta)/3}
        \norm{\kappa(r)}_{\hr{\frac32}}\,dr\\
    &\le C\paren{1+(t-t_0)^{-\eta/3}}
        \sup_{t_0\le r\le t}\norm{\kappa(r)}_{\hr{\frac32}}.
\end{aligned}
\end{equation}
The second inequality uses~\eqref{eqn:N-pointwise-bound}.

Recall that $A_3$, defined by~\eqref{eqn:orderk-symbol-factorization} with $k=3$, has symbol
$\langle\xi\rangle^3$. Write
$A_3\kappa=-4\Lambda\kappa+(A_3+4\Lambda)\kappa$.
The difference $A_3+4\Lambda$ is first order, hence lower order than
$A_3$. Its bound on $\thi{3/2}$ and the estimate for $\Lambda\kappa$
give $p_I A_3\kappa(t)\in H^{-3/2+\eta}(I)$.
Lemma~\ref{lem:A-orderk-interval}(1), with $k=3$, then yields the
supported Sobolev gain in~\eqref{eq:A-orderk-interval-est}:
\begin{equation}\label{eqn:k-est}
    \norm{\kappa(t)}_{\hr{\frac32+\eta}}
    \le C\paren{
        \norm{\Lambda\kappa(t)}_{\hi{-\frac32+\eta}}
        +\norm{\kappa(t)}_{\hr{\frac32}}}.
\end{equation}
Interpolating the uniform $\thi{3/2+\eta}$ bound on $[t_1,T]$ from
\eqref{eqn:Lam-32}--\eqref{eqn:k-est} with continuity in $\thi{3/2}$ gives
\begin{equation}\label{eqn:reconstruction-positive-time-sobolev}
    \kappa\in C\paren{[\tau,T];\thi{\frac32+\varepsilon}}.
\end{equation}
The case $\varepsilon=0$ is already part of the mild solution class.
Lemma~\ref{lem:theta-kappa} and Sobolev composition also give uniform
bounds for $\theta_\kappa$ in $H^{5/2+\eta}(I)$ and $\xk$ in
$H^{7/2+\eta}(I)$ on $[t_1,T]$. Hence the normalized curves are
uniformly bounded in $C^{3,\alpha}(\overline I)$ for $0<\alpha<\eta$.

\emph{2. Time differentiability.}
The graph norm on $D(\Lambda)=\thi{3/2}$ is equivalent to its
supported Sobolev norm, so $\kappa\in C([\tau,T];D(\Lambda))$.
We use the evolution equation to prove differentiability first in the
base space. For the semigroup argument, write
$E=H^{-3/2}(I)$, $S(t)=e^{t\Lambda}$, and
$g(t)=\p_sN(\kappa(t))$.
Choose $M>0$ with $\sup_{\tau\le t\le T}\norm{\kappa(t)}_{L^2(I)}\le M$.
Continuity in $\thi{3/2}$ gives lower-norm closeness and $C^1$ closeness
of the normalized curves. Their convex combinations retain a common
arc-chord bound by~\eqref{eqn:neighborhood-geometry}, so
\eqref{eqn:N-pointwise-difference} applies uniformly and gives
$g\in C([\tau,T];E)$.

For $\tau<t<T$ and $h>0$ with $t+h\le T$, \eqref{eqn:shifted-duhamel} gives
\begin{equation}\label{eqn:kappa-time-right-quotient}
    \frac{\kappa(t+h)-\kappa(t)}{h}
    =
    \frac{S(h)-\Id}{h}\kappa(t)
    +\frac1h\int_0^h S(h-r)g(t+r)\,dr.
\end{equation}
The first term converges to $\Lambda\kappa(t)$ in $E$ because
$\kappa(t)\in D(\Lambda)$. Strong continuity and local boundedness of
$S$ and continuity of $g$ imply that the second term tends to $g(t)$.
For the left difference quotient, write
\begin{equation}\label{eqn:kappa-time-left-quotient}
    \frac{\kappa(t)-\kappa(t-h)}{h}
    =
    \frac1h\int_0^h S(r)\Lambda\kappa(t-h)\,dr
    +\frac1h\int_0^h S(h-r)g(t-h+r)\,dr.
\end{equation}
Continuity of $\Lambda\kappa$ and $g$ in $E$ gives the same limit.
Thus $\p_t\kappa=\Lambda\kappa+g$ in $E$, with a continuous right-hand
side. At $\tau$ and $T$, the corresponding one-sided limits apply.

Since $\Lambda$ has order three on the zero extension of $\kappa$,
\begin{equation*}
    \norm{\Lambda\kappa(t)}_{\hi{-\frac32+\varepsilon}}
    \le C_\varepsilon
        \norm{\kappa(t)}_{\hr{\frac32+\varepsilon}}.
\end{equation*}
By \eqref{eqn:reconstruction-positive-time-sobolev}, $\Lambda\kappa$ is
also continuous in this space. Lemma~\ref{lem:N-est}, with
the auxiliary parameter $\delta_1=\frac12$, gives
\begin{equation*}
    \norm{g(t)}_{\hi{-\frac32+\varepsilon}}
    \le C_{\varepsilon,M,m}\norm{\kappa(t)}_{L^2(I)}
        \norm{\kappa(t)}_{\hr{\frac32}},
\end{equation*}
and~\eqref{eqn:N-pointwise-difference} gives continuity in the same space.
Consequently $\Lambda\kappa+g$ belongs to
$C([\tau,T];H^{-3/2+\varepsilon}(I))$, and
\begin{equation}\label{eqn:kappa-time-est}
    \norm{\p_t\kappa(t)}_{\hi{-\frac32+\varepsilon}}
    \le C_{\varepsilon,M,m}
        \norm{\kappa(t)}_{\hr{\frac32+\varepsilon}},
    \quad \tau\le t\le T.
\end{equation}
Integrating the identity in $E$ gives
\begin{equation}\label{eqn:kappa-time-integral}
    \kappa(t)-\kappa(\tau)
    =\int_\tau^t\paren{\Lambda\kappa(r)+g(r)}\,dr.
\end{equation}
Both sides lie in $H^{-3/2+\varepsilon}(I)$. The integral defines a
continuously differentiable primitive in this space, so
\begin{equation}\label{eqn:kappa-time-regularity}
    \kappa\in C^1\paren{[\tau,T];\hi{-\frac32+\varepsilon}},
    \quad
    \p_t\kappa=\Lambda\kappa+\p_sN(\kappa).
\end{equation}

\emph{3. Endpoint asymptotics.}
By Lemma~\ref{lem:A-orderk-interval}(2) with $k=3$, the $d^{3/2}$
endpoint expansion follows once
\begin{equation*}
    p_I A_3\kappa(t)\in H^{-1/2+\varepsilon}(I).
\end{equation*}
We therefore improve the bound for $\Lambda\kappa$ from
$H^{-3/2+\eta}(I)$ to $H^{-1/2+\varepsilon}(I)$.
Choose $\varepsilon<\rho<\eta$ to apply Lemma~\ref{lem:N-higher-est}
with the available curvature gain $\eta$ and nonlinear gain $\rho$;
the strict inequality $\rho>\varepsilon$ makes the Duhamel kernel integrable.
Equations~\eqref{eqn:Lam-32} and~\eqref{eqn:k-est} give
\begin{equation*}
    L=\sup_{t_1\le r\le T}
        \norm{\kappa(r)}_{\hr{\frac32+\eta}}<\infty.
\end{equation*}
Lemma~\ref{lem:N-higher-est}, with
$\delta=\eta$ and $\varepsilon=\rho$, gives
\begin{equation}
    N(\kappa(r))\in H^{1/2+\rho}(I),\quad t_1\le r\le T,
\end{equation}
with a uniform bound depending on $L,m$. Hence
\begin{equation}\label{eqn:pN-hi}
    \p_sN(\kappa)\in
    L^\infty\paren{[t_1,T];\hi{-\frac12+\rho}}.
\end{equation}
For $t\in[\tau,T]$, restart the mild formulation at $t_1$:
\begin{equation}
    \kappa(t)=e^{(t-t_1)\Lambda}\kappa(t_1)
        +\int_{t_1}^t e^{(t-r)\Lambda}\p_sN(\kappa(r))\,dr.
\end{equation}
For forcing in $H^{-1/2+\rho}(I)$ and target
$H^{-1/2+\varepsilon}(I)$, Proposition~\ref{prop:lambda-semigroup}
with $n=1$ gives the exponent $-1+(\rho-\varepsilon)/3$ in~\eqref{eqn:lambda-semigroup-est}:
\begin{equation}\label{eqn:Lam-12}
\begin{aligned}
    \norm{\Lambda\kappa(t)}_{\hi{-\frac12+\varepsilon}}
    &\le C(t-t_1)^{-(1+\varepsilon-\eta)/3}
        \norm{\kappa(t_1)}_{\hr{\frac32+\eta}}\\
    &\quad+C\int_{t_1}^t
        (t-r)^{-1+(\rho-\varepsilon)/3}
        \norm{\p_sN(\kappa(r))}_{\hi{-\frac12+\rho}}\,dr.
\end{aligned}
\end{equation}
For the initial term, commute $\Lambda$ with the semigroup and use
$\norm{\Lambda\kappa(t_1)}_{\hi{-3/2+\eta}}
\le C\norm{\kappa(t_1)}_{\hr{3/2+\eta}}$.
Since $t-t_1\ge\tau/2$ and the kernel is integrable,
\eqref{eqn:Lam-12} is uniform on $[\tau,T]$.

The order-one term satisfies
\begin{equation}
    \norm{(A_3+4\Lambda)\kappa(t)}_{\hi{-\frac12+\varepsilon}}
    \le C\norm{\kappa(t)}_{\hr{\frac32+\eta}}.
\end{equation}
Together with \eqref{eqn:Lam-12}, this gives
$p_I A_3\kappa(t)\in H^{-1/2+\varepsilon}(I)$.
Lemma~\ref{lem:A-orderk-interval}(2) gives~\eqref{eq:A-orderk-interval-split},
whose endpoint power is $k/2=3/2$. This is the expansion
\eqref{eqn:main-endpoint-split}, since replacing $e^{-d_\pm}d_\pm^{3/2}$
by $d_\pm^{3/2}$ changes only the remainder:
$(e^{-d_\pm}-1)d_\pm^{3/2}=O(d_\pm^{5/2})$, and
\begin{equation*}
    \chi_\pm(e^{-d_\pm}-1)d_\pm^{3/2}
    \in H^{5/2+\varepsilon}(\R),\quad 0\le\varepsilon<\frac12.
\end{equation*}
The estimate~\eqref{eq:A-orderk-interval-split-est} gives
\begin{equation}\label{eqn:split-est}
    |c_-(t)|+|c_+(t)|+\norm{\kappa_r(t)}_{\hr{\frac52+\varepsilon}}
    \le C\paren{
        \norm{\Lambda\kappa(t)}_{\hi{-\frac12+\varepsilon}}
        +\norm{\kappa(t)}_{\hr{\frac32+\eta}}}.
\end{equation}
This bound is uniform on $[\tau,T]$. Since $\tau>0$ is arbitrary, the
expansion holds at every $t\in(0,T]$.
\end{proof}

For reconstruction of the mean angle at $t=0$, we need a small positive
Sobolev gain whose singularity is integrable in time. Put
\begin{equation*}
    a=\frac{2-\delta}{3},\quad
    \beta=\frac{\delta}{4},\quad
    \gamma=\frac{\delta}{2},\quad T_0=\min\{T,1\}.
\end{equation*}
If $M\ge\norm{\kappa}_{X_T}$ and
$0<m\le\inf_{0\le t\le T}\abs{\kappa(t)}_*$, then
\begin{equation}\label{eqn:reconstruction-small-time-smoothing}
    \norm{\kappa(t)}_{\hr{\frac32+\beta}}
    \le C_{M,m}t^{-a-\beta/3},
    \quad 0<t\le T_0.
\end{equation}
Repeat the argument for~\eqref{eqn:Lam-32} with initial time $t/2$.
The $X_T$ norm gives the bound $\norm{\kappa(r)}_{\hr{3/2}}\le Mr^{-a}$.
For forcing in $H^{-3/2+\gamma}(I)$ and target $H^{-3/2+\beta}(I)$,
\eqref{eqn:lambda-semigroup-est} gives the exponent
$-1+(\gamma-\beta)/3$; with~\eqref{eqn:N-pointwise-bound}, this yields
\begin{equation*}
\begin{aligned}
    \norm{\Lambda\kappa(t)}_{\hi{-\frac32+\beta}}
    &\le Ct^{-\beta/3}\norm{\kappa(t/2)}_{\hr{\frac32}}\\
    &\quad+C\int_{t/2}^t
        (t-r)^{-1+(\gamma-\beta)/3}
        \norm{\p_sN(\kappa(r))}_{\hi{-\frac32+\gamma}}\,dr\\
    &\le C_{M,m}\paren{
        t^{-a-\beta/3}+t^{-a+(\gamma-\beta)/3}}\\
    &\le C_{M,m}t^{-a-\beta/3}.
\end{aligned}
\end{equation*}
The kernel is integrable because $\gamma>\beta$, and $r^{-a}\le
2^at^{-a}$ on the integration interval. Applying
\eqref{eqn:k-est} with exponent $\beta$ gives
\eqref{eqn:reconstruction-small-time-smoothing}, with constants also
depending on the fixed $\delta$. Since $\beta=\delta/4$ gives
$a+\beta/3<1$, the bound is integrable at $0$.

\subsection{Reconstruction of the Stokes--filament solution}
\label{subsec:reconstruction-proof}

Given $\kappa$, we solve~\eqref{eqn:ten-det} for $\sigma$ and use the force
primitive $\bm q_\kappa$ to compute the normalized boundary velocity
$\bm v_\kappa$. We use the normal component of its spatial derivative in
the tangent-angle equation. With the mean angle and centroid fixing rotation
and translation, we reconstruct $\bm X$ through
\eqref{eqn:theta-bar}--\eqref{eqn:curve-reconstruction} and set
$\lambda=\sigma-\kappa^2$. The single-layer representation in
Lemma~\ref{lem:stokes-layer-open-arc} then defines $(\bm u,p)$ for the
reconstructed filament.

\begin{proof}[Proof of Corollary~\ref{cor:main-reconstruction}]
Fix $M\ge\norm{\kappa}_{X_T}$ and
$0<m\le\inf_{0\le t\le T}\abs{\kappa(t)}_*$, and let
$\sigma=\sigma(\kappa)$ be given by
Proposition~\ref{prop:main-tension}.

\emph{Step 1. Reconstruct the angle, curve, and tension.}
For $t>0$, set
\begin{equation}\label{eqn:q-kappa-physical}
    \bm q_\kappa
    =\p_s\kappa\,\nk+\sigma\tk
    \in\thi{\frac12}.
\end{equation}
The derivative $\p_s\bm q_\kappa$ is the Stokes single-layer force
density. By~\eqref{eqn:stokeslet-splitting}, its boundary velocity in
normalized coordinates, given by~\eqref{eqn:filament-boundary-evolution}
on $\xk$, is
\begin{equation}\label{eqn:reconstruction-boundary-integral}
    \bm v_\kappa(s,t)
    =-\frac14\mc H\bm q_\kappa(s,t)
    +\frac{1}{4\pi}\int_I
        \p_\eta r(s,\eta;\kappa(t))\bm q_\kappa(\eta,t)\,d\eta.
\end{equation}
For smooth supported densities, integration by parts in
$\int_I G(\Delta\xk)\p_\eta\bm q_\kappa\,d\eta$ gives this formula.
Equations~\eqref{eqn:hilbert-bound} and~\eqref{eqn:Sop-bound} extend~\eqref{eqn:reconstruction-boundary-integral} to
$\bm q_\kappa\in\thi{1/2}$. Define also
\begin{equation}
    h_\kappa=-\nk\cdot\p_s\bm v_\kappa.
\end{equation}
The angular velocity $h_\kappa$ will determine the mean angle.
Put $a=(2-\delta)/3$, $\beta=\delta/4$, and $T_0=\min\{T,1\}$.
The $X_T$ bound supplies the basic $t^{-a}$ singularity; applying
\eqref{eqn:tension-bound}, \eqref{eqn:product-supported},
\eqref{eqn:product-restriction}, and~\eqref{eqn:Sop-bound} gives
\begin{equation}\label{eqn:reconstruction-time-bound}
    \norm{\bm q_\kappa(t)}_{\hr{\frac12}}
    +\norm{\bm v_\kappa(t)}_{\hi{\frac12}}
    +\norm{h_\kappa(t)}_{\hi{-\frac12}}
    \le C_{M,m}t^{-a},
    \quad 0<t\le T.
\end{equation}
The constants are uniform by~\eqref{eqn:inverse-uniform-bound} and~\eqref{eqn:tangent-Sobolev-bound}, which also applies to $\nk$.
For the mean angle, we use the additional gain in
\eqref{eqn:reconstruction-small-time-smoothing}.
Equation~\eqref{eqn:tension-higher-bound}, with exponent $\beta<\delta$,
gives $\norm{\sigma(t)}_{\hr{1/2+\beta}}\le C_{M,m}t^{-a}$.
Multiplication by the frame preserves $H^{1/2+\beta}$, and
the boundary integral in \eqref{eqn:reconstruction-boundary-integral}
is bounded at this index.
Consequently,
\begin{equation}\label{eqn:reconstruction-mean-bound}
\begin{aligned}
    &\norm{\bm q_\kappa(t)}_{\hr{\frac12+\beta}}
    +\norm{\bm v_\kappa(t)}_{\hi{\frac12+\beta}}
    +\norm{h_\kappa(t)}_{\hi{-\frac12+\beta}}\\
    &\hspace{3em}\le C_{M,m}t^{-a-\beta/3},
    \quad 0<t\le T_0.
\end{aligned}
\end{equation}
The zero extension of $1$ belongs to $\thi{1/2-\beta}$, so
$\dual{h_\kappa}{1}$ is a well-defined
$H^{-1/2+\beta}(I)$--$\thi{1/2-\beta}$ pairing.
Since $a+\beta/3<1$, its time integral converges at $0$
by~\eqref{eqn:reconstruction-mean-bound}.

Curvature determines the curve only up to rotation and translation;
the mean-angle and centroid ODEs restore these degrees of freedom.
Let $\bar\theta_0$ be the initial mean angle in
\eqref{eqn:reconstruction-initial-means}. Following
\eqref{eqn:theta-bar} and \eqref{eqn:reconstruction-endpoint-angle}, set
\begin{equation}\label{eqn:reconstructed-mean-angle}
    \bar\theta(t)
    =\bar\theta_0+\frac12\int_0^t\dual{h_\kappa(r)}{1}\,dr,
\end{equation}
and
\begin{equation}\label{eqn:reconstructed-angle}
\begin{aligned}
    \theta(s,t)
    &=\bar\theta(t)+J\kappa(s,t)
        -\frac12\int_I J\kappa(\eta,t)\,d\eta,\\
    \theta_*(t)
    &=\bar\theta(t)-\frac12\int_I J\kappa(\eta,t)\,d\eta.
\end{aligned}
\end{equation}
These formulas agree with~\eqref{eqn:theta-bar} and~\eqref{eqn:reconstruction-endpoint-angle}, since
$\int_I J\kappa=\int_I(1-\eta)\kappa(\eta)\,d\eta$.
For the initial data, this identity is understood by density.
The mild solution is continuous in $\thi{-1/2+\delta}$, so
Lemma~\ref{lem:theta-kappa} gives
$\theta\in C([0,T];C(\overline I))$ and $\theta_*\in C([0,T])$.
Set
\begin{equation}
    \bm\tau=\bm e_r(\theta)=Q_{\theta_*}\tk,\quad
    \bm n=Q_{\theta_*}\nk,\quad
    \bm v=Q_{\theta_*}\bm v_\kappa.
\end{equation}
The centroid prescription \eqref{eqn:X-bar} and the curve formula
\eqref{eqn:curve-reconstruction} give
\begin{equation}\label{eqn:reconstructed-centered-curve}
\begin{aligned}
    \bm Y(s,t)
    &=-\frac12\int_I(1-\eta)\bm\tau(\eta,t)\,d\eta
        +\int_{-1}^s\bm\tau(\eta,t)\,d\eta,\\
    \overline{\bm X}(t)
    &=\overline{\bm X}_0
        +\frac12\int_0^t\int_I\bm v(s,r)\,ds\,dr,\\
    \bm X(s,t)&=\overline{\bm X}(t)+\bm Y(s,t).
\end{aligned}
\end{equation}
The centroid integral converges by \eqref{eqn:reconstruction-time-bound}.
For $t>0$, set
\begin{equation}\label{eqn:reconstructed-tension}
    \lambda=\sigma-\kappa^2.
\end{equation}

\emph{Step 2. Reconstruct the Stokes fields.}
Let $(\bm u_\kappa,p_\kappa)$ be the single-layer fields of Lemma~\ref{lem:stokes-layer-open-arc}, with velocity~\eqref{eqn:stokes-single-layer}, on
$\Gamma_\kappa=\xk(I)$ with density
$\bm f_\kappa=\p_s\bm q_\kappa$.
Write
\begin{equation}
    \overline{\bm X}_\kappa(t)=\frac12\int_I\xk(s,t)\,ds,
    \quad
    \bm Y=Q_{\theta_*}\paren{\xk-\overline{\bm X}_\kappa}.
\end{equation}
The physical curve is obtained by rotating the centered normalized
curve and adding the prescribed centroid. Apply the same change of
coordinates to the fluid fields:
\begin{equation}\label{eqn:reconstructed-fluid}
\begin{aligned}
    \bm u(\bm x,t)
    &=Q_{\theta_*(t)}\bm u_\kappa\!\left(
        Q_{\theta_*(t)}^{-1}
        \paren{\bm x-\overline{\bm X}(t)}
        +\overline{\bm X}_\kappa(t),t\right),\\
    p(\bm x,t)
    &=p_\kappa\!\left(
        Q_{\theta_*(t)}^{-1}
        \paren{\bm x-\overline{\bm X}(t)}
        +\overline{\bm X}_\kappa(t),t\right).
\end{aligned}
\end{equation}
By translation invariance and rotational covariance of the Stokeslet, these
are the single-layer fields on $\Gamma(t)=\bm X(I,t)$ with density
$\bm f=\p_s\bm q$, where
\begin{equation}
    \bm q=\p_s\kappa\,\bm n+\sigma\bm\tau
        =Q_{\theta_*}\bm q_\kappa.
\end{equation}

\emph{Step 3. Establish positive-time regularity.}
Fix $0<\tau<T$ and $0<\alpha<\frac12$, and choose
$\alpha<\varepsilon<\varepsilon_1<\frac12$.
We first obtain regularity up to the endpoints for $\bm X$, $\lambda$, and the velocity
trace sufficient to interpret the endpoint and kinematic conditions
classically. Localization away from the endpoints gives $\bm X$ and
$\lambda$ one additional derivative for the interior equations.
By~\eqref{eqn:main-positive-time-regularity}, $\kappa$ is continuous in
$\thi{3/2+\varepsilon_1}$ on $[\tau,T]$.
Equation~\eqref{eqn:tension-difference} gives continuity of $\sigma$ in $\thi{1/2}$,
since the normalized arcs at nearby times are close in $C^1$ and all their
convex combinations satisfy~\eqref{eqn:neighborhood-geometry}. Equation~\eqref{eqn:tension-higher-bound}, with its parameter
$\delta=\frac12$ and the positive-time $L^2$ bound for $\kappa$, gives a
uniform $\thi{1/2+\varepsilon_1}$ bound. Interpolation gives
\begin{equation}\label{eqn:reconstruction-positive-time-tension}
    \sigma\in C\paren{[\tau,T];\thi{\frac12+\varepsilon}},
    \quad
    \bm q\in C\paren{[\tau,T];\thi{\frac12+\varepsilon}}.
\end{equation}
The same continuity holds for $\bm q_\kappa$.
Lemma~\ref{lem:theta-kappa} and Sobolev composition give
\begin{equation*}
    \theta,\bm\tau,\bm n
    \in C\paren{[\tau,T];H^{5/2+\varepsilon_1}(I)}.
\end{equation*}
Equations~\eqref{eqn:reconstructed-centered-curve} and~\eqref{eqn:reconstructed-tension}, together with Sobolev embedding, give
\begin{equation}
    \bm X\in C\paren{[\tau,T];C^{3,\alpha}(\overline I)},
    \quad
    \lambda\in C\paren{[\tau,T];C^{0,\alpha}(\overline I)}.
\end{equation}

Globally, the supported Sobolev spaces retain the endpoint singularity
in Theorem~\ref{thm:main-regularity}. On compact subintervals of $I$,
the singular terms are smooth, so~\eqref{eqn:main-endpoint-split} gives
\begin{equation}
    \kappa(\cdot,t)\in H^{\frac52+\varepsilon}_{\mathrm{loc}}(I).
\end{equation}
To improve the tension regularity in the interior, we localize its
order-one equation and invert the whole-line principal part.
Lemma~\ref{lem:comm-psH} applies to the commutators in
\eqref{eqn:Rop-def}--\eqref{eqn:F-kappa}, since
$\bm\tau,\bm n\in H^{5/2+\varepsilon_1}(I)$ and
$\p_s\kappa\,\bm n,\sigma\bm\tau\in\thi{1/2+\varepsilon}$.
Choose $\gamma\in(\varepsilon,\varepsilon_1)$. The curve belongs to
$C^{3,\gamma}(\overline I)$, so Lemma~\ref{lem:K-higher-order} with $k=3$
maps the $L^2$ densities $\p_s\kappa\,\bm n$ and $\sigma\bm\tau$ into
$H^{2+\varepsilon}(I)$. Differentiating once and multiplying by the
tangent controls the Stokes-remainder terms, as in
\eqref{eqn:R-kernel-high}--\eqref{eqn:dR-kernel-high}.
Consequently $F(\kappa),R_\kappa\sigma\in H^{1/2+\varepsilon}(I)$.
Localizing~\eqref{eqn:ten-det} with $\chi\in C_c^\infty(I)$ removes
the endpoints and gives
\begin{equation}\label{eqn:localized-tension-equation}
    \p_s\mc H_{\R}(\chi\sigma)
    =4\chi\bigl(F(\kappa)-R_\kappa\sigma\bigr)
    +[\p_s\mc H_{\R},\chi]\sigma
    \in H^{1/2+\varepsilon}(\R).
\end{equation}
Here $\p_s\mc H_{\R}$ denotes the whole-line multiplier $|\xi|$, and
$\chi(F(\kappa)-R_\kappa\sigma)$ is extended by zero. The cutoff
commutator has order zero and preserves $H^{1/2+\varepsilon}(\R)$:
the Fourier kernel
$(|\xi|-|\eta|)\widehat\chi(\xi-\eta)$ is bounded using
$\bigl||\xi|-|\eta|\bigr|\le|\xi-\eta|$ and weighted Young's inequality.
Let $A_1$ have symbol $\langle\xi\rangle$.
Since $A_1-\p_s\mc H_{\R}$ has order $-1$ and
$\chi\sigma\in H^{1/2+\varepsilon}(\R)$,
\eqref{eqn:localized-tension-equation} gives
$A_1(\chi\sigma)\in H^{1/2+\varepsilon}(\R)$ and, by inversion,
$\chi\sigma\in H^{3/2+\varepsilon}(\R)$, proving
\begin{equation}\label{eqn:interior-tension-regularity}
    \sigma(\cdot,t)\in H^{\frac32+\varepsilon}_{\mathrm{loc}}(I).
\end{equation}
Consequently,
\begin{equation}
    \bm X(\cdot,t)\in H^{\frac92+\varepsilon}_{\mathrm{loc}}(I),
    \quad
    \lambda(\cdot,t)\in H^{\frac32+\varepsilon}_{\mathrm{loc}}(I),
\end{equation}
and Sobolev embedding gives
\begin{equation}\label{eqn:main-interior-classical-regularity}
    \bm X(\cdot,t)\in C_{\mathrm{loc}}^{4,\alpha}(I),
    \quad \lambda(\cdot,t)\in C_{\mathrm{loc}}^{1,\alpha}(I).
\end{equation}
With a fixed interior cutoff, these Sobolev bounds are
uniform on $[\tau,T]$, by~\eqref{eqn:split-est} and the argument leading to~\eqref{eqn:interior-tension-regularity}. Interpolating with the positive-time continuity proved above
gives time continuity in the local H\"older spaces in
\eqref{eqn:main-interior-classical-regularity}.
Globally, $\bm q\in\thi{1/2+\varepsilon}$ and
$\bm f=\p_s \bm q\in\thi{-1/2+\varepsilon}$, while locally
$\bm q\in H^{3/2+\varepsilon}$ and $\bm f\in H^{1/2+\varepsilon}$.
Thus $\bm q\in C^{1,\alpha}_{\mathrm{loc}}(I)$ and
$\bm f\in C^{0,\alpha}_{\mathrm{loc}}(I)$.

For the fluid, Lemma~\ref{lem:stokes-layer-open-arc} supplies the
finite-energy solution and its common velocity trace.
Together with~\eqref{eqn:reconstructed-fluid}, it gives
\begin{equation}\label{eqn:reconstruction-fluid-energy}
    \bm u(\cdot,t)\in H^1_{\mathrm{loc}}(\R^2),
    \quad \grad\bm u(\cdot,t),\ p(\cdot,t)\in L^2(\R^2),
    \quad t>0.
\end{equation}
The single-layer representation in Lemma~\ref{lem:stokes-layer-open-arc} shows that $\bm u,p$ are smooth off $\overline{\Gamma(t)}$.

For the classical one-sided regularity, we use local layer-potential theory.
Along a $C^{2,\alpha}$ curve, a $C^{0,\alpha}$ single-layer density
gives velocity and pressure of class $C^{1,\alpha}$ and $C^{0,\alpha}$,
respectively, up to the curve from each side; see
\cite[Section~2.3]{Hsiao2021} and
\cite[Theorem~2.1 and Appendix~A]{KuoLaiMoriRodenberg2023Tension}.
For a compact interior subarc, cut off $\bm f$ on a slightly larger
subarc away from the free endpoints and extend the latter to a closed
$C^{2,\alpha}$ curve. The cutoff density extends by zero in $C^{0,\alpha}$,
so the cited regularity applies; the remaining density is supported
away from the smaller subarc and produces a smooth field there.
This gives one-sided $C^{1,\alpha}$ regularity for $\bm u$ and
$C^{0,\alpha}$ regularity for $p$ on every compact interior subarc.

By Lemma~\ref{lem:stokes-layer-open-arc}, the common trace of
$\bm u_\kappa$ is \eqref{eqn:reconstruction-boundary-integral}.
The whole-line Hilbert transform is bounded on $H^{1/2+\varepsilon}$.
Lemmas~\ref{lem:Sop-est} and~\ref{lem:Sop-Lip}, with input index
$1/2+\varepsilon<1$ and gain zero, give boundedness and continuity
of the remainder at the same index. The density and rotation are
continuous, so
\begin{equation}\label{eqn:reconstruction-positive-time-trace}
    \bm v=(\bm u|_\Gamma)\circ\bm X
    \in C\paren{[\tau,T];H^{\frac12+\varepsilon}(I)}
    \hookrightarrow C\paren{[\tau,T];C^{0,\alpha}(\overline I)}.
\end{equation}

In particular, $h_\kappa\in C([\tau,T];H^{-1/2+\varepsilon}(I))$.
Equations~\eqref{eqn:reconstructed-mean-angle} and~\eqref{eqn:reconstructed-centered-curve} give
$\bar\theta\in C^1([\tau,T])$ and
$\overline{\bm X}\in C^1([\tau,T])$.

To differentiate~\eqref{eqn:reconstructed-centered-curve} in time, we
write both the angle and the centered curve using the mean-zero
primitive $J_0$. It gains one spatial derivative in the negative
Sobolev scales needed here. Subtracting the spatial mean from $J$ gives,
for smooth $f$,
\begin{equation}\label{eqn:centered-primitive}
\begin{aligned}
    (J_0f)(s)
    &=(Jf)(s)-\frac12\int_I(Jf)(r)\,dr\\
    &=\int_{-1}^s f(\eta)\,d\eta
        -\frac12\int_I(1-\eta)f(\eta)\,d\eta.
\end{aligned}
\end{equation}
It extends to bounded maps
\begin{equation}\label{eqn:centered-primitive-bounds}
\begin{aligned}
    J_0:H^{-3/2+\varepsilon}(I)&\longrightarrow
        H^{-1/2+\varepsilon}(I),\\
    J_0:H^{-1/2+\varepsilon}(I)&\longrightarrow
        H^{1/2+\varepsilon}(I).
\end{aligned}
\end{equation}
To prove the first bound, consider the adjoint
\begin{equation*}
    (J_0^*g)(\eta)
    =\int_\eta^1 g(s)\,ds
        -\frac{1-\eta}{2}\int_I g(s)\,ds.
\end{equation*}
It vanishes at both endpoints and satisfies
$(J_0^*g)'=-g+\frac12\int_I g$. Thus
\begin{equation*}
    \norm{J_0^*g}_{\hr{3/2-\varepsilon}}
    \le C_\varepsilon\norm{g}_{\hi{1/2-\varepsilon}}.
\end{equation*}
The norm on the left is taken after zero extension. Since
$1<3/2-\varepsilon<3/2$, the zero endpoint values give
$J_0^*g\in\thi{3/2-\varepsilon}$.
Duality against $g\in\thi{1/2-\varepsilon}$ proves the first bound in
\eqref{eqn:centered-primitive-bounds}.
For the second bound, zero extension is bounded at the index
$-1/2+\varepsilon\in(-1/2,0)$; apply Lemma~\ref{lem:theta-kappa} and
subtract the mean. Both extensions satisfy
$\p_sJ_0f=f$ and $\dual{J_0f}{1}=0$.

Using~\eqref{eqn:centered-primitive}, we can write~\eqref{eqn:reconstructed-angle} and~\eqref{eqn:reconstructed-centered-curve} as
$\theta=\bar\theta+J_0\kappa$ and
$\bm X=\overline{\bm X}+J_0\bm\tau$.
Theorem~\ref{thm:main-regularity} and
\eqref{eqn:centered-primitive-bounds} give
\begin{equation}\label{eqn:reconstruction-angle-time}
    \p_t\theta=\p_t\bar\theta+J_0\p_t\kappa
    \in C\paren{[\tau,T];H^{-1/2+\varepsilon}(I)}.
\end{equation}
The clockwise normal convention gives $\p_\theta\bm e_r=-\bm n$.
Since $\p_t\theta$ lies only in a negative Sobolev space, a pointwise
chain rule is not available. We therefore mollify $\theta$ in time on
compact subintervals of $(0,T)$. The mollified angles converge in
$C_tH^{5/2+\varepsilon_1}$, and their derivatives converge in
$C_tH^{-1/2+\varepsilon}$. By Sobolev composition and~\eqref{eqn:product-restriction}, the products $-\bm n\,\p_t\theta$ converge in
$C_tH^{-1/2+\varepsilon}$.
Passing to the limit in the smooth chain rule yields
\begin{equation}\label{eqn:tangent-time-chain-rule}
    \p_t\bm\tau=-\bm n\,\p_t\theta.
\end{equation}
Differentiating~\eqref{eqn:reconstructed-centered-curve} gives
\begin{equation}\label{eqn:reconstruction-curve-time}
    \p_t\bm X
    =\p_t\overline{\bm X}-J_0(\bm n\,\p_t\theta)
    \in C\paren{[\tau,T];H^{1/2+\varepsilon}(I)}.
\end{equation}
The derivative is continuous on $[\tau,T]$, with endpoint values given
by the one-sided time derivatives of $\bm X$. Integrating the
distributional identity recovers $\bm X$.
Sobolev embedding therefore gives
\begin{equation}\label{eqn:main-physical-regularity}
\begin{aligned}
    \bm X&\in C\paren{[\tau,T];C^{3,\alpha}(\overline I)}
        \cap C^1\paren{[\tau,T];C^{0,\alpha}(\overline I)},\\
    \lambda&\in C\paren{[\tau,T];C^{0,\alpha}(\overline I)}.
\end{aligned}
\end{equation}

\emph{Step 4. Verify the original Stokes--filament equations.}
The identities~\eqref{eqn:reduced-velocity-tan}--\eqref{eqn:reduced-velocity-normal}
connect the reduced equation to the kinematic condition: the tangential
velocity derivative vanishes, while the normal component determines
the angle evolution. We justify them by approximation before recovering
the curve evolution.

The reconstruction gives $\p_s\theta=\kappa$ and
$\p_s\bm X=\bm\tau$, so $\abs{\p_s\bm X}=1$.
The Frenet identities give
\begin{equation}\label{eqn:reconstructed-force}
    \lambda\p_s\bm X-\p_s^3\bm X
    =\p_s\kappa\,\bm n+\sigma\bm\tau
    =\bm q.
\end{equation}

For the boundary integral $\bm v_\kappa$, the splitting~\eqref{eqn:stokeslet-splitting} gives, in
$H^{-1/2}(I)$ and $H^{-3/2}(I)$, respectively,
\begin{subequations}\label{eqn:reduced-velocity}
\begin{align}
    \tk\cdot\p_s\bm v_\kappa&=0,
        \label{eqn:reduced-velocity-tan}\\
    -\p_s\paren{\nk\cdot\p_s\bm v_\kappa}
        &=\Lambda\kappa+\p_sN(\kappa).
        \label{eqn:reduced-velocity-normal}
\end{align}
\end{subequations}
To justify these identities, first take smooth supported $\kappa$ and
$\sigma$, without imposing the tension equation. In the tangential
component of $\p_s\bm v_\kappa$, the field generated by
$\p_s(\p_s\kappa\,\nk)$ contributes $F(\kappa)$, while
\eqref{eqn:tension-tangential-id} gives the contribution
$-T_\kappa\sigma$ from $\p_s(\sigma\tk)$. For the normal component, use
$\nk\cdot\bm q_\kappa=\p_s\kappa$ in
\eqref{eqn:reconstruction-boundary-integral}. Thus the unreduced
identities are
\begin{equation}\label{eqn:unreduced-velocity-identities}
\begin{aligned}
    \tk\cdot\p_s\bm v_\kappa
    &=F(\kappa)-T_\kappa\sigma,\\
    -\p_s\paren{\nk\cdot\p_s\bm v_\kappa}
    &=\frac14\p_s^2\mc H\p_s\kappa
        +\p_s\wt N(\kappa,\sigma),
\end{aligned}
\end{equation}
where $\wt N$ is defined in \eqref{eqn:N-def}.

For $\kappa\in\thi{3/2}$ and $\sigma\in\thi{1/2}$, choose smooth
supported functions $\kappa_j\to\kappa$ and $\sigma_j\to\sigma$ in
these two spaces. For all sufficiently large $j$, the normalized curves
have a uniform arc-chord lower bound. Equations~\eqref{eqn:hilbert-bound}, \eqref{eqn:product-supported}, \eqref{eqn:Sop-bound}, and~\eqref{eqn:Sop-difference} give convergence of
$\bm q_j=(\p_s\kappa_j)\bm n_{\kappa_j}+\sigma_j\bm\tau_{\kappa_j}$
to $\bm q_\kappa$ in $\thi{1/2}$ and of the corresponding boundary
integrals in $H^{1/2}(I)$. After differentiation, the multiplication
bound~\eqref{eqn:product-restriction} controls the tangential and normal
projections in $H^{-1/2}(I)$. The operator estimates
\eqref{eqn:commutator-low-bound}, \eqref{eqn:Sop-bound}, and
\eqref{eqn:Sop-difference} give convergence of the commutator and
remainder terms there, and hence of their derivatives in $H^{-3/2}(I)$.
Thus both identities in~\eqref{eqn:unreduced-velocity-identities}
pass to the limit in the stated spaces. Imposing~\eqref{eqn:ten-det} and setting $N(\kappa)=\wt N(\kappa,\sigma(\kappa))$ in~\eqref{eqn:N-def} give
\eqref{eqn:reduced-velocity-tan}--\eqref{eqn:reduced-velocity-normal}.
Equation~\eqref{eqn:kappa-time-regularity} therefore gives, for every $t>0$,
\begin{equation}\label{eqn:kappa-velocity-form}
    \p_t\kappa
    =-\p_s\paren{\nk\cdot\p_s\bm v_\kappa}
    =\p_sh_\kappa
    \quad\text{in }H^{-3/2}(I).
\end{equation}

Since $\p_s\theta=\kappa$, \eqref{eqn:kappa-velocity-form} implies
$\p_s(\p_t\theta-h_\kappa)=0$. On $[\tau,T]$, both terms belong to
$H^{-1/2+\varepsilon}(I)$, so their means are well defined.
Equation~\eqref{eqn:reconstructed-mean-angle} gives
\begin{equation*}
    \dual{\p_t\theta-h_\kappa}{1}
    =2\p_t\bar\theta-\dual{h_\kappa}{1}=0.
\end{equation*}
The mean-angle equation therefore removes the constant of integration.
By rotational covariance and~\eqref{eqn:reduced-velocity-tan},
\begin{equation*}
    \p_t\theta=h_\kappa=-\bm n\cdot\p_s\bm v,
    \qquad \bm\tau\cdot\p_s\bm v=0.
\end{equation*}
The chain rule~\eqref{eqn:tangent-time-chain-rule} now gives
\begin{equation*}
    \p_t\bm\tau
    =-\bm n\,\p_t\theta
    =(\bm n\cdot\p_s\bm v)\bm n
    =\p_s\bm v.
\end{equation*}
Since $\p_s\bm X=\bm\tau$, this implies
$\p_s(\p_t\bm X-\bm v)=0$. The curve $\bm Y$ has zero mean,
and \eqref{eqn:reconstructed-centered-curve} prescribes
$\p_t\overline{\bm X}=\frac12\int_I\bm v$.
The centroid equation fixes precisely the remaining spatially constant
velocity: $\p_t\bm X-\bm v$ has zero mean and therefore vanishes. By
\eqref{eqn:reconstruction-positive-time-trace} and
\eqref{eqn:main-physical-regularity},
\begin{equation}
    \p_t\bm X(s,t)=\bm v(s,t)=\bm u(\bm X(s,t),t)
\end{equation}
holds pointwise for every $s\in\overline I$ and $t>0$, using the
continuous velocity trace at the endpoints.

For fixed $t>0$, Lemma~\ref{lem:stokes-layer-open-arc} gives
\eqref{eqn:intro-stokes-bulk} off $\overline{\Gamma(t)}$, pointwise
because the single-layer fields are smooth there.
The classical jump relations for the localized layer potentials, with
the convention in
\cite[Theorem~2.1 and Appendix~A]{KuoLaiMoriRodenberg2023Tension}, give
\begin{equation}
    \jump{\bm u}=0,\quad
    \jump{\Sigma\bm n}=\p_s\bm q
        =\p_s\paren{\lambda\p_s\bm X-\p_s^3\bm X}
\end{equation}
pointwise at every interior point of the filament.

\emph{Step 5. Verify endpoint conditions, initial data, far-field behavior,
and uniqueness.}
Since $\kappa\in\thi{3/2+\varepsilon}$, its distributional derivative
belongs to $\thi{1/2+\varepsilon}$. For $\varepsilon>0$, Sobolev
embedding makes the zero extensions of both $\kappa$ and $\p_s\kappa$
continuous; their support in $\overline I$ forces them to vanish at
both endpoints. Likewise, $\sigma\in\thi{1/2+\varepsilon}$ has a
continuous zero extension and zero endpoint traces. Hence
\begin{equation}
    \kappa=\p_s\kappa=\sigma=0\quad\text{on }\p I.
\end{equation}
The identities $\p_s^2\bm X=-\kappa\bm n$ and
$\lambda\p_s\bm X-\p_s^3\bm X=\p_s\kappa\,\bm n+\sigma\bm\tau$
from~\eqref{eqn:reconstructed-force} give the pointwise free-end conditions
\begin{equation}
    \p_s^2\bm X=0,\quad
    \lambda\p_s\bm X-\p_s^3\bm X=0
    \quad\text{at }s=\pm1.
\end{equation}
At the initial time, continuity of $\theta$ and the centroid, together
with \eqref{eqn:reconstruction-initial-means}, gives
\begin{equation}\label{eqn:main-X-regularity}
    \bm X\in C\paren{[0,T];C^1(\overline I)},\quad
    \bm X(\cdot,0)=\bm X_0,\quad
    \abs{\p_s\bm X}=1.
\end{equation}

Since $\p_s\bm q$ is the derivative of a supported distribution, it has
zero total force. The logarithmic far field cancels, and
Lemma~\ref{lem:stokes-layer-open-arc} gives
$\bm u(\bm x,t),p(\bm x,t)\to0$ as $\abs{\bm x}\to\infty$.

For a fixed curvature, Proposition~\ref{prop:main-tension} determines
$\sigma$. Equation~\eqref{eqn:reconstructed-mean-angle} fixes $\bar\theta$ from its prescribed
initial value, hence also $\theta$ and $\theta_*$.
The centroid equation~\eqref{eqn:reconstructed-centered-curve} fixes the
translation. If the initial angle and centroid are left unspecified,
any two reconstructions differ by a fixed rotation and translation.

Finally, Lemma~\ref{lem:stokes-layer-open-arc} and
\eqref{eqn:reconstructed-fluid} uniquely fix the fluid fields with the
prescribed far-field normalization. All components of
\eqref{eqn:intro-stokes-bulk}--\eqref{eqn:intro-kinematic} and
\eqref{eqn:intro-free-end}, together with the initial data, have now been
verified in the sense stated in Corollary~\ref{cor:main-reconstruction}.
\end{proof}


\section{Global solutions and long-time behavior}
\label{sec:global-theory}
We prove Theorem~\ref{thm:main-global}, beginning with the energy identity,
geometric control of the arc-chord constant, and rigidity of stationary
solutions. We then repeat the local fixed-point argument in an exponentially
weighted space to obtain small-data global existence and decay.
Finally, we use energy monotonicity and compactness to restart solutions
near a possible finite maximal time and prove the continuation criterion.
For global solutions, we study time translates and use zero-dissipation
rigidity to prove the energy gap and long-time alternatives.

\subsection{Energy dissipation and geometric control}
\label{subsec:energy-geometry}

\begin{proposition}[Energy dissipation and behavior at the initial time]
\label{prop:main-energy-identity}
Fix $0<\delta\le\frac12$ and $\kappa_0\in\mc U_\delta$, and let $\kappa$ be its
maximal mild solution. For $0<a<b<\tau_{\mathrm{max}}(\kappa_0)$,
the constructed solution to the Stokes--filament system satisfies
\begin{equation}\label{eqn:main-energy-positive}
    \mc E(b)
    +2\int_a^b\int_{\R^2}\abs{\grad_S\bm u(\bm x,t)}^2\,d\bm x\,dt
    =\mc E(a).
\end{equation}
If $\kappa_0\in\wt L^2(I)$, then~\eqref{eqn:main-energy-positive} extends to the initial time:
\begin{equation}\label{eqn:main-energy-initial}
    \mc E(t)
    +2\int_0^t\int_{\R^2}
        \abs{\grad_S\bm u(\bm x,r)}^2\,d\bm x\,dr
    =\mc E(0),
    \quad 0\le t<\tau_{\mathrm{max}}(\kappa_0).
\end{equation}
If $\kappa_0\notin\wt L^2(I)$, then
\begin{equation}\label{eqn:main-infinite-initial-energy}
    \lim_{t\downarrow0}\mc E(t)=\infty.
\end{equation}
\end{proposition}
\begin{proof}
\emph{Positive-time energy identity.}
Fix $0<a<b<\tau_{\mathrm{max}}(\kappa_0)$ and set
\begin{equation}
    V=\thi{\frac32},\quad H=L^2(I),\quad
    V^*=H^{-\frac32}(I).
\end{equation}
The continuity in~\eqref{eqn:XT-intro} gives $\kappa\in C([a,b];V)$, while
\eqref{eqn:main-positive-time-regularity} gives $\p_t\kappa\in L^2(a,b;V^*)$. The chain rule for a Gelfand triple shows
that $t\mapsto\mc E(t)$ is absolutely continuous and
\begin{equation}
    \frac{d}{dt}\mc E(t)
    =
    \dual{\p_t\kappa}{\kappa}_{V^*,V}
\end{equation}
for almost every $t\in(a,b)$. To identify this derivative with Stokes
dissipation, rewrite the curvature equation using the force primitive
$\bm q_\kappa$: by~\eqref{eqn:kappa-velocity-form}, distributional
integration by parts, and~\eqref{eqn:reduced-velocity-tan},
\begin{equation}\label{eqn:weak-curvature-energy}
    \dual{\p_t\kappa}{\kappa}_{V^*,V} = \dual{\nk\cdot\p_s\bm v_\kappa}{\p_s\kappa} =\dual{\p_s\bm
    v_\kappa}{\p_s\kappa\,\nk} = \dual{\p_s\bm v_\kappa} {\p_s\kappa\,\nk+\sigma\tk} =-\dual{\p_s \bm q_\kappa}{\bm
    v_\kappa}.
\end{equation}
The first three pairings on the right use the
$H^{-\frac12}(I)$--$\wt H^{\frac12}(I)$ duality. Since
$\p_s \bm q_\kappa\in\wt H^{-\frac12}(I)$, the last is the
$\wt H^{-\frac12}(I)$--$H^{\frac12}(I)$ pairing defined in
Lemma~\ref{lem:stokes-layer-open-arc}. Distributional integration by
parts with these supported densities produces no endpoint term. Applying
\eqref{eqn:weak-stokes-energy} with $\bm g=\bm q_\kappa$, and using the invariance
of the dissipation under the rotation and translation in
Corollary~\ref{cor:main-reconstruction}, gives
\begin{equation}
    \frac{d}{dt}\mc E(t)
    =
    -2\int_{\R^2}\abs{\grad_S\bm u(\bm x,t)}^2\,d\bm x.
\end{equation}
Integration over $(a,b)$ proves \eqref{eqn:main-energy-positive}.

\emph{Behavior at $t=0$ for $L^2$ data.}
If $\kappa_0\in\wt L^2(I)$, the solution constructed with
$\delta=\frac12$ agrees with the maximal solution by
Proposition~\ref{lem:restart-uniqueness}. Hence
$\kappa\in C([0,b];\wt L^2(I))$. Letting $a\downarrow0$ in
\eqref{eqn:main-energy-positive} and using strong continuity and monotone
convergence for the nonnegative dissipation proves
\eqref{eqn:main-energy-initial}.

\emph{Non-$L^2$ initial data.}
Suppose $\kappa_0\notin\wt L^2(I)$. The monotonicity in
\eqref{eqn:main-energy-positive} gives a limit in $[0,\infty]$ as
$t\downarrow0$. A finite limit would bound $\kappa(t_j)$ in $\wt L^2(I)$
for every sequence $t_j\downarrow0$, giving a weakly convergent
subsequence in $L^2(\R)$. Continuity at the initial time also gives
\begin{equation}
    \kappa(t_j)\longrightarrow\kappa_0
    \quad\text{in }\thi{-\frac12+\delta}.
\end{equation}
The distributional limits must agree, so $\kappa_0\in\wt L^2(I)$,
a contradiction.
\end{proof}

The following estimate controls chord lengths by total turning.

\begin{lemma}
\label{lem:energy-arc-chord}
Let $\kappa\in\wt L^2(I)$. For $-1\le s<\eta\le1$, set
\begin{equation}
    h=\eta-s,
    \quad
    V_{s,\eta}=\int_s^\eta\abs{\kappa(r)}\,dr.
\end{equation}
If $V_{s,\eta}<\pi$, then
\begin{equation}\label{eqn:turning-chord-bound}
    \frac{\abs{\xk(\eta)-\xk(s)}}{h}
    \ge\cos\paren{\frac{V_{s,\eta}}2}.
\end{equation}
Consequently, if $\sqrt h\norm{\kappa}_{L^2(s,\eta)}<\pi$, then
\begin{equation}\label{eqn:l2-chord-bound}
    \frac{\abs{\xk(\eta)-\xk(s)}}{h}
    \ge
    \cos\paren{
        \frac{\sqrt h}{2}\norm{\kappa}_{L^2(s,\eta)}
    }.
\end{equation}
In particular,
\begin{equation}\label{eqn:energy-arc-chord-bound}
    \frac12\norm{\kappa}_{L^2(\R)}^2<\mc E_{\mathrm{ac}}
    \quad\Longrightarrow\quad
    \abs{\kappa}_*
    \ge
    \cos\paren{\frac{\norm{\kappa}_{L^2(\R)}}{\sqrt2}}
    >0.
\end{equation}
Moreover, if a sequence $(\kappa_j)$ is bounded in $\wt L^2(I)$ and
$\abs{\kappa_j}_*\to0$, then
\begin{equation}\label{eqn:sequential-energy-barrier}
    \liminf_{j\to\infty}
    \frac12\norm{\kappa_j}_{L^2(\R)}^2
    \ge\mc E_{\mathrm{ac}}.
\end{equation}
\end{lemma}
\begin{proof}
When $V_{s,\eta}<\pi$, the tangents lie in an angular sector of width at
most $V_{s,\eta}$; projection onto its bisector bounds the chord from below.
Since $\theta_\kappa$ is absolutely continuous with
$\p_s\theta_\kappa=\kappa$, its extrema on $[s,\eta]$ satisfy
\begin{equation}
    \theta_+-\theta_-\le V_{s,\eta},
    \quad
    \theta_-=\min_{[s,\eta]}\theta_\kappa,
    \quad
    \theta_+=\max_{[s,\eta]}\theta_\kappa.
\end{equation}
If $V_{s,\eta}<\pi$, the unit vector $\bm e$ with angle
$(\theta_-+\theta_+)/2$ satisfies
$\tk(r)\cdot\bm e\ge\cos(V_{s,\eta}/2)>0$ for $r\in[s,\eta]$. Hence
\begin{equation}
    \abs{\xk(\eta)-\xk(s)}
    \ge\paren{\xk(\eta)-\xk(s)}\cdot\bm e
    =\int_s^\eta\tk(r)\cdot\bm e\,dr
    \ge h\cos\paren{\frac{V_{s,\eta}}2}.
\end{equation}
This proves~\eqref{eqn:turning-chord-bound}. The Cauchy--Schwarz bound
$V_{s,\eta}\le\sqrt h\norm{\kappa}_{L^2(s,\eta)}$ and monotonicity of cosine
give~\eqref{eqn:l2-chord-bound} under its stated hypothesis.

Suppose that
$\frac12\norm{\kappa}_{L^2(\R)}^2<\mc E_{\mathrm{ac}}=\pi^2/4$.
Since $h\le2$, every subinterval satisfies
\begin{equation}
    \frac{\sqrt h}{2}\norm{\kappa}_{L^2(s,\eta)}
    \le\frac{\norm{\kappa}_{L^2(\R)}}{\sqrt2}<\frac\pi2.
\end{equation}
Applying~\eqref{eqn:l2-chord-bound} and taking the infimum over $s<\eta$ proves
\eqref{eqn:energy-arc-chord-bound}.

Finally, if~\eqref{eqn:sequential-energy-barrier} failed, some subsequence
would satisfy $\frac12\norm{\kappa_j}_{L^2(\R)}^2\le E_0<\mc E_{\mathrm{ac}}$.
Then~\eqref{eqn:energy-arc-chord-bound} would give
$\abs{\kappa_j}_*\ge\cos\sqrt{E_0}>0$, contradicting
$\abs{\kappa_j}_*\to0$.
\end{proof}

In particular, energy below $\mc E_{\mathrm{ac}}=\pi^2/4$ excludes
arc-chord degeneration.

\begin{proposition}\label{prop:std-stt}
The only stationary mild solution of~\eqref{eqn:kap_evo} is
$\kappa_\infty=0$, which represents a straight filament.
\end{proposition}
\begin{proof}
For a stationary mild solution $\kappa$, rigidity follows from zero
dissipation and the free-end conditions.
Let $(\bm X,\bm u,p,\lambda)$ be the reconstruction associated with $\kappa$
from Corollary~\ref{cor:main-reconstruction}. Applying
\eqref{eqn:main-energy-positive} on any positive-time interval gives
$\grad_S\bm u=0$. Corollary~\ref{cor:main-reconstruction} and the far-field
normalization then give $\bm u=0$. The complement of the closure of the embedded open arc
is connected, so the bulk equation~\eqref{eqn:intro-stokes-bulk} implies that
the pressure is constant. The far-field normalization therefore gives $p=0$.
The interface condition~\eqref{eqn:intro-stokes-interface} and the force
formula~\eqref{eqn:bending-force} imply
\begin{equation}
    \p_s\paren{\lambda\p_s\bm X-\p_s^3\bm X}=0.
\end{equation}
The free-end boundary condition \eqref{eqn:intro-free-end} then gives
$\lambda\p_s\bm X-\p_s^3\bm X=0$ on $I$. Taking the normal component of
$\lambda\p_s\bm X-\p_s^3\bm X=(\lambda+\kappa^2)\bm\tau+\p_s\kappa\,\bm n$
gives $\p_s\kappa=0$ on $I$. Thus $\kappa$ is constant. The
free-end condition~\eqref{eqn:intro-free-end}, namely $\p_s^2\bm X=-\kappa\bm n=0$ on $\p I$,
then yields $\kappa=0$ on $I$.
\end{proof}

\subsection{Small-data global existence and exponential decay}
\label{subsec:small-data-proof}

We repeat the fixed-point argument of Section~\ref{sec:local-theory}
with a weight that captures both the $t^{-(2-\delta)/3}$ singularity
near zero and exponential decay $e^{\omega t}$ at large time.
Fix $0<\delta\le\frac12$ and $\omega\in(\omega_\Lambda,0)$, where
$\omega_\Lambda<0$ is the spectral bound from
Proposition~\ref{prop:lambda-semigroup}. Define the exponentially weighted
solution space
\begin{equation}\label{eqn:Y-intro}
    Y
    \coloneqq
    \left\{
    \kappa\in C\paren{[0,\infty);\thi{-\frac12+\delta}}
    \cap C\paren{(0,\infty);\thi{\frac32}}:
    \norm{\kappa}_Y<\infty
    \right\},
\end{equation}
with norm
\begin{equation}\label{eqn:Y-norm-intro}
\begin{aligned}
    \norm{\kappa}_Y
    &=\norm{\kappa}_{Y_1}+\norm{\kappa}_{Y_2},\quad
    \norm{\kappa}_{Y_1}
    =\sup_{t\ge0}e^{-t\omega}
    \norm{\kappa(t)}_{\hr{-\frac12+\delta}},\\
    \norm{\kappa}_{Y_2}
    &=\sup_{t>0}e^{-t\omega}
    \min\left\{t^{\frac{2-\delta}{3}},1\right\}
    \norm{\kappa(t)}_{\hr{\frac32}}.
\end{aligned}
\end{equation}

We write $Y^{(\delta,\omega)}$ when the parameters need to be specified.

\begin{proof}[Proof of Theorem~\ref{thm:main-global}(1)]
Fix $0<\delta\le\frac12$ and
$\omega\in(\omega_\Lambda,0)$. In $Y$, consider the closed ball of radius
$M>0$:
\begin{equation}\label{eqn:BM-intro}
    \mc B_M=\left\{\kappa\in Y:\norm{\kappa}_Y\le M\right\}.
\end{equation}

Fix $\varepsilon=\delta/2$. The constants in this proof may depend on
$\delta,\omega$, but not on time or the small radius $M$. Set
\begin{equation}
    w(t)=\min\{t^\frac{2-\delta}{3},1\}.
\end{equation}
\emph{Time-convolution bounds.}
We prove the exponentially weighted analogue~\eqref{eqn:global-duhamel-est}
of the local Duhamel estimate~\eqref{eqn:local-duhamel-est}. Set
\begin{equation}
    a=\frac{2-\delta}{3},\quad
    \beta=\frac{1+\delta-\varepsilon}{3},\quad
    \gamma=1-\frac{\varepsilon}{3}.
\end{equation}
Then $0<a,\beta,\gamma<1$, and
\begin{equation}\label{eqn:glb-conv-low}
    \sup_{t>0}\int_0^t e^{\omega\tau}
    \max\{(t-\tau)^{-\beta},1\}w(\tau)^{-1}\,d\tau <\infty,
\end{equation}
\begin{equation}\label{eqn:glb-conv-high}
    \sup_{t>0}w(t)\int_0^t e^{\omega\tau}
    \max\{(t-\tau)^{-\gamma},1\}w(\tau)^{-1}\,d\tau <\infty.
\end{equation}
For $0<t\le1$, we have $w(\tau)=\tau^a$ on $(0,t)$, so
\begin{equation}
    \int_0^t (t-\tau)^{-\beta}\tau^{-a}\,d\tau
    \le C t^{1-a-\beta}
    = Ct^{\varepsilon/3}\le C,
\end{equation}
while
\begin{equation}
    w(t)\int_0^t (t-\tau)^{-\gamma}\tau^{-a}\,d\tau
    \le C t^a t^{1-a-\gamma}
    = Ct^{\varepsilon/3}\le C.
\end{equation}
For $t>1$, split at $t/2$. On $(0,t/2)$, the kernels are uniformly
bounded and $\int_0^\infty e^{\omega\tau}w(\tau)^{-1}\,d\tau<\infty$.
On $(t/2,t)$, $w(\tau)^{-1}$ is uniformly bounded and
$e^{\omega\tau}\le e^{\omega t/2}$, with
\begin{equation}
    \int_0^{t/2}\max\{r^{-\beta},1\}\,dr
    +\int_0^{t/2}\max\{r^{-\gamma},1\}\,dr
    \le C(1+t).
\end{equation}
The late contribution is bounded by $Ce^{\omega t/2}(1+t)$, proving
\eqref{eqn:glb-conv-low}--\eqref{eqn:glb-conv-high}.

The forcing weight $e^{2\omega t}w(t)^{-1}$ matches the quadratic
nonlinearity: the low and high curvature norms contribute
$e^{\omega t}$ and $e^{\omega t}w(t)^{-1}$, respectively.
For $f\in C((0,\infty);\hi{-\frac12+\varepsilon})$,
\eqref{eqn:lambda-semigroup-est} and~\eqref{eqn:glb-conv-low}--\eqref{eqn:glb-conv-high} give
\begin{equation}\label{eqn:global-duhamel-est}
    \norm{\int_0^\cdot e^{(\cdot-\tau)\Lambda}\p_s f(\tau)\,d\tau}_Y
    \le C\sup_{t>0}e^{-2\omega t}w(t)
        \norm{f(t)}_{\hi{-\frac12+\varepsilon}},
\end{equation}
whenever the right-hand side is finite. Continuity in time follows as in the
proof of Proposition~\ref{prop:local-fixed-exponent}.

\emph{Invariance and contraction.}
For $\kappa\in\mc B_M$, \eqref{eqn:Y-norm-intro} gives
\begin{equation}
    \norm{\kappa(t)}_{\hr{-\frac12+\delta}}\le Me^{\omega t}\quad(t\ge0),
    \quad
    \norm{\kappa(t)}_{\hr{\frac32}}\le Mw(t)^{-1}e^{\omega t}\quad(t>0).
\end{equation}
Assume $M\le\min\{1,(2C_{J,\delta})^{-1}\}$. Since
$\abs{\kappa_\infty}_*=1$ for $\kappa_\infty=0$ and the small low norm
makes the normalized curves $C^1$-close to the straight filament,
\eqref{eqn:arc-chord-C1-continuity} gives
\begin{equation}\label{eqn:global-ball-geometry}
    \abs{\rho\xka{1}(t)+(1-\rho)\xka{2}(t)}_*
    \ge1-C_{J,\delta}M\ge\frac12,
    \quad 0\le\rho\le1,
\end{equation}
for all $\kappa_1,\kappa_2\in\mc B_M$ and $t\ge0$.
The inverse bound~\eqref{eqn:inverse-uniform-bound} is uniform in time
and over $\kappa\in\mc B_M$. Bounding $C_{M,1/2}$ by $C_{1,1/2}$,
we denote the resulting constants by $C$.

Equations~\eqref{eqn:N-pointwise-bound} and~\eqref{eqn:N-pointwise-difference} give, for $t>0$,
\begin{equation}\label{eqn:global-nonlinear-bounds}
    \begin{aligned}
        \norm{N(\kappa(t))}_{\hi{-\frac12+\varepsilon}}
        &\le CM^2e^{2\omega t}w(t)^{-1},\\
        \norm{N(\kappa_1(t))-N(\kappa_2(t))}_{\hi{-\frac12+\varepsilon}}
        &\le CMe^{2\omega t}w(t)^{-1}\norm{\kappa_1-\kappa_2}_Y.
    \end{aligned}
\end{equation}
Here quadratic smallness replaces the small-time factor
$T^{\varepsilon/3}$ in the local argument.
Applying~\eqref{eqn:global-duhamel-est} to~\eqref{eqn:global-nonlinear-bounds} and using the consequence of~\eqref{eqn:lambda-semigroup-est}
$\norm{e^{\cdot\Lambda}\kappa_0}_Y
\le C\norm{\kappa_0}_{\hr{-\frac12+\delta}}$ gives
\begin{equation}\label{eqn:global-map-bounds}
    \begin{aligned}
        \norm{\Psi_{\kappa_0}\kappa}_Y
        &\le C\norm{\kappa_0}_{\hr{-\frac12+\delta}}+CM^2,\\
        \norm{\Psi_{\kappa_0}\kappa_1-\Psi_{\kappa_0}\kappa_2}_Y
        &\le CM\norm{\kappa_1-\kappa_2}_Y.
    \end{aligned}
\end{equation}
Fix $C\ge1$ so that~\eqref{eqn:global-map-bounds} holds. Choose $M_*>0$
to ensure both geometry and contraction:
\begin{equation}
    M_*\le\min\{1,(2C_{J,\delta})^{-1}\},\quad CM_*\le\frac12,
\end{equation}
Set the ball radius $M=2C\norm{\kappa_0}_{\hr{-\frac12+\delta}}$ and
the corresponding initial-data threshold $M_0=M_*/(2C)$. If $\norm{\kappa_0}_{\hr{-\frac12+\delta}}\le M_0$, then
$M\le M_*$ and
\begin{equation}
    C\norm{\kappa_0}_{\hr{-\frac12+\delta}}+CM^2\le M,
    \quad CM\le\frac12.
\end{equation}
Thus $\Psi_{\kappa_0}$ is a contraction of $\mc B_M$ into itself.
The Banach fixed-point theorem gives a unique fixed point in $\mc B_M$,
which is a global mild solution by~\eqref{eqn:global-ball-geometry}.
Uniqueness in $Y$ follows from Theorem~\ref{thm:main-local} and
Proposition~\ref{lem:restart-uniqueness}. The bound
$\norm{\kappa}_Y\le M=2C\norm{\kappa_0}_{\hr{-\frac12+\delta}}$
gives~\eqref{eqn:main-global-decay}.
\end{proof}

\subsection{Continuation and asymptotic alternatives}
\label{subsec:global-behavior-proof}

Using the positive-time $L^2$ bound, we extract a limit in the local-theory
topology and restart past a proposed finite maximal time whenever the
arc-chord constant is bounded below by a positive constant along a sequence
approaching that time. For global solutions with an eventual uniform
arc-chord bound, we pass to a limit of time translates. The limiting local
solution has constant energy, hence zero Stokes dissipation, and the
zero-dissipation argument forces it to be straight. Below $\mc E_{\mathrm{ac}}$,
we apply the geometric bound and this rigidity argument to prove the energy
gap. We then prove arc-chord degeneration for positive limiting energy and,
for zero limiting energy, restart in the small-data regime to obtain
exponential decay.

\begin{proof}[Proof of Theorem~\ref{thm:main-global}(2)]
\emph{Positive-time energy bound and maximal solution.}
Fix $0<\delta\le\frac12$ and $\kappa_0\in\mc U_\delta$. The maximal
continuation from Theorem~\ref{thm:main-local} and
Proposition~\ref{lem:restart-uniqueness} satisfies
\eqref{eqn:main-local-class} on every $[0,T]$ with
$T<\tau_{\text{max}}(\kappa_0)$.

By~\eqref{eqn:main-energy-positive}, $\mc E$ is non-increasing for
positive times. Thus, for every $0<a<\tau_{\text{max}}(\kappa_0)$,
\begin{equation}\label{eqn:global-positive-time-bound}
    \sup_{a\le t<\tau_{\text{max}}(\kappa_0)}
    \norm{\kappa(t)}_{L^2(\R)}
    \le\norm{\kappa(a)}_{L^2(\R)}.
\end{equation}
The limiting energy $\mc E_*\in[0,\infty)$ therefore exists. We also use
this bound in the compactness arguments for continuation and time translates.

\emph{Finite maximal time.}
Suppose that $\tau_{\text{max}}(\kappa_0)<\infty$ and that the arc-chord
constant does not tend to zero. There are $m>0$ and times
$t_j\uparrow\tau_{\text{max}}(\kappa_0)$ such that
$\abs{\kappa(t_j)}_*\ge m$. For any fixed
$a\in(0,\tau_{\text{max}}(\kappa_0))$, estimate
\eqref{eqn:global-positive-time-bound} gives a uniform $L^2$ bound for
$\kappa(t_j)$. Fix $0<\delta_0<\frac12$ and set $s_0=-\frac12+\delta_0$.
By the compact embedding $\wt L^2(I)\hookrightarrow\thi{s_0}$, a subsequence
satisfies
\begin{equation*}
    \kappa(t_j)\longrightarrow\kappa_*
    \quad\text{in }\thi{s_0}.
\end{equation*}
By~\eqref{eqn:arc-chord-C1-continuity}, $\abs{\kappa_*}_*\ge m$, so the
limit remains admissible. By convergence in the local-theory topology,
the data $\kappa(t_j)$ lie in a single local well-posedness neighborhood of $\kappa_*$
for all sufficiently large indices. Theorem~\ref{thm:main-local} at exponent $\delta_0$ therefore gives a
common restart time $T_*>0$.
Proposition~\ref{lem:restart-uniqueness} identifies these solutions
with restarts of the maximal solution. Since
$t_j+T_*>\tau_{\text{max}}(\kappa_0)$ for large $j$, this contradicts
maximality. Hence
\begin{equation}
    \abs{\kappa(t)}_*\longrightarrow0
    \quad\text{as }t\uparrow\tau_{\text{max}}(\kappa_0).
\end{equation}
Applying~\eqref{eqn:sequential-energy-barrier} to any sequence
$t_j\uparrow\tau_{\text{max}}(\kappa_0)$, using the $L^2$ bound
\eqref{eqn:global-positive-time-bound}, gives
$\mc E_*\ge\mc E_{\mathrm{ac}}$.

\emph{Rigidity of nondegenerate time translates.}
For a global solution, we prove the implication
\eqref{eqn:uniform-geometry-zero-energy}: an eventual uniform arc-chord
lower bound forces $\mc E_*=0$. Assume
\begin{equation}\label{eqn:global-uniform-geometry}
    \liminf_{t\to\infty}\abs{\kappa(t)}_*>0.
\end{equation}
Fix an auxiliary exponent $0<\delta_0<\frac12$ and set
$s_0=-\frac12+\delta_0$. Let $t_j\to\infty$ be arbitrary. After discarding
finitely many terms, $\abs{\kappa(t_j)}_*\ge m_0$ for some $m_0>0$. The bound
\eqref{eqn:global-positive-time-bound}, weak compactness in $\wt L^2(I)$,
and the compact embedding into $\thi{s_0}$ give, after passing to a
subsequence,
\begin{equation}\label{eqn:global-translate-data}
    \kappa(t_j)\rightharpoonup\kappa_*
    \quad\text{in }\wt L^2(I),
    \quad
    \kappa(t_j)\longrightarrow\kappa_*
    \quad\text{in }\thi{s_0}.
\end{equation}
Equation~\eqref{eqn:arc-chord-C1-continuity} gives
$\abs{\kappa_*}_*\ge m_0$.

By Theorem~\ref{thm:main-local}, applied at $\kappa_*$ with exponent $\delta_0$,
the local solutions $\vph_j$ and $\vph_*$ with initial data $\kappa(t_j)$ and
$\kappa_*$ exist on a common interval $[0,T]$ for all sufficiently large $j$,
with
\begin{equation}\label{eqn:global-translate-solutions}
    \norm{\vph_j-\vph_*}_{X_T^{(\delta_0)}}\longrightarrow0.
\end{equation}
Equation~\eqref{eqn:shifted-duhamel} and
Proposition~\ref{lem:restart-uniqueness} identify the time translates:
\begin{equation}\label{eqn:time-shift-global}
    \vph_j(t)=\kappa(t_j+t),
    \quad 0\le t\le T.
\end{equation}

At each fixed positive time $t\in(0,T]$, the weighted convergence
\eqref{eqn:global-translate-solutions} gives strong convergence in
$\thi{3/2}$, hence in $L^2$. By~\eqref{eqn:time-shift-global} and
$\mc E(t_j+t)\to\mc E_*$, the limiting trajectory has constant energy:
\begin{equation}\label{eqn:global-limit-energy}
    \frac12\norm{\vph_*(t)}_{L^2(\R)}^2
    =\lim_{j\to\infty}\mc E(t_j+t)
    =\mc E_*,
    \quad 0<t\le T.
\end{equation}
Let $(\bm X_*,\bm u_*,p_*,\lambda_*)$ be the reconstruction associated with $\vph_*$.
The energy identity~\eqref{eqn:main-energy-positive} then gives,
for every $0<a<b\le T$,
\begin{equation}
    \int_a^b\int_{\R^2}
    \abs{\grad_S\bm u_*(\bm x,t)}^2\,d\bm x\,dt=0.
\end{equation}
The dissipation therefore vanishes for almost every $t\in(a,b)$. At each
such time, the zero-dissipation argument in the proof of
Proposition~\ref{prop:std-stt} gives $\vph_*(t)=0$. By continuity for positive
times,
\begin{equation}
    \vph_*(t)=0,
    \quad 0<t\le T.
\end{equation}
By~\eqref{eqn:global-limit-energy}, this proves the rigidity claim:
\begin{equation}\label{eqn:uniform-geometry-zero-energy}
    \liminf_{t\to\infty}\abs{\kappa(t)}_*>0
    \quad\Longrightarrow\quad
    \mc E_*=0.
\end{equation}

\emph{Energy gap and asymptotic alternatives.}
Suppose that the solution is global and
$\mc E_*<\mc E_{\mathrm{ac}}$. For some $t_0>0$,
$\mc E(t_0)<\mc E_{\mathrm{ac}}$. By~\eqref{eqn:main-energy-positive},
$\mc E(t)\le \mc E(t_0)$ for $t\ge t_0$. Lemma~\ref{lem:energy-arc-chord},
applied at time $t$, gives
\begin{equation}\label{eqn:eventual-energy-geometry}
    \abs{\kappa(t)}_*
    \ge\cos\sqrt{\mc E(t)}
    \ge\cos\sqrt{\mc E(t_0)}>0,
    \quad t\ge t_0.
\end{equation}
This verifies~\eqref{eqn:global-uniform-geometry}, so
\eqref{eqn:uniform-geometry-zero-energy} yields $\mc E_*=0$.
Together with the finite-time case, this proves the gap
\eqref{eqn:main-energy-gap}.

For a global solution with $\mc E_*>0$,
\eqref{eqn:uniform-geometry-zero-energy} implies
$\liminf_{t\to\infty}\abs{\kappa(t)}_*=0$, and~\eqref{eqn:main-energy-gap}
gives $\mc E_*\ge\mc E_{\mathrm{ac}}$.

If $\mc E_*=0$, the finite-time case already implies global existence, and
\begin{equation}
    \norm{\kappa(t)}_{L^2(\R)}\longrightarrow0.
\end{equation}
For all sufficiently large $t$, Lemma~\ref{lem:energy-arc-chord} gives
\begin{equation}
    1\ge\abs{\kappa(t)}_*
    \ge\cos\sqrt{\mc E(t)}.
\end{equation}
Hence $\abs{\kappa(t)}_*\to1$. The two alternatives are disjoint and
exhaust all global solutions.

\emph{Eventual exponential decay.}
Once the energy is sufficiently small, the solution enters the
small-data regime and the exponentially weighted theory applies from
that time onward. Assume $\mc E_*=0$ and fix $\omega\in(\omega_\Lambda,0)$. Choose
$T_\omega>0$ so large that $\norm{\kappa(T_\omega)}_{L^2(\R)}$ is below
the smallness threshold in Theorem~\ref{thm:main-global}(1) with
$\delta=\frac12$. Proposition~\ref{lem:restart-uniqueness}, including
consistency with respect to $\delta$, identifies the resulting global
solution with $\kappa(T_\omega+\cdot)$. Equation~\eqref{eqn:main-global-decay}
then gives the stated estimate.
Reconstructing the curve from the curvature gives exponential straightening
modulo rigid motions.
\end{proof}

\begin{proof}[Proof of Theorem~\ref{thm:main-global}(3)]
For initial data below the threshold, Lemma~\ref{lem:energy-arc-chord}
gives $\abs{\kappa_0}_*>0$, so the maximal solution exists.
Proposition~\ref{prop:main-energy-identity} gives
$\mc E(t)\le\mc E(t_0)<\mc E_{\mathrm{ac}}$ for
$t_0\le t<\tau_{\mathrm{max}}$, including $t_0=0$ for $L^2$ initial data.
By the energy gap \eqref{eqn:main-energy-gap}, $\mc E_*=0$.
Part~\emph{(2)} then gives global existence and exponential straightening.
Equations~\eqref{eqn:energy-arc-chord-bound} and~\eqref{eqn:main-energy-positive} give
\eqref{eqn:main-post-threshold-arc-chord}.
\end{proof}

\begin{proof}[Proof of Corollary~\ref{cor:main-straightening}]
Fix $\omega\in(\omega_\Lambda,0)$. Theorem~\ref{thm:main-global}(2)
gives $T>0$ such that
\begin{equation*}
    \norm{\kappa(t)}_{\hr{\frac32}}\le Ce^{\omega t},
    \quad |\kappa(t)|_*\ge\frac12,\quad t\ge T.
\end{equation*}
Write $\bm v=(\bm u|_{\Gamma})\circ\bm X$.
Equations~\eqref{eqn:inverse-uniform-bound}, \eqref{eqn:tension-bound}, \eqref{eqn:product-supported}, and~\eqref{eqn:stokes-layer-bound} give, as in~\eqref{eqn:reconstruction-time-bound},
\begin{equation}\label{eqn:reconstructed-velocity-decay}
    \norm{\bm v(t)}_{\hi{\frac12}}
    +\norm{\p_t\theta(t)}_{\hi{-\frac12}}
    \le C\norm{\bm q_\kappa(t)}_{\hr{\frac12}}
    \le C\norm{\kappa(t)}_{\hr{\frac32}}
    \le Ce^{\omega t},\quad t\ge T.
\end{equation}
The curvature bound and arc-chord lower bound make the constants
uniform in time. The centroid equation~\eqref{eqn:X-bar} gives
\begin{equation*}
    |\p_t\overline{\bm X}(t)|
    =\left|\frac12\int_I\bm v(s,t)\,ds\right|
    \le Ce^{\omega t}.
\end{equation*}
Integrating in time gives a limit $\overline{\bm X}_\infty$ with
$|\overline{\bm X}(t)-\overline{\bm X}_\infty|\le Ce^{\omega t}$.

For the orientation, use a compactly supported average because
$\p_t\theta$ is controlled in a negative Sobolev space.
Choose a nonnegative $\phi\in C_c^\infty(I)$ with $\int_I\phi=1$, and set
$a(t)=\int_I\phi(s)\theta(s,t)\,ds$. Since
$\phi\in\thi{\frac12}$, \eqref{eqn:reconstructed-velocity-decay} gives
\begin{equation*}
    |a'(t)|=|\dual{\p_t\theta(t)}{\phi}|\le Ce^{\omega t}.
\end{equation*}
Thus $a(t)$ converges to some $\theta_\infty\in\R$, with
$|a(t)-\theta_\infty|\le Ce^{\omega t}$.
Using $\p_s\theta=\kappa$ and $\int_I\phi=1$, we also have
\begin{equation*}
    \norm{\theta(\cdot,t)-a(t)}_{C(\overline I)}
    \le\norm{\kappa(t)}_{L^1(I)}\le Ce^{\omega t}.
\end{equation*}
Consequently $\theta(\cdot,t)$, and hence $\bar\theta(t)$, converges to
$\theta_\infty$ at the same rate. The identity
$\p_s\bm X=(\cos\theta,\sin\theta)$ and the convergence of the centroid
then give the stated $C^1$ estimate. By continuity of the reconstructed
angle and curve on $[0,T]$, the estimate holds for all $t\ge0$ after
increasing $C_\omega$.
\end{proof}


\appendix

\section{Proofs of auxiliary estimates}
\label{app:operator-estimates}
\label{app:divided-differences}

We give the proofs for the geometric, commutator, and boundary-integral estimates in Sections~\ref{sec:formulation} and~\ref{sec:tension}.

\subsection{Geometric estimates}
\label{app:geometry-estimates}

\begin{proof}[Proof of Lemma~\ref{tau-n-est}]
The proof reduces to composition estimates for $\bm e_r$ on $H^{s+1}(I)$.
The assumption on $s$ gives $\frac{1}{2}<s+1\le1$, so this space is an
algebra. For $f\in H^{s+1}(I)$ and $-\frac{1}{2}<s<0$, the Lipschitz
continuity of $\bm e_r$ gives
\begin{equation}
\begin{aligned}
    [\bm e_r\circ f]_{\hi{s+1}}^2
    &= \int_I\int_I
    \frac{\abs{\bm e_r(f(x))-\bm e_r(f(y))}^2}
    {\abs{x-y}^{1+2(s+1)}}\,dx\,dy \\
    &\le \int_I\int_I
    \frac{\abs{f(x)-f(y)}^2}
    {\abs{x-y}^{1+2(s+1)}}\,dx\,dy
    = [f]_{\hi{s+1}}^2.
\end{aligned}
\end{equation}
At $s=0$, the chain rule gives the corresponding estimate:
\begin{equation}
    \p_s(\bm e_r\circ f)
    =\bm e_r'(f)\p_s f,
    \quad
    \norm{\p_s(\bm e_r\circ f)}_{L^2(I)}
    \le \norm{\p_s f}_{L^2(I)}.
\end{equation}
Since $\abs{\bm e_r}\equiv 1$,
$\norm{\bm e_r\circ f}_{L^2(I)}$ is uniformly bounded. Thus
\begin{equation}\label{eqn:er-comp-Hr}
    [\bm e_r\circ f]_{\hi{s+1}}
    \le [f]_{\hi{s+1}},
    \quad
    \norm{\bm e_r\circ f}_{\hi{s+1}}
    \le C\paren{1+\norm{f}_{\hi{s+1}}}.
\end{equation}
Equation~\eqref{eqn:er-comp-Hr} also holds for $\bm e_\theta\coloneqq\bm e_r'=Q_{\frac{\pi}{2}}\bm e_r$, since
$\bm e_\theta$ is a fixed rotation of $\bm e_r$.

To estimate differences, write
\begin{equation}
    \bm e_r(f_1)-\bm e_r(f_2)
    = \paren{f_1-f_2}\bm G(f_1,f_2),
\end{equation}
where
\begin{equation}
    \bm G(f_1,f_2)
    \coloneqq \int_0^1 \bm e_\theta\paren{\eta f_1+(1-\eta)f_2}\,d\eta.
\end{equation}
Applying~\eqref{eqn:er-comp-Hr} to $\bm e_\theta$ gives
\begin{equation}
    \norm{\bm G(f_1,f_2)}_{\hi{s+1}}
    \le C\paren{1+\norm{f_1}_{\hi{s+1}}+\norm{f_2}_{\hi{s+1}}}.
\end{equation}
Since $H^{s+1}(I)$ is an algebra, it follows that
\begin{equation}\label{eqn:er-comp-diff}
\begin{aligned}
    \norm{\bm e_r(f_1)-\bm e_r(f_2)}_{\hi{s+1}}
    &\le C\paren{1+\norm{f_1}_{\hi{s+1}}+\norm{f_2}_{\hi{s+1}}}
    \norm{f_1-f_2}_{\hi{s+1}}.
\end{aligned}
\end{equation}

Equation~\eqref{eqn:theta-Hs+1} gives
\begin{equation}
    \norm{\thk}_{\hi{s+1}} \le C\norm{\kappa}_{\hr{s}},
    \quad
    \norm{\thka{1}-\thka{2}}_{\hi{s+1}}
    \le C\norm{\kappa_1-\kappa_2}_{\hr{s}}.
\end{equation}
Since $\xk'=\bm e_r(\thk)$, estimates~\eqref{eqn:er-comp-Hr}
and~\eqref{eqn:er-comp-diff} imply
\begin{equation}
    \norm{\xk'}_{\hi{s+1}}
    \le C\paren{1+\norm{\kappa}_{\hr{s}}},
    \quad
    \norm{\xka{1}'-\xka{2}'}_{\hi{s+1}}
    \le C\paren{1+M}\norm{\kappa_1-\kappa_2}_{\hr{s}}.
\end{equation}
The seminorm estimate in~\eqref{eqn:er-comp-Hr} also gives the sharper bound
\begin{equation}
    [\xk']_{\hi{s+1}}
    \le [\thk]_{\hi{s+1}}
    \le C\norm{\kappa}_{\hr{s}}.
\end{equation}

It remains to estimate the second derivative. Taking $f_1=\thk$ and $f_2=0$
in~\eqref{eqn:er-comp-diff} gives
\begin{equation}
    \norm{\xk'-\bm e_r(0)}_{\hi{s+1}}
    \le C\paren{1+\norm{\thk}_{\hi{s+1}}}\norm{\thk}_{\hi{s+1}}
    \le C\paren{1+\norm{\kappa}_{\hr{s}}}\norm{\kappa}_{\hr{s}}.
\end{equation}
Differentiating gives
\begin{equation}
    \norm{\xk''}_{\hi{s}}
    \le C \norm{\xk'-\bm e_r(0)}_{\hi{s+1}}
    \le C\paren{1+\norm{\kappa}_{\hr{s}}}\norm{\kappa}_{\hr{s}}.
\end{equation}
Applying the same differentiation estimate to the difference gives
\begin{equation}
    \norm{\xka{1}''-\xka{2}''}_{\hi{s}}
    \le C \norm{\xka{1}'-\xka{2}'}_{\hi{s+1}}
    \le C\paren{1+M}\norm{\kappa_1-\kappa_2}_{\hr{s}}.
\end{equation}
\end{proof}

\begin{proof}[Proof of Lemma~\ref{lem:Xk-holder}]
We first use the Sobolev regularity of $\kappa$ to bound
$\theta_\kappa$ in H\"older spaces, and then use
$\xk'=\bm e_r(\theta_\kappa)$. Equation~\eqref{eqn:theta-Hs+1} gives
\begin{equation}
    \norm{\theta_\kappa}_{\hi{s+1}} \le C\norm{\kappa}_{\hr{s}}.
\end{equation}
The assumptions on $s$ and $\alpha$ imply
$s+1>m+\alpha+\frac{1}{2}$. The one-dimensional Sobolev embedding gives
\begin{equation}\label{eqn:Xk-holder-higher-theta}
    \norm{\theta_\kappa}_{\ci{m,\alpha}} \le C\norm{\kappa}_{\hr{s}}.
\end{equation}

To estimate the reconstructed curve, Fa\`a di Bruno's formula shows that,
for $0\le j\le m$, each derivative
$\xk^{(j+1)}=\p_s^j(\bm e_r(\theta_\kappa))$ is a finite sum of terms of the
form
\begin{equation}
    G\paren{\theta_\kappa}\prod_{\ell=1}^j \paren{\p_s^\ell\theta_\kappa}^{\beta_\ell},
    \quad
    \sum_{\ell=1}^j \ell\beta_\ell = j,
\end{equation}
where $G$ is a smooth bounded function. Since $\ci{m,\alpha}$ is an algebra,
estimate~\eqref{eqn:Xk-holder-higher-theta} gives
\begin{equation}
    \norm{\xk'}_{\ci{m,\alpha}}
    \le C\paren{1+\norm{\theta_\kappa}_{\ci{m,\alpha}}}^{m+1}
    \le C\paren{1+\norm{\kappa}_{\hr{s}}}^{m+1}.
\end{equation}
Since $\xk(-1)=0$, it follows that
\begin{equation}
    \norm{\xk}_{\ci{m+1,\alpha}}
    \le C \norm{\xk'}_{\ci{m,\alpha}}
    \le C\paren{1+\norm{\kappa}_{\hr{s}}}^{m+1}.
\end{equation}
This proves~\eqref{eqn:Xk-holder-higher}.

When $m=0$, the fact that $\bm e_r$ is $1$-Lipschitz gives, for $x,y\in I$,
\begin{equation}
    \abs{\xk'(x)-\xk'(y)}
    = \abs{\bm e_r(\theta_\kappa(x))-\bm e_r(\theta_\kappa(y))}
    \le \abs{\theta_\kappa(x)-\theta_\kappa(y)}.
\end{equation}
By~\eqref{eqn:Xk-holder-higher-theta},
\begin{equation}
    [\xk']_{\ci{0,\alpha}} \le [\theta_\kappa]_{\ci{0,\alpha}}
    \le C\norm{\kappa}_{\hr{s}},
\end{equation}
which proves~\eqref{eqn:Xk-holder}.

For the difference estimate, let
$\theta_i=\theta_{\kappa_i}$ for $i=1,2$. By linearity of $J$ and
\eqref{eqn:Xk-holder-higher-theta},
\begin{equation}
    \norm{\theta_i}_{\ci{m,\alpha}} \le C M,
    \quad
    \norm{\theta_1-\theta_2}_{\ci{m,\alpha}}
    \le C \norm{\kappa_1-\kappa_2}_{\hr{s}}.
\end{equation}
The fundamental theorem of calculus gives
\begin{equation}
    \xka{1}'-\xka{2}'
    = \paren{\theta_1-\theta_2}\bm G(\theta_1,\theta_2),
    \quad
    \bm G(\theta_1,\theta_2)
    \coloneqq \int_0^1 \bm e_\theta\paren{\eta\theta_1+(1-\eta)\theta_2}\,d\eta.
\end{equation}
The smoothness of $\bm e_\theta$ and the algebra property of
$\ci{m,\alpha}$ imply
\begin{equation}
    \norm{\bm G(\theta_1,\theta_2)}_{\ci{m,\alpha}} \le C\paren{1+M}^{m+1}.
\end{equation}
Combining the last three estimates gives
\begin{equation}
    \norm{\xka{1}'-\xka{2}'}_{\ci{m,\alpha}}
    \le C\paren{1+M}^{m+1}\norm{\kappa_1-\kappa_2}_{\hr{s}}.
\end{equation}
Since $(\xka{1}-\xka{2})'=\xka{1}'-\xka{2}'$ and
$\xka{1}(-1)=\xka{2}(-1)=0$, it follows that
\begin{equation}
    \norm{\xka{1}-\xka{2}}_{\ci{m+1,\alpha}}
    \le C\paren{1+M}^{m+1}\norm{\kappa_1-\kappa_2}_{\hr{s}},
\end{equation}
which proves~\eqref{eqn:Xk-holder-higher-lip}.
\end{proof}

\subsection{Commutator estimates}
\label{app:commutator-estimates}

\begin{proof}[Proof of Lemma~\ref{lem:comm-psH}]
It is enough to take $g\in C_c^\infty(I)$ and then use density. Set
$T_fg\coloneqq[\p_s\mc H,f]g$. Differentiating the kernel in~\eqref{eqn:finite-hilbert-definition} gives
\begin{equation}\label{eq:comm-kernel}
    (T_fg)(s)
    =\frac{1}{\pi}\operatorname{p.v.}\int_I
    \frac{f(s)-f(\eta)}{(s-\eta)^2}\,g(\eta)\,d\eta.
\end{equation}
To separate the singular Hilbert-transform part, define
$V_0(s,\eta)=\frac{f(s)-f(\eta)}{s-\eta}$. The identity
\begin{equation}\label{eq:kernel-decomp}
    \frac{f(s)-f(\eta)}{(s-\eta)^2}
    =\p_\eta V_0(s,\eta)+\frac{\p_s f(\eta)}{s-\eta}
\end{equation}
gives
\begin{equation}\label{eq:comm-split}
    T_fg=I_1+I_2, \quad I_1(s)\coloneqq\frac{1}{\pi}\int_I \p_\eta V_0(s,\eta)g(\eta)\,d\eta, \quad
    I_2\coloneqq\mc H(\p_s f\cdot g).
\end{equation}

We first estimate the smoother term $I_1$ at two endpoints. Choose
$0<\alpha<\delta$. The embedding
$H^{\frac52+\delta}(I)\hookrightarrow C^{2,\alpha}(I)$ gives
\begin{equation}
    \abs{\p_sV_0(s,\eta)}+\abs{\p_\eta V_0(s,\eta)}
    \le C\norm{f}_{H^{\frac52+\delta}(I)}.
\end{equation}
Thus $I_1$ has a bounded kernel and
\begin{equation}
    \norm{I_1}_{L^2(I)}
    \le C\norm{f}_{H^{\frac52+\delta}(I)}\norm{g}_{L^2(\R)}.
\end{equation}

For the $H^1$ endpoint, integration by parts in $\eta$ gives
\begin{equation}
    I_1(s)=-\frac{1}{\pi}\int_I
    V_0(s,\eta)\p_\eta g(\eta)\,d\eta,
\end{equation}
where the boundary terms vanish because $g\in C_c^\infty(I)$. Since
\begin{equation}
    \abs{V_0(s,\eta)}+\abs{\p_sV_0(s,\eta)}
    \le C\norm{f}_{H^{\frac52+\delta}(I)},
\end{equation}
we obtain
\begin{equation}
    \norm{I_1}_{H^1(I)}
    \le C\norm{f}_{H^{\frac52+\delta}(I)}\norm{g}_{H^1(\R)}.
\end{equation}

Interpolation between the bounds
$I_1:L^2(I)\to L^2(I)$ and $I_1:\thi{1}\to H^1(I)$ gives
\begin{equation}\label{eqn:commutator-smooth-bound}
    \norm{I_1}_{\hi{\frac12+\varepsilon}}
    \le C\norm{f}_{H^{\frac52+\delta}(I)}
    \norm{g}_{\hr{\frac12+\varepsilon}}.
\end{equation}

For $I_2$, let $F\in H^{\frac52+\delta}(\R)$ be a bounded extension of
$f$. Because $g$ is supported in $\overline I$, the function
$(\p_sF)g$ restricts to $(\p_sf)g$ on $I$. Lemma~\ref{point_mult_sobo}
therefore gives
\begin{equation}
    \norm{(\p_sF)g}_{\hr{\frac12+\varepsilon}}
    \le C\norm{f}_{H^{\frac52+\delta}(I)}
    \norm{g}_{\hr{\frac12+\varepsilon}}.
\end{equation}
Lemma~\ref{Hf-bounded} then implies
\begin{equation}\label{eqn:commutator-hilbert-bound}
    \norm{I_2}_{H^{\frac12+\varepsilon}(I)}
    \le C\norm{f}_{H^{\frac52+\delta}(I)}
    \norm{g}_{\hr{\frac12+\varepsilon}}.
\end{equation}
Combining~\eqref{eqn:commutator-smooth-bound} and~\eqref{eqn:commutator-hilbert-bound} proves
\eqref{eq:comm-psH-bound}; density gives the result for every
$g\in\thi{\frac12+\varepsilon}$.
\end{proof}

\subsection{Divided differences and boundary integral estimates}
\label{app:boundary-integral-estimates}

We work on the interval and keep the dependence on the arc-chord constant
explicit.\par

\begin{lemma}\label{lem:diff-est}
Let $\bm Z \in \ci{1,\alpha}$ for $\alpha\in(0,1)$. Then, for
$s,\eta\in I$ with $s\ne\eta$,
\begin{align}
    &\abs{\Delta \bm Z} \le \norm{\bm Z' }_{\ci{0}} |s-\eta|, \\
    &\abs{\bm Z '(\eta) - \frac{\Delta \bm Z }{s-\eta}}, \abs{\bm Z '(s) - \frac{\Delta \bm Z }{s-\eta}} \le [\bm Z ']_{\ci{0,\alpha}} |s-\eta|^\alpha.\label{eqn:divided-difference-remainder}
\end{align}
If $h\in\R$, $s+h\in I$, and $s+h\ne\eta$, then
\begin{align}
    &|\Delta_{s,h}\Delta \bm Z | \le C\norm{\bm Z' }_{\ci{0}} \abs{h}, \\
    &\abs{\Delta_{s,h} \paren{\bm Z'(s) - \frac{\Delta \bm Z}{s-\eta}} }   \le C [\bm Z ']_{\ci{0,\alpha}}\paren{\abs{h}^\alpha+\abs{h}\abs{s-\eta}^{\alpha-1}},\\
    &\abs{\Delta_{s,h} \paren{\bm Z'(\eta) - \frac{\Delta \bm Z}{s-\eta}} }  \le C [\bm Z ']_{\ci{0,\alpha}}\abs{h}\abs{s-\eta}^{\alpha-1}.
\end{align}
In particular, if $\abs{h}\le \frac{1}{2}\abs{s-\eta}$, then
\begin{equation}\label{eqn:diff-est1}
    \abs{\Delta_{s,h} \paren{\bm Z'(s) - \frac{\Delta \bm Z}{s-\eta}} }, \abs{\Delta_{s,h} \paren{\bm Z'(\eta) - \frac{\Delta \bm Z}{s-\eta}} }  \le C [\bm Z ']_{\ci{0,\alpha}} \abs{h}^\alpha ,
\end{equation}
\end{lemma}
\begin{proof}
The fundamental theorem of calculus gives
\begin{equation}
    \abs{\Delta\bm Z} = \abs{\paren{s-\eta} \int_0^1 \bm Z'(\eta + \rho(s-\eta))\,d\rho  } \le \norm{\bm Z'}_{\ci{0}} \abs{s-\eta},
\end{equation}
\begin{equation}
    \abs{\frac{\Delta \bm Z}{s-\eta} - \bm Z'(\eta)} = \abs{ \int_0^1 \bm Z'\paren{\eta + \rho(s-\eta)} - \bm Z'(\eta)\,d\rho } \le C [\bm Z']_{\ci{0,\alpha}} \abs{s-\eta}^\alpha .
\end{equation}
\begin{equation}
    \abs{\Delta_{s,h} \Delta \bm Z} = \abs{ h \int_0^1 \bm Z'(s+\rho h) \,d\rho } \le \norm{\bm Z'}_{\ci{0}} \abs{h}.
\end{equation}
First suppose that $\abs h\le\frac12\abs{s-\eta}$. Since
\begin{equation}
    \p_s\paren{\frac{\Delta\bm Z}{s-\eta}} = \frac{1}{s-\eta}\paren{\bm Z'(s) - \frac{\Delta\bm Z}{s-\eta}},
\end{equation}
we have
\begin{equation}
\begin{aligned}
    \abs{\Delta_{s,h}\paren{\frac{\Delta\bm Z}{s-\eta}}} &= \abs{h \int_0^1\p_s\paren{\frac{\Delta\bm Z}{s-\eta}} (s+\rho h, \eta) \,d\rho } \le C[\bm Z']_{\ci{0,\alpha}}\abs{h}\abs{s-\eta}^{\alpha-1} ,
\end{aligned}
\end{equation}
Hence
\begin{equation}
\begin{aligned}
    \abs{\Delta_{s,h} \paren{\bm Z'(s) - \frac{\Delta\bm Z}{s-\eta}} } \le \abs{\Delta_{s,h}\bm Z'(s)} + \abs{\Delta_{s,h}\paren{\frac{\Delta\bm Z}{s-\eta}}} \\
    \le C [\bm Z']_{\ci{0,\alpha}}\paren{\abs{h}^\alpha+\abs{h} \abs{s-\eta}^{\alpha-1}}.
\end{aligned}
\end{equation}
\begin{equation}
    \abs{\Delta_{s,h} \paren{\bm Z'(\eta) - \frac{\Delta\bm Z}{s-\eta}} } \le\abs{\Delta_{s,h}\paren{\frac{\Delta\bm Z}{s-\eta}}} \le C [\bm Z']_{\ci{0,\alpha}}\abs{h} \abs{s-\eta}^{\alpha-1}.
\end{equation}
This proves the three finite-difference estimates in this regime, including
\eqref{eqn:diff-est1}. If
$\abs h>\frac12\abs{s-\eta}$, apply~\eqref{eqn:divided-difference-remainder}
at $(s,\eta)$ and $(s+h,\eta)$. Since
$\abs{s-\eta}<2\abs h$ and
$\abs{s+h-\eta}\le3\abs h$, both finite differences are bounded by
$C[\bm Z']_{\ci{0,\alpha}}\abs h^\alpha$. In this regime,
$\abs h^\alpha\le C\abs h\abs{s-\eta}^{\alpha-1}$, which proves the
stated general estimates.
\end{proof}

% \begin{lemma}\label{lem:knl-est}
% Let $\bm Z \in \ci{1,\alpha}$ for $\delta\in(0,1)$ with $\starnorm{\bm Z }>0$. We define
% \begin{align}
%     K_1(s,\eta) &= \frac{\Delta \bm Z  \cdot \paren{\bm Z '(t) - \frac{\Delta \bm Z }{s-\eta}}}{|\Delta \bm Z |^2}, \label{eqn:K1-def}\\
%     Q_1(s,\eta) &= \frac{\Delta \bm Z \cdot\paren{ \bm Z '(s) - \frac{\Delta \bm Z }{s-\eta}}}{|\Delta \bm Z |^2}. \label{eqn:Q1-def}
% \end{align}
% Then, it follows that
% \begin{align}
%     &|K_1(s,\eta)|, |Q_1(s,\eta)| \le \frac{[\bm Z ']_{\ci{0,\alpha}}}{\starnorm{\bm Z }} |s-\eta|^{\delta-1}, \\
%     &|\Delta_{s,h}K_1(s,\eta)|, |\Delta_{s,h}Q_1(s,\eta)|\le C\frac{\norm{\bm Z }_{\ci{1,\alpha}}^2}{\starnorm{\bm Z }^2} \paren{ |s+h-t|^{-1}|h|^\alpha +  |h| |s+h-t|^{\delta-2} },
%     % \\
%     % &|\Delta_{s,h} r(s,\eta)| \le
%     % \begin{cases}
%     %     C |s-\eta| & s\in [t-2h, t+h]\\
%     %     C|h|   & s \in \T\setminus [t-2h,t+h].
%     % \end{cases}
% \end{align}
% \end{lemma}

% \begin{proof}
% \begin{equation}
%     \begin{aligned}
%         &\Delta_{s,h} K_1(s,\eta) \\
%         &= \frac{T_{s,h}\Delta \bm Z  \cdot \paren{\bm Z '(t) - \frac{T_{s,h}\Delta \bm Z }{s+h-t}}}{|T_{s,h}\Delta \bm Z |^2} - \frac{\Delta \bm Z  \cdot \paren{\bm Z '(t) - \frac{\Delta \bm Z }{s-\eta}}}{|\Delta \bm Z |^2} \\
%         &= \frac{T_{s,h}\Delta \bm Z  \cdot \paren{\bm Z '(t) - \frac{T_{s,h}\Delta \bm Z }{s+h-t}} |\Delta \bm Z |^2- \Delta \bm Z  \cdot \paren{\bm Z '(t) - \frac{\Delta \bm Z }{s-\eta}}|T_{s,h}\Delta \bm Z |^2}{|\Delta \bm Z |^2|T_{s,h}\Delta \bm Z |^2} \\
%         &= \frac{A_1 + A_2 + A_3 + A_4}{|\Delta \bm Z |^2|T_{s,h}\Delta \bm Z |^2},
%     \end{aligned}
% \end{equation}
% where
% \begin{align}
%     &A_1 =T_{s,h}\Delta \bm Z  \cdot \paren{\bm Z '(t) - \frac{T_{s,h}\Delta \bm Z }{s+h-t}} |\Delta \bm Z |^2 - T_{s,h}\Delta \bm Z  \cdot \paren{\bm Z '(t) - \frac{\Delta \bm Z }{s-\eta}} |\Delta \bm Z |^2 \\
%     &A_2 = T_{s,h}\Delta \bm Z  \cdot \paren{\bm Z '(t) - \frac{\Delta \bm Z }{s-\eta}} |\Delta \bm Z |^2 - \Delta \bm Z  \cdot \paren{\bm Z '(t) - \frac{\Delta \bm Z }{s-\eta}} |\Delta \bm Z |^2 \\
%     &A_3 = \Delta \bm Z  \cdot \paren{\bm Z '(t) - \frac{\Delta \bm Z }{s-\eta}} |\Delta \bm Z |^2 - \Delta \bm Z  \cdot \paren{\bm Z '(t) - \frac{\Delta \bm Z }{s-\eta}} (\Delta \bm Z \cdot T_{s,h}\Delta \bm Z ) \\
%     &A_4 = \Delta \bm Z  \cdot \paren{\bm Z '(t) - \frac{\Delta \bm Z }{s-\eta}} (\Delta \bm Z \cdot T_{s,h}\Delta \bm Z ) - \Delta \bm Z  \cdot \paren{\bm Z '(t) - \frac{\Delta \bm Z }{s-\eta}}|T_{s,h}\Delta \bm Z |^2.
% \end{align}
% Applying Lemma~\ref{lem:diff-est}, they satisfy the following bounds,
% \begin{align}
%     &|A_1|\le |\Delta \bm Z |^2|T_{s,h}\Delta \bm Z | [\bm Z ']_{\ci{0,\alpha}} |h|^\alpha, \\
%     &|A_2| ,|A_3|, |A_4|\le |\Delta \bm Z |^2 [\bm Z ']_{\ci{0,\alpha}}^2|h| |s-\eta|^\alpha.
% \end{align}
% We also have,
% \begin{equation}
%     \begin{aligned}
%         &\Delta_{s,h} Q_1(s,\eta) \\
%         &= \frac{T_{s,h}\Delta \bm Z  \cdot \paren{T_{s,h}\bm Z '(s) - \frac{T_{s,h}\Delta \bm Z }{s+h-t}}}{|T_{s,h}\Delta \bm Z |^2} - \frac{\Delta \bm Z  \cdot \paren{\bm Z '(s) - \frac{\Delta \bm Z }{s-\eta}}}{|\Delta \bm Z |^2} \\
%         &= \frac{T_{s,h}\Delta \bm Z  \cdot \paren{T_{s,h}\bm Z '(s) - \frac{T_{s,h}\Delta \bm Z }{s+h-t}} |\Delta \bm Z |^2- \Delta \bm Z  \cdot \paren{\bm Z '(s) - \frac{\Delta \bm Z }{s-\eta}}|T_{s,h}\Delta \bm Z |^2}{|\Delta \bm Z |^2|T_{s,h}\Delta \bm Z |^2} \\
%         &= \frac{B_1 + B_2 + B_3 + B_4}{|\Delta \bm Z |^2|T_{s,h}\Delta \bm Z |^2}
%     \end{aligned}
% \end{equation}
% where
% \begin{align}
%     B_1 &=T_{s,h}\Delta \bm Z  \cdot \paren{T_{s,h}\bm Z '(s) - \frac{T_{s,h}\Delta \bm Z }{s+h-t}} |\Delta \bm Z |^2 - T_{s,h}\Delta \bm Z  \cdot \paren{\bm Z '(s) - \frac{\Delta \bm Z }{s-\eta}} |\Delta \bm Z |^2 \\
%     B_2 &= T_{s,h}\Delta \bm Z  \cdot \paren{\bm Z '(s) - \frac{\Delta \bm Z }{s-\eta}} |\Delta \bm Z |^2 - \Delta \bm Z  \cdot \paren{\bm Z '(s) - \frac{\Delta \bm Z }{s-\eta}} |\Delta \bm Z |^2 \\
%     B_3  &= \Delta \bm Z  \cdot \paren{\bm Z '(s) - \frac{\Delta \bm Z }{s-\eta}} |\Delta \bm Z |^2 - \Delta \bm Z  \cdot \paren{\bm Z '(s) - \frac{\Delta \bm Z }{s-\eta}} (\Delta \bm Z \cdot T_{s,h}\Delta \bm Z ) \\
%     B_4  &= \Delta \bm Z  \cdot \paren{\bm Z '(s) - \frac{\Delta \bm Z }{s-\eta}} (\Delta \bm Z \cdot T_{s,h}\Delta \bm Z ) - \Delta \bm Z  \cdot \paren{\bm Z '(s) - \frac{\Delta \bm Z }{s-\eta}}|T_{s,h}\Delta \bm Z |^2.
% \end{align}
% Similarly, we have
% \begin{align}
%     &|B_1|\le 2|\Delta \bm Z |^2|T_{s,h}\Delta \bm Z | \norm{\bm Z }_{\ci{1,\alpha}} |h|^\alpha, \\
%     &|B_2| ,|B_3|, |B_4|\le |\Delta \bm Z |^2 \norm{\bm Z }_{\ci{1,\alpha}}^2|h| |s-\eta|^\alpha.
% \end{align}
% Applying Lemma~\ref{lem:diff-est}, we obtain the results.

% \end{proof}

\begin{lemma}\label{lem-phi-ZV}
Let $V, \bm Z =(Z_1,Z_2)^T\in \ci{1,\alpha}$ for $\alpha\in(0,1)$ with
$\abs{\bm Z}_*>0$. For $s,\eta\in I$ with $s\ne\eta$, define
\begin{equation}
    \phi(s,\eta) = \frac{\Delta V}{\abs{\Delta\bm Z}}.
\end{equation}
% For $s,\eta\in I$, it follows that,
% \begin{align}
%     &\abs{\phi}\le 1,\\
%     &\abs{\p_s\phi}, \abs{\p_\eta\phi} \le C \frac{[V']_{\ci{0,\alpha}}\norm{\bm Z'}_{\ci{0}} + [\bm Z']_{\ci{0,\alpha}}\norm{V'}_{\ci{0}}}{\abs{\bm Z}_*^2}  \abs{s-\eta}^{\alpha-1},
% \end{align}
% For $s,\eta\in I$, $h\in\R$ with $s+h\in I$ and $\abs{s-\eta}\ge 2\abs{h}$,
% \begin{align}
%     &\abs{\Delta_{s,h}\phi} \le C  \frac{[V']_{\ci{0,\alpha}}\norm{\bm Z'}_{\ci{0}} + [\bm Z']_{\ci{0,\alpha}}\norm{V'}_{\ci{0}}}{\abs{\bm Z}_*^2} \abs{h} \abs{s-\eta}^{\alpha},\\
%     &\abs{\Delta_{s,h}\p_s\phi}, \abs{\Delta_{s,h}\p_\eta\phi} \le C  \frac{[V']_{\ci{0,\alpha}}\norm{\bm Z'}_{\ci{0}} + [\bm Z']_{\ci{0,\alpha}}\norm{V'}_{\ci{0}}}{\abs{\bm Z}_*^3} \norm{\bm Z'}_{\ci{0}}\abs{h}^\alpha\abs{s-\eta}^{-1} .
% \end{align}
Then
\begin{align}
    &\abs{\phi}\le C\abs{\bm Z}_*^{-1}\norm{V'}_{\ci{0,\alpha}},\label{eqn:phi-est-1}\\
    &\abs{\p_s\phi}, \abs{\p_\eta\phi} \le C \abs{\bm Z}_*^{-2}\norm{V'}_{\ci{0,\alpha}}\norm{\bm Z'}_{\ci{0,\alpha}}  \abs{s-\eta}^{\alpha-1}.\label{eqn:phi-est-2}
\end{align}
If $s,\eta,s+h\in I$, $h\in\R$, and
$\abs{s-\eta}\ge2\abs{h}$, then
\begin{align}
    &\abs{\Delta_{s,h}\phi} \le C  \abs{\bm Z}_*^{-2}\norm{V'}_{\ci{0,\alpha}}\norm{\bm Z'}_{\ci{0,\alpha}} \abs{h} \abs{s-\eta}^{\alpha-1},\label{eqn:phi-est-3}\\
    &\abs{\Delta_{s,h}\p_s\phi}, \abs{\Delta_{s,h}\p_\eta\phi} \le C\abs{\bm Z}_*^{-3}\norm{V'}_{\ci{0,\alpha}}\norm{\bm Z'}_{\ci{0,\alpha}}^2\abs{h}^\alpha\abs{s-\eta}^{-1} .\label{eqn:phi-est-4}
\end{align}
\end{lemma}

\begin{proof}
Estimate~\eqref{eqn:phi-est-1} follows directly from~\eqref{eqn:arc-chord-intro}. Differentiating with respect to $s$ gives
\begin{align}
    \p_s\phi(s,\eta) &= \frac{V'(s)}{\abs{\Delta\bm Z}} - \frac{\Delta V\paren{\Delta \bm Z\cdot\bm Z'(s)}}{\abs{\Delta\bm Z}^3} \coloneqq A + B,
\end{align}
where
\begin{align}
    A  = \frac{1}{\abs{\Delta\bm Z}} \paren{V'(s) - \frac{\Delta V}{s-\eta}} ,\quad
    B = \frac{ \Delta V\Delta\bm Z}{\abs{\Delta\bm Z}^3} \cdot\paren{\frac{\Delta \bm Z}{s-\eta} - \bm Z'(s)}.
\end{align}
By Lemma~\ref{lem:diff-est},
\begin{equation}
    \abs{A} \le \abs{\bm Z}_*^{-1}[V']_{\ci{0,\alpha}} \abs{s-\eta}^{\alpha-1},\quad \abs{B} \le \abs{\bm Z}_*^{-2}\norm{V'}_{\ci{0}} [\bm Z']_{\ci{0,\alpha}} \abs{s-\eta}^{\alpha-1}.
\end{equation}
This proves \eqref{eqn:phi-est-2}. Integrating
\eqref{eqn:phi-est-2} along the segment $s+\rho h$, $0\le\rho\le1$,
gives \eqref{eqn:phi-est-3}:
\begin{equation}
\begin{aligned}
    \abs{\Delta_{s,h}\phi} &= \abs{ h\int_0^1 \p_s \phi(s+\rho h, \eta)\,d\rho }  \\
    &\le C \abs{\bm Z}_*^{-2}\paren{[V']_{\ci{0,\alpha}}\norm{\bm Z'}_{\ci{0}} + [\bm
    Z']_{\ci{0,\alpha}}\norm{V'}_{\ci{0}}} \abs{h} \abs{s-\eta}^{\alpha-1} .
\end{aligned}
\end{equation}
To estimate the finite differences of $\p_s\phi$, expand
\begin{align}
    \Delta_{s,h} A &= \frac{1}{\abs{T_{s,h}\Delta\bm Z}} \Delta_{s,h}\paren{V'(s) - \frac{\Delta V}{s-\eta}} + \paren{V'(s) - \frac{\Delta V}{s-\eta}}\Delta_{s,h}\paren{\frac{1}{\abs{\Delta\bm Z}} },\\
    \Delta_{s,h} B &= T_{s,h}\paren{\frac{ \Delta V\Delta\bm Z}{\abs{\Delta\bm Z}^3}} \cdot\Delta_{s,h}\paren{\frac{\Delta \bm Z}{s-\eta} - \bm Z'(s)} + \Delta_{s,h}\paren{\frac{ \Delta V\Delta\bm Z}{\abs{\Delta\bm Z}^3}} \cdot \paren{\frac{\Delta \bm Z}{s-\eta} - \bm Z'(s)}
\end{align}
For $\abs{h} \le \frac{1}{2}\abs{s-\eta}$, we have
% \begin{align}
%     \abs{V'(s) - \frac{\Delta V}{s-\eta}}\le C[V']_{\ci{0,\alpha}}\abs{s-\eta}^\alpha,\\ \abs{\frac{\Delta \bm Z}{s-\eta} - \bm Z'(s)}\le C [\bm Z']_{\ci{0,\alpha}} \abs{s-\eta}^\alpha,\\
%     \abs{ \Delta_{s,h}\paren{V'(s) - \frac{\Delta V}{s-\eta}} }\le C[V']_{\ci{0,\alpha}}\abs{h}^\alpha,\\
%     \abs{\Delta_{s,h}\paren{\frac{\Delta \bm Z}{s-\eta} - \bm Z'(s)}}\le C [\bm Z']_{\ci{0,\alpha}} \abs{h}^\alpha .
% \end{align}
\begin{align}
    &\frac{1}{\abs{T_{s,h}\Delta\bm Z}} \le \abs{\bm Z}_*^{-1}\abs{s+h-\eta}^{-1} \le C\abs{\bm Z}_*^{-1} \abs{s-\eta}^{-1},\\
    &\abs{T_{s,h}\paren{\frac{ \Delta V\Delta\bm Z}{\abs{\Delta\bm Z}^3}}} \le C\abs{\bm Z}_*^{-2}\norm{V'}_{\ci{0}} \abs{s-\eta}^{-1}.
\end{align}
The remaining factors satisfy
\begin{align}
    &\abs{\partial_s\paren{\frac{1}{\abs{\Delta \bm Z}}} }=\abs{ \frac{\Delta \bm Z\cdot\p_s\bm
    Z}{\abs{\Delta\bm Z}^3} }\le \abs{\bm Z}_*^{-2}\norm{\bm Z'}_{\ci{0}}\abs{s-\eta}^{-2}, \\
    &\p_s\paren{\frac{ \Delta V\Delta\bm Z}{\abs{\Delta\bm Z}^3}} = V'(s) \frac{\Delta\bm Z}{\abs{\Delta\bm
    Z}^3} + \Delta V \paren{\frac{\bm Z'(s)}{\abs{\Delta\bm Z}^3} -3 \Delta\bm Z \frac{\Delta\bm Z\cdot\bm
    Z'(s)}{\abs{\Delta\bm Z}^5} }, \\
    &\abs{\p_s\paren{\frac{ \Delta V\Delta\bm Z}{\abs{\Delta\bm Z}^3}} } \le
    C\paren{\frac{\abs{V'(s)}}{\abs{\Delta\bm Z}^2} + \frac{\abs{\Delta V}\abs{\bm Z'(s)}}{\abs{\Delta\bm Z}^3}
    + \frac{\abs{\bm Z'}\abs{\Delta V}}{\abs{\Delta\bm Z}^3} }\notag \\
    &\quad \le C\abs{\bm Z}_*^{-3}\norm{V'}_{\ci{0}}\norm{\bm Z'}_{\ci{0}}\abs{s-\eta}^{-2}.
\end{align}
Hence
\begin{equation}
\begin{aligned}
    &\abs{\Delta_{s,h}\paren{\frac{1}{\abs{\Delta\bm Z}}}} = \abs{ h \int_0^1 \partial_s \paren{\frac{1}{\abs{\Delta \bm Z}}}(s+\rho h,\eta) \,d\rho }\\
    &\le \abs{h}\abs{\bm Z}_*^{-2}\norm{\bm Z'}_{\ci{0}}
    \int_0^1 \abs{s+\rho h-\eta}^{-2} \,d\rho  \le C \abs{\bm Z}_*^{-2}\norm{\bm Z'}_{\ci{0}} \abs{h}\abs{s-\eta}^{-2},
\end{aligned}
\end{equation}
and
\begin{equation}
\begin{aligned}
    \abs{\Delta_{s,h}\paren{\frac{ \Delta V\Delta\bm Z}{\abs{\Delta\bm Z}^3}}} &= \abs{h\int_0^1
    \p_s\paren{\frac{ \Delta V\Delta\bm Z}{\abs{\Delta\bm Z}^3}} (s+\rho h,\eta) \,d\rho }  \\
    &\le C\abs{\bm Z}_*^{-3}\norm{V'}_{\ci{0}}\norm{\bm Z'}_{\ci{0}}\abs{h}\abs{s-\eta}^{-2}.
\end{aligned}
\end{equation}
Combining these estimates with Lemma~\ref{lem:diff-est} gives
\begin{equation}
\begin{aligned}
    \abs{\Delta_{s,h}A} &\le C\abs{\bm Z}_*^{-2}\norm{\bm Z'}_{\ci{0}}[V']_{\ci{0,\alpha}} \paren{ \abs{h}^\alpha\abs{s-\eta}^{-1} + \abs{h}\abs{s-\eta}^{\alpha-2} } \\
    &\le C\abs{\bm Z}_*^{-2}\norm{\bm Z'}_{\ci{0}}[V']_{\ci{0,\alpha}}  \abs{h}^\alpha\abs{s-\eta}^{-1},
\end{aligned}
\end{equation}
and
\begin{equation}
\begin{aligned}
    \abs{\Delta_{s,h}B} \le C \abs{\bm Z}_*^{-3}\norm{V'}_{\ci{0}}\norm{\bm Z'}_{\ci{0}}[\bm Z']_{\ci{0,\alpha}} \abs{h}^\alpha  \abs{s-\eta}^{-1}  .
\end{aligned}
\end{equation}
This is \eqref{eqn:phi-est-4} for $\p_s\phi$.

The derivative with respect to $\eta$ has the analogous form
\begin{align}
    \p_\eta\phi(s,\eta) &= -\frac{V'(\eta)}{\abs{\Delta\bm Z}} + \frac{\Delta V\paren{\Delta \bm Z\cdot\bm Z'(\eta)}}{\abs{\Delta\bm Z}^3}.
\end{align}
% and
% \begin{equation}
%     \abs{V'(t) - \frac{\Delta V}{s-\eta}},\abs{\frac{\Delta \bm Z}{s-\eta} - \bm Z'(t)}\le C [\bm Z']_{\ci{0,\alpha}} \abs{s-\eta}^\alpha,
% \end{equation}
% \begin{equation}
%     \abs{ \Delta_{s,h}\paren{V'(t) - \frac{\Delta V}{s-\eta}} },\abs{\Delta_{s,h}\paren{\frac{\Delta \bm Z}{s-\eta} - \bm Z'(t)}}\le C [\bm Z']_{\ci{0,\alpha}} \abs{h}^\alpha .
% \end{equation}
The estimates for $\p_\eta\phi$ follow from the same decomposition,
with $s$ and $\eta$ interchanged.
\end{proof}

If $V$ is a component of $\bm Z$, then
$\abs{\Delta V}\le\abs{\Delta\bm Z}$. This improves the arc-chord dependence
in Lemma~\ref{lem-phi-ZV}.
\begin{lemma}\label{lem-phi-Z}
Let $\bm Z = (V, W)^T\in \ci{1,\alpha}$ for $\alpha\in(0,1)$ with
$\abs{\bm Z}_*>0$. For $s,\eta\in I$ with $s\ne\eta$, define
\begin{equation}
    \phi(s,\eta) = \frac{\Delta V}{\abs{\Delta\bm Z}}.
\end{equation}
Then
\begin{align}
    &\abs{\phi}\le 1,\\
    &\abs{\p_s\phi}, \abs{\p_\eta\phi} \le C \abs{\bm Z}_*^{-1}[\bm Z']_{\ci{0,\alpha}}  \abs{s-\eta}^{\alpha-1},
\end{align}
If $s,\eta,s+h\in I$ and $\abs{s-\eta}\ge2\abs h$, then
\begin{align}
    &\abs{\Delta_{s,h}\phi} \le C \abs{\bm Z}_*^{-1}[\bm Z']_{\ci{0,\alpha}} \abs{h} \abs{s-\eta}^{\alpha-1},\\
    &\abs{\Delta_{s,h}\p_s\phi}, \abs{\Delta_{s,h}\p_\eta\phi} \le C \abs{\bm Z}_*^{-2}\norm{\bm Z'}_{\ci{0}}[\bm Z']_{\ci{0,\alpha}} \abs{h}^\alpha\abs{s-\eta}^{-1} .
\end{align}
\end{lemma}
\begin{proof}
Refine the proof of Lemma~\ref{lem-phi-ZV} using the same decomposition as
in \eqref{eqn:phi-est-2}:
\begin{equation}
    \p_s\phi
    =
    \frac{1}{\abs{\Delta\bm Z}}
    \paren{V'(s)-\frac{\Delta V}{s-\eta}}
    +
    \frac{\Delta V\Delta\bm Z}{\abs{\Delta\bm Z}^3}
    \cdot
    \paren{\frac{\Delta\bm Z}{s-\eta}-\bm Z'(s)}.
\end{equation}
Since $V$ is a component of $\bm Z$,
$\abs{\Delta V}\le \abs{\Delta\bm Z}$.  Combining this improvement with
\eqref{eqn:divided-difference-remainder} gives
\begin{equation}
    \abs{\p_s\phi}
    \le C\abs{\bm Z}_*^{-1}[\bm Z']_{\ci{0,\alpha}}
    \abs{s-\eta}^{\alpha-1}.
\end{equation}
The formula for $\p_\eta\phi$ is the same with $s$ and $\eta$
interchanged, so the same bound holds for $\p_\eta\phi$.  The estimate
$\abs{\phi}\le 1$ follows directly from
$\abs{\Delta V}\le\abs{\Delta\bm Z}$.

For the first difference estimate, integrate the $\p_s\phi$-bound
along the segment $s+\rho h$, $0\le\rho\le1$.  Because
$\abs{s+\rho h-\eta}\ge \frac12\abs{s-\eta}$ under
$\abs{s-\eta}\ge2\abs h$,
\begin{equation}
    \abs{\Delta_{s,h}\phi}
    \le |h|\int_0^1 \abs{\p_s\phi(s+\rho h,\eta)}\,d\rho
    \le C\abs{\bm Z}_*^{-1}[\bm Z']_{\ci{0,\alpha}}
    |h|\abs{s-\eta}^{\alpha-1}.
\end{equation}
For $\Delta_{s,h}\p_s\phi$, use the decompositions of
$\Delta_{s,h}A$ and $\Delta_{s,h}B$ in the proof of
Lemma~\ref{lem-phi-ZV}. The component estimate
$\abs{\Delta V}\le\abs{\Delta\bm Z}$ again lowers the power of
$\abs{\bm Z}_*^{-1}$. By \eqref{eqn:diff-est1},
\begin{equation}
    \abs{\Delta_{s,h}\p_s\phi}
    \le C\abs{\bm Z}_*^{-2}\norm{\bm Z'}_{\ci{0}}[\bm Z']_{\ci{0,\alpha}}|h|^\alpha\abs{s-\eta}^{-1}.
\end{equation}
The same argument, with the formula for $\p_\eta\phi$, gives the last
estimate.
\end{proof}
% \begin{proof}
% First, we have
% \begin{align}
%     \p_s\phi(s,\eta) &= \frac{V'(s)}{\abs{\Delta\bm Z}} - \frac{\Delta V\paren{\Delta \bm Z\cdot\bm Z'(s)}}{\abs{\Delta\bm Z}^3} = A + B,
% \end{align}
% where
% \begin{align}
%     A  = \frac{1}{\abs{\Delta\bm Z}} \paren{V'(s) - \frac{\Delta V}{s-\eta}} ,\quad
%     B = \frac{ \Delta V\Delta\bm Z}{\abs{\Delta\bm Z}^3} \cdot\paren{\frac{\Delta \bm Z}{s-\eta} - \bm Z'(s)}.
% \end{align}
% We have
% \begin{equation}
%     \abs{A} \le \frac{ [\bm Z']_{\ci{0,\alpha}} }{\abs{\bm Z}_*} \abs{s-\eta}^{\alpha-1},\quad \abs{B} \le \frac{\abs{\Delta{\bm Z}}^2}{\abs{\Delta\bm Z}^3}  [\bm Z'] \abs{s-\eta}^\alpha \le \frac{ [\bm Z']_{\ci{0,\alpha}} }{\abs{\bm Z}_*} \abs{s-\eta}^{\alpha-1}.
% \end{equation}

% \begin{equation}
% \begin{aligned}
%     \abs{\Delta_{s,h}\phi} = \abs{ h\int_0^1 \p_s \phi(s+\eta h, t)\,d\eta } \le C \frac{[\bm Z']_{\ci{0,\alpha}}}{\abs{\bm Z}_*} \abs{h} \abs{s-\eta}^{\alpha-1} .
% \end{aligned}
% \end{equation}

% \begin{align}
%     \Delta_{s,h} A &= \frac{1}{\abs{T_{s,h}\Delta\bm Z}} \Delta_{s,h}\paren{V'(s) - \frac{\Delta V}{s-\eta}} + \paren{V'(s) - \frac{\Delta V}{s-\eta}}\Delta_{s,h}\paren{\frac{1}{\abs{\Delta\bm Z}} },\\
%     \Delta_{s,h} B &= T_{s,h}\paren{\frac{ \Delta V\Delta\bm Z}{\abs{\Delta\bm Z}^3}} \cdot\Delta_{s,h}\paren{\frac{\Delta \bm Z}{s-\eta} - \bm Z'(s)} + \Delta_{s,h}\paren{\frac{ \Delta V\Delta\bm Z}{\abs{\Delta\bm Z}^3}} \cdot \paren{\frac{\Delta \bm Z}{s-\eta} - \bm Z'(s)}
% \end{align}
% For $\abs{h} \le \frac{1}{2}\abs{s-\eta}$, we have
% \begin{equation}
%     \abs{V'(s) - \frac{\Delta V}{s-\eta}},\abs{\frac{\Delta \bm Z}{s-\eta} - \bm Z'(s)}\le C [\bm Z']_{\ci{0,\alpha}} \abs{s-\eta}^\alpha,
% \end{equation}
% \begin{equation}
%     \abs{ \Delta_{s,h}\paren{V'(s) - \frac{\Delta V}{s-\eta}} },\abs{\Delta_{s,h}\paren{\frac{\Delta \bm Z}{s-\eta} - \bm Z'(s)}}\le C [\bm Z']_{\ci{0,\alpha}} \abs{h}^\alpha .
% \end{equation}

% \begin{equation}
%     \frac{1}{\abs{T_{s,h}\Delta\bm Z}} \le \frac{1}{\abs{\bm Z}_*}\abs{s+h-t}^{-1} \le \frac{C}{\abs{\bm Z}_*} \abs{s-\eta}^{-1}.
% \end{equation}
% \begin{equation}
%     \abs{\partial_s\paren{\frac{1}{\abs{\Delta \bm Z}}} }=\abs{ \frac{\Delta \bm Z\cdot\p_s\bm Z}{\abs{\Delta\bm Z}^3} }\le \frac{\norm{\bm Z'}_{\ci{0}}}{\abs{\bm Z}_*^2}\abs{s-\eta}^{-2}
% \end{equation}

% \begin{equation}
% \begin{aligned}
%     &\abs{\Delta_{s,h}\paren{\frac{1}{\abs{\Delta\bm Z}}}} = \abs{ h \int_0^1 \partial_s \paren{\frac{1}{\Delta \bm Z}} \paren{s + \eta h}  \,d\eta }\\
%     \le \abs{h}\frac{\norm{\bm Z'}_{\ci{0}}}{\abs{\bm Z}_*^2}
%     \int_0^1 \abs{s+\eta h-t}^{-2} \,d\eta  \le C \frac{\norm{\bm Z'}_{\ci{0}}}{\abs{\bm Z}_*^2} \abs{h}\abs{s-\eta}^{-2}.
% \end{aligned}
% \end{equation}

% \begin{equation}
% \begin{aligned}
%     \abs{\Delta_{s,h}A} \le C\frac{\norm{\bm Z'}_{\ci{0}}[\bm Z']_{\ci{0,\alpha}}}{\abs{\bm Z}_*^2} \paren{ \abs{h}^\alpha\abs{s-\eta}^{-1} + \abs{h}\abs{s-\eta}^{\alpha-2} } \\
%     \le C\frac{\norm{\bm Z'}_{\ci{0}}[\bm Z']_{\ci{0,\alpha}}}{\abs{\bm Z}_*^2}  \abs{h}^\alpha\abs{s-\eta}^{-1}
% \end{aligned}
% \end{equation}

% \begin{align}
%     \abs{T_{s,h}\paren{\frac{ \Delta V\Delta\bm Z}{\abs{\Delta\bm Z}^3}}} \le \frac{C}{\abs{\bm Z}_*} \abs{s-\eta}^{-1}.
% \end{align}

% \begin{align}
%     \p_s\paren{\frac{ \Delta V\Delta\bm Z}{\abs{\Delta\bm Z}^3}} = V'(s) \frac{\Delta\bm Z}{\abs{\Delta\bm Z}^3} + \Delta V  \paren{\frac{\bm Z'(s)}{\abs{\Delta\bm Z}^3} -3 \Delta\bm Z \frac{\Delta\bm Z\cdot\bm Z'(s)}{\abs{\Delta\bm Z}^5} }
% \end{align}
% \begin{align}
%     \abs{\p_s\paren{\frac{ \Delta V\Delta\bm Z}{\abs{\Delta\bm Z}^3}} } \le C\paren{\frac{\abs{V'(s)}}{\abs{\Delta\bm Z}^2} + \frac{\abs{\Delta V}\abs{\bm Z'(s)}}{\abs{\Delta\bm Z}^3} +  \frac{\abs{\bm Z'}\abs{\Delta V}}{\abs{\Delta\bm Z}^3} }\le C \frac{\abs{\bm Z'(s)}}{\abs{\Delta\bm Z}^2} \le C\frac{\norm{\bm Z'}_{\ci{0}}}{\abs{\bm Z}_*^2}\abs{s-\eta}^{-2}.
% \end{align}
% \begin{align}
%      \Delta_{s,h}\paren{\frac{ \Delta V\Delta\bm Z}{\abs{\Delta\bm Z}^3}} = \abs{h\int_0^1  \p_s\paren{\frac{ \Delta V\Delta\bm Z}{\abs{\Delta\bm Z}^3}} (s + \eta h) \,d\eta }  \le C\frac{\norm{\bm Z'}_{\ci{0}}}{\abs{\bm Z}_*^2}\abs{h}\abs{s-\eta}^{-2}.
% \end{align}

% \begin{equation}
% \begin{aligned}
%     \abs{\Delta_{s,h}B} \le C \frac{\norm{\bm Z'}_{\ci{0}}[\bm Z']_{\ci{0,\alpha}}}{\abs{\bm Z}_*^2} \paren{  \abs{h}^\alpha  \abs{s-\eta}^{-1} + \abs{h}\abs{s-\eta}^{\alpha-2} } \le C \frac{\norm{\bm Z'}_{\ci{0}}[\bm Z']_{\ci{0,\alpha}}}{\abs{\bm Z}_*^2}\abs{h}^\alpha  \abs{s-\eta}^{-1}  .
% \end{aligned}
% \end{equation}

% Since,
% \begin{align}
%     \p_\eta\phi(s,\eta) &= -\frac{V'(t)}{\abs{\Delta\bm Z}} + \frac{\Delta V\paren{\Delta \bm Z\cdot\bm Z'(t)}}{\abs{\Delta\bm Z}^3}.
% \end{align}
% and
% \begin{equation}
%     \abs{V'(t) - \frac{\Delta V}{s-\eta}},\abs{\frac{\Delta \bm Z}{s-\eta} - \bm Z'(t)}\le C [\bm Z']_{\ci{0,\alpha}} \abs{s-\eta}^\alpha,
% \end{equation}
% \begin{equation}
%     \abs{ \Delta_{s,h}\paren{V'(t) - \frac{\Delta V}{s-\eta}} },\abs{\Delta_{s,h}\paren{\frac{\Delta \bm Z}{s-\eta} - \bm Z'(t)}}\le C [\bm Z']_{\ci{0,\alpha}} \abs{h}^\alpha .
% \end{equation}
% All the estimates above for $\p_\eta\phi$ are the same.

% \end{proof}

\begin{proof}[Proof of Lemma~\ref{LM-kernel-est}]
Differentiating the kernels gives
\begin{equation}
\begin{aligned}
    \p_\eta \ell(s,\eta;\bm Z) = \frac{1}{s-\eta} - \frac{\bm Z'(\eta)\cdot \Delta \bm Z}{|\Delta \bm Z|^2} = \frac{\Delta \bm Z}{\abs{\Delta\bm Z}^2}\cdot\paren{ \frac{\Delta \bm Z}{s-\eta} -\bm Z'(\eta) },
\end{aligned}
\end{equation}
\begin{equation}
\begin{aligned}
    \p_\eta m_{ij}(s,\eta;\bm Z) &= \p_\eta\paren{\frac{\Delta Z_i}{|\Delta\bm Z|}} \frac{\Delta Z_j}{|\Delta\bm Z|} + \frac{\Delta Z_i}{|\Delta\bm Z|}\p_\eta\paren{\frac{\Delta Z_j}{|\Delta\bm Z|}} .
\end{aligned}
\end{equation}
Lemmas~\ref{lem:diff-est} and~\ref{lem-phi-Z} immediately imply
\eqref{eqn:LM-1} and \eqref{eqn:LM-2}. For the finite-difference bounds,
expand
\begin{equation}
\begin{aligned}
    \Delta_{s,h}(\p_\eta \ell(s,\eta;\bm Z))
    &= T_{s,h}\paren{\frac{\Delta \bm Z}{\abs{\Delta\bm Z}^2}}
    \cdot\Delta_{s,h}\paren{
    \frac{\Delta \bm Z}{s-\eta} -\bm Z'(\eta) } \\
    &\quad
    + \Delta_{s,h}\paren{\frac{\Delta \bm Z}{\abs{\Delta\bm Z}^2}}
    \cdot\paren{ \frac{\Delta \bm Z}{s-\eta} -\bm Z'(\eta) }.
\end{aligned}
\end{equation}
and
\begin{equation}
\begin{aligned}
    \Delta_{s,h}(\p_\eta m_{ij}(s,\eta;\bm Z)) &= T_{s,h}\p_\eta\paren{\frac{\Delta Z_i}{|\Delta\bm Z|}} \Delta_{s,h}\paren{ \frac{\Delta Z_j}{|\Delta\bm Z|} } + \frac{\Delta Z_j}{|\Delta\bm Z|}  \Delta_{s,h}\p_\eta\paren{\frac{\Delta Z_i}{|\Delta\bm Z|}} \\
    &+ T_{s,h}\frac{\Delta Z_i}{|\Delta\bm Z|}\Delta_{s,h}\p_\eta\paren{\frac{\Delta Z_j}{|\Delta\bm Z|}} + \p_\eta\paren{\frac{\Delta Z_j}{|\Delta\bm Z|}}\Delta_{s,h}\frac{\Delta Z_i}{|\Delta\bm Z|} .
\end{aligned}
\end{equation}
Moreover,
% \begin{equation}
%     \p_s \paren{\frac{\Delta \bm Z}{\abs{\Delta\bm Z}^2}} = \frac{\bm Z'(s)}{\abs{\Delta\bm Z}^2} - 2 \frac{\Delta\bm Z\paren{\Delta\bm Z\cdot\bm Z'(s)}}{\abs{\Delta\bm Z}^4},
% \end{equation}
\begin{equation}
    \abs{\p_s \paren{\frac{\Delta \bm Z}{\abs{\Delta\bm Z}^2}}} = \abs{\frac{\bm Z'(s)}{\abs{\Delta\bm Z}^2} - 2 \frac{\Delta\bm Z\paren{\Delta\bm Z\cdot\bm Z'(s)}}{\abs{\Delta\bm Z}^4}} \le C \abs{\bm Z}_*^{-2}\norm{\bm Z'}_{\ci{0}}\abs{s-\eta}^{-2},
\end{equation}
and hence
\begin{equation}
    \abs{\Delta_{s,h} \paren{\frac{\Delta \bm Z}{\abs{\Delta\bm Z}^2}}} =\abs{h\int_0^1 \p_s \paren{\frac{\Delta \bm Z}{\abs{\Delta\bm Z}^2}}(s+\rho h,\eta) \,d\rho}\le C \abs{\bm Z}_*^{-2}\norm{\bm Z'}_{\ci{0}}\abs{h}\abs{s-\eta}^{-2}.
\end{equation}
Lemma~\ref{lem-phi-Z}, together with
$\abs h\le\frac12\abs{s-\eta}$, now gives
\begin{equation}
\begin{aligned}
    \abs{\Delta_{s,h}(\p_\eta \ell(s,\eta;\bm Z))} &\le C \abs{\bm Z}_*^{-1}[\bm Z']_{\ci{0,\alpha}} \abs{h}^\alpha\abs{s-\eta}^{-1} + C \abs{\bm Z}_*^{-2}\norm{\bm Z'}_{\ci{0}}[\bm Z']_{\ci{0,\alpha}} \abs{h} \abs{s-\eta}^{\alpha-2} \\
    &\le C \abs{\bm Z}_*^{-2}\norm{\bm Z'}_{\ci{0}}[\bm Z']_{\ci{0,\alpha}}\abs{h}^\alpha\abs{s-\eta}^{-1}.
\end{aligned}
\end{equation}
and
\begin{equation}
\begin{aligned}
    \abs{\Delta_{s,h}(\p_\eta m_{ij}(s,\eta;\bm Z))} &\le C \abs{\bm Z}_*^{-2}[\bm Z']_{\ci{0,\alpha}}^2 \abs{h}\abs{s-\eta}^{2\alpha-2} \\
    &\quad + C\abs{\bm Z}_*^{-2}\norm{\bm Z'}_{\ci{0}}[\bm Z']_{\ci{0,\alpha}} \abs{h}^\alpha\abs{s-\eta}^{-1} \\
    &\le C \abs{\bm Z}_*^{-2}\norm{\bm Z'}_{\ci{0,\alpha}}[\bm Z']_{\ci{0,\alpha}}  \abs{h}^\alpha\abs{s-\eta}^{-1}.
\end{aligned}
\end{equation}
The derivatives with respect to $s$ satisfy
\begin{equation}\label{eqn:ell-s-derivative}
\begin{aligned}
    \p_s \ell(s,\eta;\bm Z) = \frac{\bm Z'(s)\cdot \Delta \bm Z}{|\Delta \bm Z|^2} - \frac{1}{s-\eta} = \frac{\Delta \bm Z}{\abs{\Delta\bm Z}^2}\cdot\paren{ \bm Z'(s) -\frac{\Delta \bm Z}{s-\eta}  }.
\end{aligned}
\end{equation}
\begin{equation}\label{eqn:m-s-derivative}
\begin{aligned}
    \p_s m_{ij}(s,\eta;\bm Z) &= \p_s\paren{\frac{\Delta Z_i}{|\Delta\bm Z|}} \frac{\Delta Z_j}{|\Delta\bm Z|} + \frac{\Delta Z_i}{|\Delta\bm Z|}\p_s\paren{\frac{\Delta Z_j}{|\Delta\bm Z|}} .
\end{aligned}
\end{equation}
The estimates for $\p_s\ell$ and $\p_s m_{ij}$ follow from \eqref{eqn:ell-s-derivative}--\eqref{eqn:m-s-derivative} and the same bounds \eqref{eqn:diff-est1} and \eqref{eqn:phi-est-4}.
\end{proof}

% \begin{lemma}\label{YZ-est}
% \begin{equation}
%     \abs{ \p_s L_{\bm Y}(s,\eta) - \p_s \ell(s,\eta;\bm X) } \le \frac{\norm{\bm X'}_{\ci{0}}^2+\norm{\bm Y'}_{\ci{0}}^2}{\abs{\bm X}_*^2\abs{\bm Y}_*^2} \norm{\bm X'}_{\ci{0,\alpha}} \norm{\bm Y'-\bm X'}_{\ci{0,\alpha}} \abs{s-\eta}^{\alpha-1}.
% \end{equation}
% \begin{equation}
%     \abs{ \p_\eta L_{\bm Y}(s,\eta) - \p_\eta \ell(s,\eta;\bm X) } \le \frac{\norm{\bm X'}_{\ci{0}}^2+\norm{\bm Y'}_{\ci{0}}^2}{\abs{\bm X}_*^2\abs{\bm Y}_*^2} \norm{\bm X'}_{\ci{0,\alpha}} \norm{\bm Y'-\bm X'}_{\ci{0,\alpha}} \abs{s-\eta}^{\alpha-1}.
% \end{equation}
% \end{lemma}
% \begin{proof}
% Let $\bm Z = \bm Y - \bm X$,
% \begin{equation}
%     \p_s L_{\bm Y}(s,\eta) = \frac{\bm Y'(s)\cdot\Delta\bm Y}{\abs{\Delta\bm Y}^2},
% \end{equation}
% \begin{equation}
% \begin{aligned}
%     &\frac{\bm Y'(s)\cdot\Delta\bm Y}{\abs{\Delta\bm Y}^2} - \frac{\bm X'(s)\cdot\Delta\bm X}{\abs{\Delta\bm X}^2} \\
%     &= \frac{1}{\abs{\Delta\bm Y}^2\abs{\Delta\bm X}^2}\paren{\paren{\bm Y'(s)\cdot\Delta\bm Y}\abs{\Delta\bm X}^2 - \abs{\Delta\bm Y}^2\paren{\bm X'(s)\cdot\Delta\bm X}} \\
%     &= \frac{1}{\abs{\Delta\bm Y}^2\abs{\Delta\bm X}^2}\paren{\paren{\bm X'(s)\cdot\Delta\bm Z + \bm Z'(s) \cdot\Delta\bm X+\bm Z'(s)\cdot\Delta\bm Z}\abs{\Delta\bm X}^2 -  \paren{2\Delta\bm Z\cdot\Delta\bm X + \abs{\Delta\bm Z}^2}\paren{\bm X'(s)\cdot\Delta\bm X} } \\
%     &= \frac{1}{\abs{\Delta\bm Y}^2\abs{\Delta\bm X}^2}\sum_{k=1}^5 I_k.
% \end{aligned}
% \end{equation}
% where
% \begin{equation}
% \begin{aligned}
%     I_1 &= \paren{\bm X'(s) - \frac{\Delta\bm X}{s-\eta}}\cdot\Delta\bm Z \abs{\Delta \bm X}^2,\\
%     I_2 &= \paren{\bm Z'(s) - \frac{\Delta\bm Z}{s-\eta}}\cdot\Delta\bm X \abs{\Delta \bm X}^2, \\
%     I_3 &= \paren{\bm Z'(s) - \frac{\Delta\bm Z}{s-\eta}}\cdot\Delta\bm Z \abs{\Delta \bm X}^2, \\
%     I_4 &= -2 \paren{\Delta\bm Z\cdot\Delta\bm X}\paren{\bm X'(s) - \frac{\Delta\bm X}{s-\eta}}\cdot\Delta\bm X ,\\
%     I_5 &= -\abs{\Delta\bm Z}^2\paren{\bm X'(s) - \frac{\Delta\bm X}{s-\eta}}\cdot\Delta\bm X
% \end{aligned}
% \end{equation}
% Hence
% \begin{equation}
%     \abs{ \p_s L_{\bm Y}(s,\eta) - \p_s \ell(s,\eta;\bm X) } \le \frac{\norm{\bm X'}_{\ci{0}}^2+\norm{\bm Y'}_{\ci{0}}^2}{\abs{\bm X}_*^2\abs{\bm Y}_*^2} \norm{\bm X'}_{\ci{0,\alpha}} \norm{\bm Y'-\bm X'}_{\ci{0,\alpha}} \abs{s-\eta}^{\alpha-1}.
% \end{equation}
% where we have used
% \begin{equation}
%     \abs{\Delta\bm Z}\le \abs{\Delta\bm X} + \abs{\Delta\bm Y}, \quad \abs{\bm Z'}\le \abs{\bm X'} + \abs{\bm Y'}.
% \end{equation}
% \end{proof}

\begin{proof}[Proof of Lemma~\ref{lem:LM-Lip}]
We first estimate the directional derivatives of the kernels and then
integrate these estimates along the segment joining $\bm Y$ to $\bm X$.
Let $\bm Z\in C^{1,\alpha}(I)$ be a direction. Differentiation with respect
to the curve commutes with $\p_\eta$:
\begin{equation}
    \p_{\bm X}\paren{\p_\eta \ell(s,\eta;\bm X)}\bm Z
    =\p_\eta\paren{\p_{\bm X}\ell(s,\eta;\bm X)\bm Z}.
\end{equation}
The corresponding derivatives of $\ell$ and $m_{ij}$ are
\begin{align}
    \p_{\bm X}\ell(s,\eta;\bm X)\bm Z
    &=\frac{\Delta\bm X\cdot\Delta\bm Z}
    {\abs{\Delta\bm X}^2},\label{eqn:ell-curve-derivative}\\
    \p_{\bm X}m_{ij}(s,\eta;\bm X)\bm Z
    &=\frac{\Delta Z_i\Delta X_j+\Delta X_i\Delta Z_j}
    {\abs{\Delta\bm X}^2}
    -2\frac{\Delta X_i\Delta X_j
    \paren{\Delta\bm X\cdot\Delta\bm Z}}
    {\abs{\Delta\bm X}^4}.\label{eqn:m-curve-derivative}
\end{align}
Each term has the same quotient structure. For example,
\begin{equation}
    \p_\eta\paren{\frac{\Delta X_i\Delta Z_j}
    {\abs{\Delta\bm X}^2}}
    =\p_\eta\paren{\frac{\Delta X_i}{\abs{\Delta\bm X}}}
    \frac{\Delta Z_j}{\abs{\Delta\bm X}}
    +\p_\eta\paren{\frac{\Delta Z_j}{\abs{\Delta\bm X}}}
    \frac{\Delta X_i}{\abs{\Delta\bm X}}.
\end{equation}
Taking an $s$-increment and using the discrete multiplication rule gives
\begin{equation}
\begin{aligned}
    \Delta_{s,h}\p_\eta\paren{
    \frac{\Delta X_i\Delta Z_j}{\abs{\Delta\bm X}^2}}
    &=T_{s,h}\p_\eta\paren{\frac{\Delta X_i}{\abs{\Delta\bm X}}}
    \Delta_{s,h}\frac{\Delta Z_j}{\abs{\Delta\bm X}}
    +\Delta_{s,h}\p_\eta\paren{\frac{\Delta X_i}{\abs{\Delta\bm X}}}
    \frac{\Delta Z_j}{\abs{\Delta\bm X}}\\
    &\quad+T_{s,h}\p_\eta\paren{\frac{\Delta Z_j}{\abs{\Delta\bm X}}}
    \Delta_{s,h}\frac{\Delta X_i}{\abs{\Delta\bm X}}
    +\Delta_{s,h}\p_\eta\paren{\frac{\Delta Z_j}{\abs{\Delta\bm X}}}
    \frac{\Delta X_i}{\abs{\Delta\bm X}}.
\end{aligned}
\end{equation}
Lemma~\ref{lem-phi-ZV} gives
\begin{equation}
    \abs{\p_\eta\paren{\frac{\Delta X_i\Delta Z_j}
    {\abs{\Delta\bm X}^2}}}
    \le C\abs{\bm X}_*^{-2}\norm{\bm Z'}_{\ci{0,\alpha}}
    \norm{\bm X'}_{\ci{0,\alpha}}
    \abs{s-\eta}^{\alpha-1}.
\end{equation}
If $\abs h\le\frac12\abs{s-\eta}$, the same lemma and the preceding
decomposition give
\begin{equation}
    \abs{\Delta_{s,h}\p_\eta\paren{
    \frac{\Delta X_i\Delta Z_j}{\abs{\Delta\bm X}^2}}}
    \le C\abs{\bm X}_*^{-3}\norm{\bm Z'}_{\ci{0,\alpha}}
    \norm{\bm X'}_{\ci{0,\alpha}}^2
    \abs h^\alpha\abs{s-\eta}^{-1}.
\end{equation}
The same calculation applies to $\p_\eta\ell$ and
$\p_\eta m_{ij}$, using \eqref{eqn:ell-curve-derivative}--\eqref{eqn:m-curve-derivative} and
Lemma~\ref{lem-phi-ZV}. Thus, for $q=\ell,m_{ij}$,
\begin{align}
    \abs{\p_{\bm X}\paren{\p_\eta q(s,\eta;\bm X)}\bm Z}
    &\le C\abs{\bm X}_*^{-2}\norm{\bm Z'}_{\ci{0,\alpha}}
    \norm{\bm X'}_{\ci{0,\alpha}}
    \abs{s-\eta}^{\alpha-1},\label{eqn:kernel-directional-bound}\\
    \abs{\Delta_{s,h}\p_{\bm X}
    \paren{\p_\eta q(s,\eta;\bm X)}\bm Z}
    &\le C\abs{\bm X}_*^{-3}\norm{\bm Z'}_{\ci{0,\alpha}}
    \norm{\bm X'}_{\ci{0,\alpha}}^2
    \abs h^\alpha\abs{s-\eta}^{-1}.\label{eqn:kernel-directional-increment}
\end{align}
% We write
% \begin{equation}
%     (L_{\bm X}f)(s) - (L_{\bm Y}f)(s) = \int_I \p_\eta\paren{\ell(s,\eta;\bm X)-\ell(s,\eta;\bm Y)} f(\eta)\,d\eta
% \end{equation}
% \begin{equation}
%     (M_{\bm X,ij}f)(s) - (M_{\bm Y,ij}f)(s) = \int_I \p_\eta\paren{m_{ij}(s,\eta;\bm X)-m_{ij}(s,\eta;\bm Y)} f(\eta)\,d\eta
% \end{equation}
Integrating~\eqref{eqn:kernel-directional-bound} along the segment from $\bm Y$ to $\bm X$ gives:
\begin{equation}
\begin{aligned}
    \abs{\p_\eta\paren{\ell(s,\eta;\bm X)-\ell(s,\eta;\bm Y)}} &\le\int_0^1 \abs{\p_{\bm X}\p_\eta\ell
    \paren{s,\eta;\rho\bm X+(1-\rho)\bm Y} \paren{\bm X-\bm Y}}\,d\rho \\
    &\le C_{M,m} \norm{\bm X'-\bm Y'}_{\ci{0,\alpha}} \abs{s-\eta}^{\alpha-1}.
\end{aligned}
\end{equation}
The same estimate holds for $m_{ij}$:
\begin{equation}
    \abs{\p_\eta\paren{m_{ij}(s,\eta;\bm X)-m_{ij}(s,\eta;\bm Y)}}
    \le C_{M,m}
    \norm{\bm X'-\bm Y'}_{\ci{0,\alpha}}
    \abs{s-\eta}^{\alpha-1}.
\end{equation}
Integrating~\eqref{eqn:kernel-directional-increment} along the same segment shows that,
for $q=\ell,m_{ij}$ and $\abs h\le\frac12\abs{s-\eta}$,
\begin{equation}
    \abs{\Delta_{s,h}\p_\eta
    \paren{q(s,\eta;\bm X)-q(s,\eta;\bm Y)}}
    \le C_{M,m}
    \norm{\bm X'-\bm Y'}_{\ci{0,\alpha}}
    \abs h^\alpha\abs{s-\eta}^{-1}.
\end{equation}
The corresponding estimates for $\p_s\ell$ and $\p_sm_{ij}$ follow in the same
way. Equation~\eqref{eqn:kernel-mapping-bound} gives the endpoint operator bounds, and the
interpolation argument from Lemma~\ref{LM-op-est} proves
\eqref{eqn:Lop-Lip} and~\eqref{eqn:Mop-Lip}.
\end{proof}


\begin{proof}[Proof of Lemma~\ref{K-op}]
The bound~\eqref{eqn:kernel-pointwise-hypothesis} is integrable because $\alpha>0$. More precisely,
\begin{equation}
    \sup_{s\in I}\int_I\abs{k(s,\eta)}\,d\eta
    +\sup_{\eta\in I}\int_I\abs{k(s,\eta)}\,ds
    \le C a_1.
\end{equation}
The Schur test gives
\begin{equation}
    \norm{K f}_{L^2(I)}
    \le C a_1 \norm{f}_{L^2(\R)}.
\end{equation}
This proves the case $\varepsilon=0$. Suppose now that
$0<\varepsilon<\alpha$. Writing the fractional seminorm in terms of
increments gives
\begin{equation}
\begin{aligned}
    [Kf]_{\hi{\varepsilon}}^2
    &=\int_{-2}^2\frac{1}{\abs h^{1+2\varepsilon}}
    \int_{I\cap(I-h)}
    \abs{Kf(s+h)-Kf(s)}^2\,ds\,dh.
\end{aligned}
\end{equation}
For fixed $h$, another application of the Schur test gives
\begin{equation}
    \int_{I\cap(I-h)}
    \abs{Kf(s+h)-Kf(s)}^2\,ds
    \le M_1(h)M_2(h)\norm{f}_{L^2(\R)}^2,
\end{equation}
where
\begin{align}
    M_1(h)&\coloneqq
    \sup_{s\in I\cap(I-h)}\int_I
    \abs{k(s+h,\eta)-k(s,\eta)}\,d\eta,\\
    M_2(h)&\coloneqq
    \sup_{\eta\in I}\int_{I\cap(I-h)}
    \abs{k(s+h,\eta)-k(s,\eta)}\,ds.
\end{align}

For $\abs h\ge1$, \eqref{eqn:kernel-pointwise-hypothesis} directly gives the estimate below.
It therefore remains to consider $0<\abs h<1$. Split each integral into
$\abs{s-\eta}\le2\abs h$ and $\abs{s-\eta}>2\abs h$. On the near region,
\eqref{eqn:kernel-pointwise-hypothesis} gives $Ca_1\abs h^\alpha$. On the far region,
\eqref{eqn:kernel-increment-hypothesis} gives
\begin{equation}
    Ca_2\abs h^\alpha
    \int_{2\abs h}^{2}\frac{dr}{r}
    \le Ca_2\abs h^\alpha\paren{1+\abs{\log\abs h}}.
\end{equation}
The same estimates apply to both integration variables. Hence
\begin{equation}
    M_1(h)M_2(h)
    \le C\paren{a_1+a_2}^2\abs h^{2\alpha}
    \paren{1+\abs{\log\abs h}^2}.
\end{equation}
Substituting into the seminorm formula gives
\begin{equation}
\begin{aligned}
    [Kf]_{\hi{\varepsilon}}
    &\le C\paren{a_1+a_2}
    \left(\int_{-2}^2
    \abs h^{-1+2(\alpha-\varepsilon)}
    \paren{1+\abs{\log\abs h}^2}\,dh\right)^{\frac12}
    \norm{f}_{L^2(\R)} \\
    &\le C\paren{a_1+a_2}\norm{f}_{L^2(\R)},
\end{aligned}
\end{equation}
because $\varepsilon<\alpha$.
\end{proof}

\begin{proof}[Uniform trace bound in Lemma~\ref{lem:stokes-layer-open-arc}]
Fix $0<\alpha<\beta<\delta$ as in the main-text proof.
To make the constant in~\eqref{eqn:stokes-layer-bound} uniform, observe that the normalized arcs with
$\norm{\kappa}_{\hr{-\frac12+\delta}}\le M$ and $\abs{\kappa}_*\ge m$
form a precompact family in $C^{1,\alpha}$, whose limits remain regular
embedded arcs. Near each limit $\bm X_0$, a fixed closed extension
$\wt{\bm X}_0$ gives nearby extensions
$\wt{\bm X}_0+E(\bm X-\bm X_0)$, where $E$ is a bounded
$C^{1,\alpha}$ extension operator. These remain embedded and have uniform
chart, pullback, and trace bounds in a sufficiently small neighborhood.
In the whole-space variational problem
\cite[Proposition~7.1]{SayasSelgas2014}, the curve enters through the
trace operator and the normalization functionals used to fix the additive
constant in the velocity. The trace operator is controlled by these bounds;
the normalization functionals are uniformly bounded by choosing one enclosing circle in
\cite[(7.2)--(7.4)]{SayasSelgas2014}. The whole-space Korn and divergence
constants are independent of the curve. A finite cover of the compact
closure gives~\eqref{eqn:stokes-layer-bound}.
\end{proof}

\subsection{Higher-order boundary integral estimates}
\label{app:higher-order-estimates}

The next higher-order divided-difference bound controls
$\p_s N(\kappa)$.

\begin{lemma}\label{lem:K-higher-order}
Let $d\ge1$ and $k\ge2$ be integers, let
$\alpha\in(0,1)$, and let $0\le\varepsilon<\alpha$. Let
$f\in C^{k,\alpha}(I;\R^d)$, let
$U\subset\R^d$ be open, and let $\phi\in C^\infty(U)$. Define the divided
differences
\begin{equation}\label{eq:Vj-def}
    V(s,\eta) \coloneqq \frac{f(s)-f(\eta)}{s-\eta},
    \quad
    V_j(s,\eta) \coloneqq \frac{f^{(j)}(s)-f^{(j)}(\eta)}{s-\eta},
    \quad 0\le j\le k,
\end{equation}
so that $V_0 = V$. Extend $V$ continuously to the diagonal by
$V(s,s)=f'(s)$, and assume that
$\overline{V(I\times I)}\Subset U$. The integral operator $K$ is defined by
\begin{equation}\label{eq:K-def}
    (K g)(s) \coloneqq \int_I \p_\eta\phi\bigl(V(s,\eta)\bigr)\,g(\eta)\,d\eta,
    \quad s\in I.
\end{equation}
Then $K:\wt L^2(I)\to\hi{k-1+\varepsilon}$ is bounded, with
\begin{equation}\label{eq:K-bound}
    \norm{K g}_{\hi{k-1+\varepsilon}}
    \le C\,\norm{g}_{L^2(\R)},
\end{equation}
where $C$ depends on $d,k,\varepsilon,\alpha$,
$\norm{f}_{C^{k,\alpha}(I)}$, and the derivatives of $\phi$
through order $k+1$ on a fixed neighborhood of
$\overline{V(I\times I)}$.
\end{lemma}

\begin{proof}
We first derive pointwise and H\"older bounds for $V_j$.
For $0\le \ell\le k-1$ and $n=k-\ell\ge 1$, the Leibniz rule applied to
$V_\ell(s,\eta)=[f^{(\ell)}(s)-f^{(\ell)}(\eta)]/(s-\eta)$ gives
\begin{equation}\label{eq:Vl-leibniz}
    \p_s^n V_\ell(s,\eta)
    = \bigl(f^{(\ell)}(s)-f^{(\ell)}(\eta)\bigr)
      \frac{(-1)^n n!}{(s-\eta)^{n+1}}
      +\sum_{j=1}^{n}\binom{n}{j}f^{(\ell+j)}(s)
      \frac{(-1)^{n-j}(n-j)!}{(s-\eta)^{n-j+1}}.
\end{equation}
The integral Taylor remainder for $f^{(\ell)}\in\ci{n,\alpha}$ at $s$ is
\begin{equation}\label{eq:remainder}
    f^{(\ell)}(\eta)-\sum_{j=0}^{n}\frac{f^{(\ell+j)}(s)}{j!}(\eta-s)^j
    = \frac{(\eta-s)^n}{(n-1)!}
      \int_0^1(1-r)^{n-1}\bigl(f^{(k)}(s+r(\eta-s))-f^{(k)}(s)\bigr)\,dr,
\end{equation}
Combining \eqref{eq:remainder} and \eqref{eq:Vl-leibniz} gives
\begin{equation}\label{eq:Vl-formula}
    \p_s^n V_\ell(s,\eta)
    = \frac{(-1)^{n+1}\,n!}{(s-\eta)^{n+1}}
      \Bigl[\,f^{(\ell)}(\eta)
      -\sum_{j=0}^{n}\frac{f^{(\ell+j)}(s)}{j!}(\eta-s)^j\Bigr].
\end{equation}
Since $f^{(k)}\in\ci{0,\alpha}$, the integrand in \eqref{eq:remainder} is bounded by
$C\abs{r(\eta-s)}^\alpha\le C\abs{s-\eta}^\alpha$, giving
\begin{equation}\label{eq:remainder-bound}
    \Bigl\lvert f^{(\ell)}(\eta)-\sum_{j=0}^{n}\frac{f^{(\ell+j)}(s)}{j!}(\eta-s)^j
    \Bigr\rvert \le C\,\abs{s-\eta}^{n+\alpha}.
\end{equation}
Substitution into \eqref{eq:Vl-formula} gives
\begin{equation}\label{eq:Vl-ptwise}
    \abs{\p_s^{k-\ell}V_\ell(s,\eta)} \le C\,\abs{s-\eta}^{\alpha-1},
    \quad 0\le \ell\le k.
\end{equation}
For $\ell=k$, this follows directly from the
$C^{0,\alpha}$-seminorm of $f^{(k)}$.
For the H\"older bound, \eqref{eq:Vl-formula} gives
\begin{equation}\label{eq:Vl-intrep}
    \p_s^n V_\ell(s,\eta)
    = -\frac{k-\ell}{s-\eta}
      \int_0^1(1-r)^{n-1}
      \bigl(f^{(k)}(s+r(\eta-s))-f^{(k)}(s)\bigr)\,dr.
\end{equation}
For $s,\eta,s+h\in I$ with
$\abs{h}\le\tfrac{1}{2}\abs{s-\eta}$, split
$\p_s^n V_\ell(s+h,\eta)-\p_s^n V_\ell(s,\eta)$ as
\begin{align}
    &\p_s^n V_\ell(s+h,\eta)-\p_s^n V_\ell(s,\eta)\notag\\
    &= \underbrace{\Bigl(\frac{1}{s+h-\eta}-\frac{1}{s-\eta}\Bigr)
      c_n\int_0^1(1-r)^{n-1}
      \bigl(f^{(k)}(s+r(\eta-s))-f^{(k)}(s)\bigr)\,dr}_{=: I_1}
    \notag\\
    &\quad
    +\frac{c_n}{s+h-\eta}\int_0^1(1-r)^{n-1}
    \Bigl[\bigl(f^{(k)}(s+h+r(\eta-s-h))-f^{(k)}(s+h)\bigr)\notag\\
    &\hspace{10em}
    -\bigl(f^{(k)}(s+r(\eta-s))-f^{(k)}(s)\bigr)\Bigr]\,dr,
    \label{eq:Vl-split}
\end{align}
where $c_n=-(k-\ell)$. Since $\abs{s+h-\eta}\ge\tfrac{1}{2}\abs{s-\eta}$,
\begin{equation}
    \Bigl\lvert\frac{1}{s+h-\eta}-\frac{1}{s-\eta}\Bigr\rvert
    \le \frac{C\abs{h}}{\abs{s-\eta}^2},
\end{equation}
and the integral is $O(\abs{s-\eta}^\alpha)$ by $f^{(k)}\in\ci{0,\alpha}$, so
$\abs{I_1}\le C\abs{h}\abs{s-\eta}^{\alpha-2}
\le C\abs{h}^\alpha\abs{s-\eta}^{-1}$
(using $\abs{h}\le\tfrac{1}{2}\abs{s-\eta}$).
For the second integral, $\abs{s+h-\eta}^{-1}\le C\abs{s-\eta}^{-1}$.
For each $r\in[0,1]$, pairing the evaluations at $s+h+r(\eta-s-h)$ and
$s+r(\eta-s)$ gives an argument difference $h(1-r)$; the evaluations at $s+h$ and $s$
differ by $h$. H\"older continuity bounds both increments by
$C\abs{h}^\alpha$, giving a contribution
$C\abs{h}^\alpha\abs{s-\eta}^{-1}$. Thus
\begin{equation}\label{eq:Vl-holder}
    \abs{\p_s^{k-\ell}V_\ell(s+h,\eta)-\p_s^{k-\ell}V_\ell(s,\eta)}
    \le C\,\abs{h}^\alpha\,\abs{s-\eta}^{-1},
    \quad \abs{h}\le\tfrac{1}{2}\abs{s-\eta}.
\end{equation}
For $\ell=k$, writing the difference of the two quotients and using
$[f^{(k)}]_{C^{0,\alpha}}$ gives
\begin{equation}
    \abs{V_k(s+h,\eta)-V_k(s,\eta)}
    \le C\left(
        \frac{\abs h^\alpha}{\abs{s-\eta}}
        +\abs h\,\abs{s-\eta}^{\alpha-2}
    \right)
    \le C\abs h^\alpha\abs{s-\eta}^{-1},
\end{equation}
so \eqref{eq:Vl-holder} also holds in this case.

It remains to prove $K:\wt L^2(I)\to\hi{k-1+\varepsilon}$. It suffices
to show $\p_s^{k-1}(K g)\in\hi{\varepsilon}$.
Taking $k-1$ derivatives in $s$ in the sense of distributions, which is
justified by the integrable kernel bounds below, gives
\begin{equation}\label{eq:diff-under}
    \p_s^{k-1}(K g)(s)
    = \int_I k_\phi(s,\eta)\,g(\eta)\,d\eta,
    \quad
    k_\phi(s,\eta)\coloneqq \p_s^{k-1}\p_\eta\phi\bigl(V(s,\eta)\bigr).
\end{equation}
We apply Lemma~\ref{K-op} to $k_\phi$.

By the multivariable Faà di Bruno formula,
$k_\phi = \p_s^{k-1}\p_\eta\phi(V)$ is a finite sum of
terms of the form
\begin{equation}\label{eq:fdb}
    D^m\phi(V)\bigl[Q_1V,\ldots,Q_mV\bigr],
\end{equation}
where $Q_i=\p_s^{a_i}\p_\eta^{b_i}$ has positive order and
$\sum_{i=1}^m\operatorname{ord}(Q_i)=k$.
Set $D_+=\p_s+\p_\eta$. Since
$\p_\eta=D_+-\p_s$ and $D_+V_j=V_{j+1}$, every factor
$Q_iV$ reduces to
a linear combination of $\p_s^a V_b$ with $a+b=\operatorname{ord}(Q_i)$.
If $\operatorname{ord}(Q_i)<k$, then $\abs{Q_iV}$ is
bounded because the diagonal extension of $V$ belongs to
$C^{k-1}(I\times I)$.
Since the orders sum to $k$, at most one factor has order $k$; that
factor satisfies $\abs{Q_iV}\le C\abs{s-\eta}^{\alpha-1}$ by
\eqref{eq:Vl-ptwise}. All remaining factors and $D^m\phi(V)$ are bounded,
so
\begin{equation}\label{eq:L-ptwise}
    \abs{k_\phi(s,\eta)} \le C\,\abs{s-\eta}^{\alpha-1}.
\end{equation}

For the H\"older bound, let $s,\eta,s+h\in I$ satisfy
$\abs{h}\le\tfrac{1}{2}\abs{s-\eta}$. We claim that
\begin{equation}\label{eq:L-holder}
    \abs{k_\phi(s+h,\eta)-k_\phi(s,\eta)} \le C\,\abs{h}^\alpha\,\abs{s-\eta}^{-1}.
\end{equation}
Let $Q$ be a mixed derivative of total order $r<k$. The
fundamental theorem of calculus and \eqref{eq:Vl-ptwise} give
\begin{equation}\label{eq:lower-factor-holder}
    \abs{QV(s+h,\eta)-QV(s,\eta)}
    \le
    \begin{cases}
        C\abs h, & r\le k-2,\\
        C\abs h\,\abs{s-\eta}^{\alpha-1}, & r=k-1,
    \end{cases}
    \le C\abs h^\alpha\abs{s-\eta}^{-1}.
\end{equation}
Here we used
$\abs{s+\rho h-\eta}\ge\frac12\abs{s-\eta}$ for
$0\le\rho\le1$ and the boundedness of $I$.
A telescoping expansion of \eqref{eq:fdb} and
\eqref{eq:lower-factor-holder} gives \eqref{eq:L-holder} when all factors
have order less than $k$. If one factor has order $k$, it is the only
factor. Since $k\ge2$, $V$ is uniformly Lipschitz in $s$, and
\begin{equation}
    \abs{\Delta_{s,h}\!\left(D\phi(V)[QV]\right)}
    \le C\abs h\,\abs{s-\eta}^{\alpha-1}
    +C\abs h^\alpha\abs{s-\eta}^{-1}
    \le C\abs h^\alpha\abs{s-\eta}^{-1}.
\end{equation}
Hence \eqref{eq:L-holder} also holds for the top-order term.

\noindent
Lemma~\ref{K-op} and \eqref{eq:L-ptwise}--\eqref{eq:L-holder} give
\begin{equation}\label{eq:step1-bound}
    \norm{\p_s^{k-1}(K g)}_{\hi{\varepsilon}} \le C\,\norm{g}_{L^2(\R)}.
\end{equation}
Since $k\ge2$, the kernel $\p_\eta\phi(V)$ is bounded, and the
Schur test gives
$\norm{K g}_{L^2(I)}\le C\norm{g}_{L^2(\R)}$.
The one-dimensional norm equivalence
\begin{equation}
    \norm{u}_{H^{k-1+\varepsilon}(I)}
    \simeq \norm{u}_{L^2(I)}
    +\norm{\p_s^{k-1}u}_{H^\varepsilon(I)}
\end{equation}
and \eqref{eq:step1-bound} give \eqref{eq:K-bound}.
\end{proof}

\subsection{Difference estimates for the nonlinear terms}
\label{app:nonlinear-estimates}

\begin{proof}[Proof of Lemma~\ref{R-Lip}]
Set $\alpha=\frac{1}{2}\paren{\varepsilon+\delta}$, so that
$\varepsilon<\alpha<\delta$. Every interpolated curve
$\rho\xka{1}+(1-\rho)\xka{2}$ has speed at most one, so its arc-chord
constant is at most one. The assumed lower bound gives
$0<m\le 1$.
First estimate the commutator part of $R_\kappa$ in~\eqref{eqn:Rop-def}. Lemmas~\ref{lem:Xk-holder},
\ref{tau-n-est}, \ref{dsH_f_comm}, and~\ref{point_mult_sobo} give
\begin{equation}
\begin{aligned}
    &\norm{ \tka{1}\cdot[\p_s\mc H,\tka{1}]\sigma - \tka{2}\cdot[\p_s\mc H,\tka{2}]\sigma }_{\hi{-\frac{1}{2}+\varepsilon}} \\
    & = \norm{ \tka{1}\cdot[\p_s\mc H,\paren{\tka{1}-\tka{2}}]\sigma + \paren{\tka{1}-\tka{2}}\cdot[\p_s\mc H,\tka{2}]\sigma }_{\hi{-\frac{1}{2}+\varepsilon}}\\
    &\le C \norm{\tka{1}}_{\hi{\frac{1}{2}+\delta}}\paren{\norm{\p_s\tka{1}-\p_s\tka{2}}_{\hi{-\frac{1}{2}+\delta}}+ [\tka{1}-\tka{2}]_{\ci{0,\alpha}}}\norm{\sigma}_{\hr{\frac{1}{2}}} \\
    & + C\norm{\tka{1}-\tka{2}}_{\hi{\frac{1}{2}+\delta}}\paren{\norm{\p_s\tka{2}}_{\hi{-\frac{1}{2}+\delta}}+ [\tka{2}]_{\ci{0,\alpha}}} \norm{\sigma}_{\hr{\frac{1}{2}}} \\
    &\le C_{M,m} \norm{\kappa_1-\kappa_2}_{\hr{-\frac{1}{2}+\delta}}\norm{\sigma}_{\hr{\frac{1}{2}}}.
\end{aligned}
\end{equation}
For the Stokes-remainder part, write
\begin{equation}
\begin{aligned}
    &\tka{1}\cdot\p_s\int_I \p_\eta r(s,\eta;\kappa_1)\sigma(\eta)\tka{1}(\eta)\,d\eta  \\
    &\quad - \tka{2}\cdot\p_s\int_I \p_\eta r(s,\eta;\kappa_2)\sigma(\eta)\tka{2}(\eta)\,d\eta = I_1 + I_2 +
    I_3,
\end{aligned}
\end{equation}
where
\begin{align}
I_1 & = \paren{\tka{1}-\tka{2}}\cdot\p_s\int_I \p_\eta r(s,\eta;\kappa_1)\sigma(\eta)\tka{1}(\eta)\,d\eta, \\
I_2 &= \tka{2}\cdot\p_s\int_I \p_\eta \paren{r(s,\eta;\kappa_1)-r(s,\eta;\kappa_2)}\sigma(\eta)\tka{1}(\eta)\,d\eta, \\
I_3 &= \tka{2}\cdot\p_s\int_I \p_\eta r(s,\eta;\kappa_2)\sigma(\eta)\paren{\tka{1}(\eta)-\tka{2}(\eta)} \,d\eta.
\end{align}
Equations~\eqref{eqn:Sop-bound}, \eqref{eqn:product-supported}, and~\eqref{eqn:product-restriction} give
\begin{equation}
\begin{aligned}
    \norm{I_1}_{\hi{-\frac{1}{2}+\varepsilon}} &\le  \abs{\xka{1}}_*^{-2}\norm{\tka{1}}_{\ci{0,\alpha}}[\tka{1}]_{\ci{0,\alpha}} \norm{\tka{1}-\tka{2}}_{\hi{\frac{1}{2}+\delta}}\norm{\tka{1}}_{\hi{\frac{1}{2}+\delta}} \norm{\sigma}_{\hr{\frac{1}{2}}}\\
    &\le C_{M,m}M \norm{\kappa_1-\kappa_2}_{\hr{-\frac{1}{2}+\delta}}\norm{\sigma}_{\hr{\frac{1}{2}}}.
\end{aligned}
\end{equation}
Similarly,
\begin{equation}
\begin{aligned}
    \norm{I_3}_{\hi{-\frac{1}{2}+\varepsilon}} & \le  \abs{\xka{2}}_*^{-2}\norm{\tka{2}}_{\ci{0,\alpha}}[\tka{2}]_{\ci{0,\alpha}} \norm{\tka{1}-\tka{2}}_{\hi{\frac{1}{2}+\delta}}\norm{\tka{2}}_{\hi{\frac{1}{2}+\delta}} \norm{\sigma}_{\hr{\frac{1}{2}}}\\
    & \le C_{M,m}M \norm{\kappa_1-\kappa_2}_{\hr{-\frac{1}{2}+\delta}}\norm{\sigma}_{\hr{\frac{1}{2}}}.
\end{aligned}
\end{equation}
Lemma~\ref{lem:Sop-Lip} gives
\begin{equation}
\begin{aligned}
    \norm{I_2}_{\hi{-\frac{1}{2}+\varepsilon}} & \le C_{M,m} \norm{\tka{2}}_{\hi{\frac{1}{2}+\delta}} \norm{\kappa_1-\kappa_2}_{\hr{-\frac{1}{2}+\delta}} \norm{\tka{1}}_{\hi{\frac{1}{2}+\delta}} \norm{\sigma}_{\hr{\frac{1}{2}}} \\
    & \le C_{M,m} \norm{\kappa_1-\kappa_2}_{\hr{-\frac{1}{2}+\delta}} \norm{\sigma}_{\hr{\frac{1}{2}}}.
\end{aligned}
\end{equation}
Combining the commutator estimate with the bounds for $I_1,I_2,I_3$
proves the result.
\end{proof}

\begin{proof}[Proof of Lemma~\ref{F-Lip-kappa}]
Set $\alpha=\frac{1}{2}\paren{\varepsilon+\delta}$, so that
$\varepsilon<\alpha<\delta$. Each interpolated curve has speed at most one.
Thus the assumed arc-chord lower bound gives $0<m\le 1$.
Split the commutator difference as
\begin{equation}
\begin{aligned}
    [\p_s\mc H,  \tka{1}]\cdot (\p_s \kappa_1 \nka{1}) - [\p_s\mc H, \tka{2}]\cdot (\p_s \kappa_2 \nka{2}) = I_1 + I_2 + I_3,
\end{aligned}
\end{equation}
where
\begin{align}
    I_1 &= [\p_s\mc H,  \tka{1}-\tka{2}]\cdot (\p_s \kappa_1 \nka{1}) , \\
    I_2 &= [\p_s\mc H,  \tka{2}]\cdot ( (\p_s \kappa_1 -\p_s\kappa_2)\nka{1}) , \\
    I_3 &= [\p_s\mc H,  \tka{2}]\cdot (\p_s \kappa_2( \nka{1} - \nka{2})) .
\end{align}
Since $\nka{i}=Q_{\frac{\pi}{2}}^{-1}\xka{i}'$ for $i=1,2$, we have
\begin{equation}
    \norm{\nka{i}}_{\hi{\frac{1}{2}+\delta}} = \norm{\xka{i}'}_{\hi{\frac{1}{2}+\delta}},
    \quad
    \norm{\nka{1}-\nka{2}}_{\hi{\frac{1}{2}+\delta}}
    = \norm{\xka{1}'-\xka{2}'}_{\hi{\frac{1}{2}+\delta}}.
\end{equation}
Lemmas~\ref{lem:Xk-holder}, \ref{tau-n-est}, \ref{dsH_f_comm},
and~\ref{point_mult_sobo} give
\begin{equation}
\begin{aligned}
    \norm{I_1}_{\hi{-\frac{1}{2}+\varepsilon}} &\le C \paren{ [\tka{1}-\tka{2}]_{\ci{0,\alpha}} + \norm{\p_s(\tka{1}-\tka{2})}_{\hi{-\frac{1}{2}+\delta}}  } \norm{\p_s \kappa_1 \nka{1}}_{\hr{\frac{1}{2}}}\\
    % & \le C\paren{1+\norm{\kappa_2}_{\hr{-\frac{1}{2}+\delta}}} \norm{\kappa_1-\kappa_2}_{\hr{-\frac{1}{2}+\delta}}\norm{\p_s\kappa_1}_{\hr{\frac{1}{2}}} \norm{\nka{1}}_{\hi{\frac{1}{2}+\delta}} \\
    &\le C_{M,m} \norm{\kappa_1-\kappa_2}_{\hr{-\frac{1}{2}+\delta}}\norm{\kappa_1}_{\hr{\frac{3}{2}}}.
\end{aligned}
\end{equation}
\begin{equation}
\begin{aligned}
    \norm{I_2}_{\hi{-\frac{1}{2}+\varepsilon}} &\le C \paren{ [\tka{2}]_{\ci{0,\alpha}} + \norm{\p_s\tka{2}}_{\hi{-\frac{1}{2}+\delta}}  } \norm{\p_s (\kappa_1-\kappa_2) \nka{1}}_{\hr{\frac{1}{2}}}\\
    % & \le C \norm{\kappa_2}_{\hr{-\frac{1}{2}+\delta}} \norm{\p_s(\kappa_1-\kappa_2)}_{\hr{\frac{1}{2}}} \norm{\nka{1}}_{\hi{\frac{1}{2}+\delta}} \\
    &\le C_{M,m}M \norm{\kappa_1-\kappa_2}_{\hr{\frac{3}{2}}}.
\end{aligned}
\end{equation}
\begin{equation}
\begin{aligned}
    \norm{I_3}_{\hi{-\frac{1}{2}+\varepsilon}} &\le C \paren{ [\tka{2}]_{\ci{0,\alpha}} + \norm{\p_s\tka{2}}_{\hi{-\frac{1}{2}+\delta}}  } \norm{\p_s \kappa_2 (\nka{1}-\nka{2}) }_{\hr{\frac{1}{2}}}\\
    % & \le C\norm{\kappa_2}_{\hr{-\frac{1}{2}+\delta}}\norm{\p_s\kappa_2}_{\hr{\frac{1}{2}}} \norm{\nka{1} - \nka{2}}_{\hi{\frac{1}{2}+\delta}} \\
    &\le C_{M,m}M \norm{\kappa_1-\kappa_2}_{\hr{-\frac{1}{2}+\delta}}\norm{\kappa_2}_{\hr{\frac{3}{2}}}.
\end{aligned}
\end{equation}
\begin{equation}
\begin{aligned}
    \norm{\sum_{i=1}^3 I_i}_{\hi{-\frac{1}{2}+\varepsilon}}
    &\le C_{M,m}
    \norm{\kappa_1-\kappa_2}_{\hr{-\frac{1}{2}+\delta}}
    \paren{
        \norm{\kappa_1}_{\hr{\frac{3}{2}}}
        +\norm{\kappa_2}_{\hr{\frac{3}{2}}}
    } \\
    &\quad+ C_{M,m}M
    \norm{\kappa_1-\kappa_2}_{\hr{\frac{3}{2}}}.
\end{aligned}
\end{equation}
For the Stokes-remainder difference, write
\begin{equation}
\begin{aligned}
    &\tka{1}\cdot\p_s \int_I\p_\eta r(s,\eta;\kappa_1)\p_\eta\kappa_1\nka{1}(\eta) \,d\eta  \\
    &\quad - \tka{2}\cdot\p_s \int_I\p_\eta r(s,\eta;\kappa_2)\p_\eta\kappa_2\nka{2}(\eta) \,d\eta = I_1' + I_2'
    + I_3' + I_4'
\end{aligned}
\end{equation}
where
\begin{align}
    I_1' &= (\tka{1}-\tka{2})\cdot\p_s \int_I\p_\eta r(s,\eta;\kappa_1)\p_\eta\kappa_1\nka{1}(\eta) \,d\eta ,\\
    I_2' &= \tka{2}\cdot\p_s \int_I(\p_\eta r(s,\eta;\kappa_1)-\p_\eta r(s,\eta;\kappa_2))\p_\eta\kappa_1\nka{1}(\eta) \,d\eta ,\\
    I_3' &= \tka{2}\cdot\p_s \int_I\p_\eta r(s,\eta;\kappa_2)(\p_\eta\kappa_1-\p_\eta\kappa
    _2)\nka{1}(\eta) \,d\eta,\\
    I_4' &= \tka{2}\cdot\p_s \int_I\p_\eta r(s,\eta;\kappa_2)\p_\eta\kappa_2(\nka{1}-\nka{2})(\eta) \,d\eta.
\end{align}
Lemmas~\ref{lem:Sop-est} and~\ref{lem:Sop-Lip} give
\begin{equation}
\begin{aligned}
    \norm{I_1'}_{\hi{-\frac{1}{2}+\varepsilon}} &\le \norm{\tka{1}-\tka{2}}_{\hi{\frac{1}{2}+\delta}} \norm{\int_I\p_\eta r(s,\eta;\kappa_1)\p_\eta\kappa_1\nka{1}(\eta) \,d\eta}_{\hi{\frac{1}{2}+\varepsilon}} \\
    & \le C_{M,m}M \norm{\kappa_1-\kappa_2}_{\hr{-\frac{1}{2}+\delta}} \norm{\p_s \kappa_1}_{\hr{\frac{1}{2}}} \norm{\nka{1}}_{\hi{\frac{1}{2}+\delta}}\\
    & \le C_{M,m}M \norm{\kappa_1-\kappa_2}_{\hr{-\frac{1}{2}+\delta}}\norm{\kappa_1}_{\hr{\frac{3}{2}}},
\end{aligned}
\end{equation}
\begin{equation}
\begin{aligned}
    \norm{I_2'}_{\hi{-\frac{1}{2}+\varepsilon}} &\le C_{M,m}\norm{\tka{2}}_{\hi{\frac{1}{2}+\delta}} \norm{\kappa_1-\kappa_2}_{\hr{-\frac{1}{2}+\delta}}\norm{\p_s\kappa_1}_{\hr{\frac{1}{2}}} \norm{\nka{1}}_{\hi{\frac{1}{2}+\delta}} \\
    &\le C_{M,m} \norm{\kappa_1-\kappa_2}_{\hr{-\frac{1}{2}+\delta}}\norm{\kappa_1}_{\hr{\frac{3}{2}}},
\end{aligned}
\end{equation}
\begin{equation}
\begin{aligned}
    \norm{I_3'}_{\hi{-\frac{1}{2}+\varepsilon}} & \le C_{M,m}M \norm{\tka{2}}_{\hi{\frac{1}{2}+\delta}}  \norm{\p_s(\kappa_1-\kappa_2)}_{\hr{\frac{1}{2}}} \norm{\nka{1}}_{\hi{\frac{1}{2}+\delta}}\\
    &\le C_{M,m}M \norm{\kappa_1-\kappa_2}_{\hr{\frac{3}{2}}},
\end{aligned}
\end{equation}
\begin{equation}
\begin{aligned}
    \norm{I_4'}_{\hi{-\frac{1}{2}+\varepsilon}}
    &\le C_{M,m}M \norm{\tka{2}}_{\hi{\frac{1}{2}+\delta}}  \norm{\p_s\kappa_2 }_{\hr{\frac{1}{2}}} \norm{\nka{1}-\nka{2}}_{\hi{\frac{1}{2}+\delta}}\\
    &\le C_{M,m}M \norm{\kappa_1-\kappa_2}_{\hr{-\frac{1}{2}+\delta}}\norm{\kappa_2}_{\hr{\frac{3}{2}}}.
\end{aligned}
\end{equation}
\begin{equation}
\begin{aligned}
    \norm{\sum_{i=1}^4 I_i'}_{\hi{-\frac{1}{2}+\varepsilon}}
    &\le C_{M,m}
    \norm{\kappa_1-\kappa_2}_{\hr{-\frac{1}{2}+\delta}}
    \paren{
        \norm{\kappa_1}_{\hr{\frac{3}{2}}}
        +\norm{\kappa_2}_{\hr{\frac{3}{2}}}
    } \\
    &\quad+ C_{M,m}M
    \norm{\kappa_1-\kappa_2}_{\hr{\frac{3}{2}}}.
\end{aligned}
\end{equation}
The two decompositions give the estimate.
\end{proof}

\begin{proof}[Proof of~\eqref{eqn:N-pointwise-difference}]
Set $\alpha=\frac{1}{2}\paren{\varepsilon+\delta}$, so that
$\varepsilon<\alpha<\delta$. Each interpolated curve has speed at most one,
so the assumed arc-chord lower bound gives $0<m\le1$.
By~\eqref{eqn:inverse-uniform-bound}, the inverse constants $C_{\kappa_i}$
can be absorbed into $C_{M,m}$.
Whenever a frame field is multiplied by a function supported in
$\overline I$, we use a fixed bounded Sobolev extension of the frame field
to $\mathbb R$. The resulting function remains supported in $\overline I$.
Using~\eqref{eqn:N-def}, split
\begin{equation}
    N(\kappa_2) - N(\kappa_1) = \frac{1}{4} \sum_{i=1}^3 I_i +\frac{1}{4\pi} \sum_{i=1}^4 I_i'.
\end{equation}
The commutator terms are
\begin{align}
    I_1 &=[\p_s\mc H, \nka{1}-\nka{2}]\cdot\paren{\p_s\kappa_1\nka{1} +\sigma_1\tka{1}},\\
    I_2 &= [\p_s\mc H, \nka{2}]\cdot\paren{(\p_s\kappa_1-\p_s\kappa_2)\nka{1} + (\sigma_1-\sigma_2)\tka{1}},\\
    I_3 &= [\p_s\mc H, \nka{2}]\cdot\paren{\p_s\kappa_2(\nka{1}-\nka{2}) +\sigma_2(\tka{1}-\tka{2})}.
\end{align}
The Stokes-remainder terms are
\begin{align}
    I_1' &= (\nka{1}-\nka{2})\cdot\p_s\int_I \p_\eta r(s,\eta;\kappa_1) \paren{ \p_\eta\kappa_1\nka{1}+\sigma_1\tka{1} } \,d\eta, \\
    I_2' &= \nka{2}\cdot\p_s\int_I \p_\eta (r(s,\eta;\kappa_1)-r(s,\eta;\kappa_2)) \paren{ \p_\eta\kappa_1\nka{1}+\sigma_1\tka{1} } \,d\eta,\\
    I_3' &= \nka{2}\cdot\p_s\int_I \p_\eta r(s,\eta;\kappa_2) \paren{ \p_\eta(\kappa_1-\kappa_2)\nka{1}+(\sigma_1-\sigma_2)\tka{1} } \,d\eta,\\
    I_4' &= \nka{2}\cdot\p_s\int_I \p_\eta r(s,\eta;\kappa_2) \paren{ \p_\eta\kappa_2(\nka{1}-\nka{2})+\sigma_2(\tka{1}-\tka{2}) } \,d\eta.
\end{align}
Lemmas~\ref{tau-n-est} and~\ref{lem:Xk-holder} also give, for $i=1,2$,
\begin{equation}\label{eqn:N-frame-bounds}
    \norm{\nka{i}}_{\hi{\frac{1}{2}+\delta}}
    +\norm{\tka{i}}_{\hi{\frac{1}{2}+\delta}}
    \le C_{M,m},\quad
    [\nka{i}]_{\ci{0,\alpha}}
    +\norm{\p_s\nka{i}}_{\hi{-\frac{1}{2}+\delta}}
    \le C_{M,m}M,
\end{equation}
and
\begin{equation}\label{eqn:N-frame-difference}
\begin{aligned}
    &[\nka{1}-\nka{2}]_{\ci{0,\alpha}} +\norm{\p_s(\nka{1}-\nka{2})}_{\hi{-\frac{1}{2}+\delta}} \\
    &\quad +\norm{\nka{1}-\nka{2}}_{\hi{\frac{1}{2}+\delta}} +\norm{\tka{1}-\tka{2}}_{\hi{\frac{1}{2}+\delta}}
    \le C_{M,m} \norm{\kappa_1-\kappa_2}_{\hr{-\frac{1}{2}+\delta}}.
\end{aligned}
\end{equation}
Using the bounded frame extensions described above,
Lemma~\ref{point_mult_sobo} gives
\begin{equation}\label{eqn:N-density-bounds}
\begin{aligned}
    \norm{\p_s\kappa_1\nka{1}+\sigma_1\tka{1}}_{\hr{\frac{1}{2}}}
    &\le C_{M,m}
    \paren{
        \norm{\kappa_1}_{\hr{\frac{3}{2}}}
        +\norm{\sigma_1}_{\hr{\frac{1}{2}}}
    }, \\
    \norm{
        \p_s(\kappa_1-\kappa_2)\nka{1}
        +(\sigma_1-\sigma_2)\tka{1}
    }_{\hr{\frac{1}{2}}}
    &\le C_{M,m}
    \paren{
        \norm{\kappa_1-\kappa_2}_{\hr{\frac{3}{2}}}
        +\norm{\sigma_1-\sigma_2}_{\hr{\frac{1}{2}}}
    }, \\
    \norm{
        \p_s\kappa_2(\nka{1}-\nka{2})
        +\sigma_2(\tka{1}-\tka{2})
    }_{\hr{\frac{1}{2}}}
    &\le C_{M,m}
    \norm{\kappa_1-\kappa_2}_{\hr{-\frac{1}{2}+\delta}} \\
    &\quad\times
    \paren{
        \norm{\kappa_2}_{\hr{\frac{3}{2}}}
        +\norm{\sigma_2}_{\hr{\frac{1}{2}}}
    }.
\end{aligned}
\end{equation}
Applying Lemma~\ref{dsH_f_comm} to $I_1$, $I_2$, and $I_3$, and then
using~\eqref{eqn:tension-difference} and~\eqref{eqn:tension-bound}, gives
\begin{equation}
\begin{aligned}
    \norm{I_1}_{\hi{-\frac{1}{2}+\varepsilon}} &\le C_{M,m} \norm{\kappa_1-\kappa_2}_{\hr{-\frac{1}{2}+\delta}} \paren{ \norm{\kappa_1}_{\hr{\frac{3}{2}}} + \norm{\sigma_1}_{\hr{\frac{1}{2}}} }  \\
    & \le C_{M,m} \norm{\kappa_1-\kappa_2}_{\hr{-\frac{1}{2}+\delta}} \norm{\kappa_1}_{\hr{\frac{3}{2}}}
\end{aligned}
\end{equation}

\begin{equation}
\begin{aligned}
    \norm{I_2}_{\hi{-\frac{1}{2}+\varepsilon}} &\le C_{M,m}M \paren{ \norm{\kappa_1-\kappa_2}_{\hr{\frac{3}{2}}}
    +\norm{\sigma_1-\sigma_2}_{\hr{\frac{1}{2}}} } \le C_{M,m}M \norm{\kappa_1-\kappa_2}_{\hr{\frac{3}{2}}} \\
    &\quad + C_{M,m}M \norm{\kappa_1-\kappa_2}_{\hr{-\frac{1}{2}+\delta}} \paren{
    \norm{\kappa_1}_{\hr{\frac{3}{2}}} +\norm{\kappa_2}_{\hr{\frac{3}{2}}} }
\end{aligned}
\end{equation}

\begin{equation}
\begin{aligned}
    \norm{I_3}_{\hi{-\frac{1}{2}+\varepsilon}} &\le C_{M,m}M \norm{\kappa_1-\kappa_2}_{\hr{-\frac{1}{2}+\delta}} \paren{ \norm{\kappa_2}_{\hr{\frac{3}{2}}} +  \norm{\sigma_2}_{\hr{\frac{1}{2}}}  }  \\
    & \le C_{M,m}M \norm{\kappa_1-\kappa_2}_{\hr{-\frac{1}{2}+\delta}} \norm{\kappa_2}_{\hr{\frac{3}{2}}}
\end{aligned}
\end{equation}

\begin{equation}
\begin{aligned}
    \norm{\sum_{i=1}^3 I_i}_{\hi{-\frac{1}{2}+\varepsilon}} &\le C_{M,m} \norm{\kappa_1-\kappa_2}_{\hr{-\frac{1}{2}+\delta}} \paren{\norm{\kappa_1}_{\hr{\frac{3}{2}}}+\norm{\kappa_2}_{\hr{\frac{3}{2}}}} \\
    &+ C_{M,m}M \norm{\kappa_1-\kappa_2}_{\hr{\frac{3}{2}}}.
\end{aligned}
\end{equation}
For the remainder terms, we use~\eqref{eqn:N-frame-bounds}--\eqref{eqn:N-density-bounds}, the
continuity of
$\p_s:H^{\frac{1}{2}+\varepsilon}(I)\to
H^{-\frac{1}{2}+\varepsilon}(I)$, and~\eqref{eqn:Sop-bound}, \eqref{eqn:Sop-difference}, \eqref{eqn:tension-bound}, and~\eqref{eqn:tension-difference}. These bounds give
\begin{equation}
\begin{aligned}
    \norm{I_1'}_{\hi{-\frac{1}{2}+\varepsilon}} &\le C_{M,m}M \norm{\kappa_1-\kappa_2}_{\hr{-\frac{1}{2}+\delta}} \paren{ \norm{\kappa_1}_{\hr{\frac{3}{2}}} + \norm{\sigma_1}_{\hr{\frac{1}{2}}} } \\
    &\le C_{M,m}M \norm{\kappa_1-\kappa_2}_{\hr{-\frac{1}{2}+\delta}}  \norm{\kappa_1}_{\hr{\frac{3}{2}}}
\end{aligned}
\end{equation}
\begin{equation}
\begin{aligned}
    \norm{I_2'}_{\hi{-\frac{1}{2}+\varepsilon}} &\le C_{M,m} \norm{\kappa_1-\kappa_2}_{\hr{-\frac{1}{2}+\delta}}  \paren{ \norm{\kappa_1}_{\hr{\frac{3}{2}}} + \norm{\sigma_1}_{\hr{\frac{1}{2}}} } \\
    &\le C_{M,m} \norm{\kappa_1-\kappa_2}_{\hr{-\frac{1}{2}+\delta}} \norm{\kappa_1}_{\hr{\frac{3}{2}}},
\end{aligned}
\end{equation}
\begin{equation}
\begin{aligned}
    \norm{I_3'}_{\hi{-\frac{1}{2}+\varepsilon}} &\le C_{M,m}M \paren{\norm{\kappa_1-\kappa_2}_{\hr{\frac{3}{2}}}
    + \norm{\sigma_1-\sigma_2}_{\hr{\frac{1}{2}}}} \le C_{M,m}M \norm{\kappa_1-\kappa_2}_{\hr{\frac{3}{2}}} \\
    &\quad + C_{M,m}M \norm{\kappa_1-\kappa_2}_{\hr{-\frac{1}{2}+\delta}} \paren{
    \norm{\kappa_1}_{\hr{\frac{3}{2}}} +\norm{\kappa_2}_{\hr{\frac{3}{2}}} }.
\end{aligned}
\end{equation}
\begin{equation}
\begin{aligned}
    \norm{I_4'}_{\hi{-\frac{1}{2}+\varepsilon}} &\le C_{M,m}M \norm{\kappa_1-\kappa_2}_{\hr{-\frac{1}{2}+\delta}} \paren{\norm{\kappa_2}_{\hr{\frac{3}{2}}} + \norm{\sigma_2}_{\hr{\frac{1}{2}}}} \\
    &\le C_{M,m}M \norm{\kappa_1-\kappa_2}_{\hr{-\frac{1}{2}+\delta}} \norm{\kappa_2}_{\hr{\frac{3}{2}}}
\end{aligned}
\end{equation}
Hence
\begin{equation}
\begin{aligned}
    \norm{\sum_{i=1}^4 I_i'}_{\hi{-\frac{1}{2}+\varepsilon}}
    &\le C_{M,m}M
    \norm{\kappa_1-\kappa_2}_{\hr{\frac{3}{2}}} \\
    &\quad
    + C_{M,m}
    \norm{\kappa_1-\kappa_2}_{\hr{-\frac{1}{2}+\delta}}
    \paren{
        \norm{\kappa_1}_{\hr{\frac{3}{2}}}
        +\norm{\kappa_2}_{\hr{\frac{3}{2}}}
    }.
\end{aligned}
\end{equation}
Combining the commutator and remainder estimates gives
\begin{equation}
\begin{aligned}
    \norm{N(\kappa_1) - N(\kappa_2)}_{\hi{-\frac{1}{2}+\varepsilon}} &\le C_{M,m} \norm{\kappa_1-\kappa_2}_{\hr{-\frac{1}{2}+\delta}} \paren{\norm{\kappa_1}_{\hr{\frac{3}{2}}}+\norm{\kappa_2}_{\hr{\frac{3}{2}}}} \\
    & + C_{M,m}M \norm{\kappa_1-\kappa_2}_{\hr{\frac{3}{2}}}.
\end{aligned}
\end{equation}
\end{proof}


\section*{Acknowledgments}
The author thanks Yoichiro Mori for suggesting the problem and for many helpful discussions.

\bibliographystyle{plain}
\bibliography{references}
\end{document}
